% SPDX-FileCopyrightText: 2026 Will Cook
% SPDX-License-Identifier: CC-BY-4.0
%
% GENERATED by scripts/assemble_reasoning_surfaces.py from the authored files
% under paper/reasoning-parts/erdos249. Parts are inlined, never \input, so
% that the Pandoc projection cannot silently drop a section.
\documentclass[11pt]{article}
% SPDX-FileCopyrightText: 2026 Will Cook
% SPDX-License-Identifier: CC-BY-4.0
%
% Shared preamble for the two per-problem reasoning-surface papers.
% These documents are written to be read in full as working context: every
% result appears with its exact hypotheses, its evidence band, the scale at
% which it is proved, and the Lean declaration that carries it.
\usepackage{paper-house-style}
\usepackage{array}
\usepackage{adjustbox}
\usepackage{booktabs}
\usepackage{longtable}
\usepackage{enumitem}
\setlist{topsep=0.35em,itemsep=0.2em}
\usepackage{xurl}
\usepackage{xspace}
\usepackage[colorlinks=true,linkcolor=linkinternal,urlcolor=linkexternal,
  citecolor=linkinternal]{hyperref}
\hypersetup{bookmarksnumbered=true,bookmarksopen=true,bookmarksopenlevel=1,
  pdftitle={The Binary Totient Series: a detailed map of the current attack on Erdos Problem 249},
  pdfauthor={Will Cook},pdflang={en-GB}}
% Lean declaration names are full of underscores and are typeset in text mode.
% Category codes are fixed when TeX READS a line, not when a macro expands, so
% no amount of \catcode juggling inside \lean can rescue an underscore that has
% already been tokenised as a subscript. This package makes _ print as itself in
% text while leaving it as a subscript in maths, which is exactly the split we
% need. It also protects any stray underscore in authored prose.
\usepackage{underscore}

% ---------------------------------------------------------------------------
% Theorem environments.
% ---------------------------------------------------------------------------
\theoremstyle{plain}
\newtheorem{thm}{Theorem}[section]
\newtheorem{lem}[thm]{Lemma}
\newtheorem{prop}[thm]{Proposition}
\newtheorem{cor}[thm]{Corollary}
\theoremstyle{definition}
\newtheorem{defn}[thm]{Definition}
\theoremstyle{remark}
\newtheorem{rem}[thm]{Remark}
\newtheorem{obs}[thm]{Observation}

% ---------------------------------------------------------------------------
% Citation of a Lean declaration. The declaration name is the authority; the
% file:line is where it sat at the pinned formal-source checkpoint. Underscores
% are handled locally so authors may write names verbatim.
% ---------------------------------------------------------------------------
\newcommand{\PK}{Erdos249257}
\newcommand{\commit}{99f4bf47422abbd8757cbb22b50ba079d764d3a7}
\newcommand{\repobase}{https://github.com/wcook04/plectis-lean-erdos249-257/blob/\commit}
% \lean declaration file:line forms a hyperlink to the pinned public
% formal-source checkpoint. The second argument is a machine coordinate,
% never typeset: \allowbreak and spaces are stripped, an ErdosProblems/
% prefix links the ErdosProblems tree directly (otherwise the \PK tree is
% used), a trailing :N becomes a GitHub #LN anchor, :N-M and :N--M become
% range anchors #LN-LM, and comma-separated lists link the first line.
\ExplSyntaxOn
\tl_new:N \l__leanlink_arg_tl
\tl_new:N \l__leanlink_url_tl
\seq_new:N \l__leanlink_split_seq
\cs_new_protected:Npn \leanlink #1
  {
    \tl_set:Nx \l__leanlink_arg_tl { \tl_to_str:n {#1} }
    \regex_replace_all:nnN { \\allowbreak } { } \l__leanlink_arg_tl
    \regex_replace_all:nnN { \s+ } { } \l__leanlink_arg_tl
    \tl_set:Nx \l__leanlink_url_tl { \repobase }
    \exp_args:NnV \regex_match:nnTF { \A (ErdosProblems|Erdos249257)/ } \l__leanlink_arg_tl
      { }
      { \tl_put_right:Nx \l__leanlink_url_tl { / \PK } }
    \exp_args:NnV \regex_extract_once:nnNTF { (.*\.lean) :? (.*) } \l__leanlink_arg_tl \l__leanlink_split_seq
      {
        \tl_put_right:Nx \l__leanlink_url_tl { / \seq_item:Nn \l__leanlink_split_seq { 2 } }
        \tl_set:Nx \l_tmpa_tl { \seq_item:Nn \l__leanlink_split_seq { 3 } }
        \regex_replace_all:nnN { -- } { - } \l_tmpa_tl
        \exp_args:NnV \regex_extract_once:nnNTF { ([0-9]+) (- ([0-9]+))? } \l_tmpa_tl \l__leanlink_split_seq
          {
            \tl_put_right:Nx \l__leanlink_url_tl { \c_hash_str L \seq_item:Nn \l__leanlink_split_seq { 2 } }
            \str_if_eq:eeF { \seq_item:Nn \l__leanlink_split_seq { 4 } } { }
              { \tl_put_right:Nx \l__leanlink_url_tl { -L \seq_item:Nn \l__leanlink_split_seq { 4 } } }
          }
          { }
      }
      { \tl_put_right:Nx \l__leanlink_url_tl { / \tl_use:N \l__leanlink_arg_tl } }
    \href { \tl_use:N \l__leanlink_url_tl } { \textsc{Lean~source} }
  }
\newcommand{\lean}[2]{\leanlink{#2}}
\ExplSyntaxOff

% Declaration names are written \mathtt{...} in maths and, unavoidably, in
% running prose too. Rather than police every occurrence, make the command
% mode-agnostic: monospace is what is meant in both places.
\AtBeginDocument{%
  \let\origmathtt\mathtt
  \renewcommand{\mathtt}[1]{\ifmmode\origmathtt{#1}\else\texttt{#1}\fi}%
  \let\origmathrm\mathrm
  \renewcommand{\mathrm}[1]{\ifmmode\origmathrm{#1}\else\textrm{#1}\fi}%
}

% One authored part uses \m{...} as an inline-maths shorthand. Defining it is
% cheaper and safer than rewriting the prose around every occurrence.
\providecommand{\m}[1]{\ensuremath{#1}}

% ---------------------------------------------------------------------------
% Evidence band. These are never blurred: a reader must always be able to see
% what a kernel checked, what an exact computation checked, what was proved in
% ordinary mathematics, what is quoted from the literature, and what is open.
% ---------------------------------------------------------------------------
\newcommand{\ev}[1]{{\small\textsf{[#1]}}\xspace}

% Scale at which a statement is proved. This is the load-bearing axis of both
% documents: the corpus proves at fixed or bounded scale, and both open
% producers need cofinal scale.
\newcommand{\scale}[1]{{\small\textsf{scale:#1}}\xspace}

% The representation a statement is expressed in. Obstructions are relative to
% a coordinate, never to the object, so every obstruction carries one.
% \nolinkurl manipulates catcodes, so it cannot be written into the
% hyperref bookmark file. \coord appears in subsubsection titles, so it
% must supply a plain-text PDF string alongside the typeset form.
\newcommand{\coord}[1]{\texorpdfstring{{\small\nolinkurl{coord:#1}}}{coord:#1}\xspace}
\newcommand{\codeid}[1]{\nolinkurl{#1}}

% ---------------------------------------------------------------------------
% Mathematical shorthand.
% ---------------------------------------------------------------------------
\newcommand{\N}{\mathbb{N}}
\newcommand{\Npos}{\mathbb{N}_{>0}}
\newcommand{\Z}{\mathbb{Z}}
\newcommand{\Q}{\mathbb{Q}}
\newcommand{\R}{\mathbb{R}}
\newcommand{\C}{\mathbb{C}}
\newcommand{\ph}{\varphi}
\newcommand{\EB}{E}

\allowdisplaybreaks
\emergencystretch=2em

% These working records contain a small number of intentionally exhaustive
% formulas. Keep ordinary displays at their natural size, but bound any display
% wider than the classical text measure rather than letting it disappear past
% the page edge.
\long\def\[#1\]{%
  \leavevmode\par\vspace{\abovedisplayskip}%
  \noindent\makebox[\linewidth][c]{%
    \begin{adjustbox}{max width=\linewidth}%
      \(\displaystyle #1\)%
    \end{adjustbox}}%
  \par\vspace{\belowdisplayskip}%
}

\runninghead{Erd\H{o}s \#249: the binary totient series}
\title{The Binary Totient Series}
\subtitle{A detailed map of the current attack on Erd\H{o}s Problem~\#249}
\author{Will Cook}
\date{July 2026}

\begin{document}
\maketitle

\paperprovenance

\begin{abstract}
This paper gives a claim-bounded map of the current formal attack on Erd\H{o}s
Problem~249, which asks whether
$S=\sum_{n\ge1}\ph(n)/2^n$ is irrational.  The problem remains open.
The unconditional headline results are an explicit rational basis for the
full dyadic totient kernel, exact rank $2^e+1$ through every level $e\ge1$, a denominator
exclusion through approximately $7.96\times10^{34}$, and finite
lcm-diagonal certificates for every $t\le82$.  The full kernel is therefore
infinite-dimensional.  The paper also records exact equivalences between
irrationality and several cofinal certificate or tail-nonintegrality supplies.
Those equivalences re-express rather than solve the problem: no certificate at
$t=83$, unbounded certificate family, or irrationality proof is obtained.  The
contribution is the audited organization of checked results,
coordinate-specific obstructions, and open implications---not a solution, an
exhaustive reproduction of every repository theorem, or a priority claim for
every formalisation.
\end{abstract}

\section*{How to read this document}

This is a reasoning surface, not an exhaustive corpus export or a literature
survey.  It is designed to expose the principal premises, obstructions, and
open interfaces before work begins.  The machine-readable corpus remains the
current inventory when later theorem waves outrun this exposition.  Three
conventions carry the paper's claim boundary.

\emph{Evidence bands.}  Every statement is tagged.  \ev{Lean} means a proof term
was checked by the pinned Lean kernel.  \ev{Cert} means an exact finite
computation, with no floating-point decision anywhere in it.  \ev{Math} means
proved in ordinary mathematics in the sources but not formalised.  \ev{Cited}
means established in the published literature.  \ev{Open} means not proved.
These are never blurred, and a statement carrying one band is never described in
language belonging to another.

\emph{Scale.}  Every parameter-indexed statement is tagged \scale{fixed} (proved
at specific listed values), \scale{bounded} (proved below an explicit bound),
\scale{cofinal} (proved for arbitrarily large parameters), or \scale{uniform}
(proved for all parameters).  The corpus contains theorems at all four scales;
the unresolved endpoint depends on a particular cofinal certificate producer
that is not among them.  Sorting by scale keeps that distinction visible.

\emph{Coordinates.}  Every statement records the representation it is expressed
in.  An obstruction is a fact about a coordinate, not about the object, and a
wall measured in one representation may simply not exist in another; the atlas
section gives the transport maps, so any obstruction recorded here can be
re-measured elsewhere.

\emph{Frontier terminology.}  The catalogue uses \emph{producer} for an
unproved hypothesis that would supply a missing step, \emph{consumer} for a
proved implication that uses such a hypothesis, \emph{converter} for an exact
identity or reduction, and \emph{socket} for the precise place where the
missing hypothesis enters.  These are local navigation terms, not additional
mathematical objects or evidence classes.

Erd\H{o}s Problem 249 is open.  No section of this document claims
otherwise, and a reduction of the problem to another statement is recorded as a
reduction, never as progress toward a solution.
\clearpage
\tableofcontents
\clearpage
% ---- part a249_front ----
\section{The problem, and what is actually known}

Let $\varphi$ be Euler's totient function and put
\[
  S \;:=\; \sum_{n \ge 1} \frac{\varphi(n)}{2^{n}}
  \;=\; \tfrac12 + \tfrac14 + \tfrac{2}{8} + \tfrac{2}{16} + \tfrac{4}{32} + \cdots
  \;=\; 1.3676308019850223\ldots
\]
The series converges absolutely because $\varphi(n) \le n$. Erd\H{o}s Problem
\#249 asks whether $S$ is irrational. \emph{It is open.} Nothing in this
document decides it, no route recorded here is a proof, and no result here
should be read as an approach that is close to working. What this document
does is different in kind: it assembles every recorded failure, obstruction,
countermodel and dead route the programme has produced against \#249, and
classifies each one by the exact class of argument it eliminates. The claim
being made is that those failures are not independent --- they are repeated
measurements of a single obstruction, and that obstruction has a shape which
can be stated.

\begin{rem}[Literature coordinates]
The nearby literature serves three distinct roles.  Kova\v{c}--Tao's
\href{https://arxiv.org/abs/2406.17593v4}{Ahmes-series paper} maps a broader
Lambert-subseries landscape, but does not treat the coefficient series $S$.
Merca's \href{https://doi.org/10.1007/s11139-016-9856-3}{Lambert-series
factorization theorem} and the
\href{https://doi.org/10.55016/ojs/cdm.v14i1.62425}{Merca--Schmidt}
factor-pair formalism provide transform background for the
M\"obius--Mersenne coordinate.  Finally,
\href{https://arxiv.org/abs/1511.02221}{Balasubramanian--Giri--Srivastav,
Theorem~2.2} gives uniform shifted-correlation estimates for the relevant
divisor-convolution class.  It does not supply the dyadic residue small-ball
or phase anti-concentration required by $\mathrm{Sep}$.  These are contextual
or method sources, not proofs of the open certificate supply. \ev{Cited}
\end{rem}

The open question has a sharp equivalent inside the corpus, and it is the
statement everything below is measured against. Write $2^{N} S = \Phi_N + R_N$
with $\Phi_N := \sum_{n \le N} \varphi(n) 2^{N-n} \in \N$ and
$R_N := \sum_{j \ge 0} \varphi(N+1+j)/2^{j+1} \in [0,\infty)$
(\lean{Erdos249257.TotientTailPeriodKiller.totientTail}{TotientTailPeriodKiller.lean:57},
\lean{two_pow_mul_totient_series_eq}{TotientTailPeriodKiller.lean:150}).
If $S \in \Q$ then, applying Euler's theorem to the odd part of its
denominator, there are $h > 0$ and $N_0$ with $R_{N+h} - R_N \in \Z$ for every
$N \ge N_0$
(\lean{Erdos249257.eventual\_period\_of\_not\_irrational}{TotientTailPeriodKiller.lean:358},
explicit witnesses $h = \varphi(\mathrm{oddPart}(r.\mathrm{den}))$, $N_0 = v_2(r.\mathrm{den})$).

So it suffices to refute integrality cofinally on every period ray. The finite
device that refutes it is a residue certificate:
\[
  \mathrm{windowDiscrepancy}(h,N,L)
  \;:=\; \sum_{j<L}\bigl(\varphi(N{+}h{+}1{+}j) - \varphi(N{+}1{+}j)\bigr)\,2^{\,L-1-j} \in \Z,
\]
the depth-$L$ truncation of $2^{L}(R_{N+h} - R_N)$, and
\[
  \mathrm{certifiedKill}(h,N,L)
  \;:\Longleftrightarrow\;
  (N{+}h{+}L{+}2) \;<\; \mathrm{windowDiscrepancy}(h,N,L) \bmod 2^{L}
  \;<\; 2^{L} - (N{+}h{+}L{+}2),
\]
a decidable condition
(\lean{Erdos249257.TotientTailPeriodKiller.windowDiscrepancy}{TotientTailPeriodKiller.lean:66},
\lean{Erdos249257.TotientTailPeriodKiller.certifiedKill}{TotientTailPeriodKiller.lean:72}).

The radius $N{+}h{+}L{+}2$ is exactly the crude tail bound coming from
$\varphi(m) \le m$, so a certificate is sound: it forces
$R_{N+h} - R_N \notin \Z$, and in fact certificates are \emph{complete} for
non-integrality
(\lean{Erdos249257.exists\_certifiedKill\_iff\_tail\_diff\_notMem\_int}{LcmConeFlatness.lean:316}).
The entire \#249 side of the programme therefore reduces to one obligation.

\begin{defn}[The supply obligation $\mathrm{Sep}$ --- the exact open target]
\label{defn:sep}
\[
  \mathrm{Sep}
  \;:\Longleftrightarrow\;
  \forall h \ge 1\ \forall N_0\ \exists N \ge N_0\ \exists L,\;
  \mathrm{certifiedKill}(h,N,L).
\]
$\mathrm{Sep} \Rightarrow \mathrm{Irrational}(S)$
(\lean{Erdos249257.irrational\_totient\_series\_of\_certificate\_supply}{TotientTailPeriodKiller.lean:394}).
\coord{binary-digit} \scale{cofinal} \ev{Lean}
\end{defn}

Two features of $\mathrm{Sep}$ govern everything that follows. First, the
quantifier shape is $\forall\forall\exists\exists$: the corpus can and does
verify individual instances, but an instance is not the obligation. Second,
depth cannot stay bounded: from $\mathrm{certifiedKill}(h,N,L)$ one gets
$2(N{+}h{+}L{+}2) < 2^{L}$ immediately
(\lean{Erdos249257.certifiedKill\_depth\_floor}{TotientTailPeriodKiller.lean:79},
\ev{Lean}), so $L \gtrsim \log_2(N{+}h)$. A certificate is a demand that a
weighted accumulation over a window of length $\approx \log_2 N$ lands near
the centre of its dyadic arc, at full arithmetic resolution.

\subsection{The unconditional record}

The following hold with no irrationality hypothesis. They are the results a
specialist should know before reading further; each is stated at the scale at
which it is actually proved.

\begin{thm}[Denominator exclusion --- the headline unconditional fact]
\label{thm:denom}
If $S \in \Q$ then its reduced denominator exceeds
$Q_0 := 79\,639\,646\,646\,701\,375\,323\,355\,774\,875\,831\,053 \approx 7.96 \times 10^{34}$.
Equivalently, $S \ne p$ for every $p \in \Q$ with $p.\mathrm{den} \le Q_0$.
The bound is sharp for the method: $q = Q_0 + 1$ is the exact first failing
denominator, exhibited as the mediant of two explicit unimodular Farey
neighbours.
\coord{farey} \scale{bounded} \ev{Lean}
\lean{Erdos249257.tsum\_totient\_div\_pow\_two\_ne\_ratCast\_of\_den\_le\_79639646646701375323355774875831053}{CertificateKernel.lean:18384}
\lean{Erdos249257.gap\_check\_window\_1\_240\_first\_failure}{GapFareyBound.lean:225}
\end{thm}

This is obtained from a classical Stern--Brocot gap lemma
(\lean{farey\_gap}{GapFareyBound.lean:51}) applied at window $K = 240$, the
last rung of the ladder $4838 \to 2^{22} \to 2.49\times10^{17} \to Q_0$. It is
logically independent of $\mathrm{Sep}$: it uses no certificate and no period
apparatus. Its own open continuation is whether $\sup_K (b+d)(K) = \infty$;
that would close \#249 through this theorem's consumer. \ev{Open}

\begin{prop}[Finite certificate deposits]
\label{prop:deposits}
$\mathrm{certifiedKill}$ has been verified at: the $28$ diagonal instances of
the LCM pincer through $t = 64$
(\lean{Erdos249257.certifiedKill\_diagonal\_all\_imported\_through\_t64}{DiagonalPincerCertificatesT64.lean:1967},
endpoint \lean{Erdos249257.certifiedKill\_diagonal\_t64}{DiagonalPincerCertificatesT64.lean:1928});
all shifts $h \in [1,16]$ simultaneously at $(N,L) = (14,9)$, by
\texttt{decide}
(\lean{Erdos249257.certifiedKill\_all\_upto\_sixteen}{CarrySurvivorExtinction.lean:574});
and eight further period deposits at $N = 300$. Every deposit fires within a
small additive constant of the minimum depth permitted by
$\mathrm{certifiedKill\_depth\_floor}$. Note the quantifier order: this deposit
proves $\exists N \exists L\, \forall h \le 16$, whereas $\mathrm{Sep}$ needs
$\forall h\, \forall N_0\, \exists N \exists L$.
\coord{binary-digit} \scale{fixed} \ev{Lean}
\end{prop}

\begin{prop}[Cofinal positivity, and why it is exactly half a certificate]
\label{prop:sign}
The true actual-LCM tail difference is strictly positive for every $a \ge 8$,
with no irrationality hypothesis
(\lean{Erdos257PeriodNoncollapse.actualLcmTailDiff\_shift\_pos}{TotientActualLcmOrbitSign.lean:39},
\lean{Erdos257PeriodNoncollapse.actualLcmTailOrbit\_pos}{TotientActualLcmOrbitSign.lean:144}).
In the same coordinate, integrality of the orbit forces the residue to the
\emph{top edge}, exactly $2^{K} - e$
(\lean{Erdos257PeriodNoncollapse.actualLcm\_integral\_forces\_topEdgeResidue}{TotientActualLcmOrbitSign.lean:211}).
Positivity does not exclude the top edge. So the strongest cofinal fact the
corpus owns about the actual object supplies one of the two inequalities a
certificate needs and provably cannot supply the other.
\coord{actual-lcm} \scale{cofinal} \ev{Lean}
\end{prop}

\begin{prop}[Rationality forces unbounded carry rank]
\label{prop:rank}
If $S$ is rational then, for every $e$, the dyadic carry-kernel family at
level $e$ has $\Q$-rank at least $2^{e}-1$
(\lean{Erdos249257.not\_irrational\_totientSeries\_implies\_unbounded\_carryRank\_unconditional}{TotientCarryKernelRigidity.lean:300}).
The scale side is complete and uniform in $e$: the canonical family of
$2^{e}+1$ dyadic totient-kernel channels is linearly independent over $\Q$ at
every depth, via CRT plus Dirichlet
(\lean{Erdos249257.linearIndependent\_canonicalTotientKernelFamily}{TotientMahlerDefect.lean:935}),
so the full family spans an infinite-dimensional space
(\lean{Erdos249257.not\_finiteDimensional\_span\_fullTotientKernel}{TotientMahlerDefect.lean:1145}).
\coord{carry-rank} \scale{uniform} \ev{Lean}
\end{prop}

\begin{prop}[What rationality actually buys, and what it does not]
\label{prop:period-not-rank}
Rationality gives uniform eventual periodicity of the carry's dyadic sections
modulo $v$; that periodicity provably does \emph{not} promote to any $\Q$-rank
bound on the carry family
(\lean{Erdos257PeriodNoncollapse.not\_irrational\_totientSeries\_implies\_mod\_period\_and\_unbounded\_rank}{TotientTailCarryPeriod.lean:224}).
\coord{carry-rank} \scale{uniform} \ev{Lean}
\end{prop}

\begin{prop}[Exact reformulations that do not move the truth value]
\label{prop:iffs}
Several re-encodings are proved \emph{equivalent}
to $\mathrm{Irrational}(S)$, not merely sufficient for it: the period-multiple
kill supply
(\lean{periodMultipleKillSupply\_iff\_irrational}{ErdosProblems/Erdos249/PeriodMultipleEscape.lean:432}),
the base certificate supply and its lcm-diagonal normal form
(\lean{irrational\_totient\_series\_iff\_certificate\_supply}{LcmConeFlatness.lean:412};
\lean{irrational\_totient\_series\_iff\_lcm\_diagonal\_certificate\_supply}{LcmConeFlatness.lean:426}),
the directed-certificate supply and its LCM specialisation
(\lean{Erdos257PeriodNoncollapse.CofinalDirectedLcmCertificateSupply}{TotientTailCarryPeriod.lean:866}),
and the window-separated-pairs predicate
(\lean{Erdos249257.dtwWindowSeparatedPairs\_iff\_irrational\_totient\_series}{PivotAntiReconstruction.lean:1765}).
\coord{binary-digit} \scale{cofinal} \ev{Lean}
\end{prop}

\begin{prop}[A rational sequence indistinguishable from $\varphi$ by every coarse invariant]
\label{prop:parity}
There is $c : \N \to \N$ with $c(n) \le 6$ and $c(n) \le n$ for all $n$,
$c(n) \equiv \varphi(n) \pmod 2$ for \emph{every} $n$, and $c$ not eventually
periodic --- indeed for every $N, G, K$ there are $K$ explicit carry pulses
beyond $N$, pairwise separated by more than $G$ --- and yet
$\sum_n c(n)/2^{n} = 3/2 \in \Q$
(\lean{Erdos257PeriodNoncollapse.exists\_totientParity\_arbitrarilyManySeparatedCarry\_rational\_countermodel}{TotientParityCoboundaryCountermodel.lean:637};
sum at \texttt{:359}, aperiodicity at \texttt{:487}, parity match at \texttt{:194};
\texttt{\#print axioms} clean).
\coord{coefficient-word} \scale{uniform} \ev{Lean}
\end{prop}

Propositions~\ref{prop:sign}--\ref{prop:parity} are the four facts a reader
should carry into the next section: the corpus's best cofinal information is
half a certificate, its best rank information runs the wrong way, its
reformulations are equivalences rather than reductions, and every purely
qualitative property of the coefficient word is satisfied by a rational
countermodel.

\section{The wall}
\label{sec:wall}

This programme has produced many exact reformulations of \#249 and no proof.
That is the honest summary, and the rest of this section is an argument that
it is also the wrong way to read the record. The reformulations do not fail
independently. They fail in a small number of recurring ways, and when each
failure is stated as a claim about \emph{which class of argument it
eliminates} rather than as a report that something did not work, the classes
fit together into one obstruction with a describable shape. A hundred
reformulations that all die are not a hundred failures; they are a hundred
measurements of one wall, taken from different angles. The measurements are
below.

The thesis, in one sentence: \emph{every barrier in the record is a bound, and
the quantity a certificate measures is invariant under exactly the bounds that
make an argument finite.}

\subsection{The plane the barriers live on}

Fix two axes for a hypothetical proof of $\mathrm{Sep}$
(Definition~\ref{defn:sep}): the \emph{range} of indices at which it consults
$\varphi$, and the \emph{resolution} at which it consults them. Two derived
axes matter as well: the size of the proof's internal bookkeeping (carry
state, $\Q$-rank, shift-polynomial degree), and whether the target is asserted
\emph{pointwise} at a chosen index or as an \emph{average} over a block.

\begin{itemize}[leftmargin=*]
\item The quadrant [full resolution $\times$ bounded range] is closed by
  Theorem~\ref{thm:gamma} (B1): $\varphi$ may be pinned exactly on $[1,B]$ and
  the series can still be rational.
\item The complementary quadrant [coarse resolution $\times$ unbounded range]
  is closed by Proposition~\ref{prop:parity} (B7): $\varphi$'s parity may be
  pinned at \emph{every} index, together with boundedness, $c(n) \le n$, and
  arbitrarily strong aperiodicity, and the series can still be $3/2$.
\item The remaining quadrant is closed \emph{whenever the proof compresses}:
  B5 (no bounded carry state, no fixed-precision signature) and B6 (four
  independent finite truncations of the totient kernel, all with trivial
  kernel).
\item B4 closes the one construction that tried to reach the good quadrant by
  \emph{prescribing} values: residue engineering pays for amplitude in
  position, and position enters the certificate radius.
\item B2 closes the escape of retargeting; B3 records that no bounded result
  in the corpus has ever promoted.
\end{itemize}

The corpus's own results partition along these axes perfectly, which is the
first evidence that the axes are the right ones. Every unconditional
\emph{finite} deposit (Proposition~\ref{prop:deposits}, the $K=240$ Farey rung
of Theorem~\ref{thm:denom}, the actual-LCM orbits at $a = 4$ and $a = 6$) sits
at full resolution and bounded range, so B1 says extending them is evidence
forever and proof never. Every unconditional \emph{cofinal} theorem (letterwise
positivity for all $a \ge 8$, the tail bounds from $\varphi(m) \le m$,
unbounded Mersenne-shadow denominator growth) sits at coarse resolution and
unbounded range, and Proposition~\ref{prop:sign} proves in the corpus's own
coordinate that this is exactly half a certificate.

\subsection{B1: the finite-inspection barrier, as a theorem}

B1 is the one barrier that is a genuine no-go theorem about proof method, so
it is stated and proved as one. The generic vocabulary is needed first. For
$c : \N \to \N$ with $c(n) \le n$, write
$T_c := \sum_{n \ge 1} c(n)/2^{n}$,
$R^{c}_{N} := \sum_{j \ge 0} c(N{+}1{+}j)/2^{\,j+1}$,
$A_c(h,N,L) := \sum_{j<L}(c(N{+}h{+}1{+}j) - c(N{+}1{+}j))2^{\,L-1-j}$,
and let $\mathrm{Kill}_c(h,N,L)$ and $\mathrm{Sep}_c$ be
$\mathrm{certifiedKill}$ and $\mathrm{Sep}$ with $A_c$ in place of
$\mathrm{windowDiscrepancy}$. For $c = \varphi$ these are the objects of
Definition~\ref{defn:sep}.

\begin{lem}[Generic soundness]
\label{lem:gsound}
$\mathrm{Kill}_c(h,N,L) \Rightarrow R^{c}_{N+h} - R^{c}_{N} \notin \Z$.
\coord{binary-digit} \scale{uniform} \ev{Math}
\end{lem}

\begin{proof}
From $c(n) \le n$ one gets
$\lvert 2^{L}(R^{c}_{N+h} - R^{c}_{N}) - A_c(h,N,L)\rvert < N{+}h{+}L{+}2$;
this is the crude tail estimate, formalised for $\varphi$ at
\lean{Erdos249257.tail\_after\_le}{TotientTailPeriodKiller.lean:198} and used
there to prove
\lean{Erdos249257.tail\_diff\_notMem\_int\_of\_certifiedKill}{TotientTailPeriodKiller.lean:262}.
If $R^{c}_{N+h} - R^{c}_{N} = k \in \Z$ then $2^{L}k \equiv 0 \pmod{2^{L}}$
and $A_c$ lies within $N{+}h{+}L{+}2$ of $2^{L}k$, so
$A_c \bmod 2^{L}$ lies within $N{+}h{+}L{+}2$ of $0$ or of $2^{L}$,
contradicting $\mathrm{Kill}_c$.
\end{proof}

\begin{lem}[Generic tail-period law]
\label{lem:gperiod}
If $T_c = p/(2^{e}m)$ with $m$ positive and odd, and if $h\ge1$ satisfies
$m\mid 2^h-1$, then
$R^{c}_{N+h} - R^{c}_{N} \in \Z$ for every $N \ge e$.
\coord{binary-digit} \scale{cofinal} \ev{Math}
\end{lem}

\begin{proof}
$R^{c}_{N} = 2^{N}T_c - \Phi^{c}_{N}$ with $\Phi^{c}_{N} \in \Z$, so
$R^{c}_{N+h} - R^{c}_{N} = 2^{N}(2^{h}-1)T_c - (\Phi^{c}_{N+h} - \Phi^{c}_{N})$,
and $2^{N}(2^{h}-1)p/(2^{e}m) \in \Z$ because $2^{e} \mid 2^{N}$ and
$m \mid 2^{h}-1$.
\end{proof}

\begin{rem}
Lemmas~\ref{lem:gsound} and~\ref{lem:gperiod} are elementary and are stated
here in ordinary mathematics. Their $\varphi$-instances are Lean theorems
(\lean{Erdos249257.tail\_diff\_notMem\_int\_of\_certifiedKill}{TotientTailPeriodKiller.lean:262},
\lean{Erdos249257.eventual\_period\_of\_not\_irrational}{TotientTailPeriodKiller.lean:358}),
and the generic tail-period direction exists on disk inside
\lean{binaryCoeffSeries\_rational\_iff\_exists\_temperedBinaryOrbit}{GenericTailOrbitRigidity.lean:426};
a generic \emph{named} theorem of the form of Lemma~\ref{lem:gperiod} does not
exist in either tree. The composition below is mechanical but unformalised, and
is flagged \ev{Math} throughout for that reason.
\end{rem}

\begin{thm}[B1, the $\gamma$-splice: finite inspection cannot certify the supply]
\label{thm:gamma}
Let $B \ge 1$ and let $P > B$. Define $\gamma : \N \to \N$ by
\[
  \gamma(n) := \varphi(n) \ \ (n \le B), \qquad
  \gamma(n) := \begin{cases} n-1, & P \mid n \\ n, & P \nmid n \end{cases}
  \ \ (n > B).
\]
Then:
\begin{enumerate}[label=(\roman*),leftmargin=*]
\item $\gamma(n) \le n$ for all $n$, so $\gamma$ lies in the same coefficient
  class as $\varphi$;
\item $\gamma(n) = \varphi(n)$ for every $n \le B$, and consequently
  $A_\gamma(h,N,L) = A_\varphi(h,N,L)$ for every $(h,N,L)$ with
  $N + h + L \le B$;
\item $T_\gamma = D - 1/(2^{P}-1) \in \Q$, where
  $D = 2 - \sum_{n \le B}(n - \varphi(n))/2^{n}$, and the odd part of the
  reduced denominator of $T_\gamma$ is \emph{exactly} $2^{P}-1$;
\item $\mathrm{Sep}_\gamma$ is false --- indeed it fails already at $h = P$.
\end{enumerate}
Hence for every $B$ there is a coefficient sequence in the same class,
agreeing with $\varphi$ on all of $[1,B]$, whose supply obligation is false.
\coord{binary-digit} \scale{uniform} \ev{Math}
\end{thm}

\begin{proof}
(i) is immediate; (ii) holds because $A_\varphi(h,N,L)$ reads $\varphi$ only at
indices $\le N{+}h{+}L$. For (iii), since $P > B$ every multiple of $P$ exceeds
$B$, so
\[
  T_\gamma
  = \sum_{n\ge1}\frac{n}{2^{n}} - \sum_{n \le B}\frac{n - \varphi(n)}{2^{n}}
    - \sum_{k \ge 1} 2^{-kP}
  = 2 - \sum_{n \le B}\frac{n-\varphi(n)}{2^{n}} - \frac{1}{2^{P}-1}.
\]
Write $D = A/2^{B}$ with $A = 2^{B+1} - \sum_{n \le B}(n-\varphi(n))2^{\,B-n} \in \Z$.
Then $T_\gamma = \bigl(A(2^{P}-1) - 2^{B}\bigr)/\bigl(2^{B}(2^{P}-1)\bigr)$, and
$\gcd\bigl(A(2^{P}-1) - 2^{B},\, 2^{P}-1\bigr) = \gcd(2^{B}, 2^{P}-1) = 1$
because $2^{P}-1$ is odd. So no factor of $2^{P}-1$ cancels, and the odd part
of the reduced denominator is exactly $2^{P}-1$. For (iv), the odd part of the
denominator is $m = 2^{P}-1$ and $\mathrm{ord}_m(2) = P$, while
$v_2(\mathrm{den}) \le B$. By Lemma~\ref{lem:gperiod},
$R^{\gamma}_{N+P} - R^{\gamma}_{N} \in \Z$ for every $N \ge B$, so by
Lemma~\ref{lem:gsound} no $\mathrm{Kill}_\gamma(P,N,L)$ holds for any
$N \ge B$ and any $L$. Taking $N_0 = B$ refutes the inner existential of
$\mathrm{Sep}_\gamma$ at $h = P$.
\end{proof}

\begin{cor}[What B1 rules out]
\label{cor:b1}
No proof rule that is uniform over all coefficient sequences $c(n)\le n$ can
establish $\mathrm{Sep}$ from a single fixed prefix
$\{c(n):n\le B\}$: Theorem~\ref{thm:gamma} supplies a rational countermodel
with that same prefix.  This does \emph{not} invalidate an argument that uses
the fixed arithmetic sequence $\varphi$ together with compatible information
at arbitrarily large horizons; the theorem gives a different $\gamma_B$ for
each $B$, not one sequence agreeing with $\varphi$ at every $B$.
\end{cor}

\begin{rem}[Scope discipline --- what B1 does \emph{not} say]
\label{rem:b1-scope}
Theorem~\ref{thm:gamma} is a statement about proof method, not about $\varphi$.
It does not touch $S$, and it does not suggest that $\varphi$ is such a
$\gamma$; it says only that no fixed-horizon check can tell them apart. It
also does \emph{not} kill bounded-parameter certificates in general.
``Bounded parameter'' splits into four independent bounds with four different
killers: bounded index range (this theorem), bounded certificate depth $L$
(killed separately, and by a Lean theorem ---
$\mathrm{certifiedKill\_depth\_floor}$ forces $2(N{+}h{+}L{+}2) < 2^{L}$),
bounded proof state (B5), and bounded rank (B6). Conflating them would
overstate B1.
\end{rem}

The Lean-formalised sibling of the same mechanism --- a lacunary zero-valued
coboundary splice --- is Proposition~\ref{prop:parity}, which is B7. B1 and B7
are the two object-level witnesses of the wall, one along each axis.

\subsection{B2--B7}

Each barrier below is stated as a claim about a class of arguments, with its
status never blurred. \emph{Proved} means there is a theorem (Lean or ordinary
mathematics, as marked) whose content is the elimination. \emph{Observed
pattern} means an audit over the corpus: real evidence, not a theorem, and it
is never used as a premise elsewhere in this document.

\begin{prop}[B2, route-collapse: intermediate targets are not waypoints]
\label{prop:b2}
\textbf{Status: proved} (Lean, three independent mechanisms).
\begin{enumerate}[label=(\alph*),leftmargin=*]
\item \emph{Certificate completeness.} Every certificate vocabulary is an
  \emph{iff} with the underlying non-integrality, so every re-encoding of
  $\mathrm{Sep}$ is equivalent to $\mathrm{Sep}$
  (\lean{periodMultipleKillSupply\_iff\_irrational}{ErdosProblems/Erdos249/PeriodMultipleEscape.lean:432},
  forward at \texttt{:132}, converse at \texttt{:143}).
\item \emph{Sample choice.} Any weakening in which the prover may \emph{choose}
  which indices to test collapses to $\mathrm{Irrational}(S)$ outright, because
  irrationality plus doubling expansivity already supplies a two-point sample
  with the required separation
  (\lean{Erdos249257.dtwWindowSeparatedPairs\_iff\_irrational\_totient\_series}{PivotAntiReconstruction.lean:1765}).
  The mechanism is visible in the converse proof
  (\texttt{windowSeparatedPairsAt\_of\_cofinally\_scaled\_adjacent\_chord},
  \texttt{PivotAntiReconstruction.lean:1663}): it selects $T = \{N, N+1\}$ and
  two ordered pairs, and the counted-energy threshold
  $2\lvert T\rvert^{2}/5 = 8/5 \le 2\delta^{2}$ is met by $\delta = 9/10$.
\item \emph{Multi-point enrichment.} Adding vertices, higher-order differences,
  or finite shift-polynomial combinations on the same ray buys nothing, because
  integrality transports affinely along every ray $H \mapsto kH$: the four-hit
  diamond $\mathrm{Hit}(H) \wedge \mathrm{Hit}(pH) \wedge \mathrm{Hit}(qH)
  \wedge \mathrm{Hit}(pqH)$ is equivalent to $\mathrm{Hit}(H)$ alone, with no
  primality used
  (\lean{Erdos249257.fullTarget\_primeAdjunction\_diamond\_iff\_root}{FullTargetPrimeAdjunctionNoGo.lean:196}).
\end{enumerate}
\textbf{Rules out:} ``find an easier waypoint'' as a strategy class.
\textbf{Does not rule out:} targets over a \emph{fixed full block} with no
sample choice, and targets demanding a \emph{uniform quantitative margin} that
irrationality does not supply. These two exemptions are read off the collapse
proofs themselves and are exactly what the surviving routes in
\S\ref{sec:survivors} exploit.
\coord{binary-digit} \scale{cofinal} \ev{Lean}
\end{prop}

\begin{rem}
A corroborating instance: rank-$2$ second-difference certificates are sound but
measurably \emph{not} shallower than rank-$1$ --- at $(h,N) = (1,8)$ rank-$1$
fires at depth $8$ and no rank-$2$ certificate exists at depth $\le 8$
(\texttt{totient\_tail\_rank\_two\_kill\_sound\_but\_not\_shallower\_cell},
\texttt{CertificateKernel.lean:18762}, \ev{Lean}). Separately, the claim
sometimes made that the Farey growth law $\sup_K(b+d) = \infty$ is ``equivalent
in difficulty to \#249'' is an argument, not an \emph{iff}: the bound produced
is the convergent denominator of the underlying constant, which stalls at $q_0$
precisely if $S = a/q_0$. It should not be cited as an equivalence. \ev{Math}
\end{rem}

\begin{prop}[B3, no free promotion]
\label{prop:b3}
\textbf{Status: observed pattern (an audit, not a theorem).} Across the
corpus's proved results, no bounded or partial result has a proof uniform
enough to be promoted to cofinal scale by routine strengthening --- raising a
bound, widening a hypothesis, or reindexing a quantifier. Method: $56$
near-miss rows were catalogued against the four open obligations; $18$ were
flagged promotable on a statement-level read; all $18$ proof \emph{bodies} were
then opened. Verdict: $16$ conclusively not promotable with a named blocker
each, and $2$ promotable with no new mathematics but both pure restatements
that widen a target family without supplying arithmetic content. Blocker
taxonomy: hypothesis strength $\times 6$, scale-only $\times 5$,
coordinate-only $\times 4$, multiple $\times 3$. Rows with an outright
quantifier-order mismatch were judged unpromotable before the body audit and
excluded, so the audited $18$ are the \emph{most favourable} subset.
\coord{audit} \scale{n/a} \ev{Cert}
\end{prop}

\begin{obs}[The sharpest signal in the audit]
\label{obs:pointwise}
Every not-promotable \#249-supply row asks for a \emph{pointwise} fact at a
specially chosen index: one large $a$ with a top-edge residue gap, one $q$ with
a terminal-dominance inequality, one exponent past $a = 6$, one $t$ past $64$,
one sign at one LCM jump, one prime per LCM height, one $K$ past $240$. Not one
such index has ever been located unconditionally beyond the finite census. The
single audited row whose missing input is not of that shape is the
first-harmonic gap --- an \emph{average} over a block. This contrast is what
selects the leading survivor in \S\ref{sec:survivors}. It is evidence, not
proof.
\end{obs}

\begin{prop}[B4, amplitude versus radius: residue engineering is self-defeating]
\label{prop:b4}
\textbf{Status: proved} (Lean construction plus an unconditional arithmetic
margin). A certificate demands that the window residue sit at distance
$> N{+}h{+}L{+}2$ from both $0$ and $2^{L}$ modulo $2^{L}$, and the depth floor
forces $2^{L} > 2(N{+}h{+}L{+}2)$, so the required amplitude is comparable to
the modulus. Any producer that \emph{buys} that amplitude by prescribing
totient values at engineered positions pays for the prescription in position,
and position enters the radius. The corpus's strongest such construction proves
the trade is exactly self-cancelling, unconditionally and at every depth: a
two-adic pulse block lands the residue at the arc \emph{centre}
$2^{K-1} \bmod 2^{K}$ --- the ideal target --- and still cannot fire, because
its own defining congruence $p \equiv 1 + 2^{K-1} \pmod{2^{K}}$ forces
$p \ge 1 + 2^{K-1}$, so with $N = p-K$, $h = H$, $L = K$ the radius $p+H+2$
exceeds $2^{K-1}$ for every prime $p$ and every $K$.
\coord{two-adic-pulse} \scale{cofinal} \ev{Lean}
\lean{Erdos257PeriodNoncollapse.exists\_prime\_deltaTotient\_twoAdic\_pulseBlock}{TotientTwoAdicPulseBlock.lean:211}
\lean{Erdos257PeriodNoncollapse.windowDiscrepancy\_modEq\_half\_of\_twoAdic\_pulse}{TotientTwoAdicPulseBlock.lean:262}
\end{prop}

\begin{rem}[B4's own recorded repair, and its scope]
\label{rem:b4}
Two further Lean facts show that prescribing \emph{letters} is exhausted: the
full terminal dyadic staircase is unconditionally impossible --- its terminal
letter would have to be positive, strictly below a wider modulus, and divisible
by it, hence $0$
(\lean{Erdos257PeriodNoncollapse.not\_actualLcmTerminalDyadicStaircase\_of\_room}{TotientActualLcmTopEdgeStaircase.lean:251});
and the surviving \emph{punctured} staircase pins its penultimate letter to
exactly $2^{m-1}$ with $2^{m} < 2(2H{+}J{+}K{+}2)$, i.e.\ no slack at all
(\texttt{puncturedDyadicStaircase\_penultimate\_eq\_half},
\texttt{TotientActualLcmTopEdgeStaircase.lean:1187}). The repair the barrier
permits is explicit: impose the half-turn on the \emph{word} rather than on one
delta, i.e.\ ask for cofinally many $(h,N,L)$ with
$\mathrm{windowDiscrepancy}(h,N,L) \equiv 2^{L-1} \pmod{2^{L}}$ and
$2^{L-1} > N{+}h{+}L{+}2$. What is \emph{proved} is the failure of the
Dirichlet/CRT residue-engineering family by an explicit exponential margin. The
broader reading --- that only carry accumulation across the weighted word can
produce the required amplitude, since one letter obeys
$\lvert \varphi(n{+}h) - \varphi(n) \rvert < n{+}h$ while the weighted word can
reach $2^{L}$ --- is an inference from the proved instance plus the linear
growth bound, not itself a theorem on disk. \ev{Math}
\end{rem}

\begin{prop}[B5, bounded state and local signatures]
\label{prop:b5}
\textbf{Status: proved} (three Lean theorems, all problem-agnostic).
\begin{enumerate}[label=(\roman*),leftmargin=*]
\item No bounded or autonomous carry state summarising the history before
  position $m$ determines the tail after $m$: for a balanced-pulse family whose
  predecessor state is constant, no $\mathrm{decode}$ recovers the radius
  parameter, and any finite state type needs cardinality
  $\ge \lfloor m/2 \rfloor + 2$, unbounded in $m$
  (\lean{Erdos249257.balancedPulse\_no\_autonomous\_decoder}{GenericTailOrbitRigidity.lean:247}).
\item Long common suffixes erase predecessor information \emph{exactly}: two
  affine binary orbits with different seeds satisfy
  $\mathrm{orbit}_u(L) - \mathrm{orbit}_v(L) = 2^{L}(u_0 - v_0)$, so the
  endpoint residue mod $2^{L}$ is independent of the initial carry
  (\lean{Erdos249257.affineBinaryOrbit\_mod\_twoPow\_eq}{GenericTailOrbitRigidity.lean:307}).
\item Bounded local $2$-adic valuation-unit data at fixed precision excludes
  nothing: for any finite word of odd-unit symbols at fixed precision and any
  starting carry, a compatible orbit exists with every intermediate state
  centred inside its symbol's dyadic radius
  (\lean{Erdos249257.fixedPrecisionTropicalNoGo}{TropicalCurvatureCarry.lean:137}).
\end{enumerate}
\textbf{Rules out:} the finite-automaton-computes-the-expansion family
outright; any hope of distinguishing two carry histories after a long common
suffix; and any contradiction derived from a fixed-precision valuation-unit
signature. Growing precision is mandatory, not optional. None of the three
mentions $\varphi$, so all three bind \#257 identically.
\coord{carry-orbit} \scale{uniform} \ev{Lean}
\end{prop}

\begin{prop}[B6, no finite linear compression]
\label{prop:b6}
\textbf{Status: proved for four independent truncations} (Lean), with the
honest limit stated below. Every natural finite truncation of the totient
kernel has trivial kernel:
\begin{enumerate}[label=(\roman*),leftmargin=*]
\item \emph{Dyadic.} The canonical family of $2^{e}+1$ dyadic channels is
  linearly independent over $\Q$ at every depth
  (Proposition~\ref{prop:rank}), and the natural repair --- a bounded
  compressed-adjoint certificate --- is impossible
  (\lean{Erdos249257.false\_of\_compressedAdjointCertificate}{TotientMahlerDefect.lean:1183},
  structure at \texttt{:1039}).
\item \emph{M\"obius incidence.} $U_N(i,j) = \mu((i{+}1)/(j{+}1))$ when
  $(j{+}1) \mid (i{+}1)$ and $0$ otherwise is lower triangular with unit
  diagonal, so $\det U_N = 1$ for every $N$ and the jet map is injective: no
  finite incidence-quotient relation exists at any horizon
  (\lean{Erdos249257.incidenceMobius\_det\_eq\_one}{IncidenceQuotientHermitePade.lean:56},
  \texttt{mobiusCompanionJetMap\_injective} at \texttt{:77}).
\item \emph{Adjugate reconstruction.} Any finite rational row exactly isolating
  one totient value has crude two-tail cost $\ge 3$, hence never $< 1$, at any
  finite grid height, using only $\varphi(x) \le x$
  (\lean{Erdos257PeriodNoncollapse.three\_le\_totientAdjugateTailCost}{TotientTailCarryPeriod.lean:266},
  closure at \texttt{:307}).
\item \emph{Shift-polynomial.} An explicit nonzero all-horizon countermodel
  agrees with every exact whole-ray anchor and survives every finite
  commensurate LCM-cube shift polynomial, at every finite rank
  (\texttt{LcmFactorIdealPulseObstruction.lean}; module docstring, theorem
  bodies not individually re-verified in this pass, and the construction is
  explicitly synthetic --- it does not claim its compensation letters occur as
  actual totient differences).
\end{enumerate}
Additionally, every strict-subrank monomial quotient in the M\"obius--Mersenne
ladder overshoots its target by more than $1/480$, \emph{uniformly} --- every
rung $r \ge 3$ lies in $[1429/1512, 1)$ and every prefix after four atoms is
within $1/3584$ of its rung, so this is a uniform no-go rather than a census
(\lean{rankOneSubrankQuotient\_sub\_theta\_two\_gt}{ErdosProblems/Erdos249/RankOneSubrankObstruction.lean:238}).
\textbf{Honest limit:} this is four checked truncations, not a proof that no
finite-linear shortcut exists. Its correct reading is the one the source module
gives: any winning finite-linear shortcut must live in a genuinely different
coordinate or use a growing-parameter construction. The reason the barrier
bites is Proposition~\ref{prop:period-not-rank}: what rationality actually buys
is periodicity, and periodicity provably does not promote to a rank bound.
\coord{carry-rank} \scale{uniform} \ev{Lean}
\end{prop}

\begin{prop}[B7, the coarse-invariant barrier]
\label{prop:b7}
\textbf{Status: proved} (Lean, \texttt{\#print axioms} clean). This is
Proposition~\ref{prop:parity} read as an elimination.
\textbf{Rules out:} any proof of \#249 whose hypothesis set on the coefficient
word lies inside \{uniform boundedness, $c(n) \le n$, agreement with $\varphi$
mod $2$ at every index, failure of eventual periodicity in any strength up to
arbitrarily long, arbitrarily separated blocks\}. The witness satisfies every
one of those and is rational.
\textbf{Does not rule out:} arguments using actual quantitative totient size or
residue information. That complement is precisely why the corpus's standing
lesson is that only a quantitative argument in the actual-LCM or fixed-rank
style can close \#249. Any future sufficient condition stated purely in terms
of coefficient-word properties must be checked against this fixture before
being trusted, for either problem.
\coord{coefficient-word} \scale{uniform} \ev{Lean}
\end{prop}

\subsection{The shape}

Put together, the barriers say something more specific than ``the problem is
hard''. The certificate residue is a weighted accumulation over an unboundedly
long window at full arithmetic resolution, and \emph{every device that makes an
argument finite destroys precisely the quantity being measured}: bounding the
window (B1), coarsening the values (B7), compressing the state (B5),
truncating the rank (B6), or letting the prover pick the sample (B2). B4
closes the remaining escape --- buying the amplitude by construction --- with
an explicit exponential margin, and B3 records that in practice nothing
bounded has ever promoted.

There is exactly one known way to hold unbounded range and full resolution at
once without a finite bookkeeping device: stop naming an index and assert an
average instead. That is what an exponential-sum bound is. It is why the
surviving routes of \S\ref{sec:survivors} are, with the honest exceptions noted
there, the ones that never name a good index.

\section{What the wall does not block}
\label{sec:survivors}

An elimination is only worth the paper it is written on if it leaves somewhere
to stand. This section is that payoff. Five routes survive the classification
of \S\ref{sec:wall}; for each we state the exact statement it needs, which
barrier it evades, and \emph{why} the evasion is structural rather than
accidental. Two of the five are demoted explicitly, because they evade every
\emph{proved} barrier while sitting squarely inside the class that the
strongest \emph{observed} pattern (Observation~\ref{obs:pointwise}) indicts.
None of the five is a proof, and none is close to one; what has changed is that
the space of things to try is no longer large.

\subsection{Survivor 1: the first-harmonic block cancellation bound}

\begin{defn}[The open analytic obligation]
\label{defn:fh}
Write $e(x) := \exp(2\pi i x)$ and
$\mathrm{windowFirstExp}(h,N,L) := e\bigl((\mathrm{windowDiscrepancy}(h,N,L) \bmod 2^{L})/2^{L}\bigr)$.
$\mathrm{DTWFirstHarmonicNormGap}$ is the statement
\[
  \forall h \ge 1\ \forall X_0\ \exists X, L : \quad
  \max(X_0,1) \le X, \quad
  16(2X + h + L + 2) \le 2^{L}, \quad
  \Bigl\lVert \sum_{N \in [X,\,2X)} \mathrm{windowFirstExp}(h,N,L) \Bigr\rVert
  \le \tfrac{21}{25}X .
\]
It implies $\mathrm{Irrational}(S)$.
\coord{first-harmonic} \scale{cofinal} \ev{Lean}
\lean{Erdos249257.DTWFirstHarmonicNormGap}{FirstHarmonicPivot.lean:75}
\lean{Erdos249257.irrational\_totient\_series\_of\_first\_harmonic\_norm\_gap}{FirstHarmonicPivot.lean:83}
\end{defn}

The real-part form suffices and is strictly weaker: it is enough that
$\sum_{N \in [X,2X)} \mathrm{windowFirstCos}(h,N,L) \le \tfrac{9}{10}X$
(\lean{Erdos249257.exists\_certifiedKill\_of\_first\_harmonic\_gap}{FirstHarmonicGap.lean:128},
the elementary unconditional engine; the $21/25 \to 9/10$ bridge is
\lean{Erdos249257.exists\_certifiedKill\_of\_first\_harmonic\_norm\_gap}{FirstHarmonicPivot.lean:53}).
Unpacked, the required object is a constant-saving cancellation estimate for
the dyadically weighted totient-difference exponential sum
\[
  \sum_{X \le N < 2X}
  e\!\left(\frac{\sum_{j<L}\bigl(\varphi(N{+}h{+}1{+}j)-\varphi(N{+}1{+}j)\bigr)2^{\,L-1-j}}{2^{L}}\right),
  \qquad L \approx \log_2 X + O(1).
\]
\textbf{Not one instance of this bound is proved anywhere, at any $X$, $h$ or
$L$.} Both trees were searched for a theorem supplying it; only consumers
exist. \ev{Open}

\paragraph{Why it evades the wall.} This is the only route in the corpus that
evades all seven barrier classes, and the reasons are structural.

\begin{itemize}[leftmargin=*]
\item \textbf{B1.} $X$ is a free unbounded parameter and the room condition
  $16(2X{+}h{+}L{+}2) \le 2^{L}$ forces $L \gtrsim \log_2 X + 5$, so the
  estimate reads $\varphi$ at every index in $[X{+}1,\,2X{+}h{+}L]$ with both
  position and window length growing. No $\gamma$ agreeing with $\varphi$ only
  on $[1,B]$ constrains it.
\item \textbf{B2 --- the decisive point.} The proved collapse
  $\mathrm{DTWWindowSeparatedPairs} \Leftrightarrow \mathrm{Irrational}(S)$
  works \emph{because the prover may choose the sample}: the converse proof
  selects $T = \{N, N{+}1\}$ and meets the threshold with $\delta = 9/10$
  supplied by irrationality plus doubling expansivity
  (Proposition~\ref{prop:b2}(b)). Definition~\ref{defn:fh} admits no sample
  choice: the sum ranges over all of $[X,2X)$ at one common depth $L$. Since
  $\neg\,\mathrm{certifiedKill}(h,N,L)$ forces
  $\mathrm{windowFirstCos}(h,N,L) > 9/10$
  (\lean{Erdos249257.windowFirstCos\_gt\_of\_not\_certifiedKill}{FirstHarmonicGap.lean:44}),
  the gap demands that a \emph{positive proportion} of basepoints in the block
  carry certificates at a \emph{single} depth. Irrationality supplies only one
  certificate per $(h,N_0)$, at a depth that may vary with $N$. The collapse
  mechanism therefore has no purchase.
  \textbf{Honest limit:} no theorem proves
  $\mathrm{DTWFirstHarmonicNormGap}$ inequivalent to $\mathrm{Irrational}(S)$.
  What is asserted is that the one proved collapse mechanism in this lane
  demonstrably requires sample choice, which this predicate denies. Note also
  that the \emph{subset} form
  (\lean{Erdos249257.exists\_certifiedKill\_of\_first\_harmonic\_gap\_subset}{FirstHarmonicGap.lean:104})
  does permit sample choice and is therefore the collapse-exposed variant.
  Attack the full-block form.
\item \textbf{B4.} The statistic is the first additive character of the
  \emph{accumulated} weighted word, never an individual letter. Nothing is
  prescribed at any position, so the self-defeating cost that kills Dirichlet
  residue engineering does not arise. B4's own recorded repair
  (Remark~\ref{rem:b4}) is to impose the structure on the word rather than on
  one delta, which is exactly what a block exponential-sum bound does.
\item \textbf{B5, B6, B7.} No carry state is summarised and no fixed precision
  is used; the estimate is on the true residue mod $2^{L}$ with $L \to \infty$.
  A cancellation bound is not a rank statement and seeks no finite-dimensional
  relation. And the estimate uses the actual real residue, not parity,
  boundedness or aperiodicity, so neither countermodel touches it --- both
  differ from $\varphi$ at exactly the large indices where this sum lives.
\item \textbf{B3.} The audit's own verdict on this row is that the consumer
  side is \emph{finished} and already at the obligation's exact quantifier
  shape: $X$ is free, so applying the bound at $X \ge N_0$ gives
  $\exists N \ge N_0 \exists L$ directly. The missing piece is one named
  arithmetic fact, not a strengthening of anything on disk. The constants
  $9/10$, $\pi/8$ and $16$ are absolute and do not degrade with $X$, $h$ or
  $L$; there is no case analysis and no table.
\end{itemize}

\begin{rem}[Does B1 \emph{select} this route? No.]
\label{rem:b1-not-selector}
It is tempting to say the finite-inspection theorem points at analysis. It does
not. B1 says any proof must use $\varphi$ at unbounded indices --- a necessary
condition satisfied by \emph{every} open route here (the LCM diagonal as
$t \to \infty$, the actual-LCM supply for $a \ge 8$, the Farey growth law as
$K \to \infty$, the rank bound for all $e$). B1 does not discriminate.

\par\medskip

What selects this route is a sharper pair: (1) B4, proved, kills the only mechanism
the corpus ever found for reaching full resolution at unbounded range by
\emph{prescribing} values, and its own stated repair is accumulation over the
word --- which is what an exponential sum measures; and (2)
Observation~\ref{obs:pointwise}, an observed pattern and not a theorem, shows
that every single-index producer has failed to be supplied at even one large
index, while this route's missing input is a block average and therefore never
requires naming a good index.

\par\medskip

B1's real contribution is narrower but still load-bearing: it proves that the finite deposits can never be extended into a
proof, so the gap between the census and the obligation is a gap in kind, not
in degree. That is a genuine no-go about method, of the same species as a
relativization barrier. It is not a selector.
\end{rem}

\begin{rem}[Evidence in both directions]
Weak supporting evidence, observed and not proof: every landed diagonal
certificate through $t = 64$ fires within a small additive constant of the
minimum depth forced by $\mathrm{certifiedKill\_depth\_floor}$, which is the
profile a doubling-orbit equidistribution argument would produce. Countervailing
evidence, also observed: the strict-LCM-jump census records a closest central
margin of $\approx 0.000221$ of the modulus at $t = 100$, so these residues are
not robustly central and any equidistribution claim will be delicate. \ev{Cert}
\end{rem}

\subsection{Survivor 2: the four-term pivot budget}

Same lane as Survivor 1, with the analytic burden repackaged. Split the block
sum at a largest-prime pivot into centred correlation, fibre-mean,
bad-cofactor and non-supplier contributions and bound the four separately.

\begin{defn}[$\mathrm{DTWPivotResidualDecorrelation}$]
\label{defn:pivot}
For every $h > 0$ there are $s > 0$ and $\eta \in (0,1)$ such that for every
$X_0$ there are $X, L$ with $\max(X_0,1) \le X$, $h \le L - s$,
$16(2X{+}h{+}L{+}2) \le 2^{L}$, and all four of
\[
  \mathrm{Re}\,\mathrm{pivotCenteredCorrelation} \le \tfrac{14}{25}X, \quad
  \lVert \mathrm{pivotFiberMean} \rVert \le \tfrac{1}{100}X, \quad
  \lVert \mathrm{pivotBad} \rVert \le \tfrac{1}{100}X, \quad
  \lVert \mathrm{pivotNonSupplier} \rVert \le \tfrac{8}{25}X .
\]
It implies $\mathrm{Irrational}(S)$.
\coord{first-harmonic} \scale{cofinal} \ev{Lean}
\lean{Erdos249257.DTWPivotResidualDecorrelation}{FirstHarmonicPivot.lean:569}
\lean{Erdos249257.irrational\_totient\_series\_of\_pivotResidualDecorrelation}{FirstHarmonicPivot.lean:577}
\end{defn}

It inherits Survivor 1's evasions of B1, B4, B5, B6 and B7 verbatim: the four
terms are exact finite sums over canonical largest-prime supplier fibres of the
same accumulated-word phases, with no sample choice and nothing prescribed. Its
specific advantages are three.

The decomposition is an \emph{exact identity},
proved unconditionally with \texttt{\#print axioms} clean
(\lean{Erdos249257.windowFirstExp\_sum\_eq\_pivot\_decomposition}{FirstHarmonicPivot.lean:514},
budget consumer
\lean{Erdos249257.first\_harmonic\_gap\_of\_pivotBudgetAt}{FirstHarmonicPivot.lean:549}).
The one-sided budget lowers the hard requirement from a norm bound
$\lVert \cdot \rVert \le X/2$ to a real-part bound $\le 14X/25$. And the pivot
rests on genuine arithmetic rather than a sampled surrogate: on the canonical
fibre $m = 1$, $s = L-h$ the supplier set is \emph{literally} the shifted dyadic
interval of primes, proved as a membership equality
(\lean{Erdos249257.mem\_pivotFiber\_one\_overlap\_iff}{FirstHarmonicPivot.lean:423}),
with the totient factorisation at the pivot the honest
$\varphi(mp) = \varphi(m)(p-1)$ and the non-divisibility discharged by size.

Three of the four terms are then counting bookkeeping; only the first needs a
genuine correlation estimate, and only in real part.

\textbf{What it must respect.} B6 has a residual-gauge instance:
a residual-blind determinant or conditioning test cannot certify that genuine
phase reconstruction rather than a locked degenerate configuration has occurred
(\lean{Erdos249257.locked\_reconstruction\_preserves\_nonzero\_minor}{ResidualGaugeObstruction.lean:94}).
So this route must couple rows by an extra arithmetic identity, not merely
gauge-normalise columns. \textbf{Honest status:} strictly a repackaging of
Survivor 1's burden; it is listed separately only because three quarters of it
are already exact identities on disk. The module asserts no prime-distribution
or decorrelation estimate --- the socket is deliberately empty. \ev{Open}

\subsection{Survivor 3: uniform quantitative escape at cofinally many primes}

\begin{defn}[$\mathrm{DTWNaturalPrimeTailOrbitStrictGap}$]
\label{defn:primegap}
For every $h \ge 1$ and every $N_0$ there is a prime $p$ with
$\max(N_0{+}h{+}1,\, h{+}5) \le p$ and
$\mathrm{Re}\,\mathrm{tailOrbitFirstExp}(h,\,p{-}h{-}1) < \tfrac{9}{10}$,
where $\mathrm{tailOrbitFirstExp}(h,N) = e(R_{N+h} - R_N)$. Equivalently:
cofinally many primes $p$ at which the totient tail difference across the shift
$h$ stays a \emph{fixed} distance
$\ge \arccos(9/10)/2\pi \approx 0.0718$ from every integer.
It implies $\mathrm{Irrational}(S)$.
\coord{prime-pivot} \scale{cofinal} \ev{Lean}
\lean{DTWNaturalPrimeTailOrbitStrictGap}{ErdosProblems/Erdos249/TotientStrictPrimeEscape.lean:157}
\lean{irrational\_totient\_series\_of\_naturalPrimeTailOrbitStrictGap}{ErdosProblems/Erdos249/TotientStrictPrimeEscape.lean:524}
\end{defn}

It evades B1, B4, B5, B6 and B7 for the reasons given for Survivor 1: a
statement about the true real tail difference at unbounded prime positions,
with no letters prescribed, no bounded state, no rank, no coarse
coefficient-word invariant. It evades B2 by the \emph{second} of the two
exemptions in Proposition~\ref{prop:b2} --- not by denying sample choice, but by
demanding a \emph{uniform margin}. Irrationality gives non-integrality of every
tail difference and supplies no lower bound whatever on distance to $\Z$;
certificate completeness therefore cannot manufacture this predicate.

\textbf{Honest demotion.} This is a \emph{pointwise} producer: it requires
naming a good prime, and Observation~\ref{obs:pointwise} records that no
pointwise producer in this corpus has ever been supplied at even one large
index, across seven independent attempts. It evades every proved barrier while
sitting inside the class the strongest observed pattern indicts. Nothing on
disk proves it at a single prime for a single $h$. It ranks below Survivor 1.
\ev{Open}

\subsection{Survivor 4: depth-locked full-depth escape}

\begin{defn}[$\mathrm{ApFullDepthEscape}$]
\label{defn:apfde}
For every $d \ge 1$ and every $N$ there is $t \ge 1$ with
$\mathrm{certifiedKill}(td,\,N,\,td)$; unpacked,
\[
  N + 2td + 2 \;<\; \mathrm{windowDiscrepancy}(td,N,td) \bmod 2^{td}
  \;<\; 2^{td} - (N + 2td + 2).
\]
It implies $\mathrm{Irrational}(S)$.
\coord{period-ray} \scale{cofinal} \ev{Lean}
\lean{ApFullDepthEscape}{ErdosProblems/Erdos249/PeriodMultipleEscape.lean:441}
\lean{irrational\_totient\_series\_of\_apFullDepthEscape}{ErdosProblems/Erdos249/PeriodMultipleEscape.lean:450}
\end{defn}

This is the shortest fully stated open inequality the programme has produced.
Its evasion of B2 is recorded on disk rather than argued: its ambient parent
$\mathrm{PeriodMultipleKillSupply}$ is \emph{proved} equivalent to
$\mathrm{Irrational}(S)$ and $\mathrm{ApFullDepthEscape}$ implies it
(\lean{periodMultipleKillSupply\_of\_apFullDepthEscape}{ErdosProblems/Erdos249/PeriodMultipleEscape.lean:444}),
so it is at least as strong, and possibly strictly stronger --- the file's own
docstring records ``sufficient for irrationality; not known necessary''. Being
possibly strictly stronger, certificate completeness cannot collapse it back.
It evades B1 (both $N$ and $t$ unbounded), B4 (locking depth to the period
means the entire word accumulates and no letter is prescribed), B5 and B6 (no
state, no rank). Its substance is pure anti-concentration: via
\lean{windowDiscrepancy\_self\_eq\_totientBlock\_sub}{ErdosProblems/Erdos249/PeriodMultipleEscape.lean:67},
the difference of two adjacent period blocks must have central residue mod
$2^{h}$ at some multiple period $h = td$, and the room condition is automatic
once $h$ is large, so nothing but the residue's position is at stake.

\textbf{Honest demotion,} identical to Survivor 3: it is a pointwise producer
and names an index. It is listed because it is the cleanest target for
computational exploration --- with Theorem~\ref{thm:gamma} as the standing
reminder that no amount of such exploration becomes a proof. \ev{Open}

\subsection{Survivor 5: the rationality-side rank upper bound}

\begin{defn}[The coordinate-disjoint obligation]
\label{defn:rankupper}
Either: there is $C$ such that for every $c : \N \to \N$ with $c(n) \le n$ and
$T_c \notin$ the irrationals, every $v > 0$ and every tempered binary orbit
$u$ for $(c,v)$, and every $e$,
$\dim_{\Q} \mathrm{span}_{\Q}\bigl(\mathrm{canonicalCarryKernelFamily}(u,e)\bigr) \le C$;
or the same with any $g(e)$ growing strictly slower than $2^{e}-1$ in place of
$C$. Either version contradicts the proved floor of
Proposition~\ref{prop:rank} and closes \#249.
\coord{carry-rank} \scale{uniform} \ev{Open}
\end{defn}

This is the only surviving obligation that is not a residue or certificate
statement at all, so B2's certificate completeness cannot reach it --- it is
not a reformulation of $\mathrm{Sep}$ in any vocabulary. It evades B1 and B4
entirely (no window, no letters, no prescribed residues) and B5 (linear algebra
over the whole orbit, not a bounded-state summary). The scale side is finished
and uniform in $e$ (Proposition~\ref{prop:rank}).

\textbf{Honest flag, and it is severe.} B6 partly indicts this route from
inside. The most natural approach to the upper bound is dead
(\lean{Erdos249257.false\_of\_compressedAdjointCertificate}{TotientMahlerDefect.lean:1183}),
and what rationality actually buys --- uniform eventual periodicity of the
carry's dyadic sections mod $v$ --- provably does not promote to a $\Q$-rank
bound (Proposition~\ref{prop:period-not-rank}). So this survivor is genuinely
coordinate-disjoint from everything else, which is its whole value, but its
obvious approach is closed and the missing input is a rigidity theorem nobody
has stated.

\section{Reading the detailed record}

The opening note has already defined the evidence bands, scale tags,
coordinates, and catalogue terminology.  This section therefore gives only
the order of dependence needed for the longer reference material.

\paragraph{Scale and coordinate discipline.} Every direct certificate deposit
for $\mathrm{Sep}$ is fixed or bounded, although the corpus also contains
uniform and cofinal structural theorems.  Those stronger-scale results do not
produce $\mathrm{Sep}$, and Theorem~\ref{thm:gamma} shows why extending bounded
verification does not establish the cofinal statement.  The coordinate tags
are equally literal: the $\gamma$-splice concerns \coord{binary-digit}, the
rank floor \coord{carry-rank}, positivity \coord{actual-lcm}, and the Farey
exclusion \coord{farey}.  An obstruction in one coordinate does not transfer
without a proved map.

\paragraph{Order of dependence.} Read the wall in \S\ref{sec:wall} before
selecting an attack.  In particular, B4 governs prescribed totient values, B5
bounded carry states, and B2 purportedly weaker certificate conditions.  The
mathematical-ingredient catalogue that follows may then be entered by
coordinate: it records hypotheses, evidence band, scale, and Lean declaration.
The later tables give the exact implication graph, the near misses at each open
obligation, and the promotion audit summarised as B3.  These tables are
reference material; the wall and the surviving-route argument are intended to
be read linearly.
% ---- part a249_p0 ----
\section{The object, the target, and the exact record}

This part fixes the object of Erd\H{o}s Problem \#249, states exactly what is
open about it, and records every unconditional numeral, identity and
Lean-checked bound the corpus currently holds. Nothing in this part decides
irrationality. Every result below is either an identity (no hypothesis, holds
unconditionally), a finite computation (holds up to an explicitly stated
bound), or is flagged \ev{Cited}/\ev{Open} where that is the honest status.

\subsection{The object}

\begin{defn}[The Erd\H{o}s--249 constant]
\label{defn:S}
Let $\varphi$ denote Euler's totient function. Define
\[
S \;:=\; \sum_{n \ge 1} \frac{\varphi(n)}{2^{n}} \;=\; \sum_{n : \mathbb{N}} \frac{\varphi(n)}{2^{n}},
\]
the two forms coinciding because $\varphi(0)=0$; the second, $\mathbb{N}$-indexed
form is the one carried in the Lean source. The series converges absolutely
since $\varphi(n) = O(n)$. \emph{Erd\H{o}s \#249 asks whether $S$ is irrational.
This is OPEN.} Nothing in this paper decides it; every claim below is either an
unconditional identity, a finite computation with an explicitly stated range,
or is marked \ev{Cited}/\ev{Open}.
\end{defn}

\begin{rem}[Status]
No proof or disproof of $\mathrm{Irrational}(S)$ exists anywhere in the corpus,
formal or informal. What exists is: one large unconditional finite denominator
exclusion (\S\ref{ssec:farey}), several exact reformulations of $S$ that
relocate the same open question onto different coordinates without touching
its truth value (\S\ref{ssec:coprime}--\S\ref{ssec:lambert}), and a certificate
apparatus whose cofinal supply obligation is the precise open target
(recorded in full in Part~1 of this paper; only its finite, checked instances
are catalogued here, \S\ref{ssec:certtable}).
\end{rem}

\subsection{The binary-digit reduction}

\begin{prop}[Digit-shift identity]
\label{prop:shift}
For every $N : \mathbb{N}$,
\[
2^{N} \cdot S \;=\; \Phi_N + R_N,
\]
where $\Phi_N := \sum_{n \le N} \varphi(n) \cdot 2^{N-n} \in \mathbb{N}$ (the
integer prefix) and $R_N := \sum_{j \ge 0} \varphi(N+1+j)/2^{j+1}$ (the
fractional tail, \texttt{totientTail} in Lean).
\coord{binary-digit} \scale{uniform} \ev{Lean}
\lean{two_pow_mul_totient_series_eq}{TotientTailPeriodKiller.lean:150}
\end{prop}

\begin{proof}[Consequence]
Rationality and eventual periodicity of the base-2 expansion of $S$ is
\emph{exactly} a statement about the tail $R_N$: $S \in \mathbb{Q}$ forces
(by Euler's theorem applied to the odd part of the denominator) a period
$h>0$ and pre-period $N_0$ with $R_{N+h}-R_N \in \mathbb{Z}$ for all
$N \ge N_0$ (\texttt{eventual_period_of_not_irrational},
\lean{eventual_period_of_not_irrational}{TotientTailPeriodKiller.lean:358},
\ev{Lean}, explicit witness $h=\varphi(\mathrm{oddPart}(r.\mathrm{den}))$,
$N_0=v_2(r.\mathrm{den})$). Nothing in the corpus supplies the converse
cofinally; this is the certificate wall documented in full in Part~1 and is
not restated here beyond this pointer.
\end{proof}

\subsection{The unconditional denominator floor}
\label{ssec:farey}

The strongest \emph{unconditional} fact the corpus holds about $S$ is a lower
bound on its denominator should it happen to be rational. It is obtained by a
completely elementary Farey/mediant argument, entirely free of any certificate
or period apparatus, and is therefore logically independent of the open
certificate-supply obligation.

\begin{lem}[Farey gap, fully general]
\label{lem:farey}
For integers $a,b,c,d,r,s$ with $b>0$, $d>0$, $bc-ad=1$ (i.e.\ $a/b$ and
$c/d$ are unimodular Farey neighbours), and $as < rb$, $rd < cs$ (i.e.\ $r/s$
lies strictly between them): $b+d \le s$.
\coord{farey} \scale{n/a} \ev{Lean}
\lean{farey_gap}{GapFareyBound.lean:51}
\end{lem}

This lemma is a classical Stern--Brocot fact with zero totient or Mersenne
content; it is the generic engine underneath every Farey-gap bound in the
corpus, for either open problem.

\begin{prop}[The wave-17 gap certificate at window $K=240$]
\label{prop:gapwindow}
Let $V$ be the explicit committed totient residue for window $(N,K)=(1,240)$.
For every $q : \mathbb{N}$ with
\[
0 < q \;\le\; Q_0 := 79\,639\,646\,646\,701\,375\,323\,355\,774\,875\,831\,053 \;\;(\approx 7.96\times 10^{34}),
\]
\[
(q \cdot V) \bmod 2^{240} \;+\; 243\,q \;<\; 2^{240}.
\]
This bound is \emph{sharp}: $q = Q_0+1 = 79\,639\,646\,646\,701\,375\,323\,355\,774\,875\,831\,054$
is the exact first failing denominator, exhibited as the mediant of two
explicit unimodular Farey neighbours.
\coord{farey} \scale{bounded} \ev{Lean}
\lean{gap_check_window_1_240_le_79639646646701375323355774875831053}{GapFareyBound.lean:176}
\lean{gap_check_window_1_240_first_failure}{GapFareyBound.lean:225}
\end{prop}

\begin{thm}[Erd\H{o}s \#249 denominator exclusion --- the headline unconditional result]
\label{thm:denom-record}
For every $p \in \mathbb{Q}$ with reduced denominator $p.\mathrm{den} \le Q_0$,
\[
S \;\neq\; p.
\]
Equivalently: \emph{if $S$ is rational, its reduced denominator exceeds}
$Q_0 \approx 7.96 \times 10^{34}$.
\coord{farey} \scale{bounded} \ev{Lean}
\lean{tsum_totient_div_pow_two_ne_ratCast_of_den_le_79639646646701375323355774875831053}{CertificateKernel.lean:18384}
\end{thm}

\begin{rem}[What this does and does not say]
Theorem~\ref{thm:denom-record} is a complete, unconditional finite fact: no rational
of small denominator equals $S$. It says nothing about arbitrarily large
denominators and is not itself a route to irrationality. The docstring behind
Proposition~\ref{prop:gapwindow} states the honest open question this leaves:
whether $\sup_K (b+d)(K)$ (the growing analogue of $Q_0$ as the window $K$
grows) is unbounded as $K \to \infty$ --- which would close \#249 through this
theorem's consumer, entirely independently of the certificate-supply
obligation of Proposition~\ref{prop:shift}. \ev{Open}; not claimed or proved
anywhere in the corpus.
\end{rem}

\begin{rem}[Ladder of prior rungs]
$Q_0$ is the current end of an explicit sequence of increasingly wide Farey
windows computed by the same mechanism: $4838 \to 2^{22} \to 2.49\times10^{17}$
(window $K=120$) $\to Q_0$ (window $K=240$). Each rung is a finite,
independently checked instance of Lemma~\ref{lem:farey} and
Proposition~\ref{prop:gapwindow}'s pattern at a larger $K$; none is claimed to
extrapolate.
\end{rem}

\subsection{The same record, transported: the M\"obius-square and coprimality-probability forms}
\label{ssec:coprime}

The exact finite Farey record behind Theorem~\ref{thm:denom-record} is not tied to
the $\varphi(n)/2^n$ presentation of $S$; it transfers verbatim to two other
exact reformulations of the same constant, at exactly half the bound.

\begin{prop}[Fair-coin coprimality form]
\label{prop:coprime}
Let $X,Y$ be independent random variables with $\Pr(X=n)=\Pr(Y=n)=2^{-n}$
for $n \ge 1$ (independent fair-coin waiting times). Then
\[
S \;=\; \tfrac12 \;+\; \Pr\bigl(\gcd(X,Y)=1\bigr)
\;=\; \tfrac12 \;+\; \sum_{\substack{a,b\ge 1\\ \gcd(a,b)=1}} 2^{-(a+b)}.
\]
Equivalently, on the visible lattice: summing $2^{-(a+b)}$ over the half-open
coprime pairs ($a\ge1$, $b\ge0$, $\gcd(a,b)=1$) recovers $\sum_n \varphi(n)/2^n$
exactly, with no boundary correction, because the visible-point count on the
half-open antidiagonal at height $n$ equals $\varphi(n)$ for every $n$,
including $n=0,1$.
\coord{probability} \scale{n/a} \ev{Lean}
\lean{tsum_visible_coprime_pairs_eq_totient_series}{CertificateKernel.lean:18544}
\lean{totient_series_eq_half_add_visible_coprime_pairs}{CertificateKernel.lean:18557}
\end{prop}

\begin{prop}[The gcd-layer normalisation]
\label{prop:gcdlayer}
For independent fair-coin waiting times as above, $\sum_{g\ge1}\Pr(\gcd(X,Y)=g)=1$
exactly; and for every $d>0$, $\Pr(d\mid X \wedge d\mid Y) = 1/(2^d-1)^2$.
\coord{probability} \scale{uniform} \ev{Lean}
\lean{tsum_pos_coprime_inv_mersenne_eq_one}{GcdMomentCalculus.lean:349}
\lean{tsum_pos_pair_both_dvd_half_eq_inv_mersenne_sq}{GcdMomentCalculus.lean:266}
\end{prop}

\begin{thm}[Denominator exclusion for the coprimality-probability form]
\label{thm:denomcoprime}
Let
\[
Q_1 := \left\lfloor \frac{Q_0}{2} \right\rfloor = 39\,819\,823\,323\,350\,687\,661\,677\,887\,437\,915\,526.
\]
For every $a\in\mathbb{Z}$, $d\in\mathbb{N}$ with $0<d\le Q_1$: the visible
coprime-pair probability $\Pr(\gcd(X,Y)=1)$ is not equal to $a/d$.
\coord{farey} \scale{bounded} \ev{Lean}
\lean{tsum_visible_coprime_pairs_ne_int_div_of_den_le_39819823323350687661677887437915526}{CertificateKernel.lean:18572}
\end{thm}

\begin{thm}[Denominator exclusion for the M\"obius-square form]
\label{thm:denommobsq}
With $Q_1$ as in Theorem~\ref{thm:denomcoprime}: for every $a\in\mathbb{Z}$,
$d\in\mathbb{N}$ with $0<d\le Q_1$, the signed series
$T := \sum_{d\ge1} \mu(d)/(2^d-1)^2 = S - \tfrac12$ (see \S\ref{ssec:mobius})
is not equal to $a/d$.
\coord{farey} \scale{bounded} \ev{Lean}
\lean{tsum_moebius_div_two_pow_sub_one_sq_ne_int_div_of_den_le_39819823323350687661677887437915526}{CertificateKernel.lean:18487}
\end{thm}

\begin{rem}[Why the bound halves, and why this is not new information]
$Q_1 = \lfloor Q_0/2 \rfloor$ arithmetically: $Q_0$ is odd
($Q_0 = 79\,639\,646\,646\,701\,375\,323\,355\,774\,875\,831\,053$), so
$Q_0/2 = 39\,819\,823\,323\,350\,687\,661\,677\,887\,437\,915\,526.5$ and
$Q_1$ is its floor. The halving is the cost of transporting the known bound
through the affine shift
$S = \tfrac12 + T$ (resp.\ $S = \tfrac12 + \Pr(\gcd(X,Y)=1)$) on a
denominator-exclusion statement: excluding all denominators $\le Q_1$ for $T$
follows from excluding all denominators $\le Q_0$ for $S$ (a denominator-$d$
value of $T$ with $d\le Q_1$ yields a
denominator dividing $2d\le Q_0$ for $S$). Theorems~\ref{thm:denomcoprime}
and~\ref{thm:denommobsq} are therefore the \emph{same} finite Farey record as
Theorem~\ref{thm:denom-record}, transported through Proposition~\ref{prop:coprime}
and Proposition~\ref{prop:mobsq} respectively --- not independent evidence.
The converse finite implication is not asserted: subtracting $1/2$ can double
a denominator, so the $Q_1$ exclusion for $T$ alone need not recover the full
$Q_0$ exclusion for $S$.
\end{rem}

\subsection{The M\"obius--Mersenne identity and why bounded coefficients matter}
\label{ssec:mobius}

\begin{prop}[M\"obius-square reduction]
\label{prop:mobsq}
\[
S \;=\; \sum_{n\ge1}\frac{\varphi(n)}{2^n} \;=\; \frac12 \;+\; \sum_{d\ge1}\frac{\mu(d)}{(2^d-1)^2},
\]
where $\mu$ is the M\"obius function, so $\mu(d)\in\{-1,0,1\}$ for every $d$.
Consequently \emph{Erd\H{o}s \#249 $\iff$ $T:=\sum_{d\ge1}\mu(d)/(2^d-1)^2 \notin \mathbb{Q}$},
the single reduced target every other result in this subsection and
\S\ref{ssec:lambert} feeds.
\coord{mobius-mersenne} \scale{n/a} \ev{Lean}
\lean{totientSeries_eq_pnat_half_pow}{SquaredMersenneDiagonalEnclosure.lean:72}
\lean{tsum_totient_half_pow_eq_half_add_moebius_sq}{MersenneLambertLadder.lean:665}
\end{prop}

\begin{rem}[Why the bounded coefficients matter --- the Erd\H{o}s-1948 regime]
The identity of Proposition~\ref{prop:mobsq} rewrites $S$ (equivalently $T$)
as a \emph{M\"obius-twisted Lambert-squared series}: the numerator weight
$\mu(d)$ is bounded, $|\mu(d)|\le1$ for every $d$, uniformly in $d$. This
places $T$ in exactly the coefficient regime of the classical Erd\H{o}s (1948)
near-integer irrationality criterion and of the level-1 sibling identity
$L(\mu):=\sum_d \mu(d)/(2^d-1) = \tfrac12$ (rational, trivially) alongside
$L(1) = \sum_d 1/(2^d-1) = E$, the Erd\H{o}s--Borwein constant, which
\emph{is} proved irrational in this same kernel
(\lean{irrational_erdosBorwein_series}{CertificateKernel.lean:8335}, \ev{Lean}).

\par\medskip

This is the opposite regime from Proposition~\ref{prop:shift}'s
\textbf{binary-digit coordinate}, where the corresponding weight satisfies
$0\le\varphi(n)\le n$ and is unbounded.  The function $\varphi$ has average
order $6n/\pi^2$, equivalently
$\sum_{k\le x}\varphi(k)\sim 3x^2/\pi^2$, but there is no pointwise estimate
$\varphi(n)=\Theta(n)$. A
near-integer/Dirichlet-approximation argument of Erd\H{o}s-1948 shape
(formalised generically as
\lean{irrational_of_int_mul_near_int}{CertificateKernel.lean:6120} and its
base-power specialisation
\lean{irrational_of_pow_mul_near_int}{CertificateKernel.lean:6149}, both
\ev{Lean}, \coord{n/a}, fully coordinate-free) has a genuine chance of
transferring to $T$ precisely because its weight is bounded, in a way it does
not have a chance of transferring directly to the raw $\varphi(n)/2^n$ series.

\par\medskip

No such transfer is proved; \S\ref{ssec:mobius} below (cross-referenced
here, developed in Part~2 of this paper) records exactly why the transfer has
so far failed (the ``$\mu$-pollution'' obstruction) rather than merely
asserting the analogy.
\end{rem}

\subsection{The squared-Lambert gcd-moment identities}
\label{ssec:lambert}

\begin{prop}[Squared-Lambert transfer engine]
\label{prop:lambertengine}
For $w:\mathbb{N}\to\mathbb{R}$ with $|w(d)|\le d$ for all $d>0$, and $0\le r<1$:
\[
\sum_{d\ge1} w(d)\left(\frac{r^d}{1-r^d}\right)^2 \;=\; \sum_{n\ge1}\left(\sum_{e\mid n} w(e)\Bigl(\tfrac{n}{e}-1\Bigr)\right) r^n.
\]
\coord{mobius-mersenne} \scale{uniform} \ev{Lean}
\lean{tsum_lambert_linear_weight_sq_pure}{GcdMomentCalculus.lean:105}
\end{prop}

This one identity, at $r=1/2$, specialises to every squared-Lambert rung the
corpus computes; two instances are exact and directly relevant to \#249's
weight structure:

\begin{prop}[The known $\zeta_q$-rung]
\label{prop:zetaq}
\[
\sum_{d\ge1} \frac{1}{(2^d-1)^2} \;=\; \sum_{n\ge1} \frac{\sigma(n)-\tau(n)}{2^n} \;=\; \zeta_q(2)-\zeta_q(1) \text{ at } q=\tfrac12,
\]
where $\sigma$ is the sum-of-divisors function and $\tau$ the number-of-divisors
function. The \emph{identity} is machine-checked (\ev{Lean}); irrationality of
the \emph{value} $\zeta_q(2)-\zeta_q(1)$ is \ev{Cited} (Postelmans--Van Assche
$q$-Pad\'e), \emph{not} formalised in this corpus.
\coord{mobius-mersenne} \scale{n/a}
\lean{tsum_one_div_mersenne_sq_eq_sigma_sub_tau_series}{GcdMomentCalculus.lean:216}
\end{prop}

\begin{prop}[The Pillai/gcd-moment rung]
\label{prop:pillai}
\[
\sum_{d\ge1} \frac{\varphi(d)}{(2^d-1)^2} \;=\; \sum_{n\ge1} \bigl(P(n)-n\bigr)\cdot 2^{-n} \;=\; \mathbb{E}[\gcd(X,Y)],
\]
where $P(n) := \sum_{e\mid n}\varphi(e)\cdot(n/e) = (\varphi * \mathrm{Id})(n)$
is Pillai's gcd-sum function and $X,Y$ are the independent fair-coin waiting
times of Proposition~\ref{prop:coprime}. The identity itself is
machine-checked (\ev{Lean}); the value $\mathbb{E}[\gcd(X,Y)]$ is \ev{Open} ---
a cousin rung to \#249, not \#249 itself.
\coord{mobius-mersenne} \scale{n/a}
\lean{tsum_totient_div_mersenne_sq_eq_gcd_moment_series}{GcdMomentCalculus.lean:235}
\end{prop}

\begin{rem}[The level mirror]
Writing $L(f):=\sum_d f(d)/(2^d-1)$ (level 1) and $L_2(f):=\sum_d f(d)/(2^d-1)^2$
(level 2): $L(\mu)=\tfrac12$ (rational, trivial), $L(1)=E$ (Erd\H{o}s--Borwein,
irrational, \ev{Lean}), $L(\varphi)=2$ (rational); $L_2(\mu)=S-\tfrac12$
(\ev{Open}, this is \#249), $L_2(1)=\zeta_q(2)-\zeta_q(1)$ (\ev{Cited}
irrational, Proposition~\ref{prop:zetaq}), $L_2(\varphi)=\mathbb{E}[\gcd(X,Y)]$
(\ev{Open}, Proposition~\ref{prop:pillai}). At level 1 the M\"obius rung is
trivial and the $\zeta$-rung is the hard classical case; at level 2 the
$\zeta$-rung is known and the M\"obius rung \emph{is} \#249. M\"obius
projection is the one wall repeated at both levels of the ladder.
\end{rem}

\subsection{The exact certificate record: the contiguous band through $t\le82$}
\label{ssec:certtable}

The certificate apparatus itself (the predicate \texttt{certifiedKill}, its
soundness/completeness theorems, and the cofinal supply obligation that is the
actual open target) is developed in full in Part~1 of this paper. What is
recorded here is a representative historical set of machine-checked anchors,
together with the current aggregate theorem closing every scale $t\le82$.
The declaration inventory, rather than this table, is authoritative for all
individual certificate shards; no extrapolation beyond the stated band is
implied.

\paragraph{Certificate object.} For $h,N,L : \mathbb{N}$, define the window
discrepancy $\Delta_{h,N,L} := \sum_{j<L} (\varphi(N+h+1+j)-\varphi(N+1+j))\cdot 2^{L-1-j} \in \mathbb{Z}$
(\lean{windowDiscrepancy}{TotientTailPeriodKiller.lean:66}, \ev{Lean}), and
\[
\mathrm{Sep}(h,N,L) \;:\equiv\; (N+h+L+2) < \Delta_{h,N,L} \bmod 2^{L} < 2^{L}-(N+h+L+2),
\]
(\lean{certifiedKill}{TotientTailPeriodKiller.lean:72}, \ev{Lean}, decidable). By
completeness (\lean{exists_certifiedKill_iff_tail_diff_notMem_int}{LcmConeFlatness.lean:316},
\ev{Lean}), $(\exists L,\ \mathrm{Sep}(h,N,L)) \iff R_{N+h}-R_N \notin \mathbb{Z}$
for every $h,N$: an unbounded (cofinal, over $h$ and $N$) supply of
$\mathrm{Sep}$ is exactly equivalent to $\mathrm{Irrational}(S)$. The table
below records the principal $(h,N,L)$-shaped and $t$-shaped anchors; none of
them is cofinal, and none is presented as deciding \#249.

\small
\begin{longtable}{p{0.20\linewidth} p{0.40\linewidth} p{0.13\linewidth} p{0.20\linewidth}}
\textbf{Anchor} & \textbf{Exact statement of what was checked} & \textbf{Scale} & \textbf{Site} \\
\hline
\endhead
Fixed-window deposit, depth 16
&
$\mathrm{Sep}(h,12,16)$ holds for every $h \in \{1,\dots,8\}$ (by \texttt{decide},
$256$ totient values below $37$); consequently $S \ne r$ for every
$r\in\mathbb{Q}$ with $1\le h\le 8$ and $r.\mathrm{den} \mid 2^{12}(2^h-1)$.
&
\scale{fixed}
&
\lean{certifiedKill_all_small}{TotientTailPeriodKiller.lean:404}
\lean{totient_series_ne_rat_of_den_dvd}{TotientTailPeriodKiller.lean:416}
\\
Fixed-window deposit, depth 9, wider net
&
$\mathrm{Sep}(h,14,9)$ holds for every
$h\in\{1,\dots,16\}$: $S \ne r$ for every $r\in\mathbb{Q}$ with $1\le h\le16$
and $r.\mathrm{den}\mid 2^{14}(2^h-1)$.
&
\scale{fixed}
&
\lean{totient_series_ne_rat_of_den_dvd_pow_two_mul_mersenne_upto_sixteen}{CertificateKernel.lean:18890}
\\
Diagonal-pincer certificates, base 12 scales
&
$\mathrm{Sep}\bigl(H_t, H_t, D(t)\bigr)$, $H_t := \mathrm{lcm}(1,\dots,t)$,
holds (by explicit \texttt{decide}/\texttt{norm_num} on checked
\texttt{Nat.totient} values, via factored prime-power blocks with
Lucas-primality certificates) for $t \in \{1,2,3,4,5,7,8,9,11,13,16,17\}$
with certificate depths $D(t) \in \{6,5,7,7,9,14,15,14,21,22,23,26\}$
respectively (in the same order).
&
\scale{fixed}
&
\lean{certifiedKill_diagonal_all_imported}{DiagonalPincerCertificates.lean:2870}
\lean{diagonalPincerCertificateScales}{DiagonalPincerCertificates.lean:2852}
\lean{diagonalPincerKillDepth}{DiagonalPincerCertificates.lean:2854}
\\
Diagonal-pincer certificates, extended scales through $t=64$
&
The same predicate $\mathrm{Sep}(H_t,H_t,D(t))$ is additionally checked, one
sibling module per scale, at
$t \in \{19,\allowbreak23,\allowbreak25,\allowbreak27,\allowbreak29,\allowbreak31,
\allowbreak32,\allowbreak37,\allowbreak41,\allowbreak43,\allowbreak47,\allowbreak49,
\allowbreak53,\allowbreak59,\allowbreak61,\allowbreak64\}$
(16 further explicit values). Together with the base 12 scales above this is
\textbf{28 explicit values of $t$ in total}, the complete list being
$\{1,\allowbreak2,\allowbreak3,\allowbreak4,\allowbreak5,\allowbreak7,
\allowbreak8,\allowbreak9,\allowbreak11,\allowbreak13,\allowbreak16,\allowbreak17,
\allowbreak19,\allowbreak23,\allowbreak25,\allowbreak27,\allowbreak29,\allowbreak31,
\allowbreak32,\allowbreak37,\allowbreak41,\allowbreak43,\allowbreak47,\allowbreak49,
\allowbreak53,\allowbreak59,\allowbreak61,\allowbreak64\}$.
This does \emph{not} establish $\mathrm{Sep}(H_t,H_t,\cdot)$ for infinitely
many $t$; these 28 deposits are a historical strict subset of the contiguous
band in the next row.
&
\scale{fixed}
&
Per-scale Lean modules; endpoint
\lean{certifiedKill_diagonal_t64}{DiagonalPincerCertificateT64Endpoint.lean:1928}
\\
Contiguous lcm-diagonal band through $t\le82$
&
For every natural $t\le82$ there exists a depth $L$ with
$\mathrm{Sep}(H_t,H_t,L)$.  This closes every scale in the finite interval,
including plateau transfers, with no holes.  The next lcm jump is at the prime
$83$; no certificate at $t=83$ is claimed, and any such certificate must have
depth at least $125$.
&
\scale{bounded}
&
\lean{ErdosProblems.Skip.LadderT67.exists_diagonalKill_le_82}{ErdosProblems/Skip/LadderT67.lean:71264}
\lean{ErdosProblems.Skip.LadderT67.t83_depth_floor}{ErdosProblems/Skip/LadderT67.lean:71294}
\\
Farey denominator floor
&
Every $q\in\mathbb{N}$ with $0<q\le Q_0=79\,639\,646\,646\,701\,375\,323\,355\,774\,875\,831\,053$
satisfies the window-$(N,K)=(1,240)$ gap certificate; $q=Q_0+1$ is the exact
first failure.
&
\scale{bounded}
&
\lean{gap_check_window_1_240_le_79639646646701375323355774875831053}{GapFareyBound.lean:176}
\\
Coprimality-probability / M\"obius-square Farey floor
&
Every $d\in\mathbb{N}$ with $0<d\le Q_1=39\,819\,823\,323\,350\,687\,661\,677\,887\,437\,915\,526$
(exactly $\lfloor Q_0/2\rfloor$) is excluded as a denominator of
$\Pr(\gcd(X,Y)=1)$ and, separately, of $T=\sum_d\mu(d)/(2^d-1)^2$.
&
\scale{bounded}
&
\lean{tsum_visible_coprime_pairs_ne_int_div_of_den_le_39819823323350687661677887437915526}{CertificateKernel.lean:18572}
\lean{tsum_moebius_div_two_pow_sub_one_sq_ne_int_div_of_den_le_39819823323350687661677887437915526}{CertificateKernel.lean:18487}
\\
\end{longtable}
\normalsize

\begin{rem}[What this table is not]
No row above is cofinal in its indexing parameter ($h$, $t$, or the
denominator bound), and no row is claimed to extrapolate. The historical
diagonal-pincer depths are irregular and nonmonotone
($6,5,7,7,9,14,15,14,21,22,23,26,\dots$); no closed-form growth rate for
$D(t)$ is proved or conjectured in the corpus. This table records the
historical 28-scale bank through $t=64$; the later aggregate theorem for every
$t\le82$ is stated above and is not itemised row by row here.  Neither bounded
record supplies the cofinal obligation developed in Part~1.
\end{rem}
% ---- part a249_p1a ----
\section{Catalogue of mathematical ingredients}

This catalogue renders every producer and converter/identity result found for Erd\H{o}s
\#249 ($S = \sum_{n\ge1}\varphi(n)/2^n$) in the certificate-kernel lane
(\texttt{TotientTailPeriodKiller.lean}, \texttt{CertificateKernel.lean}, and the
lcm/diagonal/cone reduction family) and the M\"obius--Mersenne / squared-Lambert lane
(\texttt{GcdMomentCalculus.lean}, \texttt{RepunitMobiusNumerator.lean} through
\texttt{PrimePowerJumpDynamics.lean}, \texttt{MersenneLambertLadder.lean}). Each entry
is a self-contained premise: exact hypotheses, exact conclusion, the coordinate it was
proved in, and the Lean site. \textbf{\#249 is OPEN}; nothing below decides it. Entries
labelled \texttt{consumer}-shaped (an implication whose hypothesis is an unsupplied
cofinal predicate) are marked explicitly as conditional — the hypothesis itself carries
\ev{Open} and is never claimed as proved. Lean site paths are relative to the
public \texttt{Erdos249257/} source tree.

\subsection{Open hypotheses and candidate inputs}

Producers conclude an existence or a supply: a witnessed object, a witnessed finite
family, or an unconditional dimension/growth lower bound obtained by exhibiting
witnesses (CRT, Dirichlet primes in AP, Bertrand's postulate, explicit \texttt{decide}
computation). Sorted by scale: uniform, then bounded, then fixed; no producer in this
lane is cofinal (the cofinal-scale reduction theorems are consumer-shaped and are
catalogued as converters below, since each moves a cofinal hypothesis across a
coordinate boundary to the target conclusion).

\begin{thm}[cert:a9 --- \texttt{eventual\_period\_of\_not\_irrational}, THE TAIL-PERIOD LAW]
If $S$ is rational then a period exists: $\neg\mathrm{Irrational}(S) \to \exists h:\mathbb{N},\, 0<h \wedge \exists N_0:\mathbb{N},\, \forall N\ge N_0,\ \mathrm{totientTail}(N+h) - \mathrm{totientTail}(N) \in \mathrm{range}((\uparrow):\mathbb{Z}\to\mathbb{R})$. The witness is explicit: $h = \varphi(\mathrm{oddPart}(r.\mathrm{den}))$, $N_0 = v_2(r.\mathrm{den})$, from Euler's theorem applied to the odd part of the hypothetical denominator.
\par\noindent\ev{Lean} \scale{uniform} \coord{binary-digit} \lean{eventual\_period\_of\_not\_irrational}{TotientTailPeriodKiller.lean:358}
\end{thm}

\begin{thm}[cert:d5 --- \texttt{not\_irrational\_totientSeries\_implies\_unbounded\_carryRank\_unconditional}]
$\neg\mathrm{Irrational}(S) \to \exists v>0,\, \exists u:\mathbb{N}\to\mathbb{Z},\, \mathrm{IsTemperedBinaryOrbit}\,\varphi\,v\,u \wedge \forall e,\ 2^e-1 \le \mathrm{finrank}_{\mathbb{Q}}(\mathrm{span}(\mathrm{range}(\mathrm{canonicalCarryKernelFamily}\,u\,e)))$. Rationality of $S$ forces its associated tempered integral binary-carry orbit to have unboundedly rich dyadic-section rank; this is a second, coordinate-independent necessary condition on rationality, parallel to cert:a9 but in the carry-kernel-rank coordinate rather than the binary-digit-periodicity coordinate. It is not by itself an irrationality proof. In particular, the later countermodels rule out treating generic finite-rank shift-polynomial or compressed-adjoint observations as the missing opposite inequality; an actual-totient-specific upper bound would be a genuinely new theorem, not a surviving consequence of the present rank machinery.
\par\noindent\ev{Lean} \scale{uniform} \coord{other:carry-kernel-rank} \lean{not\_irrational\_totientSeries\_implies\_unbounded\_carryRank\_unconditional}{TotientCarryKernelRigidity.lean:300}
\end{thm}

\begin{thm}[cert:d4 --- \texttt{linearIndependent\_canonicalTotientKernelFamily}]
For every $e:\mathbb{N}$, the canonical dyadic totient-kernel family $\mathrm{canonicalTotientKernelFamily}(e) : \mathrm{TotientCanonicalIndex}(e) \to \mathbb{N}\to\mathbb{Q}$, which has exactly $2^e+1$ channels, is linearly independent over $\mathbb{Q}$. Proved unconditionally by constructing, via CRT and Dirichlet's theorem on primes in arithmetic progression (\texttt{PrimesCongruentOne}/\texttt{PrimesInAP} from Mathlib), an explicit evaluation point at which one channel becomes prime and every other channel picks up a fresh prime $\equiv 1$ modulo a large power of $2$ — a genuine witnessed producer, not merely a dimension count. Consequently $\neg\mathrm{FiniteDimensional}\,\mathbb{Q}\,(\mathrm{span}\,\mathbb{Q}\,(\mathrm{range}\,\mathrm{fullTotientKernelFamily}))$. Self-flagged: this shows the dyadic-kernel side is infinite-rank; it is not itself an irrationality proof.
\par\noindent\ev{Lean} \scale{uniform} \coord{other:dyadic-kernel-rank} \lean{card\_totientCanonicalIndex}{TotientMahlerDefect.lean:61} \lean{linearIndependent\_canonicalTotientKernelFamily}{TotientMahlerDefect.lean:935}
\end{thm}

\begin{prop}[cert:a8 --- \texttt{tail\_diff\_int\_of\_den\_dvd}]
If $S = (r:\mathbb{R})$ for $r:\mathbb{Q}$ and $r.\mathrm{den} \mid 2^N\cdot(2^h-1)$, then $\mathrm{totientTail}(N+h) - \mathrm{totientTail}(N) \in \mathrm{range}((\uparrow):\mathbb{Z}\to\mathbb{R})$: a rational value of $S$ with a denominator of this shape forces the shifted tail difference to be an integer. The sharp contrapositive engine underlying cert:a9 and cert:a6.
\par\noindent\ev{Lean} \scale{uniform} \coord{binary-digit} \lean{tail\_diff\_int\_of\_den\_dvd}{TotientTailPeriodKiller.lean:327}
\end{prop}

\begin{prop}[mob:b5 --- top fibre survives, T3, \texttt{mobiusNumerator\_gcd\_cyclotomicValue}]
For $r$ squarefree: $\gcd(|\mathrm{mobiusNumerator}(r)|,\ \mathrm{cyclotomicValue}(r)) = 1$, hence $\mathrm{cyclotomicValue}(r) \mid \mathrm{baseMobiusShadow}(r).\mathrm{den}$. The top cyclotomic channel $|\Phi_r(2)|$ is produced as a survivor: it can never cancel from the numerator and is exhibited as an unconditional divisor of the reduced denominator at every squarefree $r$.
\par\noindent\ev{Lean} \scale{uniform} \coord{cyclotomic} \lean{mobiusNumerator\_gcd\_cyclotomicValue}{CyclotomicProjectionOfShadow.lean:350} \lean{cyclotomicValue\_dvd\_baseMobiusShadow\_den}{CyclotomicProjectionOfShadow.lean:389}
\end{prop}

\begin{prop}[mob:b6 --- upper-half prime channel survival, T4]
For a scale coprime to $C = \prod_{p \in \mathrm{upperHalfPrimes}(t)} \mathrm{mersenne}(p)$ (or with prime support $\le t$), the whole finite product $C$ over the explicit upper-half prime set $\mathrm{upperHalfPrimes}(t) = \{p \text{ prime} : t/2 < p \le t\}$ divides the reduced denominator of the scaled M\"obius shadow at LCM height $t$, unscaled. A genuine finite family of surviving Mersenne-prime channels is exhibited explicitly at every $t$, not merely bounded below.
\par\noindent\ev{Lean} \scale{uniform} \coord{cyclotomic} \lean{upperHalfPrimes}{MersenneShadowCyclotomicNoncollapse.lean:914} \lean{upperHalfChannel\_product\_dvd\_den\_of\_scale\_primeFactors\_le}{MersenneShadowCyclotomicNoncollapse.lean:809} \lean{lcmHeight\_upperHalf\_product\_dvd\_den}{MersenneShadowCyclotomicNoncollapse.lean:983}
\end{prop}

\begin{prop}[mob:b7a --- denominator growth lower bound]
For $t \ge 5$ (Bertrand's postulate supplies nonemptiness of $\mathrm{upperHalfPrimes}(t)$): $2^{t/2} \le \prod_{p \in \mathrm{upperHalfPrimes}(t)} \mathrm{mersenne}(p) \le \big(\mathrm{lcmHeight}(t)\cdot\mathrm{numericMobiusShadow}(\mathrm{lcmHeight}(t))\big).\mathrm{den}$. An exponential-in-$t/2$ growth lower bound for the reduced denominator at every LCM height, produced from mob:b6's explicit surviving channel product. Denominator-only: does not by itself rule out cancellation by a foreign-defect term, hence does not by itself prove \#249.
\par\noindent\ev{Lean} \scale{uniform} \coord{mobius-mersenne} \lean{upperHalfMersenneProduct\_lower\_bound}{MersenneShadowDenominatorGrowth.lean:60}
\end{prop}

\begin{prop}[mob:b7b --- exact denominator value]
For every $t$ (no lower bound on $t$ needed for this direction): $\big(\mathrm{lcmHeight}(t)\cdot\mathrm{numericMobiusShadow}(\mathrm{lcmHeight}(t))\big).\mathrm{den} = \mathrm{mersenne}(\mathrm{lcmRadical}(t)) / \gcd\big(\mathrm{mersenne}(\mathrm{lcmRadical}(t)),\ \mathrm{lcmScale}(t)\cdot|\mathrm{oddJordanScalar}(\mathrm{lcmRadical}(t))|\big)$. A fully closed form for the reduced denominator at every scale, produced (not merely bounded) as an explicit rational function of $t$.
\par\noindent\ev{Lean} \scale{uniform} \coord{mobius-mersenne} \lean{lcmHeight\_scaledMobiusShadow\_den\_exact}{MersenneShadowDenominatorGrowth.lean:147}
\end{prop}

\begin{prop}[mob:d3 --- signed dyadic sum nonvanishing, \texttt{scaled\_dyadic\_sum\_ne\_zero}]
If a finite signed sum $\sum_{i \in s} u(i)\cdot 2^{e(m)-e(i)}$ (clearing dyadic denominators to a common exponent $e(m)$) has one index $m \in s$ with strictly maximal exponent $e(m)$ and odd coefficient $u(m)$, while every other index has strictly smaller exponent, then the cleared sum is $\equiv 1 \pmod 2$, hence nonzero. A positive finite-nonvanishing engine, produced via a rectangular Cauchy--Binet determinant expansion (\texttt{det\_mul\_rectangular}) applied to the signed-Hankel/$q$-moment construction (\texttt{hankelDet}, \texttt{truncatedMoment}); the underlying arithmetic is a $q=1/2$ Hankel-determinant framework for a possible Pad\'e route to \#249. No supply theorem is proved that such a unique-terminal configuration exists cofinally — this is a per-instance producer, not a cofinal producer.
\par\noindent\ev{Lean} \scale{bounded} \coord{p-adic} \lean{scaled\_dyadic\_sum\_ne\_zero}{SignedQMomentObstruction.lean:96} \lean{det\_mul\_rectangular}{SignedQMomentObstruction.lean:29}
\end{prop}

\begin{prop}[cert:a11 --- \texttt{certifiedKill\_all\_small}]
$\forall h \in [1,8],\ \mathrm{certifiedKill}\ h\ 12\ 16$ — a finite, fully unconditional \texttt{decide}-checked family of 8 certificate witnesses at the fixed point $N=12$, $L=16$ (256 explicit totient values below 37 evaluated). Consequence, via cert:a6: $\forall r:\mathbb{Q},\ 1\le h\le 8,\ r.\mathrm{den}\mid 2^{12}\cdot(2^h-1) \to S \ne r$.
\par\noindent\ev{Cert} \scale{fixed} \coord{other:binary-window} \lean{certifiedKill\_all\_small}{TotientTailPeriodKiller.lean:404} \lean{totient\_series\_ne\_rat\_of\_den\_dvd}{TotientTailPeriodKiller.lean:416}
\end{prop}

\begin{prop}[cert:a12 --- upto-sixteen deposit]
$\forall h \in [1,16]$, $\mathrm{certifiedKill}\ h\ 14\ 9$: a wider,
fully unconditional \texttt{decide}-checked family of certificate witnesses.
The basepoint $N=14$ yields the denominator factor $2^{14}$; the certificate
depth is $L=9$.  Consequence:
$\forall r:\mathbb{Q},\ 1\le h\le 16,\ r.\mathrm{den}\mid
2^{14}\cdot(2^h-1) \to S \ne r$.
\par\noindent\ev{Cert} \scale{fixed} \coord{other:binary-window} \lean{totient\_series\_ne\_rat\_of\_den\_dvd\_pow\_two\_mul\_mersenne\_upto\_sixteen}{CertificateKernel.lean:18890}
\end{prop}

\begin{prop}[cert:b11 --- diagonal pincer finite deposits, historical bank and current band]
Let $H_t=\operatorname{lcm}(1,\ldots,t)$ and let
\[
P(t)\quad:\Longleftrightarrow\quad \exists L,\ \mathrm{certifiedKill}(H_t,H_t,L).
\]
The checked certificate bank proves $P(t)$ at 28 explicit indices, ending at $t=64$.
For the initial indices
\[
t=1,2,3,4,5,7,8,9,11,13,16,17,\ldots,
\]
the corresponding depths begin
$6,5,7,7,9,14,15,14,21,22,23,26,\ldots$.
Each entry is a finite kernel computation on explicit totient values. The factorisations use
checked prime-power blocks and Lucas primality certificates. Thus the bank proves a finite list
of instances of $P(t)$; it does not prove that $P(t)$ holds for infinitely many $t$.
\par\noindent\ev{Cert} \scale{fixed} \coord{other:lcm-diagonal} \lean{certifiedKill\_diagonal\_all\_imported}{DiagonalPincerCertificates.lean:2870} \lean{diagonalPincerKillDepth}{DiagonalPincerCertificates.lean:2854} \lean{certifiedKill\_diagonal\_t64}{DiagonalPincerCertificateT64Endpoint.lean:1928}

The historical 28 deposits are now a strict subset of the aggregate theorem
$\forall t\le82,\ P(t)$.  That theorem closes the finite interval without
holes but supplies neither $P(83)$ nor a cofinal family.
\par\noindent\ev{Cert} \scale{bounded} \coord{other:lcm-diagonal}
\lean{ErdosProblems.Skip.LadderT67.exists_diagonalKill_le_82}{ErdosProblems/Skip/LadderT67.lean:71264}
\end{prop}

\subsection{Exact identities and reductions}

Converters and identities move a statement between coordinates without changing its
truth value: definitions, exact algebraic identities, iff-characterisations, and the
implication theorems that convert a supply hypothesis in one coordinate into the target
conclusion $\mathrm{Irrational}(S)$. This subsection carries the certificate-kernel core
(the \texttt{Sep}(h,N,L) predicate, its completeness iff, and every known reformulation
of the open supply obligation) and the M\"obius--Mersenne / squared-Lambert identity
family, including the Mersenne--Lambert five-row status ladder. Sorted by scale:
foundational (n/a) definitions and identities first, then uniform, then cofinal, then
bounded, then fixed.

\subsubsection{Foundational definitions and identities (scale n/a)}

\begin{defn}[cert:a1 --- \texttt{totientTail}]
$\mathrm{totientTail}(N) := \sum_{j\ge0}' \varphi(N+1+j)/2^{j+1}$, the fractional layer of $2^N\cdot S$; well-defined for every $N$.
\par\noindent\ev{Lean} \scale{n/a} \coord{binary-digit} \lean{totientTail}{TotientTailPeriodKiller.lean:57}
\end{defn}

\begin{defn}[cert:a3 --- \texttt{windowDiscrepancy}]
$\mathrm{windowDiscrepancy}(h,N,L) := \sum_{j<L} \big(\varphi(N+h+1+j) - \varphi(N+1+j)\big)\cdot 2^{L-1-j} \in \mathbb{Z}$, the depth-$L$ truncation of $2^L\cdot(\mathrm{totientTail}(N+h) - \mathrm{totientTail}(N))$; computable and decidable given $h,N,L$.
\par\noindent\ev{Lean} \scale{n/a} \coord{other:binary-window} \lean{windowDiscrepancy}{TotientTailPeriodKiller.lean:66}
\end{defn}

\begin{defn}[cert:a4 --- \texttt{certifiedKill}, THE Sep(h,N,L) PREDICATE]
$\mathrm{certifiedKill}(h,N,L) := (N+h+L+2:\mathbb{Z}) < \mathrm{windowDiscrepancy}(h,N,L) \bmod 2^L < 2^L - (N+h+L+2)$ — the residue of the window discrepancy modulo $2^L$ avoids the shrinking radius-$(N+h+L+2)$ neighbourhood of $0$. This is exactly the object named $\mathrm{Sep}(h,N,L)$: purely finite arithmetic, decidable, with no analytic hypothesis. This is the certificate kernel's core object; every theorem below is either a hypothesis-shape wrapping it or an unconditional finite instance of it.
\par\noindent\ev{Lean} \scale{n/a} \coord{other:binary-window} \lean{certifiedKill}{TotientTailPeriodKiller.lean:72}
\end{defn}

\begin{prop}[cert:a2 --- shift identity, \texttt{two\_pow\_mul\_totient\_series\_eq}]
$2^N \cdot \Big(\sum_{n\ge0}' \varphi(n)/2^n\Big) = \mathrm{totientPrefix}(N) + \mathrm{totientTail}(N)$, where $\mathrm{totientPrefix}(N) = \sum_{n\le N}\varphi(n)\cdot 2^{N-n} \in \mathbb{N}$ is an exact integer. Rationality/periodicity of $S$'s binary expansion reduces exactly to a statement about $\mathrm{totientTail}$. Template shape (constant $\cdot$ prefix + tail split) reusable for any $\sum f(n)/2^n$.
\par\noindent\ev{Lean} \scale{uniform} \coord{binary-digit} \lean{two\_pow\_mul\_totient\_series\_eq}{TotientTailPeriodKiller.lean:150}
\end{prop}

\begin{lem}[cert:a5 --- certificate depth floor, \texttt{certifiedKill\_depth\_floor}]
$\mathrm{certifiedKill}(h,N,L) \to 2\cdot(N+h+L+2) < 2^L$. A structural corollary of the kernel definition: the certificate depth $L$ cannot stay bounded while $N+h \to \infty$; $L$ must grow at least logarithmically with $N+h$. Necessary context for reading every fixed/bounded-scale kernel instance below.
\par\noindent\ev{Lean} \scale{uniform} \coord{other:binary-window} \lean{certifiedKill\_depth\_floor}{TotientTailPeriodKiller.lean:79}
\end{lem}

\begin{prop}[cert:a6 --- certificate soundness, the kernel's converter direction, \texttt{tail\_diff\_notMem\_int\_of\_certifiedKill}]
$\mathrm{certifiedKill}(h,N,L) \to \mathrm{totientTail}(N+h) - \mathrm{totientTail}(N) \notin \mathrm{range}((\uparrow):\mathbb{Z}\to\mathbb{R})$. This is the kernel's soundness converter: it moves a finite, decidable, binary-window fact into a real-analytic non-integrality fact. Composes with cert:a9 (tail-period law) by contradiction to kill a hypothetical rational's period.
\par\noindent\ev{Lean} \scale{uniform} \coord{other:binary-window} \lean{tail\_diff\_notMem\_int\_of\_certifiedKill}{TotientTailPeriodKiller.lean:262}
\end{prop}

\begin{thm}[cert:a7 --- certificate completeness, \texttt{exists\_certifiedKill\_iff\_tail\_diff\_notMem\_int}]
$(\exists L,\ \mathrm{certifiedKill}(h,N,L)) \iff \mathrm{totientTail}(N+h) - \mathrm{totientTail}(N) \notin \mathrm{range}((\uparrow):\mathbb{Z}\to\mathbb{R})$, for all $h,N$. Certificates are \emph{complete} receipts of non-integrality, not merely sufficient: the certificate vocabulary is dispensable, and the real target of the whole kernel is exactly the right-hand real-analytic statement. Converts the open \#249 obligation freely between "supply of certificates" language and "supply of non-integral tail differences at arbitrarily large scale" language (cert:b8 below uses the latter form). Any independent non-integrality proof from any coordinate automatically yields a certificate by this iff, and vice versa.
\par\noindent\ev{Lean} \scale{uniform} \coord{other:binary-window} \lean{exists\_certifiedKill\_iff\_tail\_diff\_notMem\_int}{LcmConeFlatness.lean:316}
\end{thm}

\begin{defn}[cert:b1 --- \texttt{periodLcm}, the universal period ray]
$\mathrm{periodLcm}(0) = 1$, $\mathrm{periodLcm}(t+1) = \mathrm{lcm}(\mathrm{periodLcm}(t), t+1)$, i.e.\ $\mathrm{periodLcm}(t) = \mathrm{lcm}(1,\dots,t)$; $t \le \mathrm{periodLcm}(t)$; every primitive period $h_0 \le t$ divides $\mathrm{periodLcm}(t)$. Pure number theory about $\mathrm{lcm}(1..t)$, zero totient content; directly reusable for \#257 or any period-search problem — "stand on the universal-period ray to remove one free parameter" is a general reduction technique.
\par\noindent\ev{Lean} \scale{n/a} \coord{other:lcm-period-ray} \lean{periodLcm}{CarrySurvivorExtinction.lean:515} \lean{le\_periodLcm}{LcmDiagonalReduction.lean:81}
\end{defn}

\begin{prop}[cert:c1 --- \texttt{farey\_gap}, the mediant lemma]
For $a,b,c,d,r,s:\mathbb{Z}$ with $b>0$, $d>0$, unimodular neighbours $bc-ad=1$, and $r$ strictly between $a/b$ and $c/d$ ($a\cdot s < r\cdot b$, $r\cdot d < c\cdot s$): $b+d \le s$. Any rational strictly between two unimodular Farey neighbours has denominator at least the sum of the neighbours' denominators. Fully problem-agnostic — pure Stern--Brocot/Farey fact, zero totient content, directly reusable for \#257.
\par\noindent\ev{Lean} \scale{n/a} \coord{farey} \lean{farey\_gap}{GapFareyBound.lean:51}
\end{prop}

\begin{prop}[cert:d1 --- \texttt{irrational\_of\_den\_mul\_abs\_sub\_tendsto\_zero}, generic Dirichlet-gap criterion]
If $u:\mathbb{N}\to\mathbb{Q}$ is eventually never equal to $x:\mathbb{R}$, and $\mathrm{den}(u(k))\cdot|x - u(k)| \to 0$, then $\mathrm{Irrational}(x)$. Classical Dirichlet-approximation irrationality criterion, zero totient/Mersenne content, already reused in this corpus to prove full-support Erd\H{o}s--Borwein irrationality (\#257-shaped). Apply to any explicit sequence of convergents to $S$ with a provable denominator$\cdot$gap $\to 0$ bound.
\par\noindent\ev{Lean} \scale{n/a} \coord{n/a} \lean{irrational\_of\_den\_mul\_abs\_sub\_tendsto\_zero}{CertificateKernel.lean:5371}
\end{prop}

\begin{prop}[cert:d2 --- \texttt{irrational\_of\_int\_mul\_near\_int}, classical near-integer criterion]
If $\forall q>0,\ \exists m,z:\mathbb{Z},\ 0 < |m\cdot\xi - z| < 1/q$, then $\mathrm{Irrational}(\xi)$. The classical Erd\H{o}s-1948-shape criterion behind digit/carry irrationality proofs. Base-power specialisation $\mathrm{irrational\_of\_pow\_mul\_near\_int}$ (witnesses of form $b^n\cdot\xi$) directly matches a digit/carry construction in any base $b$, hence directly usable for \#257's $b^n-1$ denominators. Together with cert:d1 these are the only two general-purpose irrationality criteria in the whole kernel; everything else exists to supply or substitute for their hypotheses.
\par\noindent\ev{Lean} \scale{n/a} \coord{n/a} \lean{irrational\_of\_int\_mul\_near\_int}{CertificateKernel.lean:6120} \lean{irrational\_of\_pow\_mul\_near\_int}{CertificateKernel.lean:6149}
\end{prop}

\begin{prop}[cert:d3 --- \texttt{SeparatedMinorCertificate}, linear independence from a nonzero minor]
For an indexed family $\mathrm{family}: \iota\to\mathbb{N}\to\mathbb{Q}$: a $\mathrm{SeparatedMinorCertificate}$ (an explicit finite evaluation-point assignment $\mathrm{rowIndex}:\iota\to\mathbb{N}$ with $\det(\mathrm{family}(j)(\mathrm{rowIndex}(i)))_{i,j} \ne 0$) implies $\mathrm{LinearIndependent}\ \mathbb{Q}\ \mathrm{family}$. Fully generic finite-dimensional linear algebra; $\mathrm{family}$ is a free variable, nothing here mentions totient or Mersenne structure. Reusable for any finite-rank/linear-independence obstruction in either open problem.
\par\noindent\ev{Lean} \scale{n/a} \coord{n/a} \lean{SeparatedMinorCertificate}{TotientMahlerDefect.lean:83} \lean{linearIndependent\_of\_separatedMinorCertificate}{TotientMahlerDefect.lean:91}
\end{prop}

\begin{prop}[cert:d9 --- \texttt{positive\_rational\_difference\_lower\_bound}]
For $\mathrm{pfx} < \mathrm{whole}:\mathbb{Q}$ (strict): $1/(\mathrm{whole.den}\cdot\mathrm{pfx.den}) \le (\mathrm{whole}:\mathbb{R}) - (\mathrm{pfx}:\mathbb{R})$. A positive rational difference is bounded below by the reciprocal of the product of the two \emph{actual} reduced denominators, not a displayed/guessed one. Fully general; feeds \texttt{prefixDenominator\_shell\_power\_bound\_of\_rational\_difference} and \texttt{nextSupport\_power\_bound\_of\_rational\_difference}, converting any analytic upper bound on a rational gap into a denominator-growth lower bound — plug in any tail estimate from either open problem's coordinate.
\par\noindent\ev{Lean} \scale{n/a} \coord{n/a} \lean{positive\_rational\_difference\_lower\_bound}{PrimitiveRationalGapSupply.lean:31}
\end{prop}

\begin{prop}[mob:a1a --- the M\"obius--squared-Mersenne identity]
$S = \sum_{n\ge1} \varphi(n)/2^n = \sum_{d\ge1}' \mu(d)\cdot 2^d/(2^d-1)^2 = \tfrac12 + \sum_{d\ge1}' \mu(d)/(2^d-1)^2$. Unconditional identity of convergent real series; this is the reduced target every M\"obius--Mersenne-lane result in this catalogue feeds.
\par\noindent\ev{Lean} \scale{n/a} \coord{mobius-mersenne} \lean{tsum\_totient\_half\_pow\_eq\_half\_add\_moebius\_sq}{MersenneLambertLadder.lean:665} \lean{totientSeries\_eq\_pnat\_half\_pow}{SquaredMersenneDiagonalEnclosure.lean:72} \lean{mobius\_square\_series\_eq\_partial\_add\_tail}{SquaredMersenneDiagonalEnclosure.lean:94}
\end{prop}

\begin{cor}[mob:a1b --- \#249 restated in the squared-Lambert coordinate]
Immediate from mob:a1a: $\mathrm{Irrational}(S) \iff \sum_{d\ge1}' \mu(d)/(2^d-1)^2 \notin \mathbb{Q}$. Denote the right-hand series $L_2(\mu) := S - 1/2$ (mob:a3 below). Every subsequent M\"obius-coordinate obstruction or identity in this catalogue is measured against this exact restatement, not the original totient series.
\par\noindent\ev{Lean} \scale{n/a} \coord{mobius-mersenne} \lean{tsum\_totient\_half\_pow\_eq\_half\_add\_moebius\_sq}{MersenneLambertLadder.lean:665}
\end{cor}

\begin{prop}[mob:a2 --- squared-Lambert transfer engine, \texttt{tsum\_lambert\_linear\_weight\_sq\_pure}]
For $w:\mathbb{N}\to\mathbb{R}$ with $|w(d)| \le d$ for $d>0$, and $0 \le r < 1$: $\sum_{d:\mathbb{N}^+}' w(d)\cdot(r^d/(1-r^d))^2 = \sum_{n:\mathbb{N}^+}' \Big(\sum_{e\mid n} w(e)\cdot(n/e-1)\Big)\cdot r^n$. Converts any linear-growth-bounded squared-Lambert series into a divisor-convolution power series — the single reusable brick for the whole level-2 M\"obius--Lambert ladder below. Pure Dirichlet-convolution/Lambert-series algebra, no Mersenne-specific structure; instantiate at $w = \mu, 1, \varphi$ to obtain mob:a3--a5, and directly reusable for a weighted or squared variant of \#257's series.
\par\noindent\ev{Lean} \scale{uniform} \coord{mobius-mersenne} \lean{tsum\_lambert\_linear\_weight\_sq\_pure}{GcdMomentCalculus.lean:105}
\end{prop}

\begin{defn}[mob:a3 --- $L_2(\mu)$ is exactly \#249]
$L_2(\mu) := \sum_{d:\mathbb{N}^+}' \mu(d)/(2^d-1)^2 = S - 1/2$. The open \#249 atom, restated as the M\"obius rung of the level-2 (squared) Lambert ladder; consequence of mob:a1a and mob:a2 at $w=\mu$.
\par\noindent\ev{Cited} \scale{n/a} \coord{mobius-mersenne} \lean{tsum\_totient\_half\_pow\_eq\_half\_add\_moebius\_sq}{MersenneLambertLadder.lean:665}
\end{defn}

\begin{prop}[mob:a4 --- $L_2(1)$, the known $q$-zeta anchor rung, \texttt{tsum\_one\_div\_mersenne\_sq\_eq\_sigma\_sub\_tau\_series}]
$\sum_{d:\mathbb{N}^+}' 1/(2^d-1)^2 = \sum_{n:\mathbb{N}^+}' \big(\sigma(n)-\tau(n)\big)\cdot(1/2)^n = \zeta_q(2) - \zeta_q(1)$ at $q=1/2$. The level-2 $\zeta$-rung, exactly evaluated in Lean. Irrationality of this \emph{value} is cited (Postelmans--Van Assche $q$-Pad\'e), \emph{not} formalised — only the identity itself is machine-checked. This is the "known" sibling rung mirroring mob:a3 in the level-mirror table mob:a6.
\par\noindent\ev{Lean} \scale{n/a} \coord{mobius-mersenne} \lean{tsum\_one\_div\_mersenne\_sq\_eq\_sigma\_sub\_tau\_series}{GcdMomentCalculus.lean:216}
\end{prop}

\begin{prop}[mob:a5 --- $L_2(\varphi)$, the Pillai gcd-moment rung, \texttt{tsum\_totient\_div\_mersenne\_sq\_eq\_gcd\_moment\_series}]
$\sum_{d:\mathbb{N}^+}' \varphi(d)/(2^d-1)^2 = \sum_{n:\mathbb{N}^+}' (P(n)-n)\cdot(1/2)^n$, where $P = \varphi * \mathrm{Id}$ (Pillai's gcd-sum function). Equals $\mathbb{E}[\gcd(X,Y)]$ for independent fair-coin waiting times $X,Y$ (probabilistic coordinate, mob:a7--a9). A cousin rung, not \#249 itself; status open.
\par\noindent\ev{Lean} \scale{n/a} \coord{mobius-mersenne} \lean{tsum\_totient\_div\_mersenne\_sq\_eq\_gcd\_moment\_series}{GcdMomentCalculus.lean:235}
\end{prop}

\begin{obs}[mob:a6 --- the level-mirror table]
Level 1 ($L(f) := \sum f(d)/(2^d-1)$): $L(\mu) = 1/2$ (rational, trivial); $L(1) = \mathcal{E}$, the Erd\H{o}s--Borwein constant (irrational, Erd\H{o}s 1948, machine-checked in this kernel, cert:d7c below); $L(\varphi) = 2$ (rational). Level 2 ($L_2(f) := \sum f(d)/(2^d-1)^2$): $L_2(\mu) = S - 1/2$ (mob:a3, OPEN, $=$ \#249); $L_2(1) = \zeta_q(2)-\zeta_q(1)$ (mob:a4, irrational, cited); $L_2(\varphi) = \mathbb{E}[\gcd]$ (mob:a5, open, Pillai). At level 1 the M\"obius rung is trivial and the $\zeta$-rung is hard; at level 2 the $\zeta$-rung is known and the M\"obius rung \emph{is} \#249 — M\"obius projection is the single wall at both levels. This mirror-structure phenomenon is a general observation about Dirichlet-convolution ladders, potentially informative for \#257's Mersenne-shifted sums too.
\par\noindent\ev{Lean} \scale{n/a} \coord{mobius-mersenne}
\end{obs}

\begin{prop}[mob:a7 --- gcd-divisibility factorises, \texttt{tsum\_pos\_pair\_both\_dvd\_half\_eq\_inv\_mersenne\_sq}]
For independent fair-coin waiting times $X,Y$ ($P(X=n)=2^{-n}$) and $d>0$: $P(d\mid X \wedge d\mid Y) = 1/(2^d-1)^2$. The foundation stone for reading $L_2(f)$ as $\mathbb{E}[(f*\zeta)(\gcd(X,Y))]$. Pure probability/geometric-series fact about independent geometric random variables, zero \#249-specific content; directly reusable for a two-coordinate gcd structure in \#257.
\par\noindent\ev{Lean} \scale{uniform} \coord{probability} \lean{tsum\_pos\_pair\_both\_dvd\_half\_eq\_inv\_mersenne\_sq}{GcdMomentCalculus.lean:266}
\end{prop}

\begin{prop}[mob:a8 --- reduced-direction law, \texttt{tsum\_pos\_coprime\_inv\_mersenne\_eq\_one}]
$\sum_{(a,b):\, a,b\ge1,\, \gcd(a,b)=1}' 1/(2^{a+b}-1) = 1$: the sum over exact bipartite coprime pairs. Every positive coprime pair carries a slope mass $1/(2^{a+b}-1)$, and these mass exactly one — the root cylinder $M(1,1)=1$ of the Stern--Brocot tree, mob:a9. Pure coprimality/geometric-series law, applies verbatim to any base-$b$ analogue.
\par\noindent\ev{Lean} \scale{n/a} \coord{probability} \lean{tsum\_pos\_coprime\_inv\_mersenne\_eq\_one}{GcdMomentCalculus.lean:349}
\end{prop}

\begin{prop}[mob:a9a --- Stern--Brocot cylinder recursion, \texttt{cylinderMass\_split}]
$M(a,b) := 1/\big((2^a-1)(2^b-1)\big)$ satisfies the exact telescoping identity $M(a,b) = 1/(2^{a+b}-1) + M(a+b,b) + M(a,a+b)$ for $a,b:\mathbb{N}^+$: an exact Markov measure on the Stern--Brocot mediant tree, whose closed form is a product of two independent-coordinate divisor probabilities $P(a\mid X)\cdot P(b\mid Y)$ (mob:a7).
\par\noindent\ev{Lean} \scale{uniform} \coord{other:stern-brocot} \lean{cylinderMass\_split}{GcdMomentCalculus.lean:474}
\end{prop}

\begin{prop}[mob:a9b --- depth-$d$ convergence rate, \texttt{sternBrocotDepthMass\_error} / \texttt{tendsto\_sternBrocotDepthMass}]
Children of a Stern--Brocot cylinder carry at most $2/3$ of the parent's mass; the depth-$d$ finite unfolding of mob:a9a converges to $M(a,b)$ at explicit rate $|M(a,b) - M_d(a,b)| \le (2/3)^d\cdot M(a,b)$. The "recursion limit $=$ subtree mass" identification is explicitly \emph{not} claimed beyond this rate bound (flagged as a wave-21 gap in the source). All-left cusp cylinders $M(N,1) = 1/(2^N-1)$ are near-Mersenne-reciprocal, one exponential order closer to the Erd\H{o}s--Borwein engine than raw $\varphi(n)$ coefficients — an advisory bridge candidate to Erd\H{o}s-1948-style arguments, not a proved route.
\par\noindent\ev{Lean} \scale{uniform} \coord{other:stern-brocot} \lean{sternBrocotDepthMass\_error}{GcdMomentCalculus.lean:525} \lean{tendsto\_sternBrocotDepthMass}{GcdMomentCalculus.lean:560}
\end{prop}

\begin{prop}[mob:b1 --- repunit gcd word, T1, \texttt{mobiusNumeratorPolynomial\_eq\_gcdWord}]
For $r$ squarefree, the divisor-signed polynomial $\mathrm{mobiusNumeratorPolynomial}(r) := \sum_{d\mid r} \mu(d)\cdot(r/d)\cdot\mathrm{spacedRepunit}(d,r/d)$ equals $\mathrm{gcdWord}(r)$, whose $X^k$ coefficient ($k<r$) is $\mathrm{gcdWordCoeff}(r,k) = (r/\gcd(r,k))\cdot\varphi(\gcd(r,k))$, strictly positive for $k<r$ and zero for $k\ge r$. The signed repunit numerator is exactly a positive gcd word; the classical polynomial identity underlying it (signed spaced repunit $=$ gcd word) is base-independent, though the application here is Mersenne-specific.
\par\noindent\ev{Lean} \scale{n/a} \coord{cyclotomic} \lean{mobiusNumeratorPolynomial\_eq\_gcdWord}{RepunitMobiusNumerator.lean:205} \lean{mobiusNumeratorPolynomial\_coeff}{RepunitMobiusNumerator.lean:217} \lean{mobiusNumeratorPolynomial\_coeff\_pos}{RepunitMobiusNumerator.lean:234}
\end{prop}

\begin{prop}[mob:b2 --- evaluation at $2$ recovers the integer numerator, \texttt{mobiusNumeratorPolynomial\_eval\_two}]
$\mathrm{mobiusNumeratorPolynomial}(r).\mathrm{eval}\,2 = \mathrm{mobiusNumerator}(r)$ for $r$ squarefree, where $\mathrm{mobiusNumerator}(r) := \sum_{s \subseteq \mathrm{primeFactors}(r)} (-1)^{|s|}\cdot(r/d)\cdot(\mathrm{mersenne}(r)/\mathrm{mersenne}(d))$, $d = \prod s$. Bridges the polynomial (mob:b1) world to the integer arithmetic used by every denominator-survival theorem below.
\par\noindent\ev{Lean} \scale{uniform} \coord{cyclotomic} \lean{mobiusNumeratorPolynomial\_eval\_two}{RepunitMobiusNumerator.lean:445} \lean{mobiusNumerator}{RadicalMobiusShadow.lean:101}
\end{prop}

\begin{prop}[mob:b3 --- radical shadow scale decomposition]
$\mathrm{baseMobiusShadow}(r) := \mathrm{mobiusNumerator}(r)/(2^r-1)$ (unscaled); $\mathrm{numericMobiusShadow}(H) := \mathrm{baseMobiusShadow}(\mathrm{rad}(H))/\mathrm{rad}(H)$; exactly $H\cdot\mathrm{numericMobiusShadow}(H) = (H/\mathrm{rad}(H))\cdot\mathrm{baseMobiusShadow}(\mathrm{rad}(H))$ for $H>0$. Reduced-denominator identity for the unscaled shadow: $\mathrm{baseMobiusShadow}(r).\mathrm{den} = \mathrm{mersenne}(r)/\gcd(|\mathrm{mobiusNumerator}(r)|, \mathrm{mersenne}(r))$ for $r>0$ — exact and generic, no coprimality assumed, no channel-survival hidden.
\par\noindent\ev{Lean} \scale{uniform} \coord{mobius-mersenne} \lean{scaledMobiusShadow\_eq\_radicalBase}{RadicalMobiusShadow.lean:125} \lean{baseMobiusShadow\_den}{RadicalMobiusShadow.lean:150}
\end{prop}

\begin{prop}[mob:b4 --- cyclotomic congruence, T2, \texttt{cyclotomic\_dvd\_mobiusNumeratorPolynomial\_sub}]
For $r$ squarefree, $m \mid r$: $\Phi_m \mid \big(\mathrm{mobiusNumeratorPolynomial}(r) - C(\mu(m)\cdot J_2(r/m))\big)$ in $\mathbb{Z}[X]$, where $J_2 = \mu * \mathrm{id}^2$ is the Jordan totient. Evaluated: $\mathrm{cyclotomicEval}(m) \mid \mathrm{mobiusNumerator}(r) - \mu(m)\cdot J_2(r/m)$. Every cyclotomic fibre $\Phi_m(2)$ of the numerator is congruent to an explicit constant depending only on $\mu(m)$ and the Jordan totient of the cofactor. Setting $m=r$ gives mob:b5 (top fibre survives).
\par\noindent\ev{Lean} \scale{uniform} \coord{cyclotomic} \lean{cyclotomic\_dvd\_mobiusNumeratorPolynomial\_sub}{CyclotomicProjectionOfShadow.lean:299} \lean{mobiusNumerator\_mod\_cyclotomicEval}{CyclotomicProjectionOfShadow.lean:310}
\end{prop}

\begin{prop}[mob:b8a --- prime-power jump recurrence, T5, new fibre]
Adjoining a new prime $p \nmid r$: on a genuinely new fibre $m\cdot p$ (with $m \mid r$), $\Phi_{mp} \mid \mathrm{mobiusNumeratorPolynomial}(rp) + \mathrm{expand}_p(\mathrm{mobiusNumeratorPolynomial}(r))$ — the new fibre picks up a sign-flipped $p$-expansion of the old numerator. Holds for $p$ prime, $p \nmid r$, $m \mid r$; the underlying recurrence needs no squarefreeness of $r$.
\par\noindent\ev{Lean} \scale{uniform} \coord{cyclotomic} \lean{cyclotomic\_dvd\_primeJump\_new\_fibre}{PrimePowerJumpDynamics.lean:281} \lean{cyclotomic\_dvd\_primeJump\_new\_fibre\_constant}{PrimePowerJumpDynamics.lean:347}
\end{prop}

\begin{prop}[mob:b8b --- prime-power jump recurrence, T5, old fibre]
Adjoining a new prime $p \nmid r$: on an old fibre $m$ (with $m \mid r$), $\Phi_m \mid \mathrm{mobiusNumeratorPolynomial}(rp) - C(p^2-1)\cdot\mathrm{mobiusNumeratorPolynomial}(r)$ — old fibres scale by exactly $p^2-1$. Requires $r$ squarefree (needed for the T2 constant-fibre step mob:b4, not for the underlying recurrence itself). An inductive tool for building explicit denominator/numerator values one prime at a time; a candidate base case for an induction proving mob:b7a/b7b's growth bound tight.
\par\noindent\ev{Lean} \scale{uniform} \coord{cyclotomic} \lean{cyclotomic\_dvd\_primeJump\_old\_fibre}{PrimePowerJumpDynamics.lean:310}
\end{prop}

\begin{prop}[cert:d7 --- the Mersenne--Lambert ladder identities]
Define $L(f) := \sum_{n\ge1} f(n)/(2^n-1)$ (level 1, no square). Writing $A := \varphi*\mu$ (the primitive Euler weight, $\alpha$ in prose): the Dirichlet-convolution identities $A*\zeta = \varphi$ and $\varphi*\zeta = \mathrm{Id}$ separate five exact values of $L$ along one ladder, each obtained by factoring $L(f)$ as a divisor-transform $(f*1)$ followed by binary evaluation, $L(f) = \sum_n (f*1)(m)/2^m$. $S$ sits inside this ladder as $L(A)$: this is \#249 restated in positive Erd\H{o}s--Borwein form.
\par\noindent\ev{Lean} \scale{n/a} \coord{mobius-mersenne} \lean{tsum\_totient\_div\_pow\_two\_eq\_pnat\_half\_pow}{CertificateKernel.lean:18415}
\end{prop}

\begin{longtable}{@{}l l l@{}}
\caption{The Mersenne--Lambert ladder five-row status table (docstring @ \texttt{CertificateKernel.lean:18063--18083}, body \texttt{MersenneLambertLadder.lean}).}\\
\toprule
\papertablehead
Weight $f$ & Value $L(f)$ & Status \\
\midrule
\endfirsthead
$\mu$ & $L(\mu) = 1/2$ & rational, trivial \\
$\varphi$ & $L(\varphi) = 2$ & rational \\
$\mathrm{Id}$ & $L(\mathrm{Id}) = \sum_m \sigma(m)/2^m$ & transcendental (Nesterenko 1996, \ev{Cited}, not formalised) \\
$1$ & $L(1) = \mathcal{E}$, Erd\H{o}s--Borwein constant & irrational, \ev{Lean} (\texttt{irrational\_erdosBorwein\_series}, \\
 & & \texttt{CertificateKernel.lean:8007}) — matches \#257's full-support case \\
$A = \varphi*\mu$ & $L(A) = S$ & \textbf{OPEN} — this is \#249 \\
\bottomrule
\end{longtable}

\begin{rem}[cert:d7 --- reading the ladder]
This is the single clearest bridge placing \#249 inside \#257's native Mersenne-Lambert coordinate ($\sum f(n)/(b^n-1)$, here $b=2$): if the \#257 lane's Mersenne-coordinate machinery (achievement sets, greedy orbits) can say anything about the sign or density structure of $A = \varphi*\mu$ specifically, it transfers here for free via this ladder identity. Conversely, any \#249 result about $L(A)$ phrased purely in Lambert-series terms (not totient terms) is directly \#257-lane-portable. The `evaluated' identity confirming $S$'s $\mathbb{N}$-indexed form equals its $\mathbb{N}^+$/Lambert-indexed form is the one row here checked directly against source (\lean{tsum\_totient\_div\_pow\_two\_eq\_pnat\_half\_pow}{CertificateKernel.lean:18415}); the body machinery in \texttt{MersenneLambertLadder.lean} establishing the other four rows was read only via this docstring and is flagged \ev{Cited}/\ev{Math} accordingly except where a specific declaration is named above.
\end{rem}

\begin{prop}[mob:e3 --- joint-35 cone annihilator, \texttt{oldChannel\_affine\_moment\_annihilation} / \texttt{joint35\_oldChannel\_zero}]
The four-vertex affine annihilator $q(X,Y) = XY-3X-2Y+4$ ($q(1,1)=q(3,5)=0$, $q(9,25)=152$) kills every ``old'' divisor channel exactly at any LCM height: for $d\mid H$, $d>0$, $\mathrm{transportResidueKernel}(d,15H) - 3\cdot\mathrm{transportResidueKernel}(d,3H) - 2\cdot\mathrm{transportResidueKernel}(d,5H) + 4\cdot\mathrm{transportResidueKernel}(d,H) = 0$. General form (\texttt{oldChannel\_affine\_moment\_annihilation}): any finite affine annihilator with $\sum c_i = 0$, $\sum c_i\cdot m_i = 0$ kills every old residue channel, independent of the particular $(3,5)$ choice — fully coordinate-free finite linear algebra. Sharp cone radius $19H+5L+5$ (\texttt{sharpJoint35ConeRadius}). An unbounded certificate supply built from this (not proved) would close \#249 via the same cone-flatness route as mob:e2/cert:b6.
\par\noindent\ev{Lean} \scale{uniform} \coord{other:lcm-diagonal} \lean{oldChannel\_affine\_moment\_annihilation}{JointExponentTransport.lean:36} \lean{joint35\_oldChannel\_zero}{JointExponentTransport.lean:71} \lean{sharpJoint35ConeRadius}{JointExponentTransport.lean:124}
\end{prop}

\begin{prop}[mob:f1 --- composite dilation defect identity, \texttt{supportCoeff\_mul\_eq\_add\_defect}]
For an abstract support set $A\subseteq\mathbb{N}$, $\mathrm{supportCoeff}_A(n) := \#\{a\in A: a\mid n\}$: for $a\in A$, $a>0$, $x>0$, $\mathrm{supportCoeff}_A(ax) = \mathrm{supportCoeff}_A(x) + [a\nmid x] + \mathrm{compositeDilationDefect}_A(a,x)$. On prime-only support (every element of $A$ prime) the defect is identically zero, recovering the classical $\mathrm{supportCoeff}_A(px) = \mathrm{supportCoeff}_A(x) + [p\nmid x]$. Pure divisor combinatorics over an abstract support set, no $2^n$-specific structure — a shared \#249/\#257 substrate identity: instantiate $A=\mathbb{N}$ weighted by $\varphi$ for \#249, or $A$ the chosen infinite subset for \#257's $\sum_{n\in A}1/(2^n-1)$.
\par\noindent\ev{Lean} \scale{uniform} \coord{other:support-divisor-counting} \lean{supportCoeff\_mul\_eq\_add\_defect}{CompositeDilationDefect.lean:30} \lean{compositeDilationDefect\_eq\_zero\_of\_prime\_support}{CompositeDilationDefect.lean:103} \lean{supportCoeff\_mul\_prime\_support}{CompositeDilationDefect.lean:119}
\end{prop}

\subsubsection{Cofinal-scale converters — the certificate-kernel reduction ladder}

Every entry in this block exposes a cofinal supply predicate and a checked
route to $\mathrm{Irrational}(S)$. \textbf{None of the supply predicates is
established.}  Completeness now gives registered iff theorems for the base
certificate supply and the lcm-diagonal supply; the older one-directional
consumer declarations remain useful proof components but are not the full
logical status.

\begin{thm}[cert:a10 --- certificate-supply iff, THE WALL]
$\big(\forall h:\mathbb{N},\ 0<h \to \forall N_0:\mathbb{N},\ \exists N\ge N_0,\ \exists L,\ \mathrm{certifiedKill}(h,N,L)\big) \leftrightarrow \mathrm{Irrational}(S)$. The supply side is exactly $\mathrm{Sep}(h,N,L)$ quantified as $\forall h\ge1\ \forall N_0\ge0\ \exists N\ge N_0\ \exists L$. \textbf{This supply is nowhere proved in the corpus; the iff is an exact reformulation, not progress.}
\par\noindent\ev{Lean} (equivalence proved; supply \ev{Open}) \scale{cofinal} \coord{other:binary-window} \lean{irrational\_totient\_series\_iff\_certificate\_supply}{LcmConeFlatness.lean:412}
\end{thm}

\begin{thm}[cert:b2 --- multiple-period collapse, wave-22, \texttt{irrational\_totient\_series\_of\_multiple\_certificate\_supply}]
$\big(\forall h_0>0,\ \forall N_0,\ \exists m>0,\ \exists N\ge N_0,\ \exists L,\ \mathrm{certifiedKill}(m\cdot h_0, N, L)\big) \to \mathrm{Irrational}(S)$. Weaker hypothesis than cert:a10: for every primitive period it suffices to certify \emph{some} multiple. Proof idea (\texttt{tail\_diff\_mul}, telescoping an $m\cdot h$-difference into a sum of $m$ shifted $h$-differences) is problem-agnostic; logically equivalent-strength weakening of cert:a10, still open.
\par\noindent\ev{Lean} (implication proved; hypothesis \ev{Open}) \scale{cofinal} \coord{other:lcm-period-multiple} \lean{irrational\_totient\_series\_of\_multiple\_certificate\_supply}{CarrySurvivorExtinction.lean:502}
\end{thm}

\begin{thm}[cert:b3 --- lcm-diagonal iff, the canonical single-parameter form]
$\big(\forall t_0:\mathbb{N},\ \exists t\ge t_0,\ \exists L,\ \mathrm{certifiedKill}(\mathrm{periodLcm}(t), \mathrm{periodLcm}(t), L)\big) \leftrightarrow \mathrm{Irrational}(S)$. Standing at $N=h=\mathrm{periodLcm}(t)$ beats every hypothetical rational simultaneously; conversely, pointwise completeness supplies a diagonal witness already at $t=t_0$. \textbf{This is an exact restatement with one quantified scale $t$}, not progress.  The finite floor is every $t\le82$, not a cofinal supply.
\par\noindent\ev{Lean} (equivalence proved; supply \ev{Open}) \scale{cofinal} \coord{other:lcm-diagonal} \lean{irrational\_totient\_series\_iff\_lcm\_diagonal\_certificate\_supply}{LcmConeFlatness.lean:426}
\end{thm}

\begin{lem}[cert:b4 --- lcm-window structure, \texttt{eq\_prime\_pow\_of\_not\_dvd\_periodLcm}]
For $0<j<2t$: if $j \nmid \mathrm{periodLcm}(t)$ then $\exists$ prime $p,k$ with $j=p^k \wedge t<j$ — below $2t$, every non-divisor of $\mathrm{lcm}(1..t)$ is a bare prime power exceeding $t$. Pure elementary number theory about $\mathrm{lcm}(1..t)$; structural input for anyone trying to search for a diagonal certificate — narrows where the "noise" in $\varphi$ on the window comes from.
\par\noindent\ev{Lean} \scale{uniform} \coord{other:lcm-window} \lean{eq\_prime\_pow\_of\_not\_dvd\_periodLcm}{LcmDiagonalReduction.lean:137}
\end{lem}

\begin{prop}[cert:b5 --- lcm-ray window totient factorisation, \texttt{totient\_periodLcm\_ray\_split}]
On a clean divisor $j \mid \mathrm{periodLcm}(t)$ (every prime factor of $j$ still divides $\mathrm{periodLcm}(t)/j$): $\varphi(q\cdot\mathrm{periodLcm}(t)+j) = \varphi(j)\cdot\varphi\big(q\cdot(\mathrm{periodLcm}(t)/j)+1\big)$. An exact multiplicative split of window totient values into a known local part and a cofactor forced $\equiv1$ modulo every exhausted prime — the exact algebraic handle needed to attempt an unconditional supply proof at cert:b3.
\par\noindent\ev{Lean} \scale{uniform} \coord{other:lcm-window-multiplicative} \lean{totient\_periodLcm\_ray\_split}{LcmDiagonalReduction.lean:208}
\end{prop}

\begin{thm}[cert:b6 --- lcm-cone flatness law, wave-24, \texttt{rational\_totient\_series\_forces\_lcm\_cone\_flatness}]
$\neg\mathrm{Irrational}(S) \to \exists t_1,\ \forall t\ge t_1,\ \forall q,m:\mathbb{N},\ 0<q \to \mathrm{totientTail}(q\cdot\mathrm{periodLcm}(t) + m\cdot\mathrm{periodLcm}(t)) - \mathrm{totientTail}(q\cdot\mathrm{periodLcm}(t)) \in \mathrm{range}((\uparrow):\mathbb{Z}\to\mathbb{R})$. Rationality forces one fractional constant on the \emph{entire} lcm cone $\{k\cdot\mathrm{periodLcm}(t):k\ge1\}$ at every scale $t\ge t_1$, not just the diagonal pair — a strict generalisation of cert:a9's hypothesis-generating side. Any certificate anywhere on the cone kills \#249 (cert:b7).
\par\noindent\ev{Lean} \scale{uniform} \coord{other:lcm-cone} \lean{rational\_totient\_series\_forces\_lcm\_cone\_flatness}{CertificateKernel.lean:19002}
\end{thm}

\begin{thm}[cert:b7 --- cone collapse, wave-24, annihilator umbrella, \texttt{irrational\_totient\_series\_of\_lcm\_cone\_window\_kill\_supply}]
$\big(\forall t_0:\mathbb{N},\ \exists t\ge t_0,\ \exists q,m,L:\mathbb{N},\ 0<q \wedge \mathrm{certifiedKill}(m\cdot\mathrm{periodLcm}(t),\, q\cdot\mathrm{periodLcm}(t),\, L)\big) \to \mathrm{Irrational}(S)$. One certified kill anywhere on the two-multiplier lcm cone, at arbitrarily large $t$, suffices. Diagonal (cert:b3) is the cell $q=m=1$; $q$-ray steps are $m=1$; prime-jump pairs are $(q,m)=(1,p-1)$. The widest known target still logically equivalent-in-strength to cert:a10.
\par\noindent\ev{Lean} (implication proved; hypothesis \ev{Open}) \scale{cofinal} \coord{other:lcm-cone} \lean{irrational\_totient\_series\_of\_lcm\_cone\_window\_kill\_supply}{CertificateKernel.lean:19014}
\end{thm}

\begin{cor}[cert:b8 --- pure non-integrality frontier forms, no certificate vocabulary]
Diagonal: $\big(\forall t_0,\ \exists t\ge t_0,\ \mathrm{totientTail}(2\cdot\mathrm{periodLcm}(t)) - \mathrm{totientTail}(\mathrm{periodLcm}(t)) \notin \mathrm{range}((\uparrow):\mathbb{Z}\to\mathbb{R})\big) \to \mathrm{Irrational}(S)$; cone form is the $(q,m)$ generalisation identical in shape. Via cert:a7's iff these hypotheses are \emph{exactly} equivalent to cert:b3/cert:b7 respectively. \textbf{The frontier of \#249 stripped of certificate vocabulary}: does $\mathrm{totientTail}(2H_t) - \mathrm{totientTail}(H_t) \notin \mathbb{Z}$ for infinitely many $t$, $H_t = \mathrm{lcm}(1..t)$?
\par\noindent\ev{Lean} (implication proved; hypothesis \ev{Open}) \scale{cofinal} \coord{other:real-analytic-nonintegrality} \lean{irrational\_totient\_series\_of\_lcm\_diagonal\_nonintegrality\_supply}{CertificateKernel.lean:19034} \lean{irrational\_totient\_series\_of\_lcm\_cone\_nonintegrality\_supply}{CertificateKernel.lean:19046}
\end{cor}

\begin{prop}[cert:b9a --- second-difference certificates, sound but measured not shallower, \texttt{second\_diff\_notMem\_int\_of\_certifiedRank2Kill}]
$\mathrm{certifiedRank2Kill}(h,N,L) \to \big(\mathrm{totientTail}(N+2h)-\mathrm{totientTail}(N+h)\big) - \big(\mathrm{totientTail}(N+h)-\mathrm{totientTail}(N)\big) \notin \mathrm{range}((\uparrow):\mathbb{Z}\to\mathbb{R})$. Sound second-difference non-integrality via a doubled band radius. Measured cell $(h,N)=(1,8)$: rank-1 fires at depth 8, no rank-2 certificate exists at depth $\le8$, rank-2 first fires at depth 9 — a probe over $t\le20$ finds rank-1 at least as shallow in 30/40 cells. This route is empirically \emph{not} a shortcut over rank-1; flagged do-not-re-attempt without new information.
\par\noindent\ev{Lean} \scale{uniform} \coord{other:binary-window} \lean{second\_diff\_notMem\_int\_of\_certifiedRank2Kill}{LcmConeFlatness.lean:493} (measured verdict: \lean{totient\_tail\_rank\_two\_kill\_sound\_but\_not\_shallower\_cell}{CertificateKernel.lean:19090})
\end{prop}

\begin{thm}[cert:b10a --- cone non-flatness menu refuter, wave-25, sharper than pairwise, \texttt{exists\_nonintegral\_pair\_of\_coneNonflatCert}]
For a nonempty menu $Q$ of positive vertex multipliers with the one-sided floor $\forall q\in Q,\ q\cdot H + L + 2 < 2^L$ (half the pairwise floor of $\mathrm{certifiedKill}$): $\mathrm{coneNonflatCert}(H,L,Q) \to \exists q_i,q_j \in Q,\ \mathrm{totientTail}(q_j\cdot H) - \mathrm{totientTail}(q_i\cdot H) \notin \mathrm{range}((\uparrow):\mathbb{Z}\to\mathbb{R})$. Proved by an argmin/Helly-avoidance argument over one-sided arcs: if all vertices shared one fractional part, the minimal-deep-tail vertex would be a common left endpoint of every arc, which $\mathrm{coneNonflatCert}$ denies. Information-theoretically half the depth floor of pairwise $\mathrm{certifiedKill}$; the reusable combinatorial-geometry technique transplants to any modular-residue pincer with more than two points.
\par\noindent\ev{Lean} \scale{uniform} \coord{other:lcm-cone-menu} \lean{exists\_nonintegral\_pair\_of\_coneNonflatCert}{LcmConeNonflat.lean:126}
\end{thm}

\begin{thm}[cert:b10b --- cone non-flat supply, wave-25, \texttt{irrational\_totient\_series\_of\_lcm\_cone\_nonflat\_supply}]
If $\mathrm{coneNonflatCert}$ fires (via cert:b10a) on a menu whose scale is unbounded, then $\mathrm{Irrational}(S)$. The sharpest known certificate-depth-reduced restatement of the wall; for $|Q|\ge3$ genuinely joint (menu inconsistent while every pair consistent) this is the best-known target for a search-based attempt at supplying cert:a10/cert:b3/cert:b7.
\par\noindent\ev{Lean} (implication proved; supply \ev{Open}) \scale{cofinal} \coord{other:lcm-cone-menu} \lean{irrational\_totient\_series\_of\_lcm\_cone\_nonflat\_supply}{CertificateKernel.lean:19140}
\end{thm}

\begin{prop}[cert:b12 --- survivorKill, alternative certificate via bounded carry orbit]
$\mathrm{survivorKill}(h,N,K) := \forall j < 2(N+h+1)+1,\ \exists i\le K,\ \mathrm{carryOrbit}(h,N,j-(N+h+1),i)$ escapes the strip $|\cdot|\le N+i+h+2$. Soundness: $\mathrm{survivorKill}(h,N,K) \to \mathrm{totientTail}(N+h)-\mathrm{totientTail}(N)\notin\mathrm{range}((\uparrow):\mathbb{Z}\to\mathbb{R})$ — a bounded-orbit argument, alternative to cert:a4/cert:a6's residue-window argument, exhausting all $2(N+h+1)+1$ integer box candidates and checking each provably escapes within $K$ steps. cert:a4 (window) and this (orbit) are two \emph{independently complete} coordinates for the same non-integrality target (via cert:a7's iff): the "enumerate all integer candidates in a shrinking box, verify each provably diverges" pattern is a general technique for irrationality-by-integer-orbit arguments.
\par\noindent\ev{Lean} \scale{uniform} \coord{other:carry-orbit} \lean{survivorKill}{CarrySurvivorExtinction.lean:412} \lean{tail\_diff\_notMem\_int\_of\_survivorKill}{CarrySurvivorExtinction.lean:428}
\end{prop}

\begin{thm}[mob:e1 --- first-harmonic norm-gap supply, \texttt{irrational\_totient\_series\_of\_first\_harmonic\_norm\_gap}]
Hypothesis $\mathrm{DTWFirstHarmonicNormGap} := \forall h>0,\ \forall X_0,\ \exists X\ge\max(X_0,1),\ \exists L,\ 16(2X+h+L+2)\le 2^L \wedge \big\|\sum_{N\in[X,2X)} \mathrm{windowFirstExp}(h,N,L)\big\| \le (21/25)\cdot X$. If this holds then $\mathrm{Irrational}(S)$. A genuinely cofinal first-harmonic exponential-sum (Weyl-sum) cancellation statement is sufficient for \#249; unconditional companion \texttt{exists\_certifiedKill\_of\_first\_harmonic\_gap} shows any constant-saving first-harmonic gap on one dyadic block forces a finite kill certificate (via $\cos(\pi/8)>9/10$ and a pigeonhole/averaging argument).
\par\noindent\ev{Lean} (conditional theorem proved; hypothesis \ev{Open}) \scale{cofinal} \coord{binary-digit} \lean{irrational\_totient\_series\_of\_first\_harmonic\_norm\_gap}{FirstHarmonicPivot.lean:83} \lean{exists\_certifiedKill\_of\_first\_harmonic\_gap}{FirstHarmonicGap.lean:128}
\end{thm}

\begin{thm}[mob:e2 --- prime-jump sharp-kill supply, \texttt{irrational\_totient\_series\_of\_primeJumpSharpKill\_supply}]
$\mathrm{primeJumpTailCommutator}(H,p) := D(pH) - p\cdot D(H)$, $D(H) = R_{2H}-R_H$, has an exact partial/tail split with sharp radius $3pH + (p+1)(L+2)$ (tighter than the earlier $4pH$ two-cell-disjunction radius). If $\forall t_0\, \exists t\ge t_0\, \exists p,L>0,\ \mathrm{primeJumpSharpKill}(\mathrm{periodLcm}(t), p, L)$, then $\mathrm{Irrational}(S)$. One concrete deposit is proved unconditionally: $\mathrm{primeJumpSharpKill}(12,5,15)$, kernel-\texttt{decide}d. Proof strategy: assume $S$ rational $\Rightarrow$ cone flatness (cert:b6) $\Rightarrow$ contradiction via the sharp prime-jump kill.
\par\noindent\ev{Lean} (conditional theorem + \ev{Cert} witness; supply \ev{Open}) \scale{cofinal} \coord{other:lcm-diagonal} \lean{irrational\_totient\_series\_of\_primeJumpSharpKill\_supply}{PrimeJumpWindow.lean:193} \lean{primeJumpSharpKill\_twelve\_five}{PrimeJumpWindow.lean:186}
\end{thm}
% ---- part a249_p1b ----
\subsection{Consequence theorems}

This subsection lists every premise in the corpus whose conclusion is
\emph{conditional} on an unsupplied hypothesis, with that hypothesis rendered
in display maths so an agent can recognise it as a search target. All targets
below are logically downstream of the single exact equivalence that defines
the open problem:

\begin{prop}[The exact target: \#249 = cofinal actual-orbit nonintegrality]\label{prop:NI-01}
Unconditionally,
\[
\mathrm{Irrational}\Big(\textstyle\sum_n' \varphi(n)/2^n\Big)
\iff
\forall a_0,\ \exists a \ge a_0,\ \mathtt{actualLcmTailOrbit}\ a \notin \mathrm{range}(\mathbb{Z}\to\mathbb{R}).
\]
Nothing beyond a supply of the right-hand side is needed in principle; every
other proposition in this subsection is a route to it.
\scale{n/a}\ \ev{Lean}\ \coord{mobius-mersenne}\\
\lean{irrational_totientSeries_iff_actualLcmOrbitNonintegralitySupply}{TotientActualLcmOrbitNonintegrality.lean:37}
\end{prop}

\begin{prop}[Short-window arithmetic-kill supply]\label{prop:AR-07}
If
\[
\forall a_0,\ \exists a\ge a_0,\ \exists L< 2\cdot 2^a,\quad \mathtt{LcmDiagonalArithmeticKill}(2^a,L)
\]
holds, then \#249 follows (via \mathtt{actualLcmOrbitNonintegralitySupply\_of\_shortArithmeticKill}
then Prop.~\ref{prop:NI-01}). The predicate is exactly the residue-band exclusion of
Appendix~A4/A6: $\mathtt{lcmDiagonalArithmeticWord}$ at scale $2^a$ escapes a shrinking
central arc mod $2^L$. Only two instances are proved (Prop.~\ref{prop:SK-02}'s base cases);
the cofinal supply is the open trigger of the actual-orbit batch.
\scale{cofinal}\ \ev{Open}\ \coord{mobius-mersenne}\\
\lean{PowerTwoActualLcmShortArithmeticKillSupply}{TotientActualLcmOrbitArithmetic.lean:2107}
\end{prop}

\begin{prop}[Diophantine separation supply]\label{prop:SEP-03}
If, at canonically-guarded odd ranks $q$, cofinally many $a$ satisfy
\[
\big|\,\mathtt{actualLcmTailOrbit}\ a - z\,\big| > \tfrac{1}{32} + (\text{explicit error radius}) \qquad \forall z\in\mathbb Z,
\]
then \#249 follows via the landed signed-margin producer. This restates the target as
effective irrationality-measure / anti-concentration for the actual orbit rather than
exact residue exclusion, using the explicit approximant of Prop.~\ref{prop:SEP-02-inv}.
\scale{cofinal}\ \ev{Open}\ \coord{mobius-mersenne}\\
\lean{irrational_totientSeries_of_actualLcmOrbitSeparationSupply}{TotientActualLcmOrbitSeparation.lean:305}
\end{prop}

\begin{prop}[One-sided top-edge residue-gap supply]\label{prop:TE-04}
If, for the room bound of Prop.~\ref{prop:SGN-01} ($a\ge 8$, $J+K+(a+6)<2\cdot2^a$),
\[
\mathtt{ActualLcmTopEdgeResidueGap}\ a\ J\ K\ m
\quad:\iff\quad
m\le K \ \wedge\ \big(\text{residue of }\mathtt{windowDiscrepancy}\text{ at scale }m\big)\le 2^m-(\text{room bound})
\]
holds cofinally, then \#249 follows (\mathtt{actualLcmTailDiff\_notMem\_int\_of\_topEdgeResidueGap},
sufficiency at line~1260), because Prop.~\ref{prop:SGN-01}/\ref{prop:SGN-03} already exclude the
negative-side residue independently — only the positive arc needs excluding. This is a strictly
\emph{weaker} target than the old symmetric certifiedKill band and is the genuinely easier open
target of the whole batch.
\scale{cofinal}\ \ev{Open}\ \coord{mobius-mersenne}\\
\lean{PowerTwoActualLcmTopEdgeResidueGapSupply}{TotientActualLcmTopEdgeStaircase.lean:1325}
\end{prop}

\begin{prop}[The five-link equivalence chain below the residue-gap target]\label{prop:TE-05}
Each of the following cofinal supplies is proved \emph{sufficient} for
Prop.~\ref{prop:TE-04}'s target, chained in decreasing strength:
\[
\begin{aligned}
&\mathtt{PowerTwoActualLcmTopEdgeResidueGapSupply}
 \Leftarrow \mathtt{PowerTwoAdjacentSuffixMidbandSupply}\\
&\Leftarrow \mathtt{PowerTwoOddGuardTopEdgeHalfWordBandSupply}
 \iff \mathtt{PowerTwoActualFinalTopEdgeMagnitudeSupply}\\
&\Leftarrow \mathtt{PowerTwoFlexibleActualTopEdgeMagnitudeSupply}
 \Leftarrow \mathtt{PowerTwoFlexibleActualTerminalDominanceSupply}\\
&\Leftarrow \mathtt{PowerTwoFlexibleActualTerminalCarryCorridorEscapeSupply}
\end{aligned}
\]
All six named predicates are open; none is proved. Because the chain is
implication-only-downward, proving the single \emph{weakest} link
(\texttt{...TerminalCarryCorridorEscapeSupply}) closes the entire cluster and hence \#249.
This is the true minimal remaining target of \texttt{TotientActualLcmTopEdgeStaircase}.
\scale{cofinal}\ \ev{Open}\ \coord{mobius-mersenne}\\
\lean{powerTwoActualLcmTopEdgeResidueGapSupply\_of\_adjacentSuffixMidband}{TotientActualLcmTopEdgeStaircase.lean:2108}
\end{prop}

\begin{prop}[The exact closed-form escape identity]\label{prop:TE-06}
Under integrality of the actual orbit (representative $z$) and the half-cell fit condition
$2(H+q+2)\le 4^q$ ($a\ge 8$, room bound as above),
\[
2\cdot \mathtt{actualOddHalfCenteredLift}\ a\ q \;=\; \mathtt{diagonalWindowIncrement}(2^a)(2q+2) \;-\; \mathtt{carryOrbit}\ H\ H\ z\ (2q+1)
\]
holds \emph{exactly}. Escaping either side of this named open interval (the ``terminal/carry
corridor'') is both necessary and sufficient for non-integrality at rank $q$
(\mathtt{actualLcmTailOrbit\_notMem\_int\_of\_actualTerminalCarryCorridorEscape}). This is
the sharpest fully explicit statement of ``what remains to prove'' anywhere in the batch.
\scale{bounded}\ \ev{Lean}\ \coord{mobius-mersenne}\\
\lean{two_mul_actualOddHalfCenteredLift_eq_terminal_sub_trueCarry}{TotientActualLcmTopEdgeStaircase.lean:1678}
\end{prop}

\begin{prop}[Extend the short-kill stub past $a_0=6$]\label{prop:SK-02}
Currently proved only for $a_0\le 6$:
\[
\forall a_0\le 6,\ \exists a\ L,\ a_0\le a\ \wedge\ L<2\cdot2^a\ \wedge\ \mathtt{LcmDiagonalArithmeticKill}(2^a)\ L.
\]
The module's own comment states plainly: ``the remaining endpoint gap is now precisely the
unbounded continuation beyond this finite prefix.'' Supplying a single further instance at
$a_0=7$ or $a_0=8$ already extends the stub and is independently interesting evidence toward
Prop.~\ref{prop:AR-07}'s full cofinal supply; the two base cases $a=4,6$ trace to
\mathtt{certifiedKill\_diagonal\_t16}/\mathtt{t64}.
\scale{bounded}\ \ev{Lean (finite prefix only)}\ \coord{mobius-mersenne}\\
\lean{powerTwoActualLcmShortArithmeticKillSupply_through_six}{TotientActualLcmShortKill.lean:54}
\end{prop}

\begin{prop}[Fixed-rank extremal-ordering supply]\label{prop:FR-01}
Define $\mathtt{fixedRankSecondDifference}\ H\ j := \varphi(3H+j)-2\varphi(2H+j)+\varphi(H+j)$.
If, at $H=\mathtt{periodLcm}(2^a)$ and some fixed small $j$,
\[
\mathtt{MiddleRankTotientExtremal}\ H\ j
\quad:\iff\quad
\varphi(2H+j)\ \text{is a strict min or max among}\ \{\varphi(H+j),\varphi(2H+j),\varphi(3H+j)\},
\]
then $\mathtt{fixedRankSecondDifference}\ H\ j\ne 0$, with sign matching the extremum. Notably
an \emph{ordering} suffices — no quantitative gap is required. This is a genuinely different
coordinate from the sliding LCM window: $j$ fixed and small, only three fixed ranks examined.
The remaining open step is \mathtt{PowerTwoLcmMiddleRankExtremalSupply}: does this ordering hold
cofinally in $a$.
\scale{uniform}\ \ev{Lean}\ \coord{other:fixed-rank-curvature}\\
\lean{MiddleRankTotientExtremal}{TotientFixedRankLcmAsymptotic.lean:267}
\end{prop}

\begin{prop}[The directed/LCM-specialised certificate supply]\label{prop:CP-06}
$\mathtt{directedCertifiedKill}\ h\ N\ L$ is sound \emph{and complete}:
\[
(\exists L,\ \mathtt{directedCertifiedKill}\ h\ N\ L) \iff \mathtt{totientTail}(N+h)-\mathtt{totientTail}(N)\notin \mathrm{range}(\mathbb Z\to\mathbb R),
\]
exactly, no gap; and $\mathrm{Irrational}(S)\iff \mathtt{CofinalDirectedLcmCertificateSupply}$ (the
LCM-diagonal specialisation). The asymmetric strip is a genuine finite-depth improvement over
the symmetric certificate (kills $t=3$ at depth~6, one level earlier than the symmetric one), but
the file itself notes the improvement does not turn into an independent sieve theorem: supplying
the cofinal predicate is exactly as hard as \#249 itself.
\scale{uniform}\ \ev{Lean}\ \coord{seam-integer}\\
\lean{CofinalDirectedLcmCertificateSupply}{TotientTailCarryPeriod.lean:866}
\end{prop}

\begin{prop}[Pulse-restricted survivor-search supply]\label{prop:CP-07}
Given the cofinal mod-4 pulse supply of Prop.~\ref{prop:CP-05-inv}, it suffices to kill only the
$2\bmod 4$-class candidate states (a fourfold reduction of the initial search strip) at one
cofinal arithmetic-pulse prime per putative period, to conclude \#249:
\[
\mathtt{modFourPulseSurvivorKill} \implies \mathrm{Irrational}(S).
\]
The reduction technique (use a cofinal totient-specific pulse to legally restrict the survivor
search to one residue class mod a small modulus) transfers wherever an analogous cofinal pulse
can be built on the \#257 side.
\scale{cofinal}\ \ev{Lean}\ \coord{mobius-mersenne}\\
\lean{irrational_totientSeries_of_cofinal_modFourPulseSurvivorKill}{TotientTailCarryPeriod.lean:1066}
\end{prop}

\begin{prop}[The wave-21 wall: certificate supply over all periods]\label{prop:A10}
\[
\big(\forall h\ge 1,\ \forall N_0,\ \exists N\ge N_0,\ \exists L,\ \mathtt{certifiedKill}\ h\ N\ L\big)
\iff \mathrm{Irrational}(S).
\]
This is exactly the quantifier structure $\forall h\ge1\ \forall N_0\ \exists N\ge N_0\ \exists L\ \mathrm{Sep}(h,N,L)$.
Nothing in the corpus supplies this predicate; every other reduction in Part~B
of the certificate bank is a logically equivalent-or-weaker reformulation of this same missing
supply, never independent progress on it.
\scale{cofinal}\ \ev{Lean} (equivalence; supply \ev{Open})\ \coord{mobius-mersenne}\\
\lean{irrational_totient_series_iff_certificate_supply}{LcmConeFlatness.lean:412}
\end{prop}

\begin{prop}[Diagonal collapse — one free parameter]\label{prop:B3}
\[
\big(\forall t_0,\ \exists t\ge t_0,\ \exists L,\ \mathtt{certifiedKill}(\mathtt{periodLcm}\,t)(\mathtt{periodLcm}\,t)\ L\big)
\iff \mathrm{Irrational}(S).
\]
Standing at $N=h=\mathtt{periodLcm}\,t$ beats every hypothetical rational simultaneously (any
$t\ge\max(h_0,N_0)$ serves both parameters at once), collapsing Prop.~\ref{prop:A10}'s
two unbounded parameters to one. Conversely, pointwise completeness supplies
the diagonal witness at $t=t_0$.  This is the canonical single-scale
restatement, not an advance.
\scale{cofinal}\ \ev{Lean} (equivalence; supply \ev{Open})\ \coord{mobius-mersenne}\\
\lean{irrational_totient_series_iff_lcm_diagonal_certificate_supply}{LcmConeFlatness.lean:426}
\end{prop}

\begin{prop}[Cone collapse — annihilator umbrella]\label{prop:B7}
\[
\big(\forall t_0,\ \exists t\ge t_0,\ \exists q\,m\,L,\ 0<q \wedge \mathtt{certifiedKill}(m\cdot\mathtt{periodLcm}\,t)(q\cdot\mathtt{periodLcm}\,t)\ L\big)
\implies \mathrm{Irrational}(S).
\]
One certified kill \emph{anywhere} on the two-multiplier LCM cone, at arbitrarily large $t$,
suffices; the diagonal (Prop.~\ref{prop:B3}) is the cell $q=m=1$. This is the widest known
target still logically equivalent-in-strength to Prop.~\ref{prop:A10}.
\scale{cofinal}\ \ev{Open}\ \coord{mobius-mersenne}\\
\lean{irrational_totient_series_of_lcm_cone_window_kill_supply}{CertificateKernel.lean:19014}
\end{prop}

\begin{prop}[Menu non-flatness supply — sharper than pairwise]\label{prop:B10}
For a nonempty menu $Q$ of positive vertex multipliers with one-sided floor
$\forall q\in Q,\ qH+L+2<2^L$ (\emph{half} the pairwise floor of \texttt{certifiedKill}),
\[
\mathtt{coneNonflatCert}\ H\ L\ Q \implies \exists\, q_i,q_j\in Q,\ \mathtt{totientTail}(q_jH)-\mathtt{totientTail}(q_iH)\notin \mathrm{range}(\mathbb Z\to\mathbb R).
\]
If this fires at unbounded scale over a menu with $|Q|\ge3$ genuinely joint (menu inconsistent
while every pair is separately consistent), \#249 follows via Prop.~\ref{prop:B7}. The
argmin/Helly-avoidance proof technique is a reusable combinatorial-geometry pattern for any
modular-residue pincer with more than two points.
\scale{cofinal}\ \ev{Open}\ \coord{mobius-mersenne}\\
\lean{irrational_totient_series_of_lcm_cone_nonflat_supply}{CertificateKernel.lean:19140}
\end{prop}

\begin{prop}[Unbounded Farey growth — a second independent wall]\label{prop:C2sup}
The Farey-gap denominator bound at window $K$ is currently $\sim 7.96\times10^{34}$ at $K=240$
(Prop.~\ref{prop:C2-inv}). If
\[
\sup_K\, (b+d)(K) = \infty
\]
(the bound growing without limit as $K\to\infty$), then \#249 follows via the $C3$-style
denominator-exclusion consumer, logically \emph{independently} of Prop.~\ref{prop:A10}'s
certificate-supply wall. No unboundedness claim is proved or attempted in the corpus.
\scale{cofinal}\ \ev{Open}\ \coord{other:farey-gap}\\
\lean{gap\_check\_window\_1\_240\_le\_79639646646701375323355774875831053}{GapFareyBound.lean:176}
\end{prop}

\begin{prop}[Rank lower bound; the generic compression shortcut is retired]\label{prop:D5cons}
By Prop.~\ref{prop:D4-inv}, the canonical dyadic totient-kernel family is unconditionally
$(2^e+1)$-dimensional at every level $e$. Rationality of $S$ forces an associated tempered
carry orbit with $\mathbb Q$-rank $\ge 2^e-1$ at every level (Prop.~\ref{prop:CP-02}). Formally, a theorem
\[
\text{bounding the dyadic-section rank of \emph{every} rationality-supplied tempered carry}
\]
would contradict this lower bound and close \#249. The corpus supplies no such theorem or
mechanism. More strongly, its compressed-adjoint impossibility result and its explicit
all-horizon finite-rank shift-polynomial countermodel retire the generic finite-compression
shortcut (Observation~\ref{prop:B4b-kill}). Thus the displayed upper bound is only a logically
sufficient new input, not a third live frontier and not evidence that existing rank machinery is
close to a contradiction.
\scale{cofinal}\ \ev{Open}\ \coord{other:carry-kernel-rank}\\
\lean{not_irrational_totientSeries_implies_unbounded_carryRank_unconditional}{TotientCarryKernelRigidity.lean:300}
\end{prop}

\begin{prop}[Alternative coordinate: bounded carry-orbit survivor supply]\label{prop:B12cons}
$\mathtt{survivorKill}\ h\ N\ K$ enumerates all $2(N+h+1)+1$ integer box candidates and checks
each provably diverges from the strip $|\cdot|\le N+i+h+2$ within $K$ steps; soundness gives
$\mathtt{survivorKill}\ h\ N\ K \implies \mathtt{totientTail}(N+h)-\mathtt{totientTail}(N)\notin\mathrm{range}(\mathbb Z\to\mathbb R)$,
proving the \emph{same} target as Prop.~\ref{prop:A10}/\ref{prop:NI-01} from a different
coordinate (bounded dynamics, not residue windows). Per the doctrine ``obstructions are
coordinate-relative,'' a case hard in the window coordinate may be easy here.
\scale{n/a}\ \ev{Lean}\ \coord{other:carry-orbit-dynamics}\\
\lean{tail_diff_notMem_int_of_survivorKill}{CarrySurvivorExtinction.lean:428}
\end{prop}

\subsection{Invariants}

Necessary conditions that any successful (or unsuccessful) approach must be
consistent with — facts a solution cannot contradict, whichever way \#249
eventually resolves.

\begin{prop}[Exact quotient-scale digit closed form, no cleanliness hypothesis]\label{prop:AR-04-inv}
\[
\mathtt{lcmRayArithmeticLetter}\ t\ j \;=\; \mathtt{deltaTotient}(\mathtt{periodLcm}\,t)(\mathtt{periodLcm}\,t + j) \;=\; \mathtt{diagonalWindowIncrement}\ t\ j,
\]
valid for \emph{both} divisor and non-divisor offsets $j$, with no side condition on which primes
of $j$ survive in $H/j$. Every downstream sign, positivity, or residue statement about the
diagonal word must factor through this identity; it generalises the older
$\mathtt{deltaTotient\_periodLcm\_ray\_split}$, which needed $j$'s primes to survive.
\scale{uniform}\ \ev{Lean}\ \coord{mobius-mersenne}\\
\lean{lcmRayArithmeticLetter_eq_deltaTotient}{TotientActualLcmOrbitArithmetic.lean:298}
\end{prop}

\begin{prop}[Uniform $3/4$-density retained on rough integers]\label{prop:AR-03-inv}
For $a\ge8$, $t=2^a$, every prime factor of $n>0$ exceeding $t$, and $n<2^{2\cdot2^a}$:
$n.\mathrm{primeFactors.card} < 2^a/4$, hence $(3/4)\cdot n < \varphi(n)$ over $\mathbb Q$. Any
argument bounding $\varphi$ on the short window below height $\mathtt{periodLcm}(2^a)^2$ must
respect this floor; the counting mechanism (union-bound-in-product form, then convert via
$\mathtt{Nat.totient\_eq\_mul\_prod\_factors}$) is reusable for \#257's Mersenne-rough objects
at the analogous height.
\scale{bounded}\ \ev{Lean}\ \coord{mobius-mersenne}\\
\lean{three_quarters_mul_lt_totient_of_rough_lt_two_pow}{TotientActualLcmOrbitArithmetic.lean:768}
\end{prop}

\begin{prop}[Every short-window digit sign is fixed to $+$]\label{prop:AR-05-inv}
For $a\ge8$, $0<j<2\cdot2^a$: $0<\mathtt{lcmRayArithmeticLetter}(2^a)\ j$. Every coefficient of
the actual diagonal word in the entire short window is strictly positive — divisor offsets via
a quantitative $H<4jc$ bound, foreign prime powers via predecessor descent plus
Prop.~\ref{prop:AR-03-inv}'s rough-density bracket ($H/4<c<5H/2$), exponent-one new primes via
$H/4<c$. This unconditional fact (no irrationality hypothesis) is what makes every sign theorem
in the batch — Prop.~\ref{prop:SGN-01} and the staircase-impossibility of
Prop.~\ref{prop:TE-01-killer} — provable at all.
\scale{bounded}\ \ev{Lean}\ \coord{mobius-mersenne}\\
\lean{lcmRayArithmeticLetter_pos_of_lt_two_mul}{TotientActualLcmOrbitArithmetic.lean:1703}
\end{prop}

\begin{prop}[Window word = truncated discrepancy; kill $\iff$ certificate]\label{prop:AR-06-inv}
$\mathtt{lcmDiagonalArithmeticWord}\ t\ L = \mathtt{windowDiscrepancy}(\mathtt{periodLcm}\,t)(\mathtt{periodLcm}\,t)\ L$,
and $\mathtt{LcmDiagonalArithmeticKill}\ t\ L \iff \mathtt{certifiedKill}(\mathtt{periodLcm}\,t)(\mathtt{periodLcm}\,t)\ L$.
Any residue-band existence claim about the actual diagonal word converts directly to a
$\mathtt{certifiedKill}$/non-integrality fact through this iff — it is the fixed bridge every
consumer in this subsection implicitly uses.
\scale{n/a}\ \ev{Lean}\ \coord{mobius-mersenne}\\
\lean{lcmDiagonalArithmeticKill_iff_certifiedKill}{TotientActualLcmOrbitArithmetic.lean:2045}
\end{prop}

\begin{prop}[Exact global-to-local bridge]\label{prop:SEP-01-inv}
\[
\mathtt{actualLcmTailOrbit}\ a \;=\; 2^H(2^H-1)\Big(\textstyle\sum_n' \varphi(n)/2^n\Big) - \big(\mathtt{totientPrefix}(2H)-\mathtt{totientPrefix}(H)\big),\quad H=\mathtt{periodLcm}(2^a).
\]
Every local orbit statement is secretly a statement about the global series value minus a
computable finite prefix; any rational-approximation or continued-fraction style separation
argument for \#249 must express its target through this exact affine image.
\scale{n/a}\ \ev{Lean}\ \coord{mobius-mersenne}\\
\lean{actualLcmTailOrbit_eq_scaled_totientSeries_sub_prefix}{TotientActualLcmOrbitSeparation.lean:68}
\end{prop}

\begin{prop}[Explicit computable approximant with explicit error radius]\label{prop:SEP-02-inv}
For all $a,q$: $|\mathtt{actualLcmTailOrbit}\ a - \mathtt{actualLcmRawApprox}\ a\ q| < (2H+2q+3)/2^{2q+1}$,
where $\mathtt{actualLcmRawApprox}\ a\ q := \mathtt{diagonalAdjacentSuffixRawBlock}(2^a,0,2q{+}1)/2^{2q+1}$.
Any separation-from-integers argument (Prop.~\ref{prop:SEP-03}) is forced to work with this
finite, computable block rather than the infinite tail, at this exact error rate.
\scale{uniform}\ \ev{Lean}\ \coord{mobius-mersenne}\\
\lean{abs_actualLcmTailOrbit_sub_rawApprox_lt}{TotientActualLcmOrbitSeparation.lean:141}
\end{prop}

\begin{prop}[Unconditional positive-sign corridor — a real theorem, not a supply]\label{prop:SGN-01}
For $a\ge8$, $J+(a+6)<2\cdot2^a$:
\[
0 < \mathtt{totientTail}(2H+J) - \mathtt{totientTail}(H+J),\qquad H=\mathtt{periodLcm}(2^a),
\]
with no irrationality hypothesis — the true, infinite, real translated tail difference is
strictly positive throughout almost the entire short window, proved unconditionally from
Prop.~\ref{prop:AR-05-inv} plus a directed one-sided tail bound. Any argument about the sign of
the actual orbit (not merely its residue mod $2^L$) must agree with this; specialised at $J=0$
this gives $0<\mathtt{actualLcmTailOrbit}\ a$.
\scale{bounded}\ \ev{Lean}\ \coord{mobius-mersenne}\\
\lean{actualLcmTailDiff_shift_pos}{TotientActualLcmOrbitSign.lean:39}
\end{prop}

\begin{prop}[Integrality forces the exact top-edge residue — names the obstruction]\label{prop:SGN-03}
Under the room bound of Prop.~\ref{prop:SGN-01} plus $2H+J+K+2<2^K$: integrality of the actual
orbit forces
\[
\mathtt{windowDiscrepancy}\ H\ (H{+}J)\ K \bmod 2^K \;=\; 2^K - e
\]
(the top-edge representative), provably \emph{outside} the central arc required by
$\mathtt{directedCertifiedKill}$. This documents precisely why the sign theorem alone cannot
close \#249: any contradiction needs an independent exclusion of this specific top-edge boundary
band, not a re-derivation of the carry reset. The named residue $2^K-e$ ($e>0$ small) is exactly
what Prop.~\ref{prop:TE-04} sets out to exclude.
\scale{bounded}\ \ev{Lean}\ \coord{mobius-mersenne}\\
\lean{actualLcm_integral_forces_topEdgeResidue}{TotientActualLcmOrbitSign.lean:211}
\end{prop}

\begin{prop}[Punctured staircase is pinned to the half-turn]\label{prop:TE-02-inv}
Under the room bound and the punctured-staircase hypothesis (all but the last letter vanish at
growing dyadic weight; the last letter retained below the top-edge carry band),
\[
\mathtt{lcmRayArithmeticLetter}(2^a)(J{+}K{-}1) = 2^{m-1} \quad\text{exactly, and}\quad 2^m < 2(2H+J+K+2).
\]
The modulus must be the first dyadic scale above the room bound; any extra bit of modulus makes
even the punctured route empty. Any attempt to revive a partial-staircase strategy must land
exactly here.
\scale{bounded}\ \ev{Lean}\ \coord{mobius-mersenne}\\
\lean{puncturedDyadicStaircase_penultimate_eq_half}{TotientActualLcmTopEdgeStaircase.lean:1187}
\end{prop}

\begin{prop}[Any certificate reduces to a two-bit test — universal normal form]\label{prop:TE-03-inv}
Unconditionally, for arbitrary $h,N$ (no totient content in the proof):
\[
(\exists L,\ \mathtt{certifiedKill}\ h\ N\ L) \iff \mathtt{GuardCylinderWitness}\ h\ N,
\]
where the witness is a logarithmic-depth socket $(b{+}1)$ or a two-bit mixed-guard cylinder at
scale $b=\log_2(N{+}h{+}L{+}2){+}1$. Any certificate-search algorithm or complexity bound for
either problem's tail differences must respect this compression — an unbounded-depth search is
never truly necessary once $b$ is fixed.
\scale{uniform}\ \ev{Lean}\ \coord{seam-integer}\\
\lean{exists_certifiedKill_iff_guardCylinderWitness}{TotientActualLcmTopEdgeStaircase.lean:844}
\end{prop}

\begin{prop}[Primitive kernel factor of the fixed-rank curvature]\label{prop:FR-02-inv}
On the square-root clean window $j^2\le2^a$, $a\ge4$: $2\varphi(j) \mid \mathtt{fixedRankSecondDifference}(\mathtt{periodLcm}(2^a))\ j$.
The exact local factor $2\varphi(j)$ comes from the affine-rank-3 kernel $(1,-2,1)$ combined with
the parity of three odd rough cofactors, and is provably tight — Prop.~\ref{prop:FR-03-kill}
shows no bare application of the kernel can force one more factor of $2$. Any curvature-based
attack at fixed rank must land inside this exact divisibility ceiling.
\scale{bounded}\ \ev{Lean}\ \coord{other:fixed-rank-curvature}\\
\lean{two_mul_totient_dvd_fixedRankSecondDifference}{TotientFixedRankLcmAsymptotic.lean:519}
\end{prop}

\begin{prop}[Carry displacement $\iff$ integral tail difference — central plumbing]\label{prop:CP-01-inv}
For a positive-multiplier tempered totient carry $u$ ($\mathtt{IsTemperedBinaryOrbit}\ \varphi\ v\ u$, $v>0$):
\[
(v:\mathbb Z)\mid u(N{+}k)-u(N) \iff \mathtt{totientTail}(N{+}k)-\mathtt{totientTail}(N)\in\mathrm{range}(\mathbb Z\to\mathbb R).
\]
Every statement in this batch that converts between ``carry orbit divisibility'' and ``real tail
difference is an integer'' is an instance of this identity — the underlying machinery
($\mathtt{IsTemperedBinaryOrbit}$, $\mathtt{binaryCoeffTail}$) is generic over any $f$ with
$f(n)\le n$, so it transfers to \#257's own coefficient function verbatim.
\scale{uniform}\ \ev{Lean}\ \coord{seam-integer}\\
\lean{carryShift_dvd_iff_tailDiff_mem_int}{TotientTailCarryPeriod.lean:77}
\end{prop}

\begin{prop}[What rationality does and does not buy — the exact frontier]\label{prop:CP-02}
If $\neg\mathrm{Irrational}(S)$, there exist $v>0$ and a tempered orbit $u$ such that for every
level $e$, the $\mathbb Q$-rank of the canonical carry-kernel-family span is $\ge 2^e-1$ (torsion-free
rank grows exponentially in level), \emph{yet} $u$'s dyadic sections are uniformly eventually
periodic modulo $v$ ($\mathtt{CarrySectionsEventuallyPeriodicMod}$). Rationality buys
quotient-mod-$v$ periodicity but does \emph{not} buy any bound on $\mathbb Q$-rank — this is
precisely the boundary any Mersenne-specific residue argument (mod-4 pulse, two-adic pulse) exists
to cross. Proposition~\ref{prop:D5cons} records the rank lower bound while explicitly retiring
generic finite-rank compression as a live shortcut.
\scale{uniform}\ \ev{Lean}\ \coord{other:carry-kernel-rank}\\
\lean{not_irrational_totientSeries_implies_mod_period_and_unbounded_rank}{TotientTailCarryPeriod.lean:224}
\end{prop}

\begin{prop}[Cofinal mod-4 residue pulse — genuinely totient-specific]\label{prop:CP-05-inv}
For any $h>0$ and bound $B$, there is a prime $p>B$ with $\mathtt{deltaTotient}(4h)\ p \equiv 2\pmod4$,
constructed via Dirichlet on $p\equiv(3r-H)\bmod 4r$ for a fresh prime $r\equiv1\pmod4$. Unlike the
Mersenne-primitivity killer of Prop.~\ref{prop:CP-03-kill}, this \emph{is} a genuinely
totient-specific residue supplied cofinally beyond every threshold — any argument invoking
Prop.~\ref{prop:CP-07} inherits this as its forcing mechanism. The fresh-prime CRT template is
scaled to arbitrary depth in Prop.~\ref{prop:TA-inv} below.
\scale{cofinal}\ \ev{Lean}\ \coord{other:dirichlet-crt}\\
\lean{exists_prime_deltaTotient_four_mul_mod_four_two}{TotientTailCarryPeriod.lean:384}
\end{prop}

\begin{prop}[Arbitrary-depth zero-prefix-then-pulse construction, and its transfer]\label{prop:TA-inv}
For $K\ge2$, $H>K$, $B$: cofinally many primes $p>B$ satisfy
$p\equiv1+2^{K-1}\pmod{2^K}$, $2^K\mid\varphi(p{+}H)$, and $2^K\mid\varphi(p{-}j)\wedge2^K\mid\varphi(p{-}j{+}H)$
for every $1\le j<K$ — an entire length-$(K{-}1)$ zero prefix followed by a half-turn terminal
pulse, at arbitrary two-adic depth (Prop.~\ref{prop:CP-05-inv}'s $K=2$ case generalised to every
$K$, via $K{-}1$ fresh Dirichlet primes glued by CRT). This forces
\[
\mathtt{windowDiscrepancy}\ H\ (p{-}K)\ K \equiv 2^{K-1}\pmod{2^K},
\]
and under eventual integrality transfers to $\exists z,\ (z:\mathbb R)=\mathtt{totientTail}(p{+}H)-\mathtt{totientTail}(p)\wedge z\equiv2^{K-1}\pmod{2^K}$.
Any solution attempt must be consistent with the existence of this pulse family at every depth $K$.
\scale{cofinal}\ \ev{Lean}\ \coord{p-adic}\\
\lean{exists_prime_totient_twoAdic_pulse_divisors}{TotientTwoAdicPulseBlock.lean:54}
\end{prop}

\begin{prop}[Exact Möbius-inversion tail identity]\label{prop:MP-01-inv}
$\mathtt{totientTail}\ N = \sum_d' \mu(d)\cdot 2^{d-r_d(N)}\big(q_d(N)/(2^d-1) + 1/(2^d-1)^2\big)$,
with $r_d(N)=d-N\bmod d$ (forward shift to the next multiple of $d$), $q_d(N)=\lfloor N/d\rfloor+1$ —
exact, not a definition-by-subtraction from a target. Any Möbius/Lambert-series manipulation of
$\mathtt{totientTail}$ in this batch or in \#257 factors through this identity and its fully
generic regrouping primitive $\mathtt{tsum\_lambert\_pair\_regroup\_if}$ (proved for arbitrary
weight functions $|w(d)|\le d$, $|v(m)|\le m$).
\scale{uniform}\ \ev{Lean}\ \coord{cyclotomic}\\
\lean{totientTail_eq_tsum_shiftedMobiusPulse}{TotientShiftedMobiusPulse.lean:547}
\end{prop}

\begin{prop}[The rare unconditional deposits — necessary base cases]\label{prop:SK-01-inv}
$\mathtt{LcmDiagonalArithmeticKill}(2^4)\ 23$ and $\mathtt{LcmDiagonalArithmeticKill}(2^6)\ 93$ are
proved, routed through pre-existing compressed certificates
$\mathtt{certifiedKill\_diagonal\_t16}/\mathtt{t64}$, giving $\mathtt{actualLcmTailOrbit}\ 4,6\notin\mathrm{range}(\mathbb Z\to\mathbb R)$
— the \emph{only} two unconditional actual-orbit nonintegrality facts currently proved. Any
induction or bootstrapping attempt at Prop.~\ref{prop:AR-07} must reproduce these as base cases,
and they double as a computational sanity check that the $\mathtt{lcmRayArithmeticLetter}$
machinery agrees with the older compressed-certificate route.
\scale{fixed}\ \ev{Cert}\ \coord{mobius-mersenne}\\
\lean{actualLcmTailOrbit_four_and_six_notMem_int}{TotientActualLcmShortKill.lean:35}
\end{prop}

\begin{prop}[Certificate depth floor]\label{prop:A5-inv}
$\mathtt{certifiedKill}\ h\ N\ L \implies 2(N{+}h{+}L{+}2)<2^L$. Certificate depth $L$ cannot stay
bounded while $N{+}h\to\infty$: $L$ must grow at least logarithmically. Any complexity argument
about certificate search inherits this floor, and it is exactly what Prop.~\ref{prop:TE-03-inv}
compresses to a two-bit test at the corresponding logarithmic scale.
\scale{n/a}\ \ev{Lean}\ \coord{seam-integer}\\
\lean{certifiedKill_depth_floor}{TotientTailPeriodKiller.lean:79}
\end{prop}

\begin{prop}[The tail-period law — necessary consequence of rationality]\label{prop:A9-inv}
$\neg\mathrm{Irrational}(S) \implies \exists h{:}\mathbb N,\ 0<h \wedge \exists N_0,\ \forall N\ge N_0,\ \mathtt{totientTail}(N{+}h)-\mathtt{totientTail}(N)\in\mathrm{range}(\mathbb Z\to\mathbb R)$,
with explicit witnesses $h=\varphi(\mathrm{oddPart}(r.\mathrm{den}))$, $N_0=v_2(r.\mathrm{den})$
for $S=r$. This is the ``only if'' half of the master equivalence: any rationality-refutation
strategy must eventually contradict \emph{some} period $h$ and preperiod $N_0$ produced this way.
Composes with Prop.~\ref{prop:AR-06-inv}/soundness to give Prop.~\ref{prop:A10}.
\scale{uniform}\ \ev{Lean}\ \coord{other:euler-theorem}\\
\lean{eventual_period_of_not_irrational}{TotientTailPeriodKiller.lean:358}
\end{prop}

\begin{prop}[Rationality flattens the whole LCM cone]\label{prop:B6-inv}
$\neg\mathrm{Irrational}(S) \implies \exists t_1,\ \forall t\ge t_1,\ \forall q\,m{:}\mathbb N,\ 0<q,\ \mathtt{totientTail}(q\cdot\mathtt{periodLcm}\,t + m\cdot\mathtt{periodLcm}\,t)-\mathtt{totientTail}(q\cdot\mathtt{periodLcm}\,t)\in\mathrm{range}(\mathbb Z\to\mathbb R)$.
Rationality forces ONE fractional constant on the entire cone $\{k\cdot\mathtt{periodLcm}\,t : k\ge1\}$
at every scale $t\ge t_1$, not just the diagonal pair — a strict generalisation of
Prop.~\ref{prop:A9-inv}. Any certificate anywhere on the cone therefore kills \#249
(Prop.~\ref{prop:B7}).
\scale{n/a}\ \ev{Lean}\ \coord{mobius-mersenne}\\
\lean{rational_totient_series_forces_lcm_cone_flatness}{CertificateKernel.lean:19002}
\end{prop}

\begin{prop}[Unconditional $(2^e{+}1)$-dimensional dyadic kernel — the rank floor]\label{prop:D4-inv}
The canonical dyadic-kernel family $\mathtt{canonicalTotientKernelFamily}\ e$ has exactly
$2^e+1$ channels, unconditionally linearly independent for every $e$ (CRT + Dirichlet: one channel
made prime, every other channel gets a fresh prime $\equiv1\bmod$ a large power of $2$). This is a
proved fact, needing no rationality hypothesis at all, and it is the fixed lower bound that
Prop.~\ref{prop:D5cons} would have to contradict.
\scale{n/a}\ \ev{Lean}\ \coord{other:carry-kernel-rank}\\
\lean{card_totientCanonicalIndex}{TotientMahlerDefect.lean:61}
\end{prop}

\begin{prop}[Sharp Farey rung — concrete necessary floor at $K=240$]\label{prop:C2-inv}
For all $q{:}\mathbb N$, $0<q\le79639646646701375323355774875831053$ ($\sim7.96\times10^{34}$):
$(qV)\bmod2^{240} + 243q < 2^{240}$, where $V$ is the explicit committed totient residue for
window $(N,K)=(1,240)$; this is SHARP — $q=79639646646701375323355774875831054$ is the exact
first failing denominator (an explicit Farey mediant). Any candidate rational value for $S$ must
have reduced denominator strictly larger than this bound (Prop.~\ref{prop:C3-inv} below).
\scale{fixed}\ \ev{Cert}\ \coord{other:farey-gap}\\
\lean{gap_check_window_1_240_first_failure}{GapFareyBound.lean:225}
\end{prop}

\begin{prop}[Denominator exclusion — headline unconditional consequence]\label{prop:C3-inv}
$\forall p{:}\mathbb Q,\ p.\mathrm{den}\le79639646646701375323355774875831053 \implies S\ne(p:\mathbb R)$.
If $S$ is rational, its reduced denominator exceeds $\sim7.96\times10^{34}$ — an unconditional
necessary condition on any hypothetical rational value, logically independent of the
certificate-supply wall (Prop.~\ref{prop:A10}); it excludes small denominators outright, without
needing a period at all.
\scale{fixed}\ \ev{Lean}\ \coord{other:farey-gap}\\
\lean{tsum_totient_div_pow_two_ne_ratCast_of_den_le_79639646646701375323355774875831053}{CertificateKernel.lean:18384}
\end{prop}

\begin{prop}[The two foundational irrationality engines]\label{prop:D1D2-inv}
Two, and only two, general-purpose irrationality criteria exist in the kernel, and any successful
proof of \#249 must ultimately instantiate one of them: (i) if $u{:}\mathbb N\to\mathbb Q$ is
eventually never $x$ and $\mathrm{den}(u_k)\cdot|x-u_k|\to0$, then $\mathrm{Irrational}(x)$; (ii)
if $\forall q>0,\ \exists m,z{:}\mathbb Z,\ 0<|m\xi-z|<1/q$, then $\mathrm{Irrational}(\xi)$
(with a base-power specialisation $b^n\xi$ matching any digit/carry construction in base $b$).
Every certificate, cone, or Farey argument in this catalogue exists either to supply one of these
two hypotheses or to substitute for them via the $\mathtt{certifiedKill}$ route.
\scale{n/a}\ \ev{Lean}\ \coord{other:dirichlet-approximation}\\
\lean{irrational_of_den_mul_abs_sub_tendsto_zero}{CertificateKernel.lean:5371}
\end{prop}

\begin{prop}[S sits on the Mersenne-Lambert ladder]\label{prop:D7-inv}
Writing $L(f):=\sum_{n\ge1}f(n)/(2^n-1)$: $L(\mu)=1/2$, $L(\varphi)=2$ (exactly rational,
machine-checked), $L(1)=E$ = the Erd\H os-Borwein constant, proved \textbf{irrational} in this
same kernel; $L(A)=S$ for $A=\varphi*\mu$ is exactly \#249 restated in ``positive
Erd\H os-Borwein form,'' still open; $L(\mathrm{Id})=\sum\sigma(m)/2^m$ is transcendental by
Nesterenko 1996 (cited, not formalised). Any proof of \#249 sits on this Dirichlet-convolution
ladder next to a proved irrational neighbour ($L(1)$) and a proved rational neighbour
($L(\varphi)$); this is the clearest bridge into \#257's native Lambert-series coordinate.
\scale{n/a}\ \ev{Lean, Cited}\ \coord{mobius-mersenne}\\
\lean{tsum_totient_div_pow_two_eq_pnat_half_pow}{CertificateKernel.lean:18415}
\end{prop}

\begin{prop}[Generic rational-gap adapter]\label{prop:D9-inv}
For $\mathtt{pfx}<\mathtt{whole}{:}\mathbb Q$: $1/(\mathtt{whole.den}\cdot\mathtt{pfx.den}) \le \mathtt{whole}-\mathtt{pfx}$
(as reals). Any analytic upper bound on a totient-tail gap converts, via this lemma, into a
denominator-growth lower bound extending Prop.~\ref{prop:C2-inv}'s Farey rungs without redoing the
Farey-neighbour search from scratch — a clean, fully general adapter usable on either problem's
tail estimate.
\scale{n/a}\ \ev{Lean}\ \coord{other:farey-gap}\\
\lean{positive_rational_difference_lower_bound}{PrimitiveRationalGapSupply.lean:31}
\end{prop}

\begin{prop}[Finite pincer floor — contiguous evidence through $t\le82$]\label{prop:B11-inv}
$\mathtt{certifiedKill\_diagonal\_all\_imported}$: Prop.~\ref{prop:B3}'s predicate $P\,t$ holds
for $t\in\{1,2,3,4,5,7,8,9,11,13,16,17,\dots\}$ (28 explicit cases through $t=64$), each a finite
$\mathtt{decide}/\mathtt{norm\_num}$ computation on explicit $\varphi$ values via checked
prime-power blocks with Lucas-primality certificates.  The current aggregate
theorem strengthens this historical list to $\forall t\le82,\ P\,t$, with no
holes.  It does \emph{not} establish $P\,83$ or infinitely many $t$ and hence
does not close Props.~\ref{prop:A10}/\ref{prop:B3}/\ref{prop:B7}.
\scale{bounded}\ \ev{Cert}\ \coord{mobius-mersenne}\\
\lean{certifiedKill_diagonal_all_imported}{DiagonalPincerCertificates.lean:2870}
\lean{ErdosProblems.Skip.LadderT67.exists_diagonalKill_le_82}{ErdosProblems/Skip/LadderT67.lean:71264}
\end{prop}

\begin{prop}[The fair-coin coprimality reformulation of $S$]\label{prop:C1-inv}
$\#\{(a,b): a{+}b=n,\ a>0,\ \gcd(a,b)=1\}=\varphi(n)$ for every $n$ (uniform, $\varphi(0)=0$,
$\varphi(1)=1$ included automatically). At $r=1/2$, for independent fair-coin waiting times
$X,Y$ ($P(X{=}n)=2^{-n}$, $n\ge1$): $P(\gcd(X,Y)=1) = S - 1/2$. \#249's constant is, up to the
additive constant $1/2$, a literal lattice-point probability. Any digit argument for \#249
necessarily also carries a positional meaning about visible lattice points on the addition
antidiagonal — a reformulation any proof attempt should keep in view even though it is not itself
a route to a proof.
\scale{n/a}\ \ev{Lean}\ \coord{other:visible-lattice-points}\\
\lean{tsum_pos_coprime_pair_pow}{GeometricCoprimality.lean:182}
\end{prop}

\subsection{Obstructions and countermodels}

Adversarial constructions and no-go results, kept with the mechanism that
closes each route so the machinery outlives the framing. Per the corpus's
own doctrine (project memory \texttt{feedback\_invent\_dont\_audit}): obstructions
below are coordinate-relative, and each entry names exactly which coordinate
or proof-template it excludes, not the underlying object.

\begin{obs}[Nonnegative-branch elimination for the true survivor]\label{prop:SGN-02}
If the translated actual orbit is integral at the start of Prop.~\ref{prop:SGN-01}'s positive
corridor, then for every later depth $K$ with room $J{+}K{+}(a{+}6)<2\cdot2^a$, the negated
carry-orbit trajectory $-\mathtt{carryOrbit}\ H\ (H{+}J)\ d\ K$ is negative and is an
$\mathtt{endpointSurvivor}$ at every depth in the corridor. This eliminates only the nonnegative
branch of the true survivor's sign — it is emphatically \emph{not} a nonintegrality claim;
spurious survivors of either sign may still exist. Any strategy hoping to conclude non-integrality
purely from this sign fact fails: Prop.~\ref{prop:SGN-03} names exactly what remains.
\scale{bounded}\ \ev{Lean}\ \coord{mobius-mersenne}\\
\lean{actualLcm_trueEndpointSurvivor_neg}{TotientActualLcmOrbitSign.lean:172}
\end{obs}

\begin{obs}[Total dyadic staircase is impossible — route-pruning]\label{prop:TE-01-killer}
For $a\ge8$, room bound, and modulus wide enough ($2H{+}J{+}K{+}2<2^m$):
$\mathtt{ActualLcmTerminalDyadicStaircase}\ a\ J\ K\ m$ is \textbf{false} — a terminal window
where every one of the last $m$ letters is divisible by its own growing power of two
($2^{r+1}\mid\text{letter}_r$) cannot occur. The last letter would have to be positive
(Prop.~\ref{prop:AR-05-inv}), strictly below the wide modulus, and divisible by it, forcing it to
be exactly $0$ — contradicting positivity. \textbf{Kills outright}: any strategy aiming for total
dyadic annihilation of the terminal suffix. The generic mechanism
($0<e<2^m\wedge2^m\mid e\implies e=0$) is reusable wherever a positive quantity is asked to vanish
mod a wider-than-itself modulus. The surviving route is the \emph{punctured} staircase
(Prop.~\ref{prop:TE-02-inv}) or the residue-gap producers (Prop.~\ref{prop:TE-04}).
\scale{bounded}\ \ev{Lean}\ \coord{mobius-mersenne}\\
\lean{not_actualLcmTerminalDyadicStaircase_of_room}{TotientActualLcmTopEdgeStaircase.lean:251}
\end{obs}

\begin{obs}[The bare 3-rank kernel is exactly tight — a sharpness no-go]\label{prop:FR-03-kill}
$\mathtt{fixedRankSecondDifference}(2^{n+1}\cdot2)(2^{n+1}\cdot3) = -2^{n+1}$ exactly, for all $n$
(a closed-form computed fixture). The two-adic valuation gained by Prop.~\ref{prop:FR-02-inv} is
exactly tight — the normalised $1{\times}1$ minor is odd at every depth, so the bounded-height
primitive kernel alone can \textbf{never} force one additional factor of $2$. \textbf{Kills}: any
attempt to squeeze more $2$-adic information out of the bare 3-rank curvature kernel; further
progress needs new arithmetic input, matching the file's own open target
$\mathtt{PowerTwoLcmMiddleRankExtremalSupply}$.
\scale{uniform}\ \ev{Cert}\ \coord{other:fixed-rank-curvature}\\
\lean{fixedRankSecondDifference_dyadic_fixture_exact}{TotientFixedRankLcmAsymptotic.lean:462}
\end{obs}

\begin{obs}[The parity coboundary countermodel object]\label{prop:CM-01}
$\mathtt{parityCoboundaryWeight}\ n := \mathtt{parityBaseWeight}\ n + 2\cdot\mathtt{largePowerTwoBit}\ n - 4\cdot\mathtt{largePowerTwoBit}(n{-}1)$,
where $\mathtt{parityBaseWeight}$ is the eventually-constant word $0,1,1,2,4,4,4,\dots$ and
$\mathtt{largePowerTwoBit}\ n=1$ iff $n=2^{k+3}$ (lacunary spikes starting at $8$). This is a
hand-built adversary sequence containing \emph{no} $\varphi$ in its definition at all except via
its parity. The construction technique — add a zero-valued sparse binary coboundary
$2\cdot2^{-m}-4\cdot2^{-(m+1)}=0$ at lacunary ranks to an eventually-periodic base word, destroying
periodicity while preserving the rational sum — is a completely general recipe, directly portable
to an analogous \#257 adversary matched to $1/(2^n-1)$-parity patterns.
\scale{n/a}\ \ev{Lean}\ \coord{other:binary-digit}\\
\lean{parityCoboundaryWeight}{TotientParityCoboundaryCountermodel.lean:59}
\end{obs}

\begin{obs}[What the countermodel refutes — the flagship kill of the batch]\label{prop:CM-02}
There is $c{:}\mathbb N\to\mathbb N$ with: (a) arbitrarily many, arbitrarily-separated, cofinal
explicit ``$6,0$'' carry pairs (for every $N,G,K$ a block of $K$ such pairs beyond $N$, pairwise
separated by $>G$); (b) $\forall n,\ c(n)\le6$; (c) $\forall n,\ c(n)\le n$; (d) $\forall n,\ c(n)\equiv\varphi(n)\pmod2$
(exact parity agreement with Euler's totient); (e) $c$ is \textbf{not} eventually periodic; yet
(f) $\sum_n' c(n)/2^n = 3/2$, a \emph{rational} number.

\par\medskip

\textbf{Kills}: any \#249 proof strategy
whose only inputs are uniform boundedness, the trivial growth bound $c(n)\le n$, exact
$\varphi$-parity agreement, and failure of eventual periodicity — \emph{even strengthened} to
cofinal, arbitrarily-separated, arbitrarily-long-block non-periodicity. The obstruction is
structural, not a shortage of non-periodicity witnesses: a rational coboundary can carry unlimited
\emph{visible} aperiodic structure while summing to a fixed rational.

\par\medskip

\textbf{Doctrine consequence}
(matches project memory \texttt{feedback\_erdos\_reductions\_rejected\_bank\_real\_results}): only
arguments using actual quantitative totient/Mersenne size or residue information (as in the
$\mathtt{TotientActualLcm}\ast$/$\mathtt{TotientFixedRank}\ast$ families) can possibly close \#249;
pure symbolic-word arguments cannot. This countermodel is also the ready-made stress-test fixture
for any future sufficient-condition candidate stated purely in coefficient-word terms
(boundedness/growth/parity/periodicity), on either \#249 or \#257.
\scale{cofinal}\ \ev{Cert}\ \coord{other:binary-digit}\\
\lean{exists_totientParity_arbitrarilyManySeparatedCarry_rational_countermodel}{TotientParityCoboundaryCountermodel.lean:637}
\end{obs}

\begin{obs}[Primitive Mersenne-prime factors alone force nothing]\label{prop:CP-03-kill}
For the completely-multiplicative control $c(n)=n$ (zero totient content:
$\mathtt{binaryCoeffTail}\ \mathrm{id}\ N = N{+}2$): if $q>K>0$ and $q\mid2^K-1$ (e.g. $q$ a
primitive prime factor of the homogeneous Mersenne multiplier), the shift is integral at every $N$
(value exactly $K$) but $q\nmid K$. \textbf{Kills}: ``a large primitive Mersenne prime factor
alone forces a contradiction'' as a proof strategy — this is a route-pruning warning applicable
\emph{verbatim} to \#257, whose denominators $2^n-1$ are literally the object $q\mid2^K-1$ tested
here. Any successful use of such a factor needs genuinely totient-specific residue/size input
beyond bare Mersenne primitivity — contrast with the genuinely totient-specific pulse of
Prop.~\ref{prop:CP-05-inv}.
\scale{bounded}\ \ev{Lean}\ \coord{mobius-mersenne}\\
\lean{idCoeff_mersenneModulus_does_not_annihilate_tailShift}{TotientTailCarryPeriod.lean:342}
\end{obs}

\begin{obs}[Universal no-go for absolute-adjugate coefficient reconstruction]\label{prop:CP-04-kill}
For any finite rational row $w{:}\iota\to\mathbb Q$, integer targets $x{:}\iota\to\mathbb N$ that
\emph{exactly} isolate one totient channel ($\sum_i w_i\varphi(x_i)=1$): the crude two-tail cost
$\sum_i|w_i|\cdot(2(x_i{+}1)+(x_i{+}2))$ is $\ge3$, hence never $<1$. \textbf{Kills}: the strategy
of recovering $\varphi(x)$ from $2R_{x-1}-R_x$ and bounding each tail termwise by $R_M\le M{+}2$ —
it can \emph{never} reach the strict $<1$ tail-error threshold, at any finite grid height,
independent of matrix height, row translation, determinant, or target channel. The proof only uses
$\varphi(x)\le x$, so the same triangle-inequality floor transfers verbatim to any \#257
reformulation with $c(n)\le n$.
\scale{uniform}\ \ev{Lean}\ \coord{other:adjugate-linear-algebra}\\
\lean{three_le_totientAdjugateTailCost}{TotientTailCarryPeriod.lean:266}
\end{obs}

\begin{obs}[Rank-2 second-difference certificates: sound but measured not shallower]\label{prop:B9cert-kill}
$\mathtt{certifiedRank2Kill}\ h\ N\ L \implies$ the second difference
$(\mathtt{totientTail}(N{+}2h){-}\mathtt{totientTail}(N{+}h)) - (\mathtt{totientTail}(N{+}h){-}\mathtt{totientTail}\,N) \notin\mathrm{range}(\mathbb Z\to\mathbb R)$,
sound at doubled band radius versus rank~1. Measured at $(h,N)=(1,8)$: rank-1 fires at depth~8, no
rank-2 certificate exists at depth $\le8$, rank-2 first fires at depth~9 — empirically \textbf{not}
a shortcut over rank-1 (rank-1 at least as shallow in $30/40$ probed cells for $t\le20$).
\textbf{Kills}: expecting a depth improvement from moving to second differences. Explicitly flagged
as a dead end — do not re-attempt this exact refinement without new information.
\scale{n/a}\ \ev{Cert}\ \coord{mobius-mersenne}\\
\lean{totient_tail_rank_two_kill_sound_but_not_shallower_cell}{CertificateKernel.lean:19090}
\end{obs}

\begin{obs}[Unit-gap strengthening rescues at most one lattice point]\label{prop:D8cert-kill}
$\mathtt{ReducedDenominatorUnitGapCert}\ u\ N\ K := \forall t\in[1,L],\ \neg\mathrm{Coprime}(J{+}t,u)$
(nonunits-only refinement of the Farey-gap consumer). At a prime-power reduced denominator $p^e$:
$\mathtt{ReducedDenominatorUnitGapCert}(p^e)\ N\ K \iff L{=}0 \vee (L{=}1\wedge p\mid(J{+}1))$ —
the unit-gap strengthening can rescue \textbf{at most one} additional candidate lattice point
beyond the ordinary Farey-gap certificate. \textbf{Kills}: expecting this refinement to unboundedly
strengthen Prop.~\ref{prop:C2-inv}/\ref{prop:C3-inv}'s Farey rungs — it is self-flagged in its own
docstring as ``a strict certificate-level strengthening, not a supply theorem.'' A known dead end;
do not re-attempt expecting more than $+1$ lattice point per prime power.
\scale{bounded}\ \ev{Lean}\ \coord{other:farey-gap}\\
\lean{reducedDenominatorCandidateCount_prime_pow_le_one}{PrimitiveWeightCertificate.lean:111}
\end{obs}

\begin{obs}[Synthetic all-horizon countermodel to homogeneous-factor-only proof strategies]\label{prop:B4b-kill}
There is an explicit nonzero ``all-horizon countermodel'' sequence built from multiples of
$\varphi(\mathtt{periodLcm}\,t)$ that (i) agrees with every exact whole-ray anchor
$\mathtt{deltaTotient}\ H\ (qH)=\varphi(H)$ for $2\le q<t$; (ii) stays inside the natural diagonal
bounds; and (iii) \textbf{survives every finite integer shift polynomial} — every commensurate
finite-rank LCM cube, via the normal form $P_m(E_H)\cdot\prod(E_H^n-1)$. Explicitly flagged
synthetic: it does not claim the compensation letters occur as actual totient differences.
\textbf{Kills}: the strategy of ``retain only homogeneous LCM-ray factors'' at \emph{every} finite
rank, not just rank 2 or 3. This is a representation-level exclusion (a particular factor-ideal
projection throws away information that can be adversarially reconstructed), not an exclusion of
$\varphi$ itself — an actual proof must control the fresh Möbius channel, per
$\mathtt{totient\_eq\_sum\_mobiusTotientChannel}$ in the same file.
\scale{n/a}\ \ev{Lean (docstring-sourced)}\ \coord{mobius-mersenne}\\
\lean{LcmFactorIdealPulseObstruction.lean}{LcmFactorIdealPulseObstruction.lean:1-24}
and \lean{totient\_eq\_sum\_mobiusTotientChannel}{LcmFactorIdealPulseObstruction.lean:328}
\end{obs}

\begin{obs}[Fixed-precision local valuation-unit signatures never obstruct]\label{prop:B5-kill}
For any finite word of odd-valuation-unit symbols at fixed local $2$-adic precision $u>0$, and any
starting carry state $e$, there \textbf{exists} a compatible carry orbit realising that exact word
with every intermediate state centred in its dyadic interval ($|e'|\le\mathtt{vuRadius}\ u\ \sigma$).
\textbf{Kills}: any proof strategy trying to derive a contradiction purely from ``the local
valuation-unit signature at fixed precision $u$ is incompatible with $X$'' — such signatures are
\emph{always} realisable by some carry orbit. Framed against \#249 in the docstring, but the
theorem itself carries no totient- or Mersenne-specific content: a proof needs growing precision
or extra arithmetic coupling, not a fixed-window local signature.
\scale{n/a}\ \ev{Lean}\ \coord{p-adic}\\
\lean{fixedPrecisionTropicalNoGo}{TropicalCurvatureCarry.lean:137}
\end{obs}

\begin{obs}[Square-CRT correction-suppression is independent of nonvanishing]\label{prop:B8-kill}
Square-CRT correction suppression (fixing a cofactor mod a prime square to remove a weighted
correction term on a finite horizon) is achievable, \textbf{but} suppression alone does not force
a nonzero coefficient: the smallest returned countermodel has the whole two-step finite block
vanish identically, while a separate clean witness shows it can also be nonzero — ``clean''
(correction-suppressed) is consistent with \emph{both} vanishing and nonvanishing. \textbf{Kills}:
``achieve a clean square-CRT horizon'' as a sufficient condition by itself for a nonvanishing
correction term; an additional anti-concentration or residue producer is still needed, exactly as
the module's own docstring states.
\scale{n/a}\ \ev{Lean (docstring-sourced)}\ \coord{other:square-crt-cocycle}\\
\lean{squareCRT_clean_block_can_vanish}{SquareCRTCube.lean:459}
\end{obs}

\begin{obs}[No fixed integer clears every normalised primitive-Euler coordinate]\label{prop:D6-kill}
$\forall D>0,\ \neg\forall n>0,\ \exists z{:}\mathbb Z,\ D\cdot(A(n)/n)=z$, where $A=\varphi*\mu$ is
the Mersenne-Lambert primitive weight (so $S=L(A)$, Prop.~\ref{prop:D7-inv}). Any $D$ clearing
normalised coordinates through horizon $N$ must be divisible by every odd prime $p\le N$, every
$p^2\le N$, and by $4$ once $N\ge4$ — so no fixed $D$ works for all $n$. \textbf{Kills}: any
strategy seeking a single fixed integer denominator that simultaneously integralises every
primitive-Euler coordinate $A(n)/n$ — a finite no-lift theorem in the integral
Euler/Witt-coordinate category (weaker than the general Dieudonn\'e-Dwork theorem, and explicitly
\emph{not} itself an irrationality proof). An independent-coordinate no-go, orthogonal to the
binary-window (Prop.~\ref{prop:A10}) and carry-rank (Prop.~\ref{prop:D5cons}) obstructions — a
third distinct coordinate attacking the same open problem.
\scale{n/a}\ \ev{Lean}\ \coord{mobius-mersenne}\\
\lean{no_fixed_integral_primitive_euler_index}{PrimitiveDeterminantLift.lean:169}
\end{obs}

\begin{obs}[Finite certificates only ever prove exclusion, never membership]\label{prop:B11cert-kill}
Certified-death/kill families (finite decidable checks at a given depth) are explicitly
\textbf{one-sided}: a found certificate proves exclusion, but failure to find one, or survival
through any finite probed depth, proves \emph{nothing} about membership or rationality. This is a
structural caveat recurring across every certificate family in this catalogue
($\mathtt{certifiedKill}$, $\mathtt{directedCertifiedKill}$, $\mathtt{survivorKill}$,
$\mathtt{LcmDiagonalArithmeticKill}$): none of them can ever certify membership/rationality by
finite search, only non-membership/non-integrality. \textbf{Kills}: treating an unsuccessful finite
search over any of the certificate families above as evidence toward rationality, or treating a
finite prefix of confirmed kills (Prop.~\ref{prop:SK-01-inv}, Prop.~\ref{prop:B11-inv}) as
anything more than a floor to extend.
\scale{n/a}\ \ev{Lean (structural caveat, cross-problem)}\ \coord{n/a}\\
\lean{CertifiedGreedyMersenneDeath}{GreedyAchievementSet.lean:1756}
\end{obs}
% ---- part a249_newdecls ----
\subsection{Recently added declarations}

This subsection catalogues material that entered the public record after the
rest of this paper was drafted: declarations union-merged into modules that
were already public (\texttt{SignedQMomentObstruction},
\texttt{GenericTailOrbitRigidity}, \texttt{DiagonalFreshLossBridge},
\texttt{SquaredMersenneDiagonalEnclosure},
\texttt{ActualForeignResidueProjection}, \texttt{DyadicPrefixCompression},
\texttt{RepunitMobiusNumerator}, \texttt{GeometricCoprimality},
\texttt{LcmDiagonalReduction}, \texttt{PivotAntiReconstruction}, all in
\texttt{Erdos249257/}), and the newly published per-problem tree
\texttt{ErdosProblems/Erdos249/}. Every declaration below was read directly
from source; none is inferred from a name. Several of these files contain
long chains of ``\texttt{XSupply} $\to$ irrational'' implications whose
hypothesis is an unproved cofinal predicate. Per the standing rule of this
paper, those chains are named as exactly what they are --- reductions, not
results --- and are not restated as if they narrowed the open problem. The
module \texttt{SquareCRTCube} is deliberately excluded: its private additions
did not compile even in the private tree and were reverted, so nothing from
it beyond the existing public file is cited here.

\subsubsection{The Möbius--Mersenne ladder: unconditional log-concavity}

\texttt{SignedQMomentObstruction} adds an infinite ladder
$\Theta_r := \sum_{n\ge 0} \mu(n+1)/(2^{n+1}-1)^r$ built from the same
Möbius--Mersenne atoms as the rest of the corpus, together with a full
unconditional order-two Hankel (log-concavity) theorem for every rung.

\begin{defn}
$\Theta_r := \sum_{n\ge 0}\mu(n+1)/(2^{n+1}-1)^r$, split exactly into the
first two atoms $\mathrm{TwoAtom}(r) := 1 - 3^{-r}$ and the tail
$\mathrm{TailAfterTwo}(r) := \sum_{n\ge 0}\mu(n+3)/(2^{n+3}-1)^r$.
\lean{Erdos249257.SignedQMomentObstruction.mobiusMersenneTheta}{SignedQMomentObstruction.lean:136}
\quad \lean{Erdos249257.SignedQMomentObstruction.mobiusMersenneTwoAtom}{SignedQMomentObstruction.lean:140}
\quad \lean{Erdos249257.SignedQMomentObstruction.mobiusMersenneTailAfterTwo}{SignedQMomentObstruction.lean:144}
\quad \ev{Lean} \quad \scale{uniform} \quad \coord{mobius-mersenne}.
\end{defn}

\begin{thm}[Two low rungs are exact rationals tied to \#249]
$\Theta_1 = 1/2$, and $\Theta_2 = \bigl(\sum_{n\ge 1}\varphi(n)2^{-n}\bigr) - 1/2$
--- the second rung is literally $S - 1/2$.
\lean{Erdos249257.SignedQMomentObstruction.mobiusMersenneTheta_one}{SignedQMomentObstruction.lean:197}
\quad \lean{Erdos249257.SignedQMomentObstruction.mobiusMersenneTheta_two_eq_totient_offset}{SignedQMomentObstruction.lean:214}
\quad \ev{Lean} \quad \scale{fixed} \quad \coord{mobius-mersenne}.
\end{thm}

\begin{thm}[Two-atom exact Hankel gap]
$\mathrm{TwoAtom}(r{+}1)^2 - \mathrm{TwoAtom}(r)\mathrm{TwoAtom}(r{+}2) = 4/3^{r+2}$,
hence the two-atom truncation alone is strictly log-concave at every $r$.
\lean{Erdos249257.SignedQMomentObstruction.mobiusMersenneTwoAtom_hankelGap}{SignedQMomentObstruction.lean:311}
\quad \ev{Lean} \quad \scale{uniform} \quad \coord{mobius-mersenne}.
\end{thm}

\begin{thm}[Full ladder: unconditional strict log-concavity, every rung]
For every $r\ge 1$, $\Theta_r\,\Theta_{r+2} < \Theta_{r+1}^2$: the shifted
$2\times 2$ Hankel determinant of the whole ladder is negative at every
rung, not just asymptotically. The proof propagates the exact two-atom gap
above against a geometric tail-error budget that contracts by a factor $1/4$
per rung against a gap that only contracts by $1/3$, giving an inductive
floor from rung $5$ on (kernel-checked base cases below).
\lean{Erdos249257.SignedQMomentObstruction.mobiusMersenneTheta_strict_logConcave}{SignedQMomentObstruction.lean:730}
\quad \lean{Erdos249257.SignedQMomentObstruction.mobiusMersenneTheta_strict_logConcave_of_five_le}{SignedQMomentObstruction.lean:534}
\quad \lean{Erdos249257.SignedQMomentObstruction.mobiusMersenneTheta_hankel_two_neg}{SignedQMomentObstruction.lean:740}
\quad \ev{Lean} \quad \scale{uniform} \quad \coord{mobius-mersenne}.
\end{thm}

\begin{obs}
This is a genuine unconditional analytic fact about the ladder, but it is
not by itself an irrationality obstruction or a step toward one: negative
Hankel determinants say the sequence $(\Theta_r)$ is not the moment
sequence of a positive measure, which is unrelated to whether any single
$\Theta_r$ (in particular $\Theta_2 = S - 1/2$) is rational. It is recorded
here as exactly what it is --- a structural fact about the ladder --- and
no stronger claim is made.
\end{obs}

The file also supplies the underlying algebraic infrastructure: a
rectangular Cauchy--Binet expansion derived from the Leibniz formula
(\lean{Erdos249257.SignedQMomentObstruction.det_mul_rectangular}{SignedQMomentObstruction.lean:29},
\ev{Lean}, \scale{uniform}, \coord{other:hankel-determinant}, needed because
Mathlib's pinned determinant API does not expose rectangular Cauchy--Binet
directly), and a finite unique-terminal-dyadic-exponent parity lemma showing
that clearing a common denominator by its uniquely largest power of two
preserves oddness of the numerator
(\lean{Erdos249257.SignedQMomentObstruction.scaled_dyadic_sum_odd}{SignedQMomentObstruction.lean:78},
\ev{Lean}, \scale{uniform}, \coord{other:dyadic-parity}).

\subsubsection{Generic tail-orbit rigidity: a self-labelled non-claim}

\texttt{GenericTailOrbitRigidity} is explicit in its own header that it
asserts no novelty and no priority, and contains an explicit \texttt{NON\_CLAIM}
guard against a superseded ``positive orbit'' route; it is a formal
algebra/analysis interface, reproduced here only because it entered the
public tree after the earlier parts of this paper were drafted.

\begin{thm}[T7: tempered-orbit equivalence, abstract binary series]
For any $c:\mathbb{N}\to\mathbb{N}$ with $c(n)\le n$, writing
$X_c := \sum_{n\ge 1} c(n)/2^n$, $X_c$ is rational iff there exist $v>0$ and
an integer sequence $u$ with $u(N{+}1) = 2u(N) - v\,c(N{+}1)$ and
$u(N)/2^N \to 0$; when it exists such an orbit is rigid, $u(N) = v\cdot T_c(N)$
for every $N$, where $T_c$ is the scaled tail.
\lean{Erdos249257.binaryCoeffSeries_rational_iff_exists_temperedBinaryOrbit}{GenericTailOrbitRigidity.lean:426}
\quad \lean{Erdos249257.temperedBinaryOrbit_eq_scaledTail}{GenericTailOrbitRigidity.lean:339}
\quad \ev{Lean} \quad \scale{uniform} \quad \coord{other:tail-orbit-rigidity}.
\end{thm}

\begin{prop}[Rigidity engine]
A real sequence $d$ with $d(N{+}1)=2d(N)$ and $d(N)=o(2^N)$ is identically
zero.
\lean{Erdos249257.doublingOrbit_eq_zero_of_tempered}{GenericTailOrbitRigidity.lean:315}
\quad \ev{Lean} \quad \scale{uniform} \quad \coord{other:tail-orbit-rigidity}.
\end{prop}

\begin{obs}[Finite-state no-go for successor decoders]
For fixed $m\ge 2$, the ``balanced-pulse'' family
$c_{m,r}$ (mass $r$ moved from digit $m$ to digit $m{+}1$, $0\le r\le
\lfloor(m{+}1)/2\rfloor$) has the same value and the same complete
pre-$m$ history for every $r$, yet the parameter $r$ is recovered exactly
from the first post-pulse digit. Consequently no state that identifies all
members of one such family can be decoded by any autonomous map, and the
fan-out of any correct decoder is unbounded in $m$.
\lean{Erdos249257.balancedPulse_no_autonomous_decoder}{GenericTailOrbitRigidity.lean:247}
\quad \lean{Erdos249257.balancedPulse_label_card_lower_bound}{GenericTailOrbitRigidity.lean:228}
\quad \ev{Lean} \quad \scale{uniform} \quad \coord{other:tail-orbit-rigidity}.
This is a genuine barrier of the class the brief asks for: it rules out
exactly the class of arguments that try to predict/decode the exact tail
orbit from a state depending only on the pre-pulse history, because the
family exhibited is a literal counterexample generator for any such decoder.
It says nothing about the totient-specific series itself.
\end{obs}

Two ``rational control models'' are included as honesty checks, not
irrationality statements: the pair-balanced family
$c(2k){=}k{-}a(k)$, $c(2k{+}1){=}2a(k)$ has tempered orbit for every bounded
payload $a$ and a fixed marked value $4/9$ regardless of $a$
(\lean{Erdos249257.globalBalancedCoeff_value}{GenericTailOrbitRigidity.lean:545},
\ev{Lean}, \scale{uniform}, \coord{other:tail-orbit-rigidity}); and the
identity coefficients $c(n)=n$ have the explicit multiplier-one tempered
orbit $u(N)=N+2$
(\lean{Erdos249257.idCoeff_temperedOrbit}{GenericTailOrbitRigidity.lean:564},
\ev{Lean}, \scale{fixed}, \coord{other:tail-orbit-rigidity}).

\subsubsection{Squared-Mersenne diagonal enclosure}

\texttt{SquaredMersenneDiagonalEnclosure} spends the exactly-summable
first-order Lambert term of the actual \#249 diagonal exactly, leaving only
a squared-Mersenne tail to bound.

\begin{prop}[Exact rational centre, direct]
$\mathrm{scaleDiagonalTailDifference}(H) - \mathrm{lambertProjectedDiagonal}(H,D)
= C_H \cdot \mathrm{mobiusSquareTail}(D)$ for every $H,D$ --- no residue split
and no side condition on $D$ versus $H$.
\lean{Erdos249257.SquaredMersenneDiagonalEnclosure.scaleDiagonalTailDifference_sub_lambertProjectedDiagonal}{SquaredMersenneDiagonalEnclosure.lean:138}
\quad \ev{Lean} \quad \scale{uniform} \quad \coord{other:squared-mersenne-tail}.
\end{prop}

\begin{prop}[Sharp geometric tail bound]
$|\mathrm{mobiusSquareTail}(D)| \le 4/\bigl(3(2^{D+1}-1)^2\bigr)$.
\lean{Erdos249257.SquaredMersenneDiagonalEnclosure.abs_mobiusSquareTail_le}{SquaredMersenneDiagonalEnclosure.lean:167}
\quad \ev{Lean} \quad \scale{uniform} \quad \coord{other:squared-mersenne-tail}.
\end{prop}

\begin{obs}[Directed vs. symmetric enclosure]
When the first nonzero Möbius channel past $D$ is known, its sign directs a
one-sided interval half the width of the symmetric bound
(\lean{Erdos249257.SquaredMersenneDiagonalEnclosure.mobiusSquareTail_directed_of_first_support}{SquaredMersenneDiagonalEnclosure.lean:313},
\ev{Lean}, \scale{uniform}, \coord{other:squared-mersenne-tail}), and either
enclosure composes with an integer-lattice separation hypothesis into a
consumer theorem
(\lean{Erdos249257.SquaredMersenneDiagonalEnclosure.scaleFullTarget_miss_of_lambert_projected_separation}{SquaredMersenneDiagonalEnclosure.lean:426}).
The separation hypothesis itself is exactly the remaining open content; the
enclosure is real but does not supply it.
\end{obs}

\subsubsection{Old/foreign Möbius channel split (DiagonalFreshLossBridge)}

The 4243-line addition to \texttt{DiagonalFreshLossBridge} isolates, at the
level of individual Möbius-totient channels, which channels are ``old''
(index divides the LCM height $H$) and which are ``foreign'' (index does
not), and proves exact identities for both. Most of the file's bulk beyond
what is recorded here is a long ladder of \texttt{...Supply}-conditional
``$\Rightarrow$ irrational'' implications at successively more specialised
hypotheses; those are reductions and are not restated as results.

\begin{prop}[Exact old-channel value: a positive gcd word, scaled]
For $H>0$, the sum of totient forward-differences over exactly the divisor
channels of $H$ equals
$(H/\mathrm{rad}(H))\cdot\mathrm{gcdWordCoeff}(\mathrm{rad}(H),s)$, where
$\mathrm{gcdWordCoeff}$ is the positive gcd-word coefficient of
\S\ref{repunit-gcd-word} below.
\lean{Erdos249257.DiagonalFreshLossBridge.oldMobiusIncrement_eq_scale_gcdWordCoeff}{DiagonalFreshLossBridge.lean:653}
\quad \ev{Lean} \quad \scale{uniform} \quad \coord{other:gcd-word}.
\end{prop}

\begin{prop}[Exact old/foreign split of the diagonal height increment]
$\mathrm{diagonalHeightIncrement}(H,s) = \mathrm{oldMobiusIncrement}(H,s) +
\mathrm{finiteForeignChannelIncrement}(H,s)$, a literal finite-sum identity,
not an estimate.
\lean{Erdos249257.DiagonalFreshLossBridge.diagonalHeightIncrement_eq_oldMobius_add_finiteForeign}{DiagonalFreshLossBridge.lean:829}
\quad \ev{Lean} \quad \scale{uniform} \quad \coord{other:gcd-word}.
\end{prop}

\begin{thm}[Exact doubling law for the old-channel increment]
If $H$ and $r$ are both even, $\mathrm{diagonalHeightIncrement}(2H,2r) =
2\cdot\mathrm{diagonalHeightIncrement}(H,r)$; if $H$ is even and $r$ odd, the
same doubled height increment equals $\mathrm{diagonalHeightIncrement}(H,r)$
unchanged.
\lean{Erdos249257.DiagonalFreshLossBridge.diagonalHeightIncrement_two_mul_even}{DiagonalFreshLossBridge.lean:308}
\quad \lean{Erdos249257.DiagonalFreshLossBridge.diagonalHeightIncrement_two_mul_odd}{DiagonalFreshLossBridge.lean:323}
\quad \ev{Lean} \quad \scale{uniform} \quad \coord{other:gcd-word}.
\end{thm}

\begin{obs}[Every nonzero foreign phase term is squarefree-supported]
A ``foreign'' channel $d\nmid H$ contributes a nonzero phase term to the
literal increment only if $d$ is squarefree and $d$ divides exactly one of
the two window endpoints $2H{+}s$, $H{+}s$ (never both, since $d\nmid H$
forces the two endpoints incongruent mod $d$).
\lean{Erdos249257.DiagonalFreshLossBridge.squarefree_of_foreignChannelPhaseTerm_ne_zero}{DiagonalFreshLossBridge.lean:909}
\quad \lean{Erdos249257.DiagonalFreshLossBridge.foreign_endpoint_support_disjoint}{DiagonalFreshLossBridge.lean:852}
\quad \ev{Lean} \quad \scale{uniform} \quad \coord{other:gcd-word}.
This is a real structural narrowing --- it rules out non-squarefree indices
and simultaneous double support as sources of foreign contribution --- but
it narrows a finite bookkeeping decomposition, not the analytic separation
obligation itself.
\end{obs}

\begin{prop}[Low/top echo doubling]
If a foreign channel $d$ has a lower-endpoint hit at offset $s$, its value
there is $-\mu(d)\cdot\lfloor(H{+}s)/d\rfloor$, and its top-endpoint value at
the doubled offset $2s$ is exactly twice that quantity in absolute value
(with sign flipped): $\mathrm{foreignChannelPhaseTerm}(d,H,2s) =
2\bigl(\mu(d)\lfloor(H{+}s)/d\rfloor\bigr)$.
\lean{Erdos249257.DiagonalFreshLossBridge.foreignChannelPhaseTerm_low_double_echo}{DiagonalFreshLossBridge.lean:941}
\quad \ev{Lean} \quad \scale{uniform} \quad \coord{other:gcd-word}.
\end{prop}

\subsubsection{Repunit Möbius numerator is a positive gcd word (T1)}
\label{repunit-gcd-word}

\begin{thm}[T1]
For squarefree $r$, the signed repunit numerator polynomial
$\sum_{d\mid r}\mu(d)(r/d)(1+X^d+\cdots+X^{r-d})$ equals the ``gcd word''
whose coefficient at $X^k$ ($k<r$) is $(r/\gcd(r,k))\cdot\varphi(\gcd(r,k))$
--- strictly positive for $k<r$, zero for $k\ge r$. The squarefree
hypothesis is kept explicit and is not promoted to arbitrary $r$.
\lean{Erdos249257.RepunitMobiusNumerator.mobiusNumeratorPolynomial_eq_gcdWord}{RepunitMobiusNumerator.lean:205}
\quad \lean{Erdos249257.RepunitMobiusNumerator.gcdWordCoeff}{RepunitMobiusNumerator.lean:41}
\quad \ev{Lean} \quad \scale{uniform} \quad \coord{other:gcd-word}.
\end{thm}

\begin{cor}
Evaluation at $X=2$ recovers the integral Möbius--Mersenne numerator already
owned by \texttt{RadicalMobiusShadow}, on the same squarefree boundary.
\lean{Erdos249257.RepunitMobiusNumerator.mobiusNumeratorPolynomial_eval_two}{RepunitMobiusNumerator.lean:445}
\quad \ev{Lean} \quad \scale{uniform} \quad \coord{other:gcd-word}.
\end{cor}

\subsubsection{Dyadic prefix compression: exact greedy-carry arithmetic}

\texttt{DyadicPrefixCompression} adds roughly 3200 lines of exact-arithmetic
infrastructure for a greedy half-orbit carry construction (Mersenne-weight
achievement sets); the bulk is machinery for a specific producer hypothesis.
Two pieces of genuinely reusable exact arithmetic:

\begin{defn}
For a displayed residual $p/(2L)$, the integer excess numerator above the
next dyadic point $2^{-(n+1)}$ is
$\mathrm{nextDyadicExcessIntNumerator}(p,n,L) := 2^n p - L$, chosen so the
skipped-branch comparison is an exact Diophantine inequality rather than a
real-valued phase estimate; it obeys the exact doubling recurrence
$E(p,n{+}1,L) = 2E(p,n,L) + L$.
\lean{Erdos249257.nextDyadicExcessIntNumerator}{DyadicPrefixCompression.lean:156}
\quad \lean{Erdos249257.nextDyadicExcessIntNumerator_succ}{DyadicPrefixCompression.lean:161}
\quad \ev{Lean} \quad \scale{uniform} \quad \coord{other:dyadic-carry}.
\end{defn}

\begin{prop}[Exact geometric tail bound for the Mersenne-weight remainder]
For $n\ge 2$, $\mathrm{mersenneWeightRemainder}(n) \le (4/3)(1/8)^n$, and the
corresponding infinite tail from depth $m\ge 1$ is $\le (4/21)(1/8)^m$.
\lean{Erdos249257.mersenneWeightRemainder_le_four_thirds}{DyadicPrefixCompression.lean:1162}
\quad \lean{Erdos249257.mersenneWeightRemainderTail_le_four_twentyone}{DyadicPrefixCompression.lean:1179}
\quad \ev{Lean} \quad \scale{uniform} \quad \coord{other:dyadic-carry}.
\end{prop}

\subsubsection{Pivot-fibre anti-reconstruction: an exact energy identity}

\texttt{PivotAntiReconstruction} adds roughly 1800 lines building a
finite-fibre variance/energy machinery around the ``first-harmonic''
producer hypothesis, ending in another chain of \texttt{DTW...Supply}
$\Leftrightarrow$ irrationality equivalences (again reductions, flagged as
such). The exact algebraic core underneath is reusable:

\begin{prop}[Anchor defect is a squared complex distance]
$\mathrm{firstHarmonicAnchorDefect}(h,L,T) = \sum_{N\in T}\|
\mathrm{windowFirstExp}(h,N,L) - 1\|^2$, and it equals
$2|T| - 2\sum_{N\in T}\mathrm{windowFirstCos}(h,N,L)$.
\lean{Erdos249257.TotientTailPeriodKiller.firstHarmonicAnchorDefect_eq_sum_norm_sq}{PivotAntiReconstruction.lean:562}
\quad \lean{Erdos249257.TotientTailPeriodKiller.firstHarmonicAnchorDefect_eq}{PivotAntiReconstruction.lean:550}
\quad \ev{Lean} \quad \scale{uniform} \quad \coord{other:pivot-fiber}.
\end{prop}

\begin{lem}[Separated-pairs energy floor]
For a finite family $z:T\to\mathbb{C}$ and any set of pairs $P\subseteq
T\times T$ each separated by $\ge\delta$, $|P|\cdot\delta^2 \le
\sum_{i,j\in T}\|z_i-z_j\|^2$.
\lean{Erdos249257.TotientTailPeriodKiller.card_mul_sq_le_sum_pairwise_norm_sq_of_separatedPairs}{PivotAntiReconstruction.lean:720}
\quad \ev{Lean} \quad \scale{uniform} \quad \coord{other:pivot-fiber}.
\end{lem}

\subsubsection{Actual foreign-residue projection}

\texttt{ActualForeignResidueProjection} is explicit that it is ``the proof
consumer'' for a receipt whose analytic kernel identity remains a separate
Lean validation lane; it supplies the finite bridge, not the estimate.

\begin{prop}[Explicit shadow is exactly the divisor-channel sum]
For $H>0$, $\mathrm{scaleExplicitShadow}(H) = \sum_{d\mid H}
\mathrm{residueIncrement}(d,H)$.
\lean{Erdos249257.ActualForeignResidueProjection.scaleExplicitShadow_eq_sum_divisors_residueIncrement}{ActualForeignResidueProjection.lean:338}
\quad \ev{Lean} \quad \scale{uniform} \quad \coord{other:foreign-residue}.
\end{prop}

\begin{prop}[Complement-noncancellation consumer]
If a projection's error against the true foreign defect is controlled by
the closed geometric budget $\mathrm{foreignComplementBound}(H,D)$, and the
finite rational state $\mathrm{scaleExplicitShadow}(H) +
\mathrm{projectedForeignDefect}(H,D)$ is farther from every integer than
that budget, then $S\ne \mathrm{fullTargetHit}$ at scale $H$.
\lean{Erdos249257.ActualForeignResidueProjection.scaleFullTarget_miss_of_projected_separation}{ActualForeignResidueProjection.lean:414}
\quad \ev{Lean} \quad \scale{uniform} \quad \coord{other:foreign-residue}. The
control and separation hypotheses are exactly the two remaining open
obligations; the consumer is a real theorem, not a claim they hold.
\end{prop}

\subsubsection{Geometric coprimality: a lattice coordinate, and a classical identity beside it}

\texttt{GeometricCoprimality} relocates \#249 onto the visible-lattice-point
mass of coprime pairs. Beside that bridge it formalises the classical
gcd-layering identity for the Lambert-weighted coprime sum. That identity is
an elementary rational function of its parameter, so it cannot distinguish
rational from irrational inputs; and it carries a different summand weight
from $S$, so it is not a statement about \#249.

\begin{thm}[Visible-point count is the totient, uniformly]
For every $n\in\mathbb{N}$, $\#\{(a,b): a+b=n,\ 0<a,\ \gcd(a,b)=1\} =
\varphi(n)$, with no case split at $n=0,1$.
\lean{GeometricCoprimality.card_antidiagonal_filter_pos_coprime}{GeometricCoprimality.lean:56}
\quad \ev{Lean} \quad \scale{uniform} \quad \coord{other:coprime-lattice}.
\end{thm}

\begin{prop}[Coprime-pair lattice mass bridges]
For $0\le r<1$: $\sum_{(a,b)\ \mathrm{coprime},\,a\ge 1} r^{a+b} =
\sum_n \varphi(n)r^n$; over strictly positive coprime pairs the same sum is
$\sum_n\varphi(n)r^n - r$; and layering by $\gcd = g$ recovers the full
positive-quadrant mass $\sum_g[\text{layer }g] = (r/(1-r))^2$, which equals
$1$ exactly at $r=1/2$ (two independent fair-coin waiting times have a
finite gcd almost surely).
\lean{GeometricCoprimality.tsum_coprime_pair_pow_eq_tsum_totient_mul_pow}{GeometricCoprimality.lean:120}
\quad \lean{GeometricCoprimality.tsum_pos_coprime_pair_pow}{GeometricCoprimality.lean:182}
\quad \lean{GeometricCoprimality.tsum_gcd_layer_pos_coprime_pow}{GeometricCoprimality.lean:328}
\quad \ev{Lean} \quad \scale{uniform} \quad \coord{other:coprime-lattice}.
\end{prop}

\begin{thm}[The classical visible-point Lambert identity]
For every $0\le r<1$, $\sum_{(a,b)\ \mathrm{coprime},\, a,b\ge 1}
\dfrac{r^{a+b}}{1-r^{a+b}} = \Bigl(\dfrac{r}{1-r}\Bigr)^2$, an elementary
rational function of $r$, hence rational at every rational $r$ including
$r=1/2$. This is the classical visible-point identity, and the Lean
declaration is a formalisation of it rather than a new result: writing each
pair $(A,B)$ of positive integers uniquely as $g\cdot(a,b)$ with
$\gcd(a,b)=1$ converts the quadrant sum $\sum_{A,B\ge1}r^{A+B}=(r/(1-r))^2$
into the displayed sum over visible points, which is the $\gcd$-layering
already recorded in the bridges above. Its summand weight is the Lambert
weight $r^{a+b}/(1-r^{a+b})$, not the plain weight $r^{a+b}$ under which the
same index set sums to $\sum_n\varphi(n)r^n$; the two sums share an index set
and differ in weight, and only the second is $S$.
\lean{GeometricCoprimality.tsum_pos_coprime_lambert_eq_sq}{GeometricCoprimality.lean:407}
\quad \ev{Lean} \quad \scale{uniform} \quad \coord{other:coprime-lattice}.
\end{thm}

\begin{obs}
The consequence to draw is narrow, and it is about the Lambert-weighted sum
rather than about $S$. Coprime-pair restriction, exact Stern--Brocot-type
splitting by $\gcd$, and geometric cylinder decay of the summand under the
map $n\mapsto r^n/(1-r^n)$, used together and evaluated at a rational $r$,
cannot imply irrationality of the resulting sum, because the identical
construction is provably rational at \emph{every} rational $r$ in $[0,1)$ ---
the mechanism has zero sensitivity to whether $r$ itself is rational. This
closes one route through the coprime-lattice coordinate; it does not close
the coordinate, and it is not evidence about the plain-weight sum that
actually equals $S$. Any future argument built on this specific route must
use some further feature of $r=1/2$ beyond coprimality, gcd-layering, and
geometric decay.
\end{obs}

\subsubsection{The lcm-diagonal collapse: reduced to one \texorpdfstring{$\mathbb{N}$}{N}-indexed predicate}

\texttt{LcmDiagonalReduction} (module \texttt{TotientTailPeriodKiller},
``wave 23'') removes the second free parameter from the wave-22
one-parameter reduction, producing a single decidable $\mathbb{N}$-indexed
statement equivalent to \#249. The module's own header states plainly that
the resulting supply hypothesis ``is the open content of \#249 and is NOT
claimed''; the equivalence itself is a reduction, reported as such.

\begin{obs}[Diagonal collapse --- reduction, not a result]
$\mathrm{irrational}(S)$ follows from: for every $t_0$, some $t\ge t_0$
admits a certified kill at the diagonal point $(H_t,H_t)$, $H_t :=
\mathrm{lcm}(1,\dots,t)$. This is proved as an implication only; the
antecedent is the unproved open content.
\lean{Erdos249257.TotientTailPeriodKiller.irrational_totient_series_of_lcm_diagonal_certificate_supply}{LcmDiagonalReduction.lean:112}
\quad \ev{Lean} (implication proved; antecedent \ev{Open})
\quad \scale{cofinal} \quad \coord{other:lcm-diagonal}.
\end{obs}

\begin{thm}[Window structure of the lcm ray --- unconditional]
Below $2t$, the only indices $j$ that fail to divide $H_t$ are bare prime
powers exceeding $t$; every $j\le t$ divides $H_t$ outright. On a ``clean''
divisor $j\mid H_t$ (every prime of $j$ still divides $H_t/j$) the window
totient factors exactly:
$\varphi(qH_t+j) = \varphi(j)\cdot\varphi(q(H_t/j)+1)$.
\lean{Erdos249257.TotientTailPeriodKiller.eq_prime_pow_of_not_dvd_periodLcm}{LcmDiagonalReduction.lean:137}
\quad \lean{Erdos249257.TotientTailPeriodKiller.totient_periodLcm_ray_split}{LcmDiagonalReduction.lean:208}
\quad \ev{Lean} \quad \scale{uniform} \quad \coord{other:lcm-diagonal}.
\end{thm}

\begin{obs}[Diagonal deposits]
The kernel decides $P(t)$ (a certified diagonal kill) unconditionally for
every $1\le t\le 8$ (totient arguments stay $\le 130$), at tabulated depths.
\lean{Erdos249257.TotientTailPeriodKiller.certifiedKill_periodLcm_diagonal_upto_six}{LcmDiagonalReduction.lean:257}
\quad \lean{Erdos249257.TotientTailPeriodKiller.certifiedKill_periodLcm_diagonal_at_eight}{LcmDiagonalReduction.lean:266}
\quad \ev{Cert} \quad \scale{fixed} \quad \coord{other:lcm-diagonal} --- a
finite floor, not evidence toward the cofinal supply.
\end{obs}

\subsubsection{Cyclotomic anchored kills: unconditional prime support, plus new exclusion certificates}

The newly published \texttt{ErdosProblems/Erdos249/CyclotomicAnchoredKill.lean}
(3222 lines, namespace \texttt{ErdosProblems.Erdos249.CyclotomicAnchoredKill})
discharges the abstract order-consumer producer of
\S\ref{prime-ray-curvature} completely for the concrete polynomial $X-2$, and
separately deposits new kernel-checked denominator exclusions. Its middle
$\approx\!1600$ lines are a further ladder of
\texttt{CyclotomicPrime...CarryKillSupply} $\Leftrightarrow$ irrationality
equivalences; those are reductions and are reported only as architecture,
not as narrowing progress.

\begin{thm}[Unconditional unbounded prime support for the Mersenne layer]
Every prime divisor $p$ of $2^q-1$, for prime $q$, satisfies $q\mid p-1$
(the order of $2$ mod $p$ is exactly $q$, by Fermat/Lagrange in
$(\mathbb{Z}/p)^\times$); consequently the prime divisors appearing in the
layers $\{2^n-1\}$ are unbounded, unconditionally, with no cyclotomic
resultant hypothesis left open.
\lean{ErdosProblems.Erdos249.CyclotomicAnchoredKill.prime_index_dvd_pred}{ErdosProblems/Erdos249/CyclotomicAnchoredKill.lean:33}
\quad \lean{ErdosProblems.Erdos249.CyclotomicAnchoredKill.mersenneLayer_unboundedPrimeDivisorSupply}{ErdosProblems/Erdos249/CyclotomicAnchoredKill.lean:105}
\quad \ev{Lean} \quad \scale{uniform} \quad \coord{other:cyclotomic-anchor}.
\end{thm}

\begin{thm}[Unconditional clean cyclotomic anchor existence]
For every period $h>0$ and threshold $N_0$, there exist a prime $q$ and a
prime factor $p$ of $|\Phi_{hq}(2)|$ (the binary cyclotomic layer) with $p$
coprime to $hq$, $hq\mid p-1$, and $p-1\ge N_0$. The characteristic-prime
exceptional case in the cyclotomic order decomposition is eliminated
directly, by choosing $q>2^h$ (rules out $p=q$) and $q>h$ (rules out
$p\mid h$).
\lean{ErdosProblems.Erdos249.CyclotomicAnchoredKill.exists_clean_binaryCyclotomicAnchor}{ErdosProblems/Erdos249/CyclotomicAnchoredKill.lean:1607}
\quad \lean{ErdosProblems.Erdos249.CyclotomicAnchoredKill.binaryCyclotomicLayer_prime_order_decomposition}{ErdosProblems/Erdos249/CyclotomicAnchoredKill.lean:1431}
\quad \ev{Lean} \quad \scale{uniform} \quad \coord{other:cyclotomic-anchor}.
\end{thm}

\begin{obs}
These two theorems are genuine unconditional number theory --- an exact-order
argument and a Dirichlet-plus-cyclotomic-root existence construction --- and
they close the abstract producer hypotheses of
\texttt{PrimeRayCyclotomicCurvature} (\S\ref{prime-ray-curvature}) for the
$X-2$ layer with no residual conditional. Neither result touches the
irrationality question itself; they supply arithmetic support for a
\emph{different}, still-open cofinal predicate about certified kills at
cyclotomic-anchored periods.
\end{obs}

\begin{obs}[New certified exclusions at period 30]
The kernel decides two independent certificates at the composite cyclotomic
anchor period $h=30$ (natural basepoint $N=300$, and the prime-anchored
basepoint $N=330$), yielding a genuine new denominator exclusion:
$S\ne r$ for every $r\in\mathbb{Q}$ with $r.\mathrm{den}\mid 2^{300}(2^{30}-1)$.
\lean{ErdosProblems.Erdos249.CyclotomicAnchoredKill.certifiedKill_cyclotomic_30_331_natural}{ErdosProblems/Erdos249/CyclotomicAnchoredKill.lean:3356}
\quad \lean{ErdosProblems.Erdos249.CyclotomicAnchoredKill.totient_series_ne_rat_of_den_dvd_30_300}{ErdosProblems/Erdos249/CyclotomicAnchoredKill.lean:3377}
\quad \ev{Cert} \quad \scale{fixed} \quad \coord{other:cyclotomic-anchor} --- a
finite exclusion, not evidence toward the cofinal supply.
\end{obs}

\subsubsection{Prime-ray cyclotomic curvature: the abstract order-consumer producer}
\label{prime-ray-curvature}

\begin{thm}[Bounded-degree order realisability forces unbounded prime support]
If every prime divisor of a layer $C(mq)$ has order at most $d$ in the
relevant residue group (\texttt{BoundedDegreeOrderConsumer}), then for every
fixed finite set of primes $S$, all sufficiently large prime indices $q$
avoid $S$ entirely on layer $C(mq)$; combined with a nontrivial-layer
hypothesis this forces arbitrarily large new prime divisors on cofinally
many prime indices.
\lean{ErdosProblems.Erdos249.PrimeRayCyclotomicCurvature.finitePrimeSupportEscape_of_orderConsumer}{ErdosProblems/Erdos249/PrimeRayCyclotomicCurvature.lean:165}
\quad \lean{ErdosProblems.Erdos249.PrimeRayCyclotomicCurvature.unboundedPrimeDivisorSupply_of_orderConsumer}{ErdosProblems/Erdos249/PrimeRayCyclotomicCurvature.lean:242}
\quad \ev{Lean} \quad \scale{uniform} \quad \coord{other:cyclotomic-anchor}. This
is the abstract producer instantiated unconditionally for $X-2$ in
\S\ref{prime-ray-curvature}'s sibling result above; polynomial resultant
realisability and Archimedean growth for other layers remain explicit
upstream obligations, not proved here.
\end{thm}

\subsubsection{Rank-one subrank obstruction: a uniform proved barrier}
\label{rank-one-subrank}

\texttt{RankOneSubrankObstruction} is a uniform proved barrier: it names an
entire family of candidate linear-form constructions and proves, uniformly,
that none of them can work.

\begin{thm}[Uniform rank-one no-go]
Let $\Theta_r$ be the Möbius--Mersenne ladder rung and let $t(Y,r)$ be its
first $Y$ atoms. For every $e\ge 1$ and $Y\ge 4$,
$$t(Y,e{+}2)^2/t(Y,2e{+}2) - \Theta_2 > 1/480.$$
Every rank-one monomial Schur quotient built this way overshoots $\Theta_2 =
S-1/2$ by a fixed positive margin, uniformly in $e$ and $Y$; the bound uses
only two interval facts (every rung $r\ge 3$ lies in $[1429/1512,1)$, and
every length-$\ge\!4$ prefix is within $1/3584$ of its rung).
\lean{ErdosProblems.Erdos249.RankOneSubrankObstruction.rankOneSubrankQuotient_sub_theta_two_gt}{ErdosProblems/Erdos249/RankOneSubrankObstruction.lean:238}
\quad \lean{ErdosProblems.Erdos249.RankOneSubrankObstruction.mobiusMersenneTheta_ge_alpha}{ErdosProblems/Erdos249/RankOneSubrankObstruction.lean:163}
\quad \ev{Lean} \quad \scale{uniform} \quad \coord{other:rank-one-subrank}.
\end{thm}

\begin{prop}[The gap survives positive averaging]
For any nonempty finite positively-weighted family of admissible quotients,
the weighted average still overshoots $\Theta_2$ by more than $1/480$; and
every primitive integer linear form $q\Theta_2 - p$ realising such a
quotient satisfies $|q\Theta_2 - p| > q/480$.
\lean{ErdosProblems.Erdos249.RankOneSubrankObstruction.positive_direct_sum_sub_theta_two_gt}{ErdosProblems/Erdos249/RankOneSubrankObstruction.lean:300}
\quad \lean{ErdosProblems.Erdos249.RankOneSubrankObstruction.primitive_form_abs_gt}{ErdosProblems/Erdos249/RankOneSubrankObstruction.lean:341}
\quad \ev{Lean} \quad \scale{uniform} \quad \coord{other:rank-one-subrank}.
\end{prop}

\begin{obs}
This is a proved barrier stated at the required precision: the class of
argument it rules out is exactly ``a primitive rational linear form for
$\Theta_2$ (equivalently for $S$) obtained as a rank-one strict-subrank
monomial quotient of finite Möbius--Mersenne prefixes,'' and the reason is
the explicit uniform lower bound above, not an empirical failure report.
It does not touch, and is not claimed to touch, linear forms built by any
other mechanism.
\end{obs}

\subsubsection{Period-multiple escape: the nesting identity and eight new exclusions}

\texttt{PeriodMultipleEscape} proves the period-multiple kill supply is
\emph{exactly equivalent} to irrationality of $S$ (both directions,
sufficiency by telescoping the tail-period law, necessity by certificate
completeness) --- reported here as an equivalence-class reduction, not
progress --- and separately deposits genuinely new denominator exclusions.

\begin{thm}[Nesting: the cyclotomic fan collapses onto the dyadic tower]
$Q_{a+b,N} = 2^b Q_{a,N} + Q_{b,N+a}$ (block concatenation), so the order-4
cyclotomic channel at height $h$ is literally the order-2 channel at height
$2h$: the 2--3--4 fan of channels is a nested family along the tower
$h,2h,4h,\dots$, not three unrelated moduli.
\lean{ErdosProblems.Erdos249.PeriodMultipleEscape.totientBlock_add}{ErdosProblems/Erdos249/PeriodMultipleEscape.lean:77}
\quad \lean{ErdosProblems.Erdos249.PeriodMultipleEscape.totientBlock_two_mul}{ErdosProblems/Erdos249/PeriodMultipleEscape.lean:373}
\quad \ev{Lean} \quad \scale{uniform} \quad \coord{other:lcm-period-multiple}.
\end{thm}

\begin{obs}[Supply $\Leftrightarrow$ irrational --- reduction, not a result]
\lean{ErdosProblems.Erdos249.PeriodMultipleEscape.periodMultipleKillSupply_iff_irrational}{ErdosProblems/Erdos249/PeriodMultipleEscape.lean:432}
\quad \ev{Lean} (equivalence proved; both directions of the underlying
predicate are exactly as open as \#249 itself)
\quad \scale{cofinal} \quad \coord{other:lcm-period-multiple}. The paper's
current aggregate diagonal bank certifies kills at every $H_t$ for $t\le82$;
this module's equivalence contributes no new information about $t=83$ or any
cofinal supply.
\lean{ErdosProblems.Skip.LadderT67.exists_diagonalKill_le_82}{ErdosProblems/Skip/LadderT67.lean:71264}
\end{obs}

\begin{obs}[Eight new certified denominator exclusions past the 64-smooth diagonal bank]
The kernel certifies kills, unconditionally, at the eight prime-power
periods the diagonal bank could not reach ($67,81,97,101,121,125,127,128$,
all at basepoint $N=300$), each yielding $S\ne r$ for every $r$ with
$r.\mathrm{den}\mid 2^{300}(2^h-1)$ --- new odd denominator classes,
including the Cole factors of $2^{67}-1$ and the Mersenne prime $2^{127}-1$.
One of the eight is certified at its own locked depth $L=h=67$, the concrete
instance of the sufficient (not known necessary) depth-equals-period form
$\mathrm{ApFullDepthEscape}$.
\lean{ErdosProblems.Erdos249.PeriodMultipleEscape.certifiedKill_67_300}{ErdosProblems/Erdos249/PeriodMultipleEscape.lean:471}
\quad \lean{ErdosProblems.Erdos249.PeriodMultipleEscape.certifiedKill_127_300}{ErdosProblems/Erdos249/PeriodMultipleEscape.lean:490}
\quad \lean{ErdosProblems.Erdos249.PeriodMultipleEscape.certifiedKill_67_300_fullDepth}{ErdosProblems/Erdos249/PeriodMultipleEscape.lean:498}
\quad \lean{ErdosProblems.Erdos249.PeriodMultipleEscape.totient_series_ne_rat_of_den_dvd_127_300}{ErdosProblems/Erdos249/PeriodMultipleEscape.lean:534}
\quad \ev{Cert} \quad \scale{fixed} \quad \coord{other:lcm-period-multiple}
--- eight finite exclusions, not evidence toward the cofinal supply.
\end{obs}

\subsubsection{Strict prime-orbit escape: a sharper reduction}

\texttt{TotientStrictPrimeEscape} sharpens the legacy first-harmonic
producer's threshold from a $4/5$ gap to a strict $9/10$ gap with an
adaptive truncation budget, and proves the sharper predicate still closes
\#249 through the existing singleton-certificate endpoint. This is, exactly
as its own comment states, a reduction: ``the producer itself remains
unproved.''
\lean{ErdosProblems.Erdos249.irrational_totient_series_of_naturalPrimeTailOrbitStrictGap}{ErdosProblems/Erdos249/TotientStrictPrimeEscape.lean:524}
\quad \ev{Lean} (implication proved; antecedent
\ev{Open})
\quad \scale{cofinal} \quad \coord{other:first-harmonic}.

\subsubsection{Finite Euler-sieve algebra}

\texttt{FiniteEulerSieve} records the elementary finite-stage identities
behind the squared-Möbius Euler factor, with no transcendence claim
asserted: $(1-2/p+1/p^2) = (1-1/p)^2$ and its degree-two analogue, and the
second finite difference of $1+p+\cdots+p^{e}$ recovering
$p^{e+1}(p-1)$, the prime-power totient row.
\lean{ErdosProblems.Erdos249.FiniteEulerSieve.muSqEulerFactor_one}{ErdosProblems/Erdos249/FiniteEulerSieve.lean:28}
\quad \lean{ErdosProblems.Erdos249.FiniteEulerSieve.sigmaPrimePow_secondDiff}{ErdosProblems/Erdos249/FiniteEulerSieve.lean:51}
\quad \ev{Lean} \quad \scale{uniform} \quad \coord{other:euler-sieve}.
% ---- part a249_p2 ----
\section{The scale ladder}\label{sec:249-scale-ladder}

Every statement in the Erd\H{o}s \#249 corpus that this paper draws on is tagged with a
\emph{scale}: \scale{fixed} (a finite, explicitly enumerated set of parameter values, typically
closed by \texttt{decide}/\texttt{interval\_cases}), \scale{bounded} (holds for all parameter
values on one side of a threshold, e.g.\ $\forall a\ge 8$, but the proof or the constants inside
it do not survive removing the threshold), \scale{uniform} (holds unconditionally for every value
of every free parameter, with no scale restriction at all), and \scale{cofinal} (an existential
claim of the shape $\forall N_0\,\exists N\ge N_0,\ P(N)$ — infinitely often, arbitrarily far out).

The source tables underlying this ladder also mark a handful of pure identities and converters
\texttt{n/a} when they are not indexed by any problem-scale parameter at all (they are definitions
or unconditional equivalences); for ladder purposes these are folded into \scale{uniform}, since an
unconditional statement is, if anything, stronger than a uniform one. This convention is applied
uniformly below and is stated here once rather than re-flagged on every row.

\#249's own supply obligation is \emph{exactly} cofinal. The reduction chain
(\S\ref{sec:249-scale-ladder}, and see \lean{irrational\_totientSeries\_iff\_actualLcmOrbitNonintegralitySupply}{TotientActualLcmOrbitNonintegrality.lean:37})
shows
\[
  \mathrm{Irrational}\Bigl(\sum_{n\ge 1}\varphi(n)/2^n\Bigr) \iff \forall a_0\ \exists a\ge a_0,\ \texttt{actualLcmTailOrbit}\ a \notin \mathrm{range}(\mathbb{Z}\to\mathbb{R}),
\]
an \texttt{cofinal} statement in the exponent $a$ that indexes $H=\mathrm{periodLcm}(2^a)$. Nothing
weaker in scale can close \#249; the entire question is whether the corpus's uniform, bounded, and
fixed machinery can be pushed to cofinal. The table below lists every catalogued \#249 result,
grouped by scale, with a horizontal rule separating everything that is \emph{proved} (uniform,
bounded, fixed) from the handful of statements that are themselves the \emph{cofinal target} or
consumers of it (below the rule). Quantifier prefixes are written out explicitly and exactly as
recorded against the Lean source; nothing is compressed to ``$\forall\dots$'' where the source
specifies bounds.

\begingroup
\begin{landscape}
\fontsize{8}{9.8}\selectfont
\setlength{\tabcolsep}{3pt}
\begin{longtable}{@{}p{0.22\linewidth}p{0.18\linewidth}p{0.08\linewidth}p{0.34\linewidth}p{0.12\linewidth}@{}}
\caption{The \#249 scale ladder. Every row above the rule is proved; every row below it is the
cofinal target or a direct consumer of it.}\\
\toprule
\papertablehead
\textbf{Result} & \textbf{Lean site} & \textbf{Scale} & \textbf{Quantifier prefix} & \textbf{Coordinate}\\
\midrule
\endfirsthead
\toprule
\papertablehead
\textbf{Result} & \textbf{Lean site} & \textbf{Scale} & \textbf{Quantifier prefix} & \textbf{Coordinate}\\
\midrule
\endhead
\midrule \multicolumn{5}{r}{\emph{continued on next page}} \\
\endfoot
\bottomrule
\endlastfoot

\papertablehead\multicolumn{5}{l}{\emph{-- \scale{uniform} (includes source-tagged \texttt{n/a} unconditional identities/converters) --}}\\
totientTail well-defined & \lean{totientTail}{TotientTailPeriodKiller.lean:57} & uniform & $\forall N$ & binary-digit\\
$2^N\!\cdot\! S$ split & \lean{two\_pow\_mul\_totient\_series\_eq}{TotientTailPeriodKiller.lean:150} & uniform & $\forall N$ & binary-digit\\
windowDiscrepancy (def) & \lean{windowDiscrepancy}{TotientTailPeriodKiller.lean:66} & uniform & $\forall h\,\forall N\,\forall L$ & other-binary-window\\
certifiedKill Sep$(h,N,L)$ (def) & \lean{certifiedKill}{TotientTailPeriodKiller.lean:72} & uniform & $\forall h\,\forall N\,\forall L$ & other-binary-window\\
certificate depth floor & \lean{certifiedKill\_depth\_floor}{TotientTailPeriodKiller.lean:79} & uniform & $\forall h\,\forall N\,\forall L$, given \texttt{certifiedKill}$\,h\,N\,L$ & other-binary-window\\
kill engine soundness & \lean{tail\_diff\_notMem\_int\_of\_certifiedKill}{TotientTailPeriodKiller.lean:262} & uniform & $\forall h\,\forall N\,\forall L$, given \texttt{certifiedKill} & other-binary-window\\
certificates are complete receipts & \lean{exists\_certifiedKill\_iff\_tail\_diff\_notMem\_int}{LcmConeFlatness.lean:316} & uniform & $\forall h\,\forall N$ & other-binary-window\\
forced integrality under den.\ divisibility & \lean{tail\_diff\_int\_of\_den\_dvd}{TotientTailPeriodKiller.lean:327} & uniform & $\forall N\,\forall h\,\forall r{:}\mathbb{Q}$, given $S=r$, $r.\mathrm{den}\mid 2^N(2^h{-}1)$ & binary-digit\\
tail-period law & \lean{eventual\_period\_of\_not\_irrational}{TotientTailPeriodKiller.lean:358} & uniform & $\lnot\mathrm{Irrational}\,S\to\exists h{>}0\,\exists N_0\,\forall N{\ge}N_0$ & binary-digit\\
periodLcm basic facts & \lean{periodLcm}{CarrySurvivorExtinction.lean:515} & uniform & $\forall t$ & other-lcm-period-ray\\
non-divisor$<2t$ is a bare prime power & \lean{eq\_prime\_pow\_of\_not\_dvd\_periodLcm}{LcmDiagonalReduction.lean:137} & uniform & $\forall t\,\forall j\,(0{<}j{<}2t)$, given $j\nmid\mathrm{periodLcm}\,t$ & other-lcm-window\\
LCM-ray multiplicative split & \lean{totient\_periodLcm\_ray\_split}{LcmDiagonalReduction.lean:208} & uniform & $\forall t\,\forall j\,\forall q$, given $j\mid\mathrm{periodLcm}\,t$, clean-divisor & other-lcm-window-multiplicative\\
rationality flattens the whole LCM cone & \lean{rational\_totient\_series\_forces\_lcm\_cone\_flatness}{CertificateKernel.lean:19002} & uniform & $\lnot\mathrm{Irrational}\,S\to\exists t_1\,\forall t{\ge}t_1\,\forall q{>}0\,\forall m$ & other-lcm-cone\\
second-difference kill engine & \lean{second\_diff\_notMem\_int\_of\_certifiedRank2Kill}{LcmConeFlatness.lean:493} & uniform & $\forall h\,\forall N\,\forall L$, given \texttt{certifiedRank2Kill} & other-binary-window\\
cone-menu nonintegral pair & \lean{exists\_nonintegral\_pair\_of\_coneNonflatCert}{LcmConeNonflat.lean:126} & uniform & $\forall H\,\forall L\,\forall Q{\ne}\emptyset$, given \texttt{coneNonflatCert} & other-lcm-cone-menu\\
survivorKill engine (carry-orbit route) & \lean{tail\_diff\_notMem\_int\_of\_survivorKill}{CarrySurvivorExtinction.lean:428} & uniform & $\forall h\,\forall N\,\forall K$, given \texttt{survivorKill} & other-carry-orbit\\
Farey neighbour denominator law & \lean{farey\_gap}{GapFareyBound.lean:51} & uniform & $\forall a\,b\,c\,d\,r\,s{:}\mathbb{Z}$, $b{>}0,d{>}0,bc{-}ad{=}1$, unimodular-between & farey\\
Dirichlet near-integer criterion & \lean{irrational\_of\_den\_mul\_abs\_sub\_tendsto\_zero}{CertificateKernel.lean:5371} & uniform & $\forall u{:}\mathbb{N}\to\mathbb{Q}\,\forall x{:}\mathbb{R}$ & n/a (generic)\\
Erd\H{o}s 1948 near-integer criterion & \lean{irrational\_of\_int\_mul\_near\_int}{CertificateKernel.lean:6120} & uniform & $\forall \xi{:}\mathbb{R}$, $\forall q{>}0\,\exists m\,z\,0{<}|m\xi{-}z|{<}1/q$ & n/a (generic)\\
finite-minor $\Rightarrow$ linear independence & \lean{linearIndependent\_of\_separatedMinorCertificate}{TotientMahlerDefect.lean:91} & uniform & $\forall\,\mathrm{family}{:}\iota\to\mathbb{N}\to\mathbb{Q}$, given SeparatedMinorCertificate & n/a (generic)\\
dyadic totient-kernel rank $=2^e{+}1$ & \lean{linearIndependent\_canonicalTotientKernelFamily}{TotientMahlerDefect.lean:935} & uniform & $\forall e$ & other-dyadic-kernel-rank\\
rationality forces unbounded carry-kernel rank & \lean{not\_irrational\_totientSeries\_implies\_unbounded\_carryRank\_unconditional}{TotientCarryKernelRigidity.lean:300} & uniform & $\lnot\mathrm{Irrational}\,S\to\exists v{>}0\,\exists u\,\forall e$ & other-carry-kernel-rank\\
no fixed $D$ clears all primitive Euler jets & \lean{no\_fixed\_integral\_primitive\_euler\_index}{PrimitiveDeterminantLift.lean:169} & uniform & $\forall D{>}0$ & mobius-mersenne\\
Mersenne–Lambert ladder value identities & \lean{tsum\_totient\_div\_pow\_two\_eq\_pnat\_half\_pow}{CertificateKernel.lean:18415} & uniform & value identities, $L(\mu){=}\tfrac12$, $L(\varphi){=}2$, $L(1){=}$Erd\H{o}s--Borwein & mobius-mersenne\\
prime-power reduced-denom.\ unit-gap ceiling & \lean{reducedDenominatorUnitGapCert\_prime\_pow\_iff}{PrimitiveWeightCertificate.lean:54} & uniform & $\forall p\ \text{prime}\,\forall e{>}0$ & farey\\
positive rational-difference lower bound & \lean{positive\_rational\_difference\_lower\_bound}{PrimitiveRationalGapSupply.lean:31} & uniform & $\forall\ \text{whole}\,\text{pfx}{:}\mathbb{Q},\ \text{pfx}{<}\text{whole}$ & n/a (generic)\\
M\"obius-square identity ($S=\tfrac12+\Sigma$) & \lean{tsum\_totient\_half\_pow\_eq\_half\_add\_moebius\_sq}{MersenneLambertLadder.lean:665} & uniform & unconditional & mobius-mersenne\\
squared-Lambert transfer engine & \lean{tsum\_lambert\_linear\_weight\_sq\_pure}{GcdMomentCalculus.lean:105} & uniform & $\forall w{:}\mathbb{N}\to\mathbb{R}\,(|w(d)|{\le}d)\,\forall r\in[0,1)$ & mobius-mersenne\\
$L_2(\mu)=S-\tfrac12$ & prose:L2\_of\_mu\_is\_249 & uniform & unconditional & mobius-mersenne\\
$L_2(1)=\zeta_q(2)-\zeta_q(1)$ identity & \lean{tsum\_one\_div\_mersenne\_sq\_eq\_sigma\_sub\_tau\_series}{GcdMomentCalculus.lean:216} & uniform & unconditional & mobius-mersenne\\
$L_2(\varphi)$ = gcd moment & \lean{tsum\_totient\_div\_mersenne\_sq\_eq\_gcd\_moment\_series}{GcdMomentCalculus.lean:235} & uniform & unconditional & mobius-mersenne/probability\\
level-mirror table & prose:level\_mirror\_table & uniform & unconditional & mobius-mersenne\\
gcd-divisibility factorizes & \lean{tsum\_pos\_pair\_both\_dvd\_half\_eq\_inv\_mersenne\_sq}{GcdMomentCalculus.lean:266} & uniform & $\forall d{:}\mathbb{N},\,d{>}0$ & probability\\
reduced-direction law $\Sigma=1$ & \lean{tsum\_pos\_coprime\_inv\_mersenne\_eq\_one}{GcdMomentCalculus.lean:349} & uniform & unconditional (sum over all coprime pairs) & probability\\
Stern--Brocot cylinder recursion & \lean{cylinderMass\_split}{GcdMomentCalculus.lean:474} & uniform & $\forall a\,b{:}\mathbb{N}^+\,\forall\text{depth}\,d$ & other-stern-brocot\\
repunit gcd-word T1 & \lean{mobiusNumeratorPolynomial\_eq\_gcdWord}{RepunitMobiusNumerator.lean:205} & uniform & $\forall r\,(\mathrm{Squarefree}\,r)\,\forall k{<}r$ & cyclotomic\\
eval-at-2 recovers numerator & \lean{mobiusNumeratorPolynomial\_eval\_two}{RepunitMobiusNumerator.lean:445} & uniform & $\forall r\,(\mathrm{Squarefree}\,r)$ & cyclotomic\\
radical-shadow scale decomposition & \lean{scaledMobiusShadow\_eq\_radicalBase}{RadicalMobiusShadow.lean:125} & uniform & $\forall H{>}0\,\forall r{>}0$ & mobius-mersenne\\
cyclotomic congruence T2 & \lean{cyclotomic\_dvd\_mobiusNumeratorPolynomial\_sub}{CyclotomicProjectionOfShadow.lean:299} & uniform & $\forall r\,(\mathrm{Squarefree}\,r)\,\forall m\,(m{\mid}r)$ & cyclotomic\\
top-fibre survives T3 & \lean{mobiusNumerator\_gcd\_cyclotomicValue}{CyclotomicProjectionOfShadow.lean:350} & uniform & $\forall r\,(\mathrm{Squarefree}\,r)$ & cyclotomic\\
upper-half prime channel survival T4 & \lean{upperHalfChannel\_product\_dvd\_den\_of\_scale\_primeFactors\_le}{MersenneShadowCyclotomicNoncollapse.lean:809} & uniform & $\forall t\,\forall p\ \text{prime}\in(t/2,t]$, scale-coprime-to-$C$ & cyclotomic\\
denominator lower bound / exact value & \lean{lcmHeight\_scaledMobiusShadow\_den\_exact}{MersenneShadowDenominatorGrowth.lean:147} & uniform & $\forall t$ (exact-value); $t{\ge}5$ (lower bound) & mobius-mersenne\\
prime-power jump recurrence T5 & \lean{cyclotomic\_dvd\_primeJump\_new\_fibre}{PrimePowerJumpDynamics.lean:281} & uniform & $\forall p\ \text{prime}\,(p{\nmid}r)\,\forall m\,(m{\mid}r)$ & cyclotomic\\
squared-Mersenne diagonal tail enclosure & \lean{mobius\_square\_series\_eq\_partial\_add\_tail}{SquaredMersenneDiagonalEnclosure.lean:94} & uniform & $\forall D{:}\mathbb{N}$ & mobius-mersenne\\
$(3,5)$ joint annihilator & \lean{joint35\_oldChannel\_zero}{JointExponentTransport.lean:71} & uniform & $\forall d{>}0\,(d{\mid}H)$; general: finite affine annihilator & other-lcm-diagonal\\
composite-dilation defect identity & \lean{supportCoeff\_mul\_eq\_add\_defect}{CompositeDilationDefect.lean:30} & uniform & $\forall a{\in}A\,(a{>}0)\,\forall x{>}0$ & other-divisor-support-coeff\\
sublogarithmic zero-window T11 & \lean{supportCoeffZeroWindow\_length\_le\_eps\_logb\_add}{SublogDivisorCoverage.lean:392} & uniform & $\forall \varepsilon{>}0\,\exists B$, given rationality of the base-2 support series & other-support-divisor-counting\\
campbell shift $\leftrightarrow$ Mersenne endpoint (shared namespace, \#257-flavored) & \lean{shiftWindowZero\_iff\_half\_mem\_mersenneAchievementSet}{CampbellShiftSynchronization.lean:488} & uniform & iff, no extra parameter & other-greedy-mersenne-achievement\\
totient overlap-factor multiplicativity & \lean{totient\_mul\_eq\_overlapFactor\_mul}{TotientActualLcmOrbitArithmetic.lean:43} & uniform & $\forall x{>}0$ ($j$ fixed) & other-totient-arithmetic\\
totient relative-Euler-product split & \lean{totient\_mul\_eq\_totient\_mul\_relativeEulerProduct}{TotientActualLcmOrbitArithmetic.lean:100} & uniform & $\forall j\,x{>}0$ & other-totient-arithmetic\\
lcmRayArithmeticLetter $=$ deltaTotient $=$ diagonalWindowIncrement & \lean{lcmRayArithmeticLetter\_eq\_deltaTotient}{TotientActualLcmOrbitArithmetic.lean:298} & uniform & $\forall t\,\forall j$ & mobius-mersenne\\
lcmDiagonalArithmeticWord $=$ windowDiscrepancy & \lean{lcmDiagonalArithmeticWord\_eq\_windowDiscrepancy}{TotientActualLcmOrbitArithmetic.lean:1998} & uniform & $\forall t\,\forall L$ & mobius-mersenne\\
\#249 $\iff$ cofinal actual-orbit nonintegrality & \lean{irrational\_totientSeries\_iff\_actualLcmOrbitNonintegralitySupply}{TotientActualLcmOrbitNonintegrality.lean:37} & uniform & iff & mobius-mersenne\\
actualLcmTailOrbit $=$ scaled series $-$ prefix & \lean{actualLcmTailOrbit\_eq\_scaled\_totientSeries\_sub\_prefix}{TotientActualLcmOrbitSeparation.lean:68} & uniform & $\forall a$ & mobius-mersenne\\
explicit finite-block approximation & \lean{abs\_actualLcmTailOrbit\_sub\_rawApprox\_lt}{TotientActualLcmOrbitSeparation.lean:141} & uniform & $\forall a\,\forall q$ & mobius-mersenne\\
guard-cylinder normal form & \lean{exists\_certifiedKill\_iff\_guardCylinderWitness}{TotientActualLcmTopEdgeStaircase.lean:844} & uniform & $\forall h\,N,\ \text{iff}$ & seam-integer\\
fixedRankSecondDifference sign under extremality & \lean{fixedRankSecondDifference}{TotientFixedRankLcmAsymptotic.lean:256} & uniform & $\forall H\,j$, given MiddleRankTotientExtremal & other-fixed-rank-curvature\\
dyadic fixture exactness (kernel-valuation sharpness) & \lean{fixedRankSecondDifference\_dyadic\_fixture\_exact}{TotientFixedRankLcmAsymptotic.lean:462} & uniform & $\forall n$ (closed-form) & other-fixed-rank-curvature\\
parityCoboundaryWeight (def) & \lean{parityCoboundaryWeight}{TotientParityCoboundaryCountermodel.lean:59} & uniform & $\forall n$ (definition) & binary-digit\\
totientTail $=$ positive foreign-residue kernel sum & \lean{totientTail\_eq\_tsum\_positiveForeignResidueKernel}{TotientShiftedMobiusForeignBridge.lean:61} & uniform & $\forall N$ & cyclotomic\\
totientTail $=$ shifted M\"obius-pulse sum & \lean{totientTail\_eq\_tsum\_shiftedMobiusPulse}{TotientShiftedMobiusPulse.lean:547} & uniform & $\forall N$ & cyclotomic\\
Lambert double-sum regroup engine & \lean{tsum\_lambert\_pair\_regroup\_if}{TotientShiftedMobiusPulse.lean:124} & uniform & $\forall w\,v{:}\mathbb{N}\to\mathbb{R}\,(|w(d)|{\le}d,\,|v(m)|{\le}m)$ & cyclotomic\\
carryShift-integrality iff & \lean{carryShift\_dvd\_iff\_tailDiff\_mem\_int}{TotientTailCarryPeriod.lean:77} & uniform & $\forall N\,k$, given tempered orbit & seam-integer\\
rationality $\Rightarrow$ periodicity \& unbounded rank & \lean{not\_irrational\_totientSeries\_implies\_mod\_period\_and\_unbounded\_rank}{TotientTailCarryPeriod.lean:224} & uniform & $\lnot\mathrm{Irrational}\,S\to\exists v\,u\,\forall e$ & seam-integer\\
adjugate tail-cost floor $\ge 3$ & \lean{three\_le\_totientAdjugateTailCost}{TotientTailCarryPeriod.lean:266} & uniform & $\forall\ \text{finite}\ w{:}\iota\to\mathbb{Q}\,x{:}\iota\to\mathbb{N}$, given $\sum w_i\varphi(x_i){=}1$ & other-adjugate-linear-algebra\\
directedCertifiedKill exact iff & \lean{directedCertifiedKill}{TotientTailCarryPeriod.lean:766} & uniform & $\forall h\,N\,L$ & seam-integer\\
tempered orbit rigidity iff & \lean{binaryCoeffSeries\_rational\_iff\_exists\_temperedBinaryOrbit}{GenericTailOrbitRigidity.lean:426} & uniform & $\forall c{:}\mathbb{N}\to\mathbb{N}\,(c(n){\le}n)$, iff & binary-digit\\
doubling-orbit rigidity & \lean{doublingOrbit\_eq\_zero\_of\_tempered}{GenericTailOrbitRigidity.lean:315} & uniform & $\forall d{:}\mathbb{N}\to\mathbb{R}$, doubling $+$ tempered $\to d{\equiv}0$ & other-abstract-recursion\\
finite-state no-go (unbounded state) & \lean{balancedPulse\_no\_autonomous\_decoder}{GenericTailOrbitRigidity.lean:247} & uniform & $\forall m\,\forall\mathrm{State}$, collapsing-state $\to$ no decoder & binary-digit\\
fixed-depth affine reset & \lean{affineBinaryOrbit\_mod\_twoPow\_eq}{GenericTailOrbitRigidity.lean:307} & uniform & $\forall a{:}\mathbb{N}\to\mathbb{Z}\,\forall u_0\,v_0\,\forall L$ & binary-digit\\
signed dichotomy engine (full-block certs) & \lean{int\_coeff\_series\_irrational\_or\_bpow\_mul\_eq\_intCast\_of\_full\_block\_certificates}{CertificateKernel.lean:13934} & uniform & $\forall b{\ge}2\,\forall c\,\forall q$, given full-block certificate at $q$ & binary-digit\\
Lambert double-sum regroup (dyadic forward-difference calc.) & \lean{dvd\_iteratedForwardDifference}{LcmFactorIdealPulseObstruction.lean:211} & uniform & $\forall f,c{:}\mathbb{N}\to\mathbb{Z}\,\forall\ \text{shift-lists}$ & binary-digit\\
fair-coin coprimality bridge & \lean{tsum\_coprime\_pair\_pow\_eq\_tsum\_totient\_mul\_pow}{GeometricCoprimality.lean:120} & uniform & $\forall n$; $\forall r\in[0,1)$ & probability\\
linear-descender rigidity & \lean{linearDescender\_eq\_smul\_eval}{AdelicHeightObstruction.lean:120} & uniform & $\forall V,W\,\forall\mathrm{ev}\,\forall L$, $\ker(\mathrm{ev}){\le}\ker(L)$ & p-adic\\
rational denominator survival law & \lean{divisor\_dvd\_divInt\_den}{RationalDenominatorSurvival.lean:17} & uniform & $\forall D{>}0\,\forall m{\mid}D\,\forall a$ & p-adic\\
scalar-localization complement-dvd (adelic height tax) & \lean{scalarLocalization\_complement\_dvd}{AdelicHeightObstruction.lean:23} & uniform & $\forall x{:}\mathbb{Q}\,\forall H\,(H{\mid}x.\mathrm{den})\,\forall c\,((cx).\mathrm{den}{\mid}H)$ & p-adic\\
signed Hankel terminal-parity engine & \lean{scaled\_dyadic\_sum\_ne\_zero}{SignedQMomentObstruction.lean:96} & bounded & $\exists m{\in}s$ (strictly maximal, odd coeff), $\forall i{\ne}m\,e(i){<}e(m)$ & p-adic\\
residual gauge obstruction (locked minors) & \lean{det\_residualMonomialMatrix}{ResidualGaugeObstruction.lean:47} & uniform & $\forall d\,\forall e{:}\mathrm{Fin}\,d{\to}\mathbb{N}\,\forall z{:}\mathrm{Fin}\,d{\to}\mathbb{C}$ ($z$ nonzero) & other-first-harmonic-phase\\
incidence-quotient no compression & \lean{incidenceMobius\_det\_eq\_one}{IncidenceQuotientHermitePade.lean:56} & uniform & $\forall N{:}\mathbb{N}$ & mobius-mersenne\\

\papertablehead\multicolumn{5}{l}{\emph{-- \scale{bounded} --}}\\
Farey window $K{=}240$ denominator bound & \lean{gap\_check\_window\_1\_240\_le\_79639646646701375323355774875831053}{GapFareyBound.lean:176} & bounded &
$\forall q{:}\mathbb{N},\,0{<}q{\le}$\newline
$79{,}639{,}646{,}646{,}701{,}$\newline
$375{,}323{,}355{,}774{,}875{,}831{,}053$ & farey\\
denominator lower bound $S{\ne}p/q$, $q{\le}Q_0$ & \lean{tsum\_totient\_div\_pow\_two\_ne\_ratCast\_of\_den\_le\_79639646646701375323355774875831053}{CertificateKernel.lean:18384} & bounded & $\forall p{:}\mathbb{Q},\,p.\mathrm{den}{\le}Q_0$ & farey\\
foreign-residue tail-limit converter & \lean{controlledForeignProjection\_of\_tail\_limit}{LambertDiagonalEnclosure.lean:31} & bounded & $\forall H\,D$ with $2H{\le}D$, given tail-limit convergence & mobius-mersenne\\
fixed-precision tropical no-go & \lean{fixedPrecisionTropicalNoGo}{TropicalCurvatureCarry.lean:137} & bounded & $\forall u{>}0$ (fixed precision), $\forall$ finite carry word & p-adic\\
$\mu$-clean-Pad\'e feasibility wall & prose:mu\_pollution\_qpade\_wall & bounded & claimed $\forall d$; machine-audited only $d{\le}12$ & mobius-mersenne (q-Pad\'e)\\
iterated pullback / bounded-$\Omega$ vanishing & \lean{iteratedDilationDifference\_eq\_zero\_of\_cardFactors\_le}{SupportDilationDifferences.lean:250} & bounded & given $\Omega(a){\le}K\,\forall a{\in}A$ & other-support-divisor-counting\\
rough-density $\Rightarrow$ totient $\ge\tfrac34 n$ & \lean{three\_quarters\_mul\_lt\_totient\_of\_rough\_lt\_two\_pow}{TotientActualLcmOrbitArithmetic.lean:768} & bounded & $\forall a{\ge}8\,\forall n{>}0$ ($n$ is $t$-rough, $t{=}2^a$, $n{<}2^{2\cdot2^a}$) & mobius-mersenne\\
short-window sign fixed positive & \lean{lcmRayArithmeticLetter\_pos\_of\_lt\_two\_mul}{TotientActualLcmOrbitArithmetic.lean:1703} & bounded & $\forall a{\ge}8\,\forall j\,(0{<}j{<}2\cdot2^a)$ & mobius-mersenne\\
unconditional lower-half positivity & \lean{actualLcmTailDiff\_shift\_pos}{TotientActualLcmOrbitSign.lean:39} & bounded$^\dagger$ & $\forall a{\ge}8\,\forall J\,(J{+}(a{+}6){<}2\cdot2^a)$, no rationality hypothesis & mobius-mersenne\\
top-edge residue survivor negative & \lean{actualLcm\_trueEndpointSurvivor\_neg}{TotientActualLcmOrbitSign.lean:172} & bounded & $\forall a{\ge}8\,\forall J\,K$ (room bound) & mobius-mersenne\\
integrality forces top-edge residue & \lean{actualLcm\_integral\_forces\_topEdgeResidue}{TotientActualLcmOrbitSign.lean:211} & bounded & $\forall a{\ge}8\,\forall J\,K$ (room, $2H{+}J{+}K{+}2{<}2^K$) & mobius-mersenne\\
short-kill supply through $a{\le}6$ & \lean{powerTwoActualLcmShortArithmeticKillSupply\_through\_six}{TotientActualLcmShortKill.lean:54} & bounded & $\forall a_0{\le}6\,\exists a\,L$ & mobius-mersenne\\
terminal dyadic staircase impossible with room & \lean{not\_actualLcmTerminalDyadicStaircase\_of\_room}{TotientActualLcmTopEdgeStaircase.lean:251} & bounded & $\forall a{\ge}8\,\forall J\,K\,m$ (room, modulus-wide) & mobius-mersenne\\
punctured staircase penultimate pin & \lean{puncturedDyadicStaircase\_penultimate\_eq\_half}{TotientActualLcmTopEdgeStaircase.lean:1187} & bounded & $\forall a{\ge}8\,\forall J\,K\,m$ (room $+$ punctured hyp.) & mobius-mersenne\\
top-edge residue-gap $\Rightarrow$ nonintegral & \lean{actualLcmTailDiff\_notMem\_int\_of\_topEdgeResidueGap}{TotientActualLcmTopEdgeStaircase.lean:1260} & bounded & $\forall a{\ge}8\,\forall J\,K\,m$ (room, one-sided) & mobius-mersenne\\
odd-rank centred-lift $=$ terminal $-$ true carry & \lean{two\_mul\_actualOddHalfCenteredLift\_eq\_terminal\_sub\_trueCarry}{TotientActualLcmTopEdgeStaircase.lean:1678} & bounded & $\forall a{\ge}8\,\forall q$ (room, fit bound $2(H{+}q{+}2){\le}4^q$) & mobius-mersenne\\
prime-power kernel valuation floor & \lean{two\_mul\_totient\_dvd\_fixedRankSecondDifference}{TotientFixedRankLcmAsymptotic.lean:519} & bounded & $\forall a{\ge}4\,\forall j{>}0\,(j^2{\le}2^a)$ & other-fixed-rank-curvature\\
finite rational-separation $\Rightarrow$ $\lnot$ hit (\#257-flavored) & \lean{renormalizedResidueAgreement\_of\_twice\_le}{TotientShiftedMobiusForeignBridge.lean:338} & bounded & $\forall H{>}0\,\forall D{\ge}2H$, given finite separation & mobius-mersenne\\
homogeneous Mersenne multiplier does not annihilate & \lean{idCoeff\_mersenneModulus\_does\_not\_annihilate\_tailShift}{TotientTailCarryPeriod.lean:342} & bounded & $\forall K,q\,(K{>}0,K{<}q,q{\mid}2^K{-}1)$ & mobius-mersenne\\
numerator-gap $\Rightarrow$ $\lnot$ hit (\#257-flavored) & \lean{scaleFullTarget\_miss\_of\_lambert\_projected\_num\_gap}{LambertProjectedNumeratorGap.lean:23} & bounded & $\forall H,D$ fixed, given uniform numerator-gap bound & mobius-mersenne\\
periodic-weight Lambert dichotomy & \lean{irrational\_or\_bpow\_mul\_eq\_intCast\_intWeightedErdosSeries\_periodic}{CertificateKernel.lean:14175} & bounded & $\forall b{\ge}2\,\forall m{>}0\,\forall w$ ($w(n{+}m){=}w(n)$) & mobius-mersenne\\
non-negative periodic $\Rightarrow$ irrational (unconditional) & \lean{irrational\_intWeightedErdosSeries\_periodic\_of\_coeff\_nonneg\_of\_frequently\_ne\_zero}{CertificateKernel.lean:14583} & bounded & $\forall b{\ge}2\,\forall m{>}0$, $w$ periodic, $\ge0$, frequently ${\ne}0$ & mobius-mersenne\\

\papertablehead\multicolumn{5}{l}{\emph{-- \scale{fixed} --}}\\
all $h\in[1,8]$ kill certificate at depth 16, $N{=}12$ & \lean{certifiedKill\_all\_small}{TotientTailPeriodKiller.lean:404} & fixed & $\forall h\in[1,8]$, \texttt{certifiedKill}\,$h$\,12\,16 (by \texttt{decide}) & other-binary-window\\
$h\in[1,16]$, exponent 14 kill table & \lean{totient\_series\_ne\_rat\_of\_den\_dvd\_pow\_two\_mul\_mersenne\_upto\_sixteen}{CertificateKernel.lean:18890} & fixed & $\forall h\in[1,16]$, denominator exponent 14 (by \texttt{decide}) & other-binary-window\\
rank-2 kill not shallower at $(h,N){=}(1,8)$ & \lean{totient\_tail\_rank\_two\_kill\_sound\_but\_not\_shallower\_cell}{CertificateKernel.lean:19090} & fixed & measured over $t{\le}20$ (30/40 cells) & other-binary-window\\
28-point diagonal pincer table through $t{=}64$ & \lean{certifiedKill\_diagonal\_all\_imported}{DiagonalPincerCertificates.lean:2870} & fixed & $t\in\{1,2,3,4,5,7,8,9,11,13,16,17,\dots\}$ (28 explicit values) & other-lcm-diagonal\\
contiguous diagonal certificate band through $t{\le}82$ & \lean{ErdosProblems.Skip.LadderT67.exists_diagonalKill_le_82}{ErdosProblems/Skip/LadderT67.lean:71264} & bounded & $\forall t{\le}82\,\exists L$; no $t{=}83$ certificate claimed & other-lcm-diagonal\\
short kill at $a{=}4$ and $a{=}6$ & \lean{lcmDiagonalArithmeticKill\_two\_pow\_four}{TotientActualLcmShortKill.lean:18} & fixed & $a{=}4$, $a{=}6$ (concrete, kernel-checked) & mobius-mersenne\\
square-CRT clean block: vanishing / nonvanishing witnesses & \lean{squareCRT\_clean\_block\_can\_vanish}{SquareCRTCube.lean:459} & fixed & witness $n{=}52$ (vanishes); witness $n{=}27$ (nonzero) & other-crt-dyadic-residue\\
echo-vs-height two-front wall & prose:echo\_versus\_height\_two\_front\_wall & fixed & empirical over tested families only (not universal) & mobius-mersenne (q-Pad\'e, quantitative)\\
period-4 zeroset escape witness & \lean{intWeightedCoeff\_periodFourSignWeight\_eq\_zero\_of\_mod\_four\_eq\_three}{CertificateKernel.lean:14226} & fixed & $\forall n,\,n{\equiv}3\!\!\pmod 4$ (fixed weight $w$) & mobius-mersenne\\

\midrule
\papertablehead\multicolumn{5}{c}{\emph{--- frontier: everything below this line is the cofinal target or a direct consumer of it ---}}\\
\midrule

\papertablehead\multicolumn{5}{l}{\emph{-- \scale{cofinal} (OPEN) --}}\\
\textbf{THE WALL} — certificate supply $\Rightarrow$ Irrational $S$ & \lean{irrational\_totient\_series\_of\_certificate\_supply}{TotientTailPeriodKiller.lean:394} & \textbf{cofinal} & $\forall h{\ge}1\,\forall N_0{\ge}0\,\exists N{\ge}N_0\,\exists L,\ \texttt{certifiedKill}\,h\,N\,L$ & other-binary-window\\
multiple-certificate supply (weakened) & \lean{irrational\_totient\_series\_of\_multiple\_certificate\_supply}{CarrySurvivorExtinction.lean:502} & cofinal & $\forall h_0{>}0\,\forall N_0\,\exists m{>}0\,\exists N{\ge}N_0\,\exists L$ & other-lcm-period-multiple\\
LCM-diagonal one-parameter restatement & \lean{irrational\_totient\_series\_of\_lcm\_diagonal\_certificate\_supply}{LcmDiagonalReduction.lean:112} & cofinal & $\forall t_0\,\exists t{\ge}t_0\,\exists L$ & other-lcm-diagonal\\
LCM-cone window-kill supply (widest target) & \lean{irrational\_totient\_series\_of\_lcm\_cone\_window\_kill\_supply}{CertificateKernel.lean:19014} & cofinal & $\forall t_0\,\exists t{\ge}t_0\,\exists q{>}0\,\exists m\,\exists L$ & other-lcm-cone\\
LCM diagonal/cone nonintegrality restatement & \lean{irrational\_totient\_series\_of\_lcm\_diagonal\_nonintegrality\_supply}{CertificateKernel.lean:19034} & cofinal & $\forall t_0\,\exists t{\ge}t_0$, $\texttt{totientTail}(2H_t){-}\texttt{totientTail}(H_t)\notin\mathbb{Z}$ & other-real-analytic-nonintegrality\\
LCM cone-nonflat supply (sharpest depth-reduced form) & \lean{irrational\_totient\_series\_of\_lcm\_cone\_nonflat\_supply}{CertificateKernel.lean:19140} & cofinal & $\exists$ unbounded-scale menu $Q$ with \texttt{coneNonflatCert} firing & other-lcm-cone-menu\\
first-harmonic norm-gap supply & \lean{irrational\_totient\_series\_of\_first\_harmonic\_norm\_gap}{FirstHarmonicPivot.lean:83} & cofinal & $\forall h{>}0\,\forall X_0\,\exists X{\ge}\max(X_0,1)\,\exists L$, room $16(2X{+}h{+}L{+}2){\le}2^L$ & binary-digit\\
prime-jump sharp-kill supply & \lean{irrational\_totient\_series\_of\_primeJumpSharpKill\_supply}{PrimeJumpWindow.lean:193} & cofinal & $\forall t_0\,\exists t{\ge}t_0\,\exists p,L{>}0$ & binary-digit\\
$a$-exponent short-arithmetic kill supply (the open trigger) & \lean{PowerTwoActualLcmShortArithmeticKillSupply}{TotientActualLcmOrbitArithmetic.lean:2107} & cofinal & $\forall a_0\,\exists a\,L,\ a_0{\le}a\land L{<}2\cdot2^a$ & mobius-mersenne\\
actual-orbit separation supply & \lean{irrational\_totientSeries\_of\_actualLcmOrbitSeparationSupply}{TotientActualLcmOrbitSeparation.lean:305} & cofinal & cofinal $1/32{+}$error separation at guarded odd ranks & mobius-mersenne\\
top-edge-residue-gap sufficiency chain (5 links, weakest first) & \lean{irrational\_totientSeries\_of\_flexibleActualTerminalDominanceSupply}{TotientActualLcmTopEdgeStaircase.lean:2221} & cofinal & $\forall a_0\,\exists a{\ge}a_0\,\exists\langle\text{supply predicate}\rangle$ & mobius-mersenne\\
rational countermodel refutes \{bound,parity,aperiodicity\} & \lean{exists\_totientParity\_arbitrarilyManySeparatedCarry\_rational\_countermodel}{TotientParityCoboundaryCountermodel.lean:637} & cofinal & $\exists c{:}\mathbb{N}\to\mathbb{N}$, $c{\le}6$, $c{\le}n$, $c{\equiv}\varphi\!\!\pmod2$, arbitrarily-separated blocks & binary-digit\\
mod-4 pulse cofinal supply (via CRT+Dirichlet) & \lean{exists\_prime\_deltaTotient\_four\_mul\_mod\_four\_two}{TotientTailCarryPeriod.lean:384} & cofinal & $\forall h{>}0\,\forall B\,\exists\text{prime}\,p{>}B$ & mobius-mersenne\\
mod-4 pulse survivor kill (fourfold reduction) & \lean{irrational\_totientSeries\_of\_cofinal\_modFourPulseSurvivorKill}{TotientTailCarryPeriod.lean:1066} & cofinal & given CP-05's cofinal supply & mobius-mersenne\\
two-adic pulse at arbitrary depth $K$ & \lean{exists\_prime\_totient\_twoAdic\_pulse\_divisors}{TotientTwoAdicPulseBlock.lean:54} & cofinal & $\forall K{\ge}2\,\forall H{>}K\,\forall B\,\exists\text{cofinally many primes}\,p{>}B$ & p-adic\\
two-adic pulse transfer (needs eventual integrality) & \lean{eventual\_integral\_tailDiff\_has\_cofinal\_twoAdic\_half\_pulse}{TotientTwoAdicPulseBlock.lean:327} & cofinal & $\forall K{\ge}2$, given TA-01 pulse $+$ eventual integrality & p-adic\\
\end{longtable}
\end{landscape}
\endgroup

\noindent $^\dagger$ The source normal-form table tags \texttt{actualLcmTailDiff\_shift\_pos} \scale{bounded}
because the window parameter $J$ is finite-range relative to $2\cdot2^a$; the exponent $a$ itself is
completely unbounded ($\forall a\ge 8$), and this row's constants ($4$, $8$, $32$) do not degrade with
$a$. It is flagged here, not silently re-tagged, because it is exactly the kind of row an automated
promotion sweep must not misclassify — see the audit below.

\subsection{Checked negative result: no free promotion from bounded to cofinal}

A ladder like the one above invites an obvious question: is any \scale{bounded} result actually
uniform in disguise, so that removing its threshold is a routine generalisation rather than new
mathematics? This was checked directly rather than assumed. An audit read the Lean proof body
behind every bounded-scale result in the \#249 corpus — eighteen candidates in total, spanning
the near-miss list against the \#249-supply obligation catalogued in the interface index (fourteen
rows, each independently checked against the exact cofinal target it is closest to) together with
the remaining bounded/fixed rows in the certificate, M\"obius--Lambert, and actual-orbit banks
above — and found \textbf{no} bounded result whose proof is uniform enough in its threshold
parameter to promote to cofinal or uniform for free. Every candidate failed for one of three
reasons, and the failure mode is worth recording because it is a checked negative result, not an
absence of search:

\begin{enumerate}[leftmargin=*]
\item \textbf{Finite certificate tables.} Results such as
\lean{certifiedKill\_diagonal\_all\_imported}{DiagonalPincerCertificates.lean:2870} (28 explicit
values of $t$ through $t{=}64$, now a strict subset of the aggregate
$\forall t\le82$ band),
\lean{certifiedKill\_all\_small}{TotientTailPeriodKiller.lean:404} ($h\in[1,8]$ at fixed depth
16), and
\lean{totient\_series\_ne\_rat\_of\_den\_dvd\_pow\_two\_mul\_mersenne\_upto\_sixteen}{CertificateKernel.lean:18890}
($h\in[1,16]$ at exponent 14) are each a finite \emph{list} of independently verified rows, not a
single argument instantiated at a free parameter. There is no uniform proof underneath to strip
the bound from; each new row is a fresh finite computation.
\item \textbf{\texttt{interval\_cases}/\texttt{decide} over a bounded range.} Results such as
\lean{three\_quarters\_mul\_lt\_totient\_of\_rough\_lt\_two\_pow}{TotientActualLcmOrbitArithmetic.lean:768}
($a{\ge}8$) and the entire
\lean{actualLcmTailDiff\_shift\_pos}{TotientActualLcmOrbitSign.lean:39}/top-edge-staircase family
(room conditions of the shape $J{+}(a{+}6){<}2\cdot2^a$) are proved by case-splitting a residue or
a divisor structure that is only exhaustively enumerable inside the stated range; the mathematical
content genuinely narrows as the range widens, so \texttt{decide} cannot simply be re-run at a
larger bound without an exponential blow-up in the search space it certifies.
\item \textbf{Constants that degrade with the parameter.} The clearest instance is the
echo-versus-height wall (prose:echo\_versus\_height\_two\_front\_wall,
\S5e--5f of the ambitious-modular-route note): height-optimal ladder families pay quadratic
height $P(1)\sim 2^{K^2/2}$ matching
\lean{lcmHeight\_scaledMobiusShadow\_den\_exact}{MersenneShadowDenominatorGrowth.lean:147}'s own
growth rate, while height-minimal integer-relation minimizers regrow the echo term $E_R$ to match
— for every tested family, the two costs cannot both be driven to zero as the scale parameter
grows. The interface index records the same phenomenon for
\texttt{HalfRung}$(J)$-shaped truncation rungs: the certificate window $B(J)$ is driven by
$L_J=\mathrm{lcm}(2,\dots,J)$, which grows super-exponentially in $J$, so the per-$J$ constant
degrades with the very parameter the cofinal claim needs to range over — ``the hallmark of a proof
that does not survive bound removal'' (interface index, row \texttt{d-3a/d-3b}).
\end{enumerate}

None of the fourteen near-miss rows against the \#249-supply obligation is closed by a scale
argument alone either: reading each one's own recorded \texttt{exact\_mismatch} shows the residual
gap is a missing \emph{arithmetic input} (a constant-saving Weyl-sum bound for
\texttt{e1-companion}; a one-sided top-edge residue gap for \texttt{TE-04}/\texttt{TE-05-weakest}), never a
generalisable proof technique sitting one \texttt{omega} call away from cofinal. The one row where
the source tag itself is imprecise — \texttt{SGN-01}, marked \scale{bounded} even though its
underlying exponent $a$ is unbounded, see $^\dagger$ above — is precisely the kind of false
positive the audit was built to catch, and even there what remains open is the \emph{other} half of
the certificate band (the top-edge residue), not a scale defect in SGN-01 itself.

\section{Coordinate atlas and transport maps}

\#249 is not proved or disproved in a single representation. The corpus expresses
$S=\sum_{n\ge1}\varphi(n)/2^n$ in at least eight distinct coordinates, and the central
methodological lesson of the whole programme — flagged repeatedly in the source banks as
AP17/AP18, ``obstructions are coordinate-relative'' — is that a mechanism which kills a proof
strategy in one coordinate can be silent, or even actively helpful, in another. This section lists
every coordinate used, gives at least one exact transport map out of it into another coordinate
with its Lean site, and then makes the coordinate-relativity claim precise for the one pair where
it matters most: binary-digit versus M\"obius--Mersenne.

\begin{description}[leftmargin=*,style=nextline]

\item[binary-digit] $S$ as the literal base-2 expansion object: \texttt{totientTail},
\texttt{windowDiscrepancy}, \texttt{certifiedKill}. \emph{Transport to mobius-mersenne:} the
core identity
\[
  S=\sum_{n\ge1}\frac{\varphi(n)}{2^n}=\frac12+\sum_{d\ge1}\frac{\mu(d)}{(2^d-1)^2}
\]
proved by \lean{tsum\_totient\_half\_pow\_eq\_half\_add\_moebius\_sq}{MersenneLambertLadder.lean:665}
(consumed at \lean{mobius\_square\_series\_eq\_partial\_add\_tail}{SquaredMersenneDiagonalEnclosure.lean:94})
turns every binary-digit statement about $S$ into a statement about a squared M\"obius--Mersenne
Lambert series, and back.

\item[mobius-mersenne] $S$ as $\tfrac12+\sum_d\mu(d)/(2^d-1)^2$, the Lambert ladder
$L(w)=\sum_n w(n)/(2^n-1)$ at $w=\mu,1,\varphi$, and the whole cyclotomic denominator-survival
chain built on it. \emph{Transport to cyclotomic:} evaluating the M\"obius-numerator polynomial at
$x=2$,
\lean{mobiusNumeratorPolynomial\_eval\_two}{RepunitMobiusNumerator.lean:445} together with the
definition \lean{mobiusNumerator}{RadicalMobiusShadow.lean:101}, moves the squared-denominator
question into finite cyclotomic-polynomial arithmetic over $\mathbb{Z}$.

\item[probability/coprime-pair] $S-\tfrac12=\Pr(\gcd(X,Y)=1)$ for independent fair-coin waiting
times $X,Y$. \emph{Transport from mobius-mersenne:} the gcd-divisibility factorisation
\lean{tsum\_pos\_pair\_both\_dvd\_half\_eq\_inv\_mersenne\_sq}{GcdMomentCalculus.lean:266}
($\Pr(d\mid X\wedge d\mid Y)=1/(2^d-1)^2$) together with the reduced-direction law
\lean{tsum\_pos\_coprime\_inv\_mersenne\_eq\_one}{GcdMomentCalculus.lean:349} converts the
squared-Lambert sum into a probability mass, and the independently-built
\lean{tsum\_coprime\_pair\_pow\_eq\_tsum\_totient\_mul\_pow}{GeometricCoprimality.lean:120} gives
the same bridge from the totient generating function directly.

\item[cyclotomic] The T1--T6 chain: \texttt{mobiusNumeratorPolynomial}, its coefficients (an exact
divisor-counting formula, \lean{mobiusNumeratorPolynomial\_coeff}{RepunitMobiusNumerator.lean:217}),
cyclotomic-value survival
(\lean{cyclotomicValue\_dvd\_baseMobiusShadow\_den}{CyclotomicProjectionOfShadow.lean:389}), and
prime-power jump dynamics. \emph{Transport to mobius-mersenne (denominator growth):} the upper-half
prime-channel survival theorem
\lean{upperHalfChannel\_product\_dvd\_den\_of\_scale\_primeFactors\_le}{MersenneShadowCyclotomicNoncollapse.lean:809}
feeds directly into the exact denominator-growth formula
\lean{lcmHeight\_scaledMobiusShadow\_den\_exact}{MersenneShadowDenominatorGrowth.lean:147} in the
Möbius--Mersenne coordinate.

\item[p-adic] Fixed and growing 2-adic valuation-unit dynamics: the adelic height/denominator-tax
law, the tropical fixed-precision no-go, the signed-Hankel terminal-parity engine, the two-adic
pulse-block construction. \emph{Transport to farey:} the scalar-localisation complement-dvd law
\lean{scalarLocalization\_complement\_dvd}{AdelicHeightObstruction.lean:23} together with the
Mersenne-height corollary \lean{positiveRat\_mersenne\_height}{AdelicHeightObstruction.lean:103}
($2^r\mid|x.\mathrm{num}|\wedge x<2/(2^n-1)\Rightarrow 2^r(2^n-1)<2\cdot x.\mathrm{den}$) is exactly
a denominator lower bound of the same shape the Farey coordinate proves directly.

\item[farey] Unimodular denominator exclusion:
\lean{farey\_gap}{GapFareyBound.lean:51} and the committed $K{=}240$ window certificate
\lean{gap\_check\_window\_1\_240\_le\_79639646646701375323355774875831053}{GapFareyBound.lean:176}
giving $S\ne p/q$ for every $q\le Q_0\approx7.96\times10^{34}$. \emph{Transport to binary-digit:}
the window bound is built directly from a committed $2^{240}$-scale totient residue
$V$ for the window $(N,K)=(1,240)$ — a binary-digit object — via
\lean{gap\_check\_window\_1\_240\_first\_failure}{GapFareyBound.lean:225}, which shows the bound
is sharp at exactly the mediant $b_K+d_K$.

\item[lcm-orbit] The \texttt{periodLcm}/LCM-diagonal/LCM-cone family:
\lean{periodLcm}{CarrySurvivorExtinction.lean:515},
\lean{eq\_prime\_pow\_of\_not\_dvd\_periodLcm}{LcmDiagonalReduction.lean:137}, the multiplicative ray
split, and the cone-flatness/cone-nonflat producers. \emph{Transport to binary-digit:} the identity
\lean{lcmRayArithmeticLetter\_eq\_deltaTotient}{TotientActualLcmOrbitArithmetic.lean:298} and
\lean{diagonalWindowIncrement\_eq\_lcmRayArithmeticLetter}{TotientActualLcmOrbitArithmetic.lean:310}
identifies the LCM-ray letter literally with a \texttt{deltaTotient}/\texttt{windowDiscrepancy}
value, so every LCM-orbit statement is, term for term, a binary-digit statement in disguise.

\item[first-harmonic/exponential-sum] Weyl-sum block cancellation over the totient window
discrepancy: \lean{irrational\_totient\_series\_of\_first\_harmonic\_norm\_gap}{FirstHarmonicPivot.lean:83}
and the unconditional companion at
\lean{exists\_certifiedKill\_of\_first\_harmonic\_gap}{FirstHarmonicGap.lean:128} (any
constant-saving cancellation, via $\cos(\pi/8)>9/10$ plus pigeonhole, forces a finite kill
certificate). \emph{Transport to binary-digit:} by construction the exponential sum is literally
$\sum_N \cos\bigl(2\pi\cdot(\texttt{windowDiscrepancy}\ h\ N\ L \bmod 2^L)/2^L\bigr)$ — a first-harmonic
statistic \emph{of} the binary-digit residue, so a norm-gap in this coordinate is definitionally a
statement about \texttt{windowDiscrepancy}.

\end{description}

\subsection{Obstructions are coordinate-relative: the binary-digit vs.\ M\"obius--Mersenne case}

The clearest instance of coordinate-relative obstruction in the whole \#249 corpus is the fate of
the classical Erd\H{o}s (1948) near-integer digit method. That method needs, in essence, a finite
automaton: a bounded amount of state carried forward from one digit block to the next, so that the
value of a far-away digit block can be recovered from a bounded summary of everything before it. In
the binary-digit coordinate this is exactly what fails, and it fails as a proved Lean theorem, not
a heuristic remark: \lean{balancedPulse\_no\_autonomous\_decoder}{GenericTailOrbitRigidity.lean:247}
shows that for a balanced-pulse family of radius $m$, any finite \texttt{State} type collapsing the
family to a single autonomous summary must satisfy $|\mathrm{State}|\ge\lfloor m/2\rfloor+2$ —
unbounded in $m$ — so no finite-state decoder can recover the shift $r$ from its collapsed state
(\ev{Lean}, \scale{uniform}, coordinate \texttt{binary-digit}).

This is the formal shadow of the
elementary facts that $0\le\varphi(n)\le n$, that $\varphi$ is unbounded, and
that $\varphi$ has average order $\tfrac6{\pi^2}n$, equivalently
$\sum_{k\le x}\varphi(k)\sim\tfrac3{\pi^2}x^2$ (\ev{Math}),
rather than any false pointwise estimate $\varphi(n)=\Theta(n)$: the digit-window discrepancy
\texttt{windowDiscrepancy}$\,h\,N\,L$ genuinely needs $\Omega(m)$ bits of state to track as $N$
grows, which is exactly why \#249's own binary-digit engine is forced into a \emph{cofinal
certificate supply} (\lean{irrational\_totient\_series\_of\_certificate\_supply}{TotientTailPeriodKiller.lean:394})
rather than a single finite-state argument: the same shape of no-go recurs at
\lean{det\_residualMonomialMatrix}{ResidualGaugeObstruction.lean:47} (locked gauge defeats
residual-blind rank certificates for the first-harmonic pivot) and at
\lean{linearIndependent\_canonicalTotientKernelFamily}{TotientMahlerDefect.lean:935} (the dyadic
totient-kernel span is genuinely infinite-rank, not compressible to any fixed dimension).

The M\"obius--Mersenne coordinate does not carry this obstruction, because it is built on a
coefficient sequence, $\mu(d)\in\{-1,0,1\}$, that is bounded rather than growing with $d$. The
squared-Lambert transfer engine
\lean{tsum\_lambert\_linear\_weight\_sq\_pure}{GcdMomentCalculus.lean:105} takes any weight with
$|w(d)|\le d$ — $\mu$ trivially qualifies with room to spare — and converts it into a
divisor-convolution power series for free; no finite-state decoder is ever needed because the
relevant recursion (the cyclotomic T1--T6 chain: \lean{mobiusNumeratorPolynomial\_eq\_gcdWord}{RepunitMobiusNumerator.lean:205}
through \lean{cyclotomic\_dvd\_primeJump\_new\_fibre}{PrimePowerJumpDynamics.lean:281}) is a finite,
exact polynomial identity at every squarefree radical $r$, uniform in $r$, with no growth-driven
state explosion.

This is precisely why the corpus's strongest \emph{unconditional} denominator
result — the exact-value and lower-bound formulas of
\lean{lcmHeight\_scaledMobiusShadow\_den\_exact}{MersenneShadowDenominatorGrowth.lean:147} and
\lean{upperHalfMersenneProduct\_lower\_bound}{MersenneShadowDenominatorGrowth.lean:60} — lives
entirely in the M\"obius--Mersenne coordinate and is scale-\texttt{uniform}, while the sharpest
\emph{conditional} route in the binary-digit coordinate is scale-\texttt{cofinal} and unclosed. The
two coordinates are not equivalent representations of the same difficulty: the wall each one hits
is a fact about that coordinate's own state-growth, not about \#249 itself, which is exactly the
AP17/AP18 discipline this corpus was built to enforce — before reporting an obstruction, name the
coordinate it was measured in, and test whether a re-representation removes it.
% ---- part a249_p3 ----
\section{Frontier index: what stands between the corpus and a proof}

This section is the join layer of the paper. Parts I--II inventory what the corpus proves; this
part inventories, obligation by obligation, exactly what is \emph{missing} to turn each proved
reduction into a proof of Erd\H{o}s \#249, in a form usable as a premise without opening the Lean
tree. Erd\H{o}s \#249 asks whether $S = \sum_{n\ge 0} \varphi(n)/2^n$ is irrational. It is OPEN.
Nothing in this section decides it. Every row below either (a) states a proved theorem exactly, or
(b) states an unproved but well-posed proposition (an \emph{open producer} or \emph{open supply})
together with the exact proved consumer that would turn it into irrationality of $S$.

\subsection{Reading this index}

Every entry follows the same discipline. \textbf{Yields}: the exact statement proved, with
quantifiers in their proved order. \textbf{Exact mismatch}: precisely where the proved statement
falls short of the open obligation it is compared against --- never ``it doesn't work,'' always the
named axis of shortfall. \textbf{What would close it}: the exact remaining proposition, stated in
full, whose proof (via the cited Lean consumer, already on disk) proves $\mathrm{Irrational}(S)$.

Each numbered claim carries three tags. \scale{fixed} means proved at finitely many explicit
numerals with no argument in the growing parameter; \scale{bounded} means proved on a bounded range
of a parameter that is itself unbounded elsewhere in the statement; \scale{uniform} means proved for
\emph{every} value of the relevant parameter (the strongest positive tag short of matching the
open target); \scale{cofinal} means the statement itself has the shape \m{\forall a_0\ \exists a\ge
a_0}, i.e.\ it is exactly the quantifier shape the open obligation needs; \scale{n/a} applies to
structural non-existence results with no growing parameter. The evidence tags are
\ev{Lean}, \ev{Cert}, \ev{Math}, \ev{Cited}, \ev{Open}, used exactly as defined in Part I.

The canonical open obligation for \#249, proved in Part II to be sufficient
(\lean{irrational_totient_series_of_certificate_supply}{TotientTailPeriodKiller.lean:394}), is the
\textbf{249-supply} obligation:
\[
  \forall h \ge 1,\ \forall N_0,\ \exists N \ge N_0,\ \exists L,\ \mathtt{certifiedKill}\ h\ N\ L .
\]
Fourteen near-miss rows are catalogued against it below: one headline row (the sharpest reduction
in the corpus, given its own subsection) and thirteen further rows, each attacking the same
obligation from a different coordinate.

\subsection{The headline: \#249 reduces to a single unsupplied exponential-sum bound}
\label{ssec:headline}

The following two definitions, verified against the live Lean tree in this session, are the
sharpest statement of the \#249 frontier that the corpus contains.

\begin{defn}[Window discrepancy]
\lean{windowDiscrepancy}{TotientTailPeriodKiller.lean:66}. For $h,N,L\in\mathbb N$,
\[
  A_{h,N,L} \;=\; \sum_{j=0}^{L-1} \bigl(\varphi(N+h+1+j) - \varphi(N+1+j)\bigr)\cdot 2^{\,L-1-j}
  \ \in \mathbb Z .
\]
This is the depth-$L$ truncation of $2^L\cdot(R_{N+h}-R_N)$, where
$R_N = \sum_{j\ge 1}\varphi(N+j)/2^j$ is the local totient tail
(\lean{totientTail}{TotientTailPeriodKiller.lean:57}).
\end{defn}

\begin{defn}[Certified kill]
\lean{certifiedKill}{TotientTailPeriodKiller.lean:72}.
\[
  \mathtt{certifiedKill}\ h\ N\ L \;:\Leftrightarrow\;
  (N{+}h{+}L{+}2) < A_{h,N,L}\bmod 2^L < 2^L - (N{+}h{+}L{+}2).
\]
That is: the residue of $A_{h,N,L}$ modulo $2^L$ avoids the radius-$(N{+}h{+}L{+}2)$ neighbourhood
of $0$. \lean{certifiedKill_depth_floor}{TotientTailPeriodKiller.lean:79} records the necessary
room condition this forces: $2(N{+}h{+}L{+}2) < 2^L$.
\end{defn}

\begin{defn}[First-harmonic real and complex characters]
\lean{windowFirstCos}{FirstHarmonicGap.lean:27} and
\lean{windowFirstExp}{FirstHarmonicPivot.lean:38}:
\[
  \mathtt{windowFirstCos}\ h\ N\ L = \cos\!\Bigl(2\pi\,\tfrac{A_{h,N,L}\bmod 2^L}{2^L}\Bigr),\qquad
  \mathtt{windowFirstExp}\ h\ N\ L = \exp\!\Bigl(i\,2\pi\,\tfrac{A_{h,N,L}\bmod 2^L}{2^L}\Bigr).
\]
\lean{windowFirstExp_re}{FirstHarmonicPivot.lean:42} records $\mathrm{Re}(\mathtt{windowFirstExp})
= \mathtt{windowFirstCos}$ and \lean{norm_windowFirstExp}{FirstHarmonicPivot.lean:48} that
$\|\mathtt{windowFirstExp}\,h\,N\,L\| = 1$: this is literally the first additive character of the
endpoint discrepancy modulo $2^L$, a standard exponential-sum object.
\end{defn}

\begin{thm}[Real-part certificate consumer]
\label{thm:hgap-real}
\lean{exists_certifiedKill_of_first_harmonic_gap}{FirstHarmonicGap.lean:128}. \ev{Lean}.
\scale{uniform}. \coord{first-harmonic}.
For all $h,X,L$ with $0<X$ and the room condition $16(2X{+}h{+}L{+}2)\le 2^L$, if
\[
  \sum_{N=X}^{2X-1} \mathtt{windowFirstCos}\ h\ N\ L \;\le\; \tfrac{9}{10}\,X ,
\]
then $\exists N\in[X,2X)$ with $\mathtt{certifiedKill}\ h\ N\ L$.
\end{thm}

\begin{thm}[Subset consumer --- no density or partition hypothesis]
\label{thm:hgap-subset}
\lean{exists_certifiedKill_of_first_harmonic_gap_subset}{FirstHarmonicGap.lean:104}.
\ev{Lean}. \scale{uniform}. \coord{first-harmonic}.
For any nonempty finite $T\subseteq\mathbb N$ with $T\subset[0,2X)$ and the same room condition, if
\[
  \sum_{N\in T} \mathtt{windowFirstCos}\ h\ N\ L \;\le\; \tfrac{9}{10}\,|T| ,
\]
then $\exists N\in T$ with $\mathtt{certifiedKill}\ h\ N\ L$. This is strictly stronger than
Theorem~\ref{thm:hgap-real}: $T$ can be \emph{any} explicitly chosen nonempty finite subset of the
dyadic block, not the whole block, and the proof (an averaging pigeonhole, \S below) never uses
that $T$ has positive density or comes from a partition.
\end{thm}

\begin{thm}[Complex norm consumer and the open producer]
\label{thm:hgap-norm}
\lean{exists_certifiedKill_of_first_harmonic_norm_gap}{FirstHarmonicPivot.lean:53} and
\lean{irrational_totient_series_of_first_harmonic_norm_gap}{FirstHarmonicPivot.lean:83}.
\ev{Lean} for both implications; the hypothesis \m{DTWFirstHarmonicNormGap} is \ev{Open}.
The complex norm bound is strictly stronger than the real-part bound
(\m{|z|\ge\mathrm{Re}(z)}, and $21/25 < 9/10$ absorbs the slack), so it composes through
Theorem~\ref{thm:hgap-real} to the same certificate. Define
\[
  \mathtt{DTWFirstHarmonicNormGap} :\Leftrightarrow\;
  \forall h{>}0\ \forall X_0\ \exists X, L,\ \max(X_0,1)\le X\ \wedge\ 16(2X{+}h{+}L{+}2)\le 2^L\ \wedge
\]
\[
  \Bigl\|\ \sum_{N=X}^{2X-1}\mathtt{windowFirstExp}\ h\ N\ L\ \Bigr\| \;\le\; \tfrac{21}{25}\,X .
\]
\coord{first-harmonic}. \scale{cofinal} (this is the open target itself). Then
\[
  \mathtt{DTWFirstHarmonicNormGap} \;\Longrightarrow\; \mathrm{Irrational}\Bigl(\sum_{n\ge 0}
  \tfrac{\varphi(n)}{2^n}\Bigr),
\]
proved in full, with no further gap, by chaining
Theorem~\ref{thm:hgap-norm}$\to$Theorem~\ref{thm:hgap-real}$\to$
\lean{irrational_totient_series_of_certificate_supply}{TotientTailPeriodKiller.lean:394}.
\end{thm}

\begin{obs}[The exact negative]
\label{obs:no-instance}
No instance of \m{\mathtt{DTWFirstHarmonicNormGap}}'s hypothesis, nor of \m{hgap} in
Theorems~\ref{thm:hgap-real}--\ref{thm:hgap-subset}, is proved anywhere in the corpus, at any $h$,
$X$, $L$. \ev{Open} is exact, not conservative rounding. The two consumer theorems that
\emph{use} the subset form (Theorem~\ref{thm:hgap-subset}) are conditional consumers, not a
supply: \lean{AdjacentPhaseSeparation.lean}{AdjacentPhaseSeparation.lean:267} instantiates $T$ as
the explicit two-element set $\{N,M\}$ (\texttt{apply exists\_certifiedKill\_of\_first\_harmonic\_gap\_subset (\{N, M\} : Finset Nat)}). The two call sites in
\lean{PivotAntiReconstruction.lean}{PivotAntiReconstruction.lean:882} and
\lean{PivotAntiReconstruction.lean}{PivotAntiReconstruction.lean:1572} use, respectively, an
arbitrary supplied finite set $T$ and the explicit singleton $\{N\}$; neither call site supplies
such a set or discharges its harmonic-gap hypothesis. There is no theorem anywhere in the corpus
that discharges $hgap$ at a single instance, let alone cofinally.
\end{obs}

\begin{rem}[Why this is the sharpest reduction, and why it is tractable]
This reduction converts the open producer for \#249 from certificate-supply language --- ``exhibit
a certified kill $(N,L)$'' --- into a \emph{constant-saving cancellation estimate for the first
additive character of the totient window discrepancy over a dyadic block}: exactly the shape of a
classical exponential-sum bound in analytic number theory (a Weyl-type estimate for
$\sum_N e(\theta_N)$). Three features make this genuinely the narrowest gap in the corpus.

First, the saving required is a \emph{constant}, not $o(X)$: $21/25$ (complex norm) or $9/10$
(real part), fixed independent of $h$, $X$, $L$. Any nontrivial cancellation estimate --- even one
far weaker than square-root cancellation --- would suffice; this is not asking for GRH-strength
input.

Second, the room condition $16(2X{+}h{+}L{+}2)\le 2^L$ forces only $L \gtrsim \log_2 X + 5$, which
is satisfiable for every $X$ by simply taking $L$ large enough (the depth floor grows logarithmically
in $X$, not polynomially); no scale obstruction of the kind documented elsewhere in this paper
(e.g.\ the $\sqrt{2\log_2 N}$ silent-position defect for \#257, Part~I) stands in the way of the
producer's own room requirement.

Third, the subset form (Theorem~\ref{thm:hgap-subset}) removes even the requirement that the
cancellation happen on the whole dyadic block or a positive-density subset of it: any single
explicitly named finite $T$ with the averaged bound would do. In particular a supplier-fibre or
sparse-subsequence argument --- e.g.\ picking $N$ to be one less than a prime, as the exact pivot
algebra in \lean{FirstHarmonicPivot.lean}{FirstHarmonicPivot.lean} already sets up via
\texttt{pivotArgument}, \texttt{pivotPrime}, \texttt{pivotSupplierBases} --- is a legitimate route
to this obligation and does not need to control the discrepancy on a full interval.
\end{rem}

\subsection{Row-by-row: the remaining near misses against 249-supply}

The headline (\S\ref{ssec:headline}) is row \textbf{e1-companion} of the source interface catalogue
in full. The remaining thirteen rows below attack the same 249-supply obligation from independent
coordinates in the Lean tree: the actual-LCM sign corridor, the top-edge staircase, the raw
rational approximant, the short-window arithmetic word, the diagonal pincer certificate bank, the
$h$-uniform single-window certificate, the Farey/continued-fraction denominator growth law, the
2-adic pulse block, the LCM-jump slack scalar, the prime-jump commutator, the M\"obius--Mersenne
denominator channel, and the carry-kernel rank. None of them is closed; each is recorded with its
exact remaining content.

\subsubsection{SGN-01 --- the positive-sign half of the certified-kill band (\coord{actual-lcm-sign})}

\begin{defn}
The \emph{actual LCM tail orbit} at exponent $a$ is
$\mathtt{actualLcmTailOrbit}\ a = R_{2H+2H} - R_{H+... }$, precisely
$\mathrm{totientTail}(2H)-\mathrm{totientTail}(H)$ where $H = \mathtt{periodLcm}(2^a)$
(\lean{actualLcmTailOrbit_pos}{TotientActualLcmOrbitSign.lean:144}).
\end{defn}

\begin{thm}
\lean{actualLcmTailDiff_shift_pos}{TotientActualLcmOrbitSign.lean:39} and
\lean{actualLcmTailOrbit_pos}{TotientActualLcmOrbitSign.lean:144}. \ev{Lean}. \scale{uniform}.
For every $a\ge 8$ and every $J$ with $J+(a{+}6) < 2\cdot 2^a$,
\[
  0 \;<\; \mathrm{totientTail}(2H{+}J) - \mathrm{totientTail}(H{+}J), \qquad H=\mathtt{periodLcm}(2^a).
\]
Hence $0 < \mathtt{actualLcmTailOrbit}\ a$ for every $a\ge 8$. The proof is genuinely uniform in
$a$: it rests on two facts proved for all $a\ge 8$ by a two-case structural split (divisor letter
vs.\ foreign prime power), with absolute constants $4,8,32$ and no lookup table ---
\lean{lcmRayArithmeticLetter_pos_of_lt_two_mul}{TotientActualLcmOrbitArithmetic.lean:1703} and
\lean{periodLcm_lt_eight_mul_t_mul_lcmRayArithmeticLetter}{TotientActualLcmOrbitArithmetic.lean:1906}.
\end{thm}

\begin{rem}[Exactly half the certificate band, and where the other half goes]
$\mathtt{certifiedKill}$ needs the residue to avoid a \emph{symmetric} radius-$(N{+}h{+}L{+}2)$
neighbourhood of $0$ on both sides. SGN-01 gives positivity only. Its companion
\lean{actualLcm_integral_forces_topEdgeResidue}{TotientActualLcmOrbitSign.lean:211} proves that,
under integrality, the consequence is not a kill but the opposite: the residue is forced to the
\emph{top edge}, exactly $2^K - e$ where $e$ is the true carry orbit (\lean{carryOrbit}{CarrySurvivorExtinction.lean:387}). So the
lower half of the certified-kill band is discharged unconditionally and cofinally by SGN-01; the
upper half is untouched, and what remains open is a genuinely different, one-sided statement about
the top edge --- row TE-04 below. The corpus's own normal-form tables tag this row
\scale{bounded}; that tag is misleading, since only the window offset $J$ is bounded relative to
the height, and the height $2\cdot 2^a$ itself is unbounded: the correct tag is \scale{uniform}.
\end{rem}

\subsubsection{TE-04 --- the one-sided top-edge staircase (\coord{actual-lcm-top-edge})}

This is the sharpest \emph{landed weakening} of the 249-supply obligation in the corpus: the
symmetric certified-kill band is replaced by a single upper inequality, halving the information
the open producer must supply.

\begin{thm}[Consumers, fully proved]
\lean{actualLcmTailDiff_notMem_int_of_topEdgeResidueGap}{TotientActualLcmTopEdgeStaircase.lean:1260},
\lean{actualLcmTailOrbit_notMem_int_of_topEdgeResidueGap}{TotientActualLcmTopEdgeStaircase.lean:1314},
\lean{irrational_totientSeries_of_topEdgeResidueGapSupply}{TotientActualLcmTopEdgeStaircase.lean:2184}.
\ev{Lean}. \scale{uniform}. For every $a\ge 8$ and every $J,K,m$ inside the sign corridor
($J{+}K{+}(a{+}6) < 2\cdot 2^a$), a one-sided residue gap at precision $m\le K$ ---
room $2H{+}J{+}K{+}2 < 2^m$ and $A_{H,H+J,K}\bmod 2^m \le 2^m - (2H{+}J{+}K{+}2)$ --- already
forces $\mathrm{totientTail}(2H{+}J) - \mathrm{totientTail}(H{+}J)\notin\mathbb Z$. No lower margin
at all is demanded. The proof chain is complete: the theorem holds for every $a\ge 8$, and
\emph{\m{\mathtt{PowerTwoActualLcmTopEdgeResidueGapSupply}} $\Rightarrow$ Irrational $S$} is proved.
\end{thm}

\begin{defn}[The open producer]
\lean{PowerTwoActualLcmTopEdgeResidueGapSupply}{TotientActualLcmTopEdgeStaircase.lean:1325}.
\ev{Open}. \scale{cofinal}.
\[
  \forall a_0\ \exists a,K,m,\ a_0\le a\ \wedge\ 8\le a\ \wedge\ K{+}(a{+}6)<2\cdot 2^a\ \wedge\
  \mathtt{ActualLcmTopEdgeResidueGap}\ a\ 0\ K\ m .
\]
No unconditional statement anywhere in the corpus locates the residue of the diagonal word
$A_{H,H,K}$ (where $H=\mathtt{periodLcm}(2^a)$) below the top strip at even one large $a$.
\lean{actualLcmTopEdgeResidueGap_iff_terminal}{TotientActualLcmTopEdgeStaircase.lean:909} further
reduces the condition to the last $m$ arithmetic letters of the word alone.
\end{defn}

\begin{prop}[The five strictly weaker sufficient links --- proving any one closes \#249]
\label{prop:te-chain}
All five are proved sufficient for
\lean{PowerTwoActualLcmTopEdgeResidueGapSupply}{TotientActualLcmTopEdgeStaircase.lean:1325} in the
same module, and each is \emph{strictly weaker} in the sense that it drops a hypothesis, widens a
band, or asks only for magnitude rather than a signed inequality.
\begin{enumerate}[leftmargin=*]
  \item \lean{PowerTwoAdjacentSuffixMidbandSupply}{TotientActualLcmTopEdgeStaircase.lean:1334}
    (\ev{Open}, \scale{cofinal}):
    \[
      \forall a_0\ \exists a,m,\ a_0\le a \wedge 8\le a \wedge m{+}1{+}(a{+}6)<2\cdot 2^a \wedge
      2H{+}m{+}3 < 2^m \wedge
    \]
    \[
      2H{+}m{+}2 \le \mathtt{diagonalAdjacentSuffixResidue}(2^a)\,0\,m \le 2^m - (2H{+}m{+}2)
    \]
    (a two-sided band on the adjacent-suffix residue directly, one candidate depth $m$).
  \item \lean{PowerTwoOddGuardTopEdgeHalfWordBandSupply}{TotientActualLcmTopEdgeStaircase.lean:1349}
    (\ev{Open}, \scale{cofinal}): at the odd guarded depth $2q{+}1$, a two-sided band of half-width
    $H{+}q{+}2$ on the half-word residue mod $4^q$ --- ``substantially weaker than the older fixed
    $1/32$ central band.''
  \item \lean{PowerTwoActualFinalTopEdgeMagnitudeSupply}{TotientActualLcmTopEdgeStaircase.lean:1471}
    (\ev{Open}, \scale{cofinal}), \emph{proved equivalent} to the previous one via
    \lean{topEdgeHalfWordBandSupply_iff_actualFinalCenteredMagnitudeSupply}{TotientActualLcmTopEdgeStaircase.lean:1479}:
    \[
      \forall a_0\ \exists a,q,\ \max(14,a_0)\le a \wedge \mathtt{oddGuardedCanonicalAdjacentSuffixDepth}(2^a)=2q{+}1
      \wedge H{+}q{+}2 \le |\mathtt{actualOddHalfCenteredLift}\ a\ q| .
    \]
  \item \lean{PowerTwoFlexibleActualTopEdgeMagnitudeSupply}{TotientActualLcmTopEdgeStaircase.lean:1927}
    (\ev{Open}, \scale{cofinal}): the same magnitude bound with the depth restriction relaxed from
    ``canonical guarded'' to any odd $2q{+}1$ satisfying the half-cell fit
    $2(H{+}q{+}2)\le 4^q$ and the sign-corridor room $2q{+}2{+}(a{+}6)<2\cdot 2^a$.
  \item \lean{PowerTwoFlexibleActualTerminalDominanceSupply}{TotientActualLcmTopEdgeStaircase.lean:1908}
    --- the weakest of all five; treated separately as row TE-05-weakest below.
\end{enumerate}
Proving \emph{any one} of these five closes Erd\H{o}s \#249; the corpus has already proved all the
implications from each to
\lean{PowerTwoActualLcmTopEdgeResidueGapSupply}{TotientActualLcmTopEdgeStaircase.lean:1325} (chain:
midband $\to$ residue-gap at \lean{}{TotientActualLcmTopEdgeStaircase.lean:2109}; half-word-band
$\to$ midband at \lean{}{TotientActualLcmTopEdgeStaircase.lean:2058}; final-magnitude
$\Leftrightarrow$ half-word-band at \lean{}{TotientActualLcmTopEdgeStaircase.lean:1479};
flexible-magnitude $\to$ midband at \lean{}{TotientActualLcmTopEdgeStaircase.lean:2008};
final-magnitude $\to$ flexible-magnitude at \lean{}{TotientActualLcmTopEdgeStaircase.lean:1993}).
\end{prop}

\subsubsection{TE-05-weakest --- one scalar inequality at cofinally many odd ranks (\coord{actual-lcm-top-edge})}

This is the weakest landed link in the whole \#249 sufficiency lattice: a single comparison
between one terminal arithmetic letter and twice a centred lift.

\begin{defn}
\lean{PowerTwoFlexibleActualTerminalDominanceSupply}{TotientActualLcmTopEdgeStaircase.lean:1908}.
\ev{Open}. \scale{cofinal}. Writing $H=\mathtt{periodLcm}(2^a)$,
\[
  \forall a_0\ \exists a,q,\ a_0\le a \wedge 8\le a \wedge 2q{+}2{+}(a{+}6)<2\cdot 2^a \wedge
  2(H{+}q{+}2)\le 4^q \wedge
\]
\[
  \mathtt{diagonalWindowIncrement}(2^a)(2q{+}2) \;\le\; 2\cdot \mathtt{actualOddHalfCenteredLift}\ a\ q .
\]
\end{defn}

\begin{thm}[Consumer]
\lean{actualLcmTailOrbit_notMem_int_of_actualTerminalDominance}{TotientActualLcmTopEdgeStaircase.lean:1880}
and
\lean{irrational_totientSeries_of_flexibleActualTerminalDominanceSupply}{TotientActualLcmTopEdgeStaircase.lean:2221}.
\ev{Lean}. \scale{uniform}. The dominance hypothesis at a single odd rank already excludes
integrality of \m{\mathtt{actualLcmTailOrbit}\ a}, and the supply predicate composes to
\m{\mathrm{Irrational}(S)}.
\end{thm}

\begin{obs}[The exact identity, and the internal tension it exposes]
\label{obs:staircase-tension}
\lean{two_mul_actualOddHalfCenteredLift_eq_terminal_sub_trueCarry}{TotientActualLcmTopEdgeStaircase.lean:1678}.
\ev{Lean}. \scale{uniform}. For $a\ge 8$, room $2q{+}1{+}1{+}(a{+}6)<2\cdot2^a$, half-cell fit
$2(H{+}q{+}2)\le 4^q$, and any integral representative $z$ with $z = \mathtt{actualLcmTailOrbit}\ a$:
\[
  2\cdot \mathtt{actualOddHalfCenteredLift}\ a\ q
  \;=\; \mathtt{diagonalWindowIncrement}(2^a)(2q{+}2) - \mathtt{carryOrbit}\ H\ H\ z\ (2q{+}1) .
\]
This is an \emph{exact equality}, not a bound. Comparing it against the dominance inequality above,
\emph{the dominance inequality is literally equivalent to
$\mathtt{carryOrbit}\,H\,H\,z\,(2q{+}1)\le 0$}.

\par\medskip

But SGN-02
(\lean{actualLcm_trueEndpointSurvivor_neg}{TotientActualLcmOrbitSign.lean:172}, \ev{Lean},
\scale{uniform}) proves that under integrality the true carry orbit is \emph{strictly positive}
throughout this exact corridor: for $a\ge 8$ and $J{+}K{+}(a{+}6)<2\cdot 2^a$,
\[
  0 < \mathtt{carryOrbit}\ H\ (H{+}J)\ d\ K
\]
whenever $d$ is the real-valued representative of the translated tail difference.

\par\medskip

So the two branches of the corridor-escape disjunction that
\lean{actualLcmTailOrbit_notMem_int_of_actualTerminalCarryCorridorEscape}{TotientActualLcmTopEdgeStaircase.lean:1824}
offers are \emph{not} symmetric: SGN-02 has already eliminated the branch that the identity most
naturally supplies (dominance, $\mathtt{carryOrbit}\le 0$), and the surviving branch is the
\emph{lower} escape
\[
  2\cdot \mathtt{actualOddHalfCenteredLift}\ a\ q \;\le\;
  \mathtt{diagonalWindowIncrement}(2^a)(2q{+}2) - (2H{+}2q{+}3) .
\]

\par\medskip

This is stated precisely so the reader can check it: it is neither a refutation of the dominance
route (SGN-02 constrains the sign of $\mathtt{carryOrbit}$, it does not itself bound
$\mathtt{diagonalWindowIncrement}$ or $\mathtt{actualOddHalfCenteredLift}$ against $0$, so dominance
is not shown \emph{false}, only shown to entail a carry-orbit sign the census never violates) nor a
proof (no theorem excludes dominance outright). It is the precise reason the finite census keeps
landing just inside the corridor rather than outside it, and it identifies the lower-escape branch
as the coordinate where a producer is actually needed.
\end{obs}

\begin{prop}[What would close it]
Either (i) the lower-escape branch cofinally --- as displayed just above --- or (ii) the two-sided
magnitude form \lean{PowerTwoFlexibleActualTopEdgeMagnitudeSupply}{TotientActualLcmTopEdgeStaircase.lean:1927}
(item 4 of Proposition~\ref{prop:te-chain}), which asks only $H{+}q{+}2 \le
|\mathtt{actualOddHalfCenteredLift}\ a\ q|$ and which the file proves
(\lean{flexibleActualTerminalCarryCorridorEscapeSupply_of_magnitude}{TotientActualLcmTopEdgeStaircase.lean:1938})
implies corridor escape via a clean sign split (positive branch escapes above the terminal letter,
negative branch escapes below the directed bound), sidestepping the SGN-02 tension entirely because
it does not commit to a sign for the centred lift.
\end{prop}

\subsubsection{SEP-02 --- separation from an explicit rational approximant (\coord{raw-approximant})}

\begin{thm}
\lean{abs_actualLcmTailOrbit_sub_rawApprox_lt}{TotientActualLcmOrbitSeparation.lean:141} and
\lean{irrational_totientSeries_of_actualLcmOrbitSeparationSupply}{TotientActualLcmOrbitSeparation.lean:305}.
\ev{Lean}. \scale{uniform}. Unconditionally, for every $a$ and $q$,
\[
  \bigl|\, \mathtt{actualLcmTailOrbit}\ a - \mathtt{actualLcmRawApprox}\ a\ q \,\bigr|
  \;<\; \frac{4H + 2(2q{+}1) + 4}{2^{2q+2}} ,
\]
where $\mathtt{actualLcmRawApprox}\ a\ q$ is an explicit finite computable rational block. The
error radius is uniform in $a$ and $q$ and decays like $H/4^q$: the analytic tail is fully
discharged.
\end{thm}

\begin{rem}
Consequently a cofinal $1/32$-separation of the raw approximant from every integer suffices for
\#249. The corpus verifies the analogous kill \emph{unconditionally} at exactly two exponents,
$a=4$ and $a=6$
(\lean{actualLcmTailOrbit_four_and_six_notMem_int}{TotientActualLcmShortKill.lean:35}, \ev{Cert},
\scale{fixed}, via \lean{certifiedKill_diagonal_t16}{DiagonalPincerCertificates.lean:2842} and
\lean{certifiedKill_diagonal_t64}{DiagonalPincerCertificateT64Endpoint.lean:1928}). Nothing beyond
$a=6$ is proved. The depth $q$ is not free: the supply predicate pins $2q{+}1 =
\mathtt{oddGuardedCanonicalAdjacentSuffixDepth}(2^a)$, so per scale $a$ there is exactly one
admissible depth, not a search over depths.
\end{rem}

\begin{defn}[What would close it]
$\mathtt{PowerTwoActualLcmOrbitSeparationSupply}$: \ev{Open}, \scale{cofinal}.
\[
  \forall a_0\ \exists a\ge\max(2,a_0)\ \exists q,\
  \mathtt{oddGuardedCanonicalAdjacentSuffixDepth}(2^a)=2q{+}1 \wedge
\]
\[
  \forall z\in\mathbb Z,\ \tfrac1{32} + \mathtt{actualLcmRawErrorRadius}\ a\ q \;\le\;
  |\mathtt{actualLcmTailOrbit}\ a - z| .
\]
Since the error radius is explicit, this reduces to a distance-to-nearest-integer lower bound for
one explicitly computable rational per scale --- the cleanest reduction of the \#249 obligation to
a purely finite, computable question in the whole corpus.
\end{defn}

\subsubsection{SK-02 --- the short-window arithmetic kill, verbatim, truncated at $a_0\le 6$ (\coord{short-window-arithmetic})}

\begin{thm}
\lean{powerTwoActualLcmShortArithmeticKillSupply_through_six}{TotientActualLcmShortKill.lean:54}
and
\lean{irrational_totientSeries_of_shortArithmeticKillSupply}{TotientActualLcmShortKill.lean:63}.
\ev{Cert} for the witness, \ev{Lean} for the consumer. \scale{bounded}.
\[
  \forall a_0\le 6,\ \exists a,L,\ a_0\le a \wedge L < 2\cdot 2^a \wedge
  \mathtt{LcmDiagonalArithmeticKill}(2^a)\,L
\]
--- literally the open cofinal supply predicate below with its universal quantifier truncated at
$a_0\le 6$. This is the purest \scale{fixed}-vs-\scale{cofinal} row in the corpus: the proof term is
$\langle$\texttt{6, 93, ha0, by norm\_num, lcmDiagonalArithmeticKill\_two\_pow\_six}$\rangle$ --- a single
hard-coded witness $(a,L)=(6,93)$, itself discharged from
\lean{certifiedKill_diagonal_t64}{DiagonalPincerCertificateT64Endpoint.lean:1928} (a
\texttt{norm\_num} evaluation over an explicit table of $\varphi$ values). There is no argument in
$a$ whatsoever; removing the bound $a_0\le 6$ requires an entirely new proof, not a re-run.
\end{thm}

\begin{defn}[What would close it]
$\mathtt{PowerTwoActualLcmShortArithmeticKillSupply}$: \ev{Open}, \scale{cofinal}.
\[
  \forall a_0\ \exists a,L,\ a_0\le a \wedge L < 2\cdot 2^a \wedge
  \mathtt{LcmDiagonalArithmeticKill}(2^a)\,L .
\]
The short-window restriction $L<2\cdot 2^a$ buys extra structure --- every non-divisor offset in
that window is a bare prime power, by \texttt{eq\_prime\_pow\_of\_not\_dvd\_periodLcm} --- so this
is a better-equipped target than the raw SEP-02 supply even though it is formally a stronger
statement (it implies certified-kill directly, without the separation-and-round step).
\end{defn}

\subsubsection{b11 --- the diagonal pincer certificate bank (\coord{diagonal-pincer})}

\begin{thm}
\lean{certifiedKill_diagonal_all_imported}{DiagonalPincerCertificates.lean:2870}, with
\lean{diagonalPincerCertificateScales}{DiagonalPincerCertificates.lean:2852} and
\lean{diagonalPincerKillDepth}{DiagonalPincerCertificates.lean:2854}. \ev{Cert}. \scale{fixed}.
\[
  \forall t\in\{1,2,3,4,5,7,8,9,11,13,16,17\},\ \exists L,\
  \mathtt{certifiedKill}\ (\mathtt{periodLcm}\ t)\ (\mathtt{periodLcm}\ t)\ L
\]
at depths $\{6,5,7,7,9,14,15,14,21,22,23,26\}$ respectively, extended by the separate T19\ldots T64
modules to 28 historical values through $t=64$ (the $t=64$ endpoint is
\lean{certifiedKill_diagonal_t64}{DiagonalPincerCertificateT64Endpoint.lean:1928}).
The later aggregate theorem closes every scale $t\le82$ with no holes
(\lean{ErdosProblems.Skip.LadderT67.exists_diagonalKill_le_82}{ErdosProblems/Skip/LadderT67.lean:71264}).
Every witness
is a \texttt{norm\_num} evaluation over a hard-coded block of $\varphi$ values --- e.g.\ the $t=17$
certificate lists $\varphi(12252241),\ldots,\varphi(24504506)$ explicitly --- so nothing in the
proof is a function of $t$.
\end{thm}

\begin{rem}[A stronger signal in the depth table than the theorem states]
\lean{certifiedKill_depth_floor}{TotientTailPeriodKiller.lean:79} forces $2(2H_t{+}L{+}2)<2^L$,
i.e.\ $L \gtrsim \log_2(4\cdot\mathtt{periodLcm}\ t)$ at any certified depth. Comparing that floor
against $\mathtt{diagonalPincerKillDepth}$ shows every landed certificate fires within a small
additive constant of the theoretical minimum depth, at every tested scale --- a numerically
supported (\ev{Cert}, not \ev{Math}) anti-concentration observation, not a theorem.
\end{rem}

\begin{prop}[What would close it]
\[
  \exists C\ \forall t_0\ \exists t\ge t_0\ \exists L\le \log_2(4\cdot\mathtt{periodLcm}\ t)+C,\quad
  \mathtt{certifiedKill}\ (\mathtt{periodLcm}\ t)\ (\mathtt{periodLcm}\ t)\ L .
\]
This is an anti-concentration statement about the diagonal word's residue at the first admissible
depth --- exactly the shape a doubling-orbit equidistribution argument would produce.
\end{prop}

\subsubsection{a12 --- one window, sixteen periods, wrong-way quantifiers (\coord{h-uniform-certificate})}

\begin{thm}
\lean{certifiedKill_all_upto_sixteen}{CarrySurvivorExtinction.lean:574} and
\lean{totient_series_ne_rat_of_den_dvd_pow_two_mul_mersenne_upto_sixteen}{CertificateKernel.lean:18890}
(shallower sibling \lean{certifiedKill_all_small}{TotientTailPeriodKiller.lean:404}). \ev{Cert}.
\scale{fixed}.
\[
  \forall h\in[1,16],\quad \mathtt{certifiedKill}\ h\ 14\ 9
\]
--- a single position $N=14$ and a single depth $L=9$ that simultaneously certify every period $h$
up to $16$. Consequence: $S\ne r$ for every rational $r$ with $r.\mathrm{den}\mid 2^{14}(2^h{-}1)$,
$1\le h\le 16$. Proof is \texttt{decide} on a 9-bit window; no part of it is a function of $N$ or of
the $h$-range.
\end{thm}

\begin{rem}[Quantifier inversion --- the only $h$-uniform row in the corpus]
The 249-supply obligation is $\forall h\ \forall N_0\ \exists N\ge N_0\ \exists L$. This row proves
$\exists N\ \exists L\ \forall h\in[1,16]$: the quantifiers are inverted, and the $h$-range is
finite. The inversion helps in one direction (one $(N,L)$ covers a whole block of periods, more
than the obligation asks) and hurts in the other ($N$ is pinned at $14$, not cofinal). This is the
only row in the corpus with $h$-uniformity, and it is the reason a scale-uniform version would be
unusually strong.
\end{rem}

\begin{prop}[What would close it]
A scale-uniform version of the same shape: $\exists f:\mathbb N\to\mathbb N$ with $f(N)\to\infty$
such that $\forall N_0\ \exists N\ge N_0\ \exists L$ with $\mathtt{certifiedKill}\ h\ N\ L$ for all
$h\le f(N)$. This implies the obligation immediately (fix $h$, take $N_0$ large enough that
$f(N)\ge h$) and, unlike the obligation, is a statement about one window at a time rather than a
per-period search.
\end{prop}

\subsubsection{c3 --- the Farey denominator floor and its stalled growth law (\coord{farey-window})}

\begin{thm}
\lean{tsum_totient_div_pow_two_ne_ratCast_of_den_le_79639646646701375323355774875831053}{CertificateKernel.lean:18384},
built from the classical mediant lemma \lean{farey_gap}{GapFareyBound.lean:51} and the window
sharpness certificate \lean{gap_check_window_1_240_le_79639646646701375323355774875831053}{GapFareyBound.lean:176}
/ \lean{gap_check_window_1_240_first_failure}{GapFareyBound.lean:225}. \ev{Lean} for the general
mediant lemma; \ev{Cert} for the $K=240$ window instantiation. \scale{fixed}. If $S$ is rational,
its reduced denominator exceeds $7.9639646646701375323355774875831053\times 10^{34}$ (the exact
numeral in the declaration name). This is the strongest unconditional statement about $S$ in the
corpus, logically independent of the certificate-supply reduction.
\end{thm}

\begin{rem}[Why re-running the same window buys nothing]
The cofinal upgrade of this statement --- $S\ne p/q$ for every bound $q$, not just $q =
7.96\times10^{34}$ --- is literally $\mathrm{Irrational}(S)$, so the target coordinate is right.
What blocks promotion is proved on disk:
\lean{gap_check_window_1_240_first_failure}{GapFareyBound.lean:225} establishes that the $K=240$
window bound is \emph{sharp}, at exactly the mediant $b{+}d$, so re-running the same argument at
the same window buys nothing further. Each new $K$ needs (i) a freshly committed $2^K$-scale
totient residue $V_K$ and (ii) fresh continued-fraction convergents of $V_K'/2^K$, both hard-coded
numerals in the current proofs. The growth of the bound $b_K{+}d_K$ is governed by the convergent
denominators of the underlying constant.  Rationality would force eventual stalling, so unbounded
growth would prove \#249; however, no converse reduction or proved logical equivalence is known.
The sentence is therefore a diagnosis of why this fixed-window method reaches the original
difficulty, not an iff theorem.
\end{rem}

\begin{obs}[A closed sub-route]
\lean{ReducedDenominatorUnitGapCert}{PrimitiveWeightCertificate.lean:30} and
\lean{PrimitiveReducedDenominatorUnitGapSupply}{PrimitiveWeightCertificate.lean:47}, the corpus's
own attempt at strengthening this window refinement via a unit-modulo-odd-part criterion, prove a
hard ceiling in
\lean{reducedDenominatorUnitGapCert_prime_pow_iff}{PrimitiveWeightCertificate.lean:54} and
\lean{reducedDenominatorCandidateCount_prime_pow_le_one}{PrimitiveWeightCertificate.lean:111}
(\ev{Lean}, \scale{n/a}): at a prime-power denominator, the strengthening rescues at most \emph{one}
extra lattice point. That route is closed --- it cannot be the source of a growth law.
\end{obs}

\begin{prop}[What would close it]
A proved growth law for the window family: a function $g$ with $g(K)\to\infty$ and a proof that for
every $K$ the $(N{=}1,K)$ gap check passes for all $q\le g(K)$. Equivalently, a lower bound on the
convergent denominators of the totient-window constant, uniform in $K$.
\end{prop}

\subsubsection{TA-01/TA-02 --- the two-adic pulse block reaches the arc centre and fails by an exponential (\coord{two-adic-pulse})}

\begin{thm}
\lean{exists_prime_totient_twoAdic_pulse_divisors}{TotientTwoAdicPulseBlock.lean:54},
\lean{exists_prime_deltaTotient_twoAdic_pulseBlock}{TotientTwoAdicPulseBlock.lean:211},
\lean{windowDiscrepancy_modEq_half_of_twoAdic_pulse}{TotientTwoAdicPulseBlock.lean:262},
\lean{eventual_integral_tailDiff_has_cofinal_twoAdic_half_pulse}{TotientTwoAdicPulseBlock.lean:327}.
\ev{Lean}. \scale{uniform} (the theorem is unconditional and holds for \emph{every} $K$, $H$, $B$,
not merely cofinally many). For every $K\ge 2$, every $H>K$, and every bound $B$, there are primes
$p>B$ with a length-$(K{-}1)$ zero prefix and a terminal half-turn, giving
\[
  A_{H,p-K,K} \equiv 2^{K-1} \pmod{2^K}.
\]
Under eventual integrality this transfers to an integer $z$ with $(z:\mathbb R) =
\mathrm{totientTail}(p{+}H) - \mathrm{totientTail}(p)$ and $z\equiv 2^{K-1}\pmod{2^K}$.
\end{thm}

\begin{rem}[Lands at the exact centre of the arc, fails by an exponential --- and why]
This lands the residue at the exact centre of the certified-kill arc: $2^{K-1}$ is maximally far
from $0$ modulo $2^K$, precisely the coordinate the obligation wants. Taking $N{=}p{-}K$,
$h{=}H$, $L{=}K$ gives radius $N{+}h{+}L{+}2 = p{+}H{+}2$, so $\mathtt{certifiedKill}$ requires
$2^{K-1} > p{+}H{+}2$. But $z\equiv 2^{K-1}\pmod{2^K}$ forces $|z|\ge 2^{K-1}$, while the directed
tail bound forces $|z| < p{+}H{+}2$; and the construction's own congruence
$p\equiv 1{+}2^{K-1}\pmod{2^K}$ (equivalently $v_2(p{-}1)=K{-}1$) forces $p\ge 1{+}2^{K-1}$. So
$p{+}H{+}2 > 2^{K-1}$ always, and no choice of $p$ rescues it. The quantifier order is inverted:
the theorem gives, for fixed $K$, cofinally many \emph{large} $p$; the certificate needs a $p$
\emph{small} relative to $2^{K-1}$, impossible for a single letter since $|\varphi(N{+}H)-\varphi(N)|
< N{+}H$. Depth-promotion is not the issue: it has already been performed, generalizing an
exponent-2 construction to every $K$ uniformly by CRT-gluing over a
$\mathrm{Unit}\oplus\mathrm{Fin}(K{-}1)\oplus\mathrm{Fin}(K{-}1)$ family.
\end{rem}

\begin{prop}[What would close it]
The same half-turn residue realised by an \emph{accumulated} window rather than a single terminal
letter: cofinally many $(h,N,L)$ with $A_{h,N,L}\equiv 2^{L-1}\pmod{2^L}$ and
$2^{L-1} > N{+}h{+}L{+}2$. Only the weighted sum $\sum \Delta\varphi\cdot 2^{L-1-j}$ can carry
magnitude $2^L$; the per-letter pulse provably cannot. Concretely: a pulse-block construction whose
zero prefix and half-turn are imposed on the \emph{word}, not on one delta.
\end{prop}

\subsubsection{d-3a/d-3b --- cofinal jump positions, one unproved slack sign (\coord{lcm-jump-slack})}

\begin{thm}[Cofinal producer, fully proved]
\lean{exists_periodLcm_strict_jump_ge}{DiagonalFreshLossBridge.lean:2635}. \ev{Lean}.
\scale{cofinal} (proved, not open): $\forall t_0\ \exists t\ge t_0$ with
$\mathtt{periodLcm}\ t < \mathtt{periodLcm}(t{+}1)$, witnessed by $p{-}1$ for any prime $p>t_0$.
The power-of-two specialization
\lean{CanonicalAdjacentSuffixPowerTwoPostJumpSlackSupply}{DiagonalFreshLossBridge.lean:2693} shows
the positions $t=2^a{-}1$ already suffice, so no position search is needed at all.
\end{thm}

\begin{defn}[The slack scalar and the open supply]
\lean{canonicalAdjacentSuffixCentralSlack}{DiagonalFreshLossBridge.lean:2651}:
\[
  \mathtt{canonicalAdjacentSuffixCentralSlack}\ t = \min\bigl(d-2^{m-5},\ (2^m-2^{m-5})-d\bigr),
\]
with $m$ the canonical adjacent-suffix depth and $d$ the residue at that depth. The open
producers, both \ev{Open} \scale{cofinal}, are
\lean{CanonicalAdjacentSuffixPostJumpSlackSupply}{DiagonalFreshLossBridge.lean:2686} (any strict
jump) and \lean{CanonicalAdjacentSuffixPowerTwoPostJumpSlackSupply}{DiagonalFreshLossBridge.lean:2693}
(power-of-two positions), each asking $0\le \mathtt{canonicalAdjacentSuffixCentralSlack}(t{+}1)$
cofinally, and both compose to
$\mathrm{Irrational}(S)$ at
\lean{irrational_totientSeries_of_canonicalAdjacentSuffixPostJumpSlackSupply}{DiagonalFreshLossBridge.lean:2941}
and
\lean{irrational_totientSeries_of_canonicalAdjacentSuffixPowerTwoPostJumpSlackSupply}{DiagonalFreshLossBridge.lean:2948}.
\end{defn}

\begin{rem}[Census evidence only]
\ev{Cert}, \scale{fixed}. All $40$ strict jumps up to endpoint $113$ pass, closest margin
$\approx 0.000221$ of the modulus at $t=100$; the power-of-two endpoints $4,8,16,32$ pass with
slack fractions $0.41,0.036,0.40,0.19$. Census, not theorem. The structural handle nobody has used:
$\mathtt{periodLcm\_succ\_eq\_prime\_mul\_of\_strict\_jump}$
(\lean{periodLcm_succ_eq_prime_mul_of_strict_jump}{DiagonalFreshLossBridge.lean:2332}) says a
strict jump at $t$ means $t{+}1=p^k$ and $\mathtt{periodLcm}(t{+}1) = p\cdot\mathtt{periodLcm}(t)$
--- so the required statement is a transfer lemma for how the adjacent-suffix residue at height $H$
moves under $H\mapsto pH$.
\lean{canonicalAdjacentSuffixCentralSlack_eq_of_periodLcm_eq}{DiagonalFreshLossBridge.lean:2658}
already proves the slack is constant across each LCM plateau, so only jump indices matter.
\end{rem}

\begin{prop}[What would close it]
$0\le \mathtt{canonicalAdjacentSuffixCentralSlack}(2^a)$ for cofinally many $a$.
\end{prop}

\subsubsection{e2 --- the prime-jump commutator, one fixed witness (\coord{prime-jump-commutator})}

\begin{thm}
\lean{primeJumpSharpRadius}{PrimeJumpWindow.lean:37} and
\lean{primeJumpTailCommutator_notMem_int_of_sharpKill}{PrimeJumpWindow.lean:178}. \ev{Lean}.
\scale{uniform}. Uniform consumer for all $H,p,L$: a residue of the four-vertex commutator
$J(H,p)$ outside the sharp radius $3pH{+}(p{+}1)(L{+}2)$ forces $J(H,p)\notin\mathbb Z$. This is
tighter than the earlier $4pH$ two-cell disjunction.
\end{thm}

\begin{thm}[The one witness]
\lean{primeJumpSharpKill_twelve_five}{PrimeJumpWindow.lean:186}. \ev{Cert}. \scale{fixed}.
$\mathtt{primeJumpSharpKill}\ 12\ 5\ 15$ (i.e.\ $H=\mathtt{periodLcm}\ 4=12$, $p=5$, $L=15$),
proved by \texttt{decide} on an explicit 15-bit window with no dependence on $t$ or $p$; only a
single instance exists, at $t=4$, the smallest nontrivial height. The full endpoint composition is
\lean{irrational_totient_series_of_primeJumpSharpKill_supply}{PrimeJumpWindow.lean:193}.
\end{thm}

\begin{rem}
The supply hypothesis $\forall t_0\ \exists t\ge t_0\ \exists p,L{>}0$ with
$\mathtt{primeJumpSharpKill}(\mathtt{periodLcm}\ t)\ p\ L$ is a genuinely cheaper target than the
raw SEP-02 supply: it asks for one fresh prime $p$ per LCM height rather than a full central-arc
certificate. Because the consumer route is
\texttt{rational\_totient\_series\_forces\_lcm\_cone\_flatness}$\to$contradiction, the natural
attack is to pick $p$ as the fresh prime introduced at the next strict LCM jump (linking this row
to d-3a), where the commutator's old-channel contributions are annihilated by
\lean{oldChannel_affine_moment_annihilation}{JointExponentTransport.lean:36}.
\end{rem}

\subsubsection{b7 --- the M\"obius--Mersenne denominator channel: unbounded, but off-coordinate (\coord{mobius-mersenne})}

\begin{thm}
\lean{upperHalfMersenneProduct_lower_bound}{MersenneShadowDenominatorGrowth.lean:60},
\lean{lcmHeight_scaledMobiusShadow_den_lower_bound}{MersenneShadowDenominatorGrowth.lean:85},
\lean{exists_upperHalf_mersenne_channel_dvd_den_with_bounds}{MersenneShadowDenominatorGrowth.lean:99}.
\ev{Lean}. \scale{uniform}. For every $t\ge 5$, uniformly in $t$:
\[
  2^{t/2} \;\le\; \prod_{p\in \mathrm{upperHalfPrimes}(t)} \mathtt{mersenne}(p) \;\le\;
  \mathrm{den}\bigl(\mathtt{lcmHeight}(t)\cdot \mathtt{numericMobiusShadow}(\mathtt{lcmHeight}(t))\bigr).
\]
An unbounded denominator lower bound at every LCM height, with an individual surviving channel
isolated ($2^{t/2}\le \mathtt{mersenne}(p) < 2^t$). The proof is genuinely uniform in $t$
(Bertrand's postulate via \texttt{upperHalfPrimes\_nonempty}, plus $2^{p-1}\le 2^p{-}1$ factorwise
--- no table, no constant degrading with $t$).
\end{thm}

\begin{rem}[The corpus's only proved unbounded-growth quantity in this coordinate, but the wrong coordinate for \#249]
This lives in a different coordinate from the 249-supply obligation: it bounds the reduced
denominator of the \emph{scaled} shadow at LCM height $t$, and there is no landed transport from a
denominator lower bound to $\mathtt{certifiedKill}$ or to a totient-tail non-integrality. Two
further honest caveats are recorded in the module's own docstring: it does not rule out
cancellation by the foreign-defect term, and the growth rate $2^{t/2}$ is far below the scale
$\mathtt{lcmHeight}(t)\approx 2^{1.44t}$ against which the enclosure consumers measure error. The
existing consumers built on it
(\texttt{scaleFullTarget\_miss\_of\_lambert\_projected\_separation},
\texttt{scaleFullTarget\_miss\_of\_lambert\_projected\_num\_gap}) are \#257-facing
(\texttt{ScaleFullTargetHit}), not \#249-facing.
\end{rem}

\begin{prop}[What would close it]
Either (i) a Dirichlet-criterion bridge in the shape of
\lean{irrational_of_den_mul_abs_sub_tendsto_zero}{CertificateKernel.lean:5371}: a sequence of
rationals $u_t$ with $u_t\ne S$ and $\mathrm{den}(u_t)\cdot|S-u_t|\to 0$, which requires the
approximation error to beat the denominator, not merely the denominator to grow; or (ii) a
transport map from ``channel of size $\ge 2^{t/2}$ survives in the denominator'' to ``windowDiscrepancy
residue avoids the arc,'' which does not exist in the corpus. Route (i) needs the growth rate
raised from $2^{t/2}$ to beat the shadow's own truncation error.
\end{prop}

\subsubsection{d4/d5 --- the parallel rank obligation, and a corrected pair of citations (\coord{carry-rank})}

This row runs a second, independent open obligation in parallel to 249-supply, in the carry-kernel
\emph{rank} coordinate rather than the residue coordinate. \textbf{Correction}: the source index
attributed the two theorems below to swapped files; direct reads this session confirm the correct
attribution is as given.

\begin{thm}
\lean{linearIndependent_canonicalTotientKernelFamily}{TotientMahlerDefect.lean:935}. \ev{Lean}.
\scale{uniform}. Unconditional and uniform in $e$: the dyadic totient-kernel family is linearly
independent at every depth $e$ (dimension exactly $2^e{+}1$, via CRT and Dirichlet on primes in
arithmetic progression).
\end{thm}

\begin{thm}
\lean{not_irrational_totientSeries_implies_unbounded_carryRank_unconditional}{TotientCarryKernelRigidity.lean:300}.
\ev{Lean}. \scale{uniform}. If $S$ is rational there is a tempered integral carry orbit $u$ with
\[
  2^e - 1 \;\le\; \mathrm{finrank}_{\mathbb Q}\,\mathrm{span}(\mathtt{canonicalCarryKernelFamily}\ u\ e)
  \qquad\text{for every } e .
\]
\end{thm}

\begin{rem}[The scale side is finished; the missing direction is proved dead on one route]
So: rationality forces \emph{unbounded} carry-kernel rank, and the $2^e{-}1$ floor holds for all
$e$ with a uniform proof. What is missing is the opposite inequality --- a rationality-side rank
upper bound, which would contradict the floor and close \#249. The corpus proves one natural route
to it is dead:
\lean{false_of_compressedAdjointCertificate}{TotientMahlerDefect.lean:1183} (\ev{Lean},
\scale{n/a}) shows no
\lean{CompressedAdjointCertificate}{TotientMahlerDefect.lean:1171} ($Q\cdot v\cdot A = $ boundary
with $|\mathrm{boundary}| < Q\cdot v$ and $A\ne 0$) can exist, and
\texttt{not\_finiteDimensional\_span\_fullTotientKernel} (immediately above it in the same file)
shows the full family spans an infinite-dimensional space. Cross-checked against a separate row,
$\mathtt{lcm\_factorIdeal\_finiteRank\_shiftAlgebra\_not\_sufficient}$: a single explicit
countermodel defeats every finite-rank shift-polynomial observation simultaneously, so finite-rank
amplification per se is not the missing rigidity.
\end{rem}

\begin{prop}[What would close it]
A rationality-side rank upper bound: $\exists C$ such that any tempered integral binary carry
orbit for a rational $\mathtt{binaryCoeffSeries}$ has
$\mathrm{finrank}_{\mathbb Q}\,\mathrm{span}(\mathtt{canonicalCarryKernelFamily}\ u\ e)\le C$ for
all $e$, or any bound growing slower than $2^e{-}1$.
\texttt{not\_irrational\_totientSeries\_implies\_mod\_period\_and\_unbounded\_rank}
(\texttt{TotientTailCarryPeriod.lean:224}) pins the obstacle precisely: rationality buys uniform
eventual periodicity of $u$'s dyadic sections mod $v$, and that periodicity provably does
\emph{not} promote to a $\mathbb Q$-rank bound without extra arithmetic input.
\end{prop}

\subsection{Synthesis: where the frontier actually is}

Collecting the fourteen rows, three facts stand out as loci for future effort, stated with the same
precision as the rows themselves rather than as summary rhetoric.

First, the single closest approach to 249-supply is the first-harmonic reduction of
\S\ref{ssec:headline}: a constant-saving ($21/25$ or $9/10$) exponential-sum cancellation estimate,
with a satisfiable room condition and, via the subset consumer
(Theorem~\ref{thm:hgap-subset}), no requirement that the saving hold on a full interval or a
positive-density subset. This is the only row in the index whose remaining content is recognisably
a single classical estimate rather than a compound arithmetic-geometric statement.

Second, the top-edge staircase (TE-04 through TE-05-weakest, \coord{actual-lcm-top-edge}) is the
only row where the corpus has both halved the obligation (one-sided instead of symmetric) and
exposed, via the exact identity of Observation~\ref{obs:staircase-tension}, precisely which of two
logically symmetric escape branches survives the sign machinery already proved elsewhere in the
same file family. That five strictly weaker equivalent or sufficient forms are already proved
inter-derivable (Proposition~\ref{prop:te-chain}) means effort here is not fragmented across
restatements: proving any one closes the others automatically.

Third, the d4/d5 rank obligation is structurally independent of the residue-coordinate rows above:
it is a second, self-contained sufficient condition for \#249 (an upper rank bound contradicting
the proved $2^e{-}1$ lower bound), in a coordinate (carry-kernel rank) where a natural finite-rank
strengthening is already proved impossible. A rank upper bound, if found, would not need to route
through $\mathtt{certifiedKill}$ at all.

No row in this index is a solution, a partial solution in the sense of resolved cases, or evidence
that \#249 is likely true or false. Every open producer listed is exactly as open as the original
Erd\H{o}s--Borwein question it reduces to; what the index adds is the exact remaining content,
verified against the live Lean tree, of each reduction.
% ---- part a249_p4 ----
\section{Closed routes, with the mechanism that closed them}

This section catalogues every no-go, countermodel, obstruction, and refuted
proof strategy found in the \#249/\#257 corpus that bears on Erd\H{o}s \#249
(\(S := \sum_{n\ge1}\varphi(n)/2^n\) irrational). Erd\H{o}s \#249 is
\ev{Open}; nothing in this section decides it. What follows is a machinery
inventory: for each dead route, the exact statement that closed it, the exact
scope of the closure, and the exact machinery salvaged from the wreckage.
The organising warning, stated once here and applicable to every item below:
several of these ``obstructions'' are statements about a
\emph{representation} of the problem (a proof strategy, a certificate
family, a coordinate system) and not about the object \(S\) itself. Confusing
the two is the single most common error this catalogue exists to prevent.

\subsection{Meta-level: finite inspection cannot certify the supply (the
\texorpdfstring{\(\gamma\)}{gamma}-splice remark)}

\textbf{(a) Route as conceived.} Every unconditional route to \#249 in this
corpus bottoms out in a \emph{certificate-supply} obligation of the shape
\(\forall h\ge1\,\forall N_0\,\exists N\ge N_0\,\exists L,\ \mathrm{Sep}(h,N,L)\)
(\lean{Erdos249257.irrational\_totient\_series\_iff\_certificate\_supply}{LcmConeFlatness.lean:412}, the exact
open-obligation form quoted throughout the existing manuscript). A natural
hope is that verifying \(\mathrm{Sep}\) — or its diagonal specialisation — at
every scale up to some large, explicit bound \(B\) is evidence that the
supply holds, and that pushing \(B\) far enough would eventually amount to a
proof.

\textbf{(b) Exact mechanism that closed it.} The manuscript's own remark
(\ev{Math}, prose-only, not formalised in Lean:
\emph{\S5.4, ``finite inspection cannot establish the supply''}) gives an
explicit splice construction. Take any coefficient stream \(\gamma:\mathbb
N\to\mathbb N\) that agrees with \(\varphi\) on every index \(\le B\), and
alter \(\gamma\) at one residue class beyond \(B\) so that the resulting
binary series \(\sum\gamma(n)/2^n\) is forced \emph{rational}. Such a
\(\gamma\) exists for every \(B\) (this is the same coboundary-splice
technique that produces \S\ref{sec:parity-countermodel}'s
\texttt{parityCoboundaryWeight} witness, generalised: insert a lacunary
zero-valued coboundary edit \(2/2^m - 4/2^{m+1} = 0\) at some \(m>B\)). Because
\(\gamma\) agrees with \(\varphi\) on every index \(\le B\), \(\gamma\) passes
every finite certificate that only inspects indices \(\le B\) — in
particular every instance of \(\mathrm{Sep}(h,N,L)\) with \(N+L\le B\) that
\(\varphi\) itself passes or fails. Yet \(\sum\gamma(n)/2^n\in\mathbb Q\) by
construction.

\textbf{(c) Precise scope.} This is a statement about \emph{proof method},
not about \(\varphi\): it does not touch \(S\) at all. What it excludes is
the inference ``\(\mathrm{Sep}\) verified up to bound \(B\), for arbitrarily
large \(B\), therefore \(\mathrm{Sep}\) holds cofinally.'' The scope is
exactly \(\scale{bounded}\) versus \(\scale{cofinal}\): any finite-B
verification, however large, is compatible with both outcomes, because a
rational stream can be built to survive any \emph{fixed} inspection horizon.
It says nothing about whether \(\varphi\) itself is such a \(\gamma\) — only
that no finite-horizon check can distinguish \(\varphi\) from a \(\gamma\)
that is.

\textbf{(d) Salvage.} The splice construction is exactly the mechanism
underlying \S\ref{sec:parity-countermodel}'s Lean-formalised countermodel
below, so its \emph{content} is not lost even though this specific remark is
unformalised: everywhere a proof strategy in this corpus proposes to certify
a supply by exhaustive finite search (the diagonal deposits at
\(t\in\{1,\dots,64\}\), \lean{Erdos249257.certifiedKill\_diagonal\_all\_imported\_through\_t64}{DiagonalPincerCertificatesT64.lean:1967}; the row-200{,}000 sqrt-escape
census on the \#257 side), this remark is the standing reason such a search,
however large, is evidence and not proof. \ev{Math} \scale{n/a}
\coord{binary-digit}.

\subsection{The parity-aperiodicity countermodel
(\texttt{TotientParityCoboundaryCountermodel})}
\label{sec:parity-countermodel}

\textbf{(a) Route as conceived.} A natural strengthening of \#249 attempts:
show that any coefficient word \(c:\mathbb N\to\mathbb N\) that (i) is
uniformly bounded, (ii) satisfies the trivial linear-growth bound
\(c(n)\le n\), (iii) matches \(\varphi(n)\bmod 2\) exactly, and (iv) is not
eventually periodic — even strengthened to \emph{cofinally, arbitrarily
separated, arbitrarily long blocks} of non-periodicity — must have irrational
binary series \(\sum c(n)/2^n\). This is the natural target for anyone
trying to use only \(\varphi\)'s coarse parity/aperiodicity profile.

\textbf{(b) Exact mechanism.} The construction
\(\mathtt{parityCoboundaryWeight}(n) := \mathtt{parityBaseWeight}(n) +
2\cdot\mathtt{largePowerTwoBit}(n) - 4\cdot\mathtt{largePowerTwoBit}(n-1)\),
where \(\mathtt{parityBaseWeight}\) is the eventually-constant word
\(0,1,1,2,4,4,4,\dots\) and \(\mathtt{largePowerTwoBit}(n)=1\) iff
\(n=2^{k+3}\), is a lacunary zero-valued coboundary splice on top of an
eventually-constant base. The strongest of five landed theorems,
\begin{prop}[Totient-parity arbitrarily-many-separated-carry rational
countermodel]
There exists \(c:\mathbb N\to\mathbb N\) such that: \(c(n)\le6\) for all
\(n\); \(c(n)\le n\); \(c(n)\equiv\varphi(n)\pmod2\) for every \(n\); for
every \(N,G,K\) there is a block of \(K\) explicit \((6,0)\) carry-pulse
pairs beyond \(N\), each pair separated by more than \(G\); and
\(\sum_n c(n)/2^n = 3/2\).
\end{prop}
\lean{Erdos257PeriodNoncollapse.exists\_totientParity\_arbitrarilyManySeparatedCarry\_rational\_countermodel}{TotientParityCoboundaryCountermodel.lean:637}
(sum computation \lean{tsum\_parityCoboundaryWeight\_eq\_three\_halves}{TotientParityCoboundaryCountermodel.lean:359}; non-periodicity
\lean{parityCoboundaryWeight\_not\_eventually\_periodic}{TotientParityCoboundaryCountermodel.lean:487}; exact parity match
\lean{parityCoboundaryWeight\_mod\_two\_eq\_totient}{TotientParityCoboundaryCountermodel.lean:194}). \ev{Lean} \scale{cofinal}
\coord{binary-digit}.

\textbf{(c) Precise scope.} This is an \emph{object}-level existence witness
(\(c\ne\varphi\), only \(c\equiv\varphi\pmod2\)) that closes a
\emph{proof-route}: it proves that the hypothesis set \(\{\)bounded, linear
growth, \(\varphi\)-parity match, non-eventual-periodicity\(\}\), no matter
how strongly the non-periodicity clause is strengthened (up to the
arbitrarily-separated arbitrarily-numerous form above), can never entail
irrationality of the associated binary series, because \(c\) satisfies every
hypothesis and is rational. It says nothing about \(\varphi\) itself — the
witness sequence \(c\) is a different, hand-built sequence. Any future
sufficient-condition candidate stated purely in terms of \(\{\)coefficient
boundedness, growth, parity, periodicity\(\}\) must be checked against this
countermodel before being trusted, for \emph{either} \#249 or \#257.

\textbf{(d) Salvage.} The lacunary coboundary-splice technique itself —
inserting zero-valued edits \(2\cdot2^{-m}-4\cdot2^{-(m+1)}=0\) at sparse
ranks to destroy periodicity while preserving the rational sum — is a fully
general recipe for building rational-valued, non-eventually-periodic,
bounded coefficient sequences matching \emph{any} parity template, and is
literally the mechanism instantiated informally in
Theorem~\ref{thm:gamma}'s \(\gamma\)-splice construction. It is also a
ready-made stress-test fixture for any future parity-based sufficient
condition proposed anywhere in the corpus. Confirms the corpus-wide lesson
(project memory \texttt{feedback\_erdos\_reductions\_rejected\_bank\_real\_results}):
only arguments using \emph{actual} quantitative totient/Mersenne size or
residue information — as in the \texttt{TotientActualLcm*},
\texttt{TotientFixedRank*} families below — can possibly close \#249; pure
symbolic-word arguments cannot.

\subsection{Two scoped \#249 no-go countermodels: square-CRT correction
suppression is ambiguous}

\textbf{(a) Route as conceived.} The totient prime-dilation cocycle
\(\varphi(pm)=(p-1)\varphi(m)+[p\mid m]\varphi(m)\) carries a correction term
whenever \(p\mid m\). A natural strategy imposes a square congruence
\(n\equiv A+pa\pmod{p^2}\) with a ``clean'' shift \(h\) (\(p\nmid a+h\)) to
suppress that correction on a finite horizon, then hopes that
correction-suppression (``cleanliness'') is itself enough to guarantee that
the resulting finite-difference cube is \emph{nonzero} — i.e. that cleaning
away the correction term automatically leaves behind a genuine, nonvanishing
arithmetic signal usable in a curvature argument.

\textbf{(b) Exact mechanism that closed it.} Two explicit, opposite
finite computations refute the hoped-for entailment simultaneously:
\begin{itemize}[nosep]
\item \emph{Vanishing witness.} The smallest returned countermodel has the
\emph{entire} two-step finite block \(\mathtt{actualOneCubeCoeff}\) equal to
zero at \(n=52\), \(r_0=13\), \(r_1=13/18\)
(\lean{squareCRT\_clean\_block\_can\_vanish}{SquareCRTCube.lean:459}, kernel-\texttt{decide}d;
cleanliness witnessed by \lean{squareCRT\_vanishing\_countermodel\_is\_clean}{SquareCRTCube.lean:468}).
\item \emph{Nonvanishing witness.} A separate, nearby clean configuration at
\(n=27\) has the same block \emph{nonzero}
(\lean{squareCRT\_clean\_block\_can\_be\_nonzero}{SquareCRTCube.lean:477}; cleanliness witnessed by
\lean{squareCRT\_nonzero\_countermodel\_is\_clean}{SquareCRTCube.lean:485}).
\end{itemize}
Both witnesses are exhibited alongside the positive achievability result
\lean{exists\_squareCRT\_clean\_horizon}{SquareCRTCube.lean:327} (a bounded common base satisfying prescribed
CRT anchors/residues does exist). \ev{Cert} (theorem bodies not
independently re-verified beyond signatures and docstring — flagged
\texttt{DOCSTRING-SOURCED} in the source bank) \scale{fixed}
\coord{other:square-crt-cocycle}.

\textbf{(c) Precise scope.} These are two \#249-scoped countermodels — the
coordinate is the totient prime-dilation cocycle specifically, not a
problem-agnostic phenomenon. Together they show that
\emph{correction-suppression is logically independent of the
sign/nonvanishing question it was hoped to settle}: ``clean'' (i.e.
correction-free) is consistent with \emph{both} outcomes. This is a
statement about a specific finite-difference construction (the
representation), not about \(\varphi\): it does not exclude a nonvanishing
finite-difference cube in general, only the inference from cleanliness
\emph{alone} to nonvanishing.

\textbf{(d) Salvage.} The general cocycle identity
\(\varphi(pm)=(p-1)\varphi(m)+[p\mid m]\varphi(m)\)
(\lean{totient\_mul\_prime\_cocycle}{SquareCRTCube.lean:38}) and the CRT-witness apparatus
(\lean{exists\_squareCRT\_base}{SquareCRTCube.lean:230}, \lean{exists\_large\_squareCRT\_base}{SquareCRTCube.lean:263}) remain fully reusable for
\emph{any} residue-forcing construction, in any coordinate. More valuably,
the abstract cube-cancellation lemma extracted from the same file,
\lean{alternating\_powerset\_sum\_eq\_zero\_of\_insert\_invariant}{SquareCRTCube.lean:363} — ``if inserting one
coordinate never changes the value of \(f\), the alternating powerset sum of
\(f\) vanishes'' — is completely divorced from totient/CRT content and is
directly reusable for any Hilbert-cube or higher-difference argument on
either problem. What remains as the honest open gap: an independent
anti-concentration or residue-producer argument for why the clean residual
does \emph{not} vanish, which this module explicitly does not supply.

\subsection{The adelic height obstruction}
\label{sec:adelic-height}

\textbf{(a) Route as conceived.} Several proof strategies attempt to
\emph{clear} an unwanted denominator factor from a rational value \(x\) by
multiplying by a small integer coefficient \(c\), hoping to discard the
complementary denominator information entirely (e.g. localising a diagonal
tail difference to one LCM ray channel while dropping the rest).

\textbf{(b) Exact mechanism.}
\begin{lem}[Scalar-localisation complement divisibility]
For \(x:\mathbb Q\), \(c:\mathbb Z\), \(H:\mathbb N\): if \(H\mid x.\mathrm{den}\)
and \((c\cdot x).\mathrm{den}\mid H\), then \(x.\mathrm{den}/H \mid
c.\mathrm{natAbs}\).
\end{lem}
\lean{Erdos249257.AdelicHeightObstruction.scalarLocalization\_complement\_dvd}{AdelicHeightObstruction.lean:23}. The equality
strengthening, \(\exists t:\mathbb Z,\ H\cdot c\cdot x = t\cdot x.\mathrm{num}\), is
\lean{Erdos249257.AdelicHeightObstruction.scalarLocalization\_integer\_eq\_mul\_num}{AdelicHeightObstruction.lean:56}. The
Mersenne-scale corollary — for positive \(x\), if \(2^r\mid x.\mathrm{num.natAbs}\)
and \(x<2/(2^n-1)\), then \(2^r(2^n-1)<2\cdot x.\mathrm{den}\) — is
\lean{Erdos249257.AdelicHeightObstruction.positiveRat\_mersenne\_height}{AdelicHeightObstruction.lean:103} (generic form
\lean{positiveRat\_numDivisor\_mul\_lt\_two\_mul\_den}{AdelicHeightObstruction.lean:77}).

A companion rigidity fact
rules out smuggling the discarded information back in through a linear
channel: if a linear map \(\Lambda:V\to W\) factors through a surjective
evaluation \(\mathrm{ev}\) (\(\ker\mathrm{ev}\le\ker\Lambda\)), then
\(\Lambda\) is \emph{exactly} \(\mathrm{ev}(\cdot)\bullet w_0\) for a single
fixed \(w_0\) —
\lean{Erdos249257.AdelicHeightObstruction.linearDescender\_eq\_smul\_eval}{AdelicHeightObstruction.lean:120}.

The upstream
primitive these import,
\lean{Erdos249257.RationalDenominatorSurvival.divisor\_dvd\_divInt\_den}{RationalDenominatorSurvival.lean:17} and
\lean{survivingDivisor\_dvd\_scaled\_divInt\_den}{RationalDenominatorSurvival.lean:38}, gives the exact
survival/cancellation law: a divisor \(m\mid D\) of a displayed denominator
survives reduction of \(a/D\) iff \(m\) is coprime to \(a\); after scaling the
numerator by \(h\), a surviving coprime divisor \(C\mid D\) shrinks to
exactly \(C/\gcd(C,h)\), never further. \ev{Lean} \scale{uniform}
\coord{other:rational-height}.

\textbf{(c) Precise scope.} This is pure \(\mathbb Q\)-arithmetic with zero
Mersenne or totient content — it is not a statement about \(S\) at all, but
about \emph{any} rational-height bookkeeping argument. Its exact content:
scalar denominator-clearing never \emph{erases} the complementary
denominator; it moves that complement into the coefficient's (Archimedean)
size. Any construction that tries to shrink a displayed denominator down to
one surviving channel pays for it in coefficient growth. This directly
explains \emph{why} denominator-compression strategies are structurally hard:
the Mersenne-shadow denominator lower bound
\(2^{t/2}\le\prod_{p\in\mathrm{upperHalfPrimes}\,t}\mathrm{mersenne}(p)\)
(\lean{upperHalfMersenneProduct\_lower\_bound}{MersenneShadowDenominatorGrowth.lean:60}) forces a
quadratic-in-scale coefficient cost via this lemma if one tries to localise
to a single channel.

\textbf{(d) Salvage.} Fully problem-agnostic and directly reusable for \#257
or any other Lambert-type denominator-survival argument as-is. It is the
general mechanism underlying the corpus's Farey-gap denominator-exclusion
family (Part C of the certificate bank,
\lean{tsum\_totient\_div\_pow\_two\_ne\_ratCast\_of\_den\_le\_79639646646701375323355774875831053}{CertificateKernel.lean:18384}, the strongest
current unconditional statement about \(S\): if \(S\) is rational its
reduced denominator exceeds \(79{,}639{,}646{,}646{,}701{,}375{,}323{,}355{,}774{,}875{,}831{,}053\approx7.96\times10^{34}\)).

\subsection{The signed \(q\)-moment machinery — a positive tool, not an
obstruction}

\textbf{(a) Route as conceived.} A ``signed Hankel determinant'' /
Hermite--Pad\'e style route to \#249 would exhibit a finite signed linear
combination of dyadically-cleared totient terms and show it is provably
nonzero, giving a nonvanishing certificate for a first-harmonic or
Möbius-companion construction.

\textbf{(b) Exact mechanism (a nonvanishing \emph{engine}, not a no-go).}
\begin{lem}[Unique-terminal parity nonvanishing]
If a finite signed sum \(\sum_{i\in s}u(i)\cdot2^{e(m)-e(i)}\), clearing
dyadic denominators to a common exponent \(e(m)\), has one index \(m\in s\)
with strictly maximal exponent \(e(m)\) and odd coefficient \(u(m)\), while
every other index has strictly smaller exponent, then the whole cleared sum
is \(\equiv1\pmod2\), hence nonzero.
\end{lem}
\lean{scaled\_dyadic\_sum\_ne\_zero}{SignedQMomentObstruction.lean:96}, built from a fully generic rectangular
Cauchy--Binet-style determinant-of-product expansion
\lean{det\_mul\_rectangular}{SignedQMomentObstruction.lean:29} (definitions
\lean{truncatedMoment}{SignedQMomentObstruction.lean:108} and
\lean{hankelDet}{SignedQMomentObstruction.lean:117}). \ev{Lean} \scale{bounded}
\coord{p-adic}.

\textbf{(c) Precise scope — the point this section exists to make.} Despite
living in a file named alongside the corpus's ``obstruction'' modules, this
is explicitly \emph{not} phrased as a no-go in its own docstring: it is the
finite nonvanishing engine that a Hankel--Padé construction for \#249 would
\emph{consume}. The parity lemma and the Cauchy--Binet expansion use no
totient or Mersenne structure at all — pure linear algebra, fully
problem-agnostic, directly transplantable to any base-\(q\) dyadic-clearing
argument including \#257's series family. What is missing, and what
prevents this from closing anything, is a \emph{supply} theorem: no result
in the corpus proves that a unique-terminal configuration of this shape
exists cofinally for the totient-derived family. The obstruction, such as it
is, is not in this lemma but in the absence of its hypothesis's cofinal
supply — exactly the same shape of gap as \S\ref{sec:mahler-defect}'s Mahler
defect below.

\textbf{(d) Salvage.} The entire lemma is reusable machinery, not wreckage —
nothing here needs salvaging because nothing here failed. It stands ready
for whichever future construction can supply the missing cofinal
unique-terminal-configuration existence.

\subsection{The residual-gauge obstruction}

\textbf{(a) Route as conceived.} A ``first-harmonic pivot'' strategy for
\#249 (via exponential/character-sum cancellation,
\lean{FirstHarmonicPivot.lean}{FirstHarmonicPivot.lean}) might try to certify a nonzero determinant or
full-rank minor of a residual-weighted monomial matrix as sufficient
evidence that a genuine phase reconstruction (as opposed to a degenerate
locked configuration) has occurred.

\textbf{(b) Exact mechanism.}
\begin{prop}[Locked reconstruction preserves a nonzero minor]
For a monomial matrix \(\mathrm{residualMonomialMatrix}(e,r,z)(i,j) =
r(i,j)\cdot z(j)^{e(i)}\) with all residual entries \(z(j)\ne0\) and
\(e(i)=1\) for a distinguished row \(i\): \(\det(\mathrm{phasePowerMatrix})
\ne0 \implies \det(\mathrm{residualMonomialMatrix})\ne0\), \emph{and} the
distinguished row is identically \(1\) regardless.
\end{prop}
\lean{Erdos249257.locked\_reconstruction\_preserves\_nonzero\_minor}{ResidualGaugeObstruction.lean:94}
(row-dependent relaxation \lean{rowResidualMonomialMatrix\_reconstructs\_phasePowerMatrix}{ResidualGaugeObstruction.lean:124}; the
column-gauge identity underlying it,
\(\det(\mathrm{residualMonomialMatrix}) =
\det(\mathrm{phasePowerMatrix})\cdot\prod_j W(j)\), is
\lean{det\_residualMonomialMatrix}{ResidualGaugeObstruction.lean:47}; a row-dependent relaxation still fails at
\lean{rowReconstruction\_zero\_row\_one}{ResidualGaugeObstruction.lean:145}). \ev{Lean} \scale{uniform}
\coord{other:first-harmonic-phase}.

\textbf{(c) Precise scope.} Pure linear algebra over \(\mathbb C\), fully
generic in dimension \(d\) and exponent function \(e\), zero \#249-specific
content — this is a statement about \emph{residual-weighted monomial
matrices as a representation}, not about the totient object. It rules out
exactly one class of certificate: a residual-blind rank/determinant/
conditioning test, used alone, cannot distinguish genuine phase
reconstruction from an explicit ``locked'' degenerate configuration in which
a residual weighting scales every column by \(W(j)=z(j)^{-1}\), collapsing
the exponent-one row to the constant row \((1,\dots,1)\) while the
determinant stays nonzero whenever the base (Vandermonde-shaped)
\(\mathrm{phasePowerMatrix}\) determinant is nonzero. It does \emph{not}
rule out determinants carrying an additional arithmetic coupling identity
between rows.

\textbf{(d) Salvage.} A hard requirement for any future first-harmonic
determinant construction (Part E of the certificate bank,
\lean{irrational\_totient\_series\_of\_first\_harmonic\_norm\_gap}{FirstHarmonicPivot.lean:83}): the construction \emph{must}
couple rows by an extra arithmetic identity, not merely gauge-normalise
columns. The diagonal-gauge-action technique itself is reusable for any
determinant-based nonvanishing strategy in any coordinate.

\subsection{Cone flatness and cone nonflatness}

\textbf{(a) Route as conceived.} The diagonal certificate-supply obligation
(\lean{Erdos249257.irrational\_totient\_series\_iff\_lcm\_diagonal\_certificate\_supply}{LcmConeFlatness.lean:426}) concerns only the pair \((H,2H)\) at LCM
height \(H=\mathrm{periodLcm}(t)\). Two natural strengthenings: (i) show
rationality of \(S\) is even more constrained than a single diagonal pair —
it flattens an \emph{entire cone} of ratios; (ii) conversely, exhibit
\emph{nonflatness} on a wider ``menu'' of cone vertices as a cheaper
sufficient certificate than a single pairwise \(\mathrm{certifiedKill}\).

\textbf{(b) Exact mechanism — flatness (necessary consequence of
rationality).}
\begin{prop}[Wave-24 LCM-cone flatness]
\(\neg\mathrm{Irrational}(S) \implies \exists t_1,\ \forall t\ge t_1,\
\forall q,m:\mathbb N,\ 0<q \implies \mathrm{totientTail}(q\cdot
\mathrm{periodLcm}(t)+m\cdot\mathrm{periodLcm}(t)) -
\mathrm{totientTail}(q\cdot\mathrm{periodLcm}(t)) \in
\mathrm{range}((\uparrow):\mathbb Z\to\mathbb R)\).
\end{prop}
\lean{rational\_totient\_series\_forces\_lcm\_cone\_flatness}{CertificateKernel.lean:19002} (body
\texttt{LcmConeFlatness.lean}). \ev{Lean} \scale{uniform} (for all
sufficiently large \(t\), \emph{given} rationality) \coord{other:lcm-ray-totient}.
Rationality forces one fractional constant on the entire LCM cone
\(\{k\cdot\mathrm{periodLcm}(t):k\ge1\}\) at every scale \(t\ge t_1\), not
merely the single diagonal pair — the umbrella collapse theorem
\lean{irrational\_totient\_series\_of\_lcm\_cone\_window\_kill\_supply}{CertificateKernel.lean:19014} shows any
\(\mathrm{certifiedKill}\) \emph{anywhere} on the two-multiplier cone, at
arbitrarily large \(t\), forces \(\mathrm{Irrational}(S)\).

\textbf{(b\('\)) Exact mechanism — nonflatness (a sharper sufficient
certificate).}
\begin{prop}[Wave-25 cone-nonflat menu refuter]
For a nonempty menu \(Q\) of positive vertex multipliers with a
\emph{one-sided} floor \(\forall q\in Q,\ q\cdot H+L+2<2^L\) (half the
pairwise floor of \(\mathrm{certifiedKill}\)): if \(\mathrm{coneNonflatCert}\,
H\,L\,Q\) fires (an argmin/Helly-avoidance argument: the vertices cannot all
share one fractional part, else the minimal-deep-tail vertex would be a
common left endpoint of every arc), then \(\exists q_i,q_j\in Q,\
\mathrm{totientTail}(q_j H)-\mathrm{totientTail}(q_i H)\notin
\mathrm{range}((\uparrow):\mathbb Z\to\mathbb R)\).
\end{prop}
\lean{exists\_nonintegral\_pair\_of\_coneNonflatCert}{LcmConeNonflat.lean:126} (supply theorem
\lean{irrational\_totient\_series\_of\_lcm\_cone\_nonflat\_supply}{CertificateKernel.lean:19140}). \ev{Lean}
\scale{bounded} (proved for a given finite menu \(Q\); an
\emph{unbounded}-scale menu is the open sufficient target)
\coord{other:lcm-ray-totient}.

\textbf{(c) Precise scope.} Flatness is a genuine \emph{necessary
consequence} of rationality — it is a statement about what rationality of
\(S\) would force, conditional on rationality, not an unconditional fact
about \(S\). Nonflatness is an unconditional \emph{sufficient} certificate
mechanism, information-theoretically half the depth floor of the pairwise
\(\mathrm{certifiedKill}\) (one-sided versus two-sided radius charging), but
it has been proved only for individual finite menus \(Q\); no menu of
unbounded scale has been supplied.

Neither direction decides \#249: flatness
gives a contrapositive route (exhibit \emph{any} cone nonintegrality at
arbitrarily large \(t\) and rationality is refuted) and nonflatness sharpens
the certificate needed, but the underlying producer obligation — a
nonintegral tail difference at cofinally many scales — is unchanged.

Note
also the coordinate warning from the manuscript's own Appendix C
(\lean{rational\_totient\_series\_forces\_lcm\_cone\_flatness}{LcmConeFlatness.lean:90}): cone flatness is a
\emph{necessary} consequence, and only the sharper diagonal-pincer /
full-target-avoidance normal forms
(\lean{diagonal\_int\_iff\_foreignDiagonalDefect\_hits\_fullTarget}{DiagonalPincerDecomposition.lean:215},
\lean{irrational\_totientSeries\_of\_full\_target\_avoidance\_supply}{DiagonalPincerDecomposition.lean:290}) are an exact
\(\mathrm{iff}\) with irrationality.

\textbf{(d) Salvage.} The argmin/Helly-avoidance combinatorial-geometry
technique behind cone-nonflatness is a reusable pattern for any modular-residue
pincer argument with more than two points, in any coordinate. The rank-2
second-difference certificate family built on the same window apparatus was
independently tested and found \emph{not} to be a shortcut over rank 1
(measured, not merely conjectured: at \((h,N)=(1,8)\), rank-1 fires at depth
8 while no rank-2 certificate exists at depth \(\le8\);
\lean{totient\_tail\_rank\_two\_kill\_sound\_but\_not\_shallower\_cell}{CertificateKernel.lean:19090}) — an
explicitly flagged dead end not to re-attempt without new information.

\subsection{The Mahler defect: dyadic totient-kernel finite rank per level,
infinite rank overall}
\label{sec:mahler-defect}

\textbf{(a) Route as conceived.} A finite-linear-algebra shortcut to \#249:
find a bounded-depth linear relation among the dyadic totient-kernel
channels \(n\mapsto\varphi(2^jn+r)\) that would compress the infinite
family to a finite-dimensional space, from which a rationality-forcing
contradiction (or a genuine rank obstruction) might be extracted.

\textbf{(b) Exact mechanism.}
\begin{thm}[Infinite dyadic totient-kernel rank]
For every depth \(e\ge0\), the canonical family of \(2^e+1\) dyadic
totient-kernel channels (\(\mathtt{card\_totientCanonicalIndex}\),
\lean{card\_totientCanonicalIndex}{TotientMahlerDefect.lean:61}) is linearly independent over
\(\mathbb Q\), proved via a \texttt{SeparatedMinorCertificate} (an explicit
finite evaluation-point assignment with nonzero determinant, forced by CRT +
Dirichlet's theorem on primes in arithmetic progressions — one channel made
prime, every other channel forced through a fresh prime \(\equiv1\bmod\) a
large power of \(2\)).
\end{thm}
\lean{linearIndependent\_canonicalTotientKernelFamily}{TotientMahlerDefect.lean:935}; consequently the full
infinite family spans an infinite-dimensional \(\mathbb Q\)-space,
\lean{not\_finiteDimensional\_span\_fullTotientKernel}{TotientMahlerDefect.lean:1145}. A companion
impossibility result closes the natural repair attempt directly: a bounded
\emph{compressed-adjoint certificate} — a triple \((Q,A,\mathrm{boundary})\)
with \(Q\cdot v\cdot A=\mathrm{boundary}\), \(|\mathrm{boundary}|<Q\cdot v\),
\(A\ne0\) — is provably impossible,
\lean{false\_of\_compressedAdjointCertificate}{TotientMahlerDefect.lean:1183}. \ev{Lean}
\scale{uniform} (holds for every \(e\); existence side is unconditional, via
Mathlib's CRT + primes-in-AP machinery) \coord{binary-digit} (dyadic
totient-kernel — \emph{not} the Möbius coordinate).

\paragraph{Integral coordinates and the complete relation module.}
The independence statement also determines the arithmetic of the channel
lattice.  Fix $e\ge1$, write $F_{j,r}(n)=\varphi(2^j n+r)$, and let
$I_e=\{(j,r):0\le j\le e,\ 0\le r<2^j\}$.  Retain $F_{0,0}$,
$F_{1,0}$, and the odd-residue channels at levels $1,\ldots,e$; call
this retained index set $J_e$.  The totient identities reduce every omitted
channel to $F_i=a_iF_{j(i)}$ for a retained channel $j(i)$ and an integer
$a_i\ge0$.  Thus the retained family generates the lattice
$L_e=\sum_{i\in I_e}\mathbb Z F_i$.  Its rational independence gives
integer independence as well: every element of $L_e$ has unique integer
coordinates in these $2^e+1$ retained channels.

To describe \emph{all} integral relations, let $E_i$ denote the formal
coordinate vector in $\mathbb Z^{I_e}$ and use the evaluation map
\[\adjustbox{max width=\linewidth}{$\displaystyle
 \operatorname{ev}_e:\mathbb Z^{I_e}\longrightarrow L_e,
 \qquad (c_i)\longmapsto\sum_{i\in I_e}c_iF_i.
$}\]
For each omitted index put $R_i=E_i-a_iE_{j(i)}$.  Then
\[\adjustbox{max width=\linewidth}{$\displaystyle
 \ker(\operatorname{ev}_e)
   =\bigoplus_{i\in I_e\mathbin{\backslash} J_e}\mathbb Z R_i,
 \qquad
 \operatorname{rank}_{\mathbb Z}\ker(\operatorname{ev}_e)=2^e-2.
$}\]
Indeed, subtracting $\sum_{i\notin J_e}c_iR_i$ from a relation removes
all omitted coordinates.  The remainder is a relation among retained
channels, so independence makes it zero.  Conversely, the coefficient of
$E_i$ in $\sum_{h\notin J_e}b_hR_h$ is exactly $b_i$, proving uniqueness.
The decisive feature is the unit pivot $1$ in each omitted coordinate
(or $-1$ if a relation is oriented oppositely).  Elimination uses no division;
it therefore identifies the integral relation lattice, not merely its
rational span.  The depth-two example in the short note illustrates these
two-term rows.

The \href{https://github.com/wcook04/plectis-erdos/blob/25ef6245d15a47548c6926369ae8f1a0f0a14a80/ErdosProblems/Erdos249/PaperCompleteR8/KernelRelationBasis.lean#L398}{integral normal form} and the \href{https://github.com/wcook04/plectis-erdos/blob/25ef6245d15a47548c6926369ae8f1a0f0a14a80/ErdosProblems/Erdos249/PaperCompleteR8/FullKernelAssemblies.lean#L156}{full dyadic rational basis} are kernel-checked in Lean.
The formal statement gives the same construction for every integer base
$k\ge2$, with relation rank $\sum_{j=1}^{e-1}k^j$.  The full dyadic
assembly separately gives a basis for the infinite rational span and for
its finitely supported rational relations.  These finite and infinite
statements should not be confused with a claim about infinite sums of
relations.

\textbf{(c) Precise scope — coordinate-relative, stated explicitly by the
source module.} This does \emph{not} show irrationality of \(S\). It proves
the dyadic-kernel \emph{side} is infinite-rank; the module's own docstring
names the missing input as ``a rationality-side finite-rank compression, or
equivalent contradiction.''

The companion necessary-consequence-of-rationality
theorem, \(\neg\mathrm{Irrational}(S) \implies \exists v>0,\exists U:\mathbb
N\to\mathbb Z,\ \mathrm{IsTemperedBinaryOrbit}(\varphi,v,U)\wedge\forall
e,\ 2^e-1\le\mathrm{finrank}_{\mathbb Q}\,\mathrm{span}(\mathrm{range}(
\mathrm{canonicalCarryKernelFamily}(U,e)))\)
(\lean{not\_irrational\_totientSeries\_implies\_unbounded\_carryRank\_unconditional}{TotientCarryKernelRigidity.lean:300}),
shows the SCALE side is finished — the \(2^e-1\) floor holds for all \(e\)
with a uniform proof — but the opposite inequality (a rationality-side rank
\emph{upper} bound) is exactly what is missing, and its absence is the
entire open content here.

A \emph{third}, independent coordinate exhibits
the same coordinate-relative phenomenon: the Möbius-incidence companion
matrix \(U_N(i,j)=\mu((i+1)/(j+1))\) if \((j+1)\mid(i+1)\) else \(0\) is
lower triangular with diagonal \(1\), hence \(\det U_N=1\) for every \(N\),
hence its jet-evaluation map is injective — a finite ``incidence quotient''
compression is impossible at any finite horizon in the Möbius coordinate too
(\lean{incidenceMobius\_det\_eq\_one}{IncidenceQuotientHermitePade.lean:56}). Both natural finite
truncations — dyadic-residue and Möbius-incidence — turn out to have no
nontrivial kernel, in two completely different coordinates. That is strong
evidence any winning finite-linear-algebra shortcut, if one exists, must
live in a genuinely third coordinate or use an infinite/growing-parameter
construction, not evidence that no shortcut exists at all.

\textbf{(d) Salvage.} \lean{SeparatedMinorCertificate}{TotientMahlerDefect.lean:83} and
\lean{linearIndependent\_of\_separatedMinorCertificate}{TotientMahlerDefect.lean:91} are fully generic finite-rank
linear-independence infrastructure — the indexed family is a free variable,
reusable for any finite-rank obstruction argument on either problem (e.g.
testing independence of Mersenne-indexed sequences for \#257 via an explicit
nonzero minor). The CRT+Dirichlet evaluation-forcing technique (make one
channel prime, force every other channel through a fresh arithmetic-progression
prime) is a reusable construction pattern independent of the totient
specifics.

\subsection{Further closed routes, catalogued for completeness}

The following additional no-gos and route-pruning results were read in full
and are recorded here in compressed form; each follows the same (a)/(b)/(c)/(d)
discipline as above but is presented as a table row for space.

\begin{landscape}
\small
\begin{longtable}{@{}p{0.20\linewidth}p{0.38\linewidth}p{0.36\linewidth}@{}}
\toprule
\papertablehead
\textbf{Route} & \textbf{Mechanism \& site} & \textbf{Scope / salvage} \\
\midrule
\endfirsthead
\toprule
\papertablehead
\textbf{Route} & \textbf{Mechanism \& site} & \textbf{Scope / salvage} \\
\midrule
\endhead
\bottomrule
\endlastfoot
Full terminal dyadic staircase (annihilate every suffix letter mod growing
powers of 2) &
\(a\ge8\), room bound \(\Rightarrow\) \texttt{ActualLcmTerminalDyadicStaircase}
is false: the terminal letter would have to be positive (by short-window
positivity) and strictly below a wider-than-itself modulus while divisible
by it, forcing it to \(0\), contradicting positivity.
\lean{not\_actualLcmTerminalDyadicStaircase\_of\_room}{TotientActualLcmTopEdgeStaircase.lean:251}.
\ev{Lean} \scale{uniform} \coord{binary-digit} &
Kills the FULL terminal-staircase proof strategy outright and
unconditionally; the \emph{punctured} staircase (all but the last letter
vanish) survives and is pinned exactly to the half-turn value \(2^{m-1}\) at
the penultimate letter
(\lean{puncturedDyadicStaircase\_penultimate\_eq\_half}{TotientActualLcmTopEdgeStaircase.lean:1187}). The
generic mechanism (\(0<e<2^m\wedge2^m\mid e\Rightarrow e=0\)) is reusable
anywhere a positive quantity is asked to vanish mod a wider modulus. \\
LCM factor-ideal retention (keep only the homogeneous ``factor-only'' part
of an LCM-ray decomposition) &
An explicit nonzero all-horizon countermodel, built from multiples of
\(\varphi(\mathrm{periodLcm}(t))\), agrees with every exact whole-ray anchor
\(\Delta\varphi(H,qH)=\varphi(H)\) for \(2\le q<t\) and survives every finite
commensurate LCM-cube shift-polynomial, at every finite rank.
\lean{LcmFactorIdealPulseObstruction.lean:1--24}{LcmFactorIdealPulseObstruction.lean:1-24} (module docstring;
theorem bodies \ev{Cert}, DOCSTRING-SOURCED). \scale{uniform}
\coord{other:lcm-ray-totient} &
This is a \emph{representation}-level exclusion, not an object-level one: it
shows the homogeneous factor-ideal projection discards fresh Möbius-channel
information that can be adversarially reconstructed. Explicitly flagged
SYNTHETIC — no claim that the compensation letters occur as actual totient
differences. Salvage: a future factor-ideal argument must control the
\emph{fresh} channel via
\lean{totient\_eq\_sum\_mobiusTotientChannel}{LcmFactorIdealPulseObstruction.lean:328}, not just the old one. \\
Two-prime ``full-target diamond'' (four simultaneous ray hits certify more
than one) &
The four-hit diamond
\(\mathrm{Hit}(H)\wedge\mathrm{Hit}(pH)\wedge\mathrm{Hit}(qH)\wedge
\mathrm{Hit}(pqH)\) is logically equivalent to \(\mathrm{Hit}(H)\) alone, no
primality of \(p,q\) needed — hits transport multiplicatively for free
because \(\mathrm{Hit}(H)\iff\mathrm{IsIntegralValue}(\text{scaleDiagonalTailDifference}(H))\)
and integrality transports affinely along every ray \(H\mapsto kH\).
\lean{Erdos249257.fullTarget\_primeAdjunction\_diamond\_iff\_root}{FullTargetPrimeAdjunctionNoGo.lean:196}.
\ev{Lean} \scale{uniform} \coord{other:mobius-mersenne} &
The 4-condition diamond carries zero information beyond its base point — a
``curvature-zero collapse'' of a multi-point certificate. Any future
positive-holonomy argument needs an independently-defined projection of the
foreign state \emph{plus} control of the discarded complement; a bare
2/4-point diamond cannot supply new information beyond its own scale. \\
Fixed-precision local valuation-unit signature (2-adic ``tropical curvature
carry'' attack) &
For any finite word of odd-unit-part 2-adic symbols at fixed precision
\(u>0\) and any starting carry state \(e\), there exists a compatible carry
orbit realising the word with every intermediate state centred inside its
symbol's dyadic radius.
\lean{Erdos249257.TotientTailPeriodKiller.fixedPrecisionTropicalNoGo}{TropicalCurvatureCarry.lean:137}.
\ev{Lean} \scale{bounded} \coord{p-adic} &
Bounded LOCAL valuation-unit data at FIXED precision can never exclude all
finite centred carry completions — no obstruction is derivable from a
fixed-window local signature alone; growing precision, or extra arithmetic
coupling, is required. Pure 2-adic dynamics, zero totient content, reusable
against any similarly-shaped fixed-window carry-certificate attack. \\
Fixed-depth affine carry reset (distinguishing two carry histories after a
long common tail) &
For the affine binary orbit \(\mathrm{orbit}(0)=u_0\),
\(\mathrm{orbit}(n{+}1)=2\,\mathrm{orbit}(n)-a(n{+}1)\) with common forcing
word \(a\): \(\mathrm{orbit}_u(L)-\mathrm{orbit}_v(L)=2^L(u_0-v_0)\) exactly.
\lean{Erdos249257.affineBinaryOrbit\_mod\_twoPow\_eq}{GenericTailOrbitRigidity.lean:307}. \ev{Lean}
\scale{uniform} \coord{binary-digit} &
After \(L\) common steps the endpoint residue mod \(2^L\) is
\emph{independent} of the initial carry — long common tails erase
predecessor information exactly, not approximately. Directly relevant to any
argument hoping to distinguish two carry histories from a shared long
suffix; the terminal forcing word alone determines the residue. \\
Bounded finite-state / autonomous successor decoder for the binary carry &
For a balanced-pulse family at location \(m\) whose predecessor \texttt{State}
is constant across the whole family, no \(\mathrm{decode}:\mathrm{State}\to
\mathbb N\) recovers the radius parameter; any finite \texttt{Fintype State}
needs \(\mathrm{card}\ge\lfloor m/2\rfloor+2\), unbounded in \(m\).
\lean{Erdos249257.balancedPulse\_no\_autonomous\_decoder}{GenericTailOrbitRigidity.lean:247}. \ev{Lean}
\scale{uniform} \coord{binary-digit} &
Rules out, unconditionally, ANY proof strategy that defines a bounded or
autonomous ``carry state'' summarising pre-\(m\) history and claims it
exactly determines the post-\(m\) tail — independent of whether the
coefficients are \(\varphi\) or a Möbius-support indicator. A hard stop
against finite-automaton-computes-the-expansion style attacks on either
problem. \\
Period-4 sign weight as a route to A7's ``frequently nonzero'' hypothesis for
free &
The divisor coefficient of the period-4 sign weight
(\(w(a){=}1\) if \(a{\equiv}1\), \(-1\) if \(a{\equiv}3\), else \(0\), mod
4) vanishes identically at every \(n\equiv3\pmod4\), via the divisor-pairing
involution \(d\mapsto n/d\).
\lean{Erdos249257.intWeightedCoeff\_periodFourSignWeight\_eq\_zero\_of\_mod\_four\_eq\_three}{CertificateKernel.lean:14226}.
\ev{Lean} \scale{fixed} \coord{mobius-mersenne} &
Witnesses that Dirichlet-convolution coefficients CAN vanish identically on
an arithmetic progression, so the ``frequently nonzero'' hypothesis in the
nonnegative-periodic irrationality closure
(\lean{irrational\_intWeightedErdosSeries\_periodic\_of\_coeff\_nonneg\_of\_frequently\_ne\_zero}{CertificateKernel.lean:14583})
is not free — it is why the corpus lands a dichotomy for general periodic
signed weights rather than an unconditional supply theorem. \\
Certified finite-depth ``death'' certificates as a route to membership proofs
(both problems) &
\(\mathtt{CertifiedGreedyMersenneDeath}(x,\mathrm{level},\mathrm{lookahead})\)
is a decidable, finite-depth certificate proving non-membership only
(\lean{greedyMersenneDeath\_not\_mem}{GreedyAchievementSet.lean:1747}; witness
\lean{three\_fourths\_certifiedGreedyMersenneDeath}{GreedyAchievementSet.lean:1778}, \(3/4\) certified dead at
level 1). \ev{Cert} \scale{bounded} \coord{greedy-orbit} &
Structural one-sidedness, recurring across the WHOLE certified-kill family
(\lean{TotientTailPeriodKiller.lean}{TotientTailPeriodKiller.lean}, \lean{LcmConeFlatness.lean}{LcmConeFlatness.lean}, this one):
certified-death \(\Rightarrow\) non-membership, but absence of a found
certificate, or survival through any finite depth, proves \emph{nothing}
about membership or rationality. \\
Reconstructing one totient value from a two-tail absolute-value linear
combination &
For any finite \(w:\iota\to\mathbb Q\), \(x:\iota\to\mathbb N\) with
\(\sum w_i\varphi(x_i)=1\) (exact isolation), the crude two-tail cost
\(\sum|w_i|(2(x_i+1)+(x_i+2))\ge3\), using only \(\varphi(x)\le x\).
\lean{Erdos257PeriodNoncollapse.three\_le\_totientAdjugateTailCost}{TotientTailCarryPeriod.lean:266}
(closure \lean{not\_totientAdjugateTailCost\_lt\_one}{TotientTailCarryPeriod.lean:307}). \ev{Lean}
\scale{uniform} \coord{other:totient-adjugate-linear-algebra} &
The strategy needs cost \(<1\) and gets \(\ge3\) at ANY finite grid height —
kills the entire ``isolate one coefficient via absolute-value two-tail
bookkeeping'' family before it starts. The proof mechanism is problem-agnostic
(any \(c(n)\le n\) sequence gets the same floor by the triangle inequality)
and transfers verbatim to a \#257 reformulation. \\
Homogeneous Mersenne-primitive prime factor alone forces a contradiction &
For the completely-multiplicative control \(c(n)=n\) (zero totient content):
if \(K<q\), \(q\mid2^K-1\), the corresponding tail shift is integral at every
\(N\) yet \(q\nmid K\). \lean{idCoeff\_mersenneModulus\_does\_not\_annihilate\_tailShift}{TotientTailCarryPeriod.lean:342}.
\ev{Lean} \scale{n/a} \coord{other:mersenne-modulus} &
A primitive prime factor of the homogeneous multiplier \(2^K-1\) alone
supplies NO contradiction from tail integrality. A direct cross-problem
warning: \#257's denominators \(2^n-1\) are literally this \(q\mid2^K-1\)
object, so ``a large primitive Mersenne prime factor alone forces a
contradiction'' should be checked against this lemma before being attempted
on either problem. \\
Prime-power reduced-denominator unit-gap strengthening (rescuing extra Farey
lattice points) &
At a prime-power reduced denominator \(p^e\), the unit-gap refinement can
rescue at most ONE additional candidate lattice point beyond the ordinary
gap certificate. \lean{reducedDenominatorCandidateCount\_prime\_pow\_le\_one}{PrimitiveWeightCertificate.lean:111}
(iff-characterisation \lean{reducedDenominatorUnitGapCert\_prime\_pow\_iff}{PrimitiveWeightCertificate.lean:54}). \ev{Lean}
\scale{uniform} \coord{other:farey-gap} &
A formal ceiling on how much this specific refinement can ever buy — a
known dead end for extending the \(\approx7.96\times10^{34}\) Farey rung
(\S\ref{sec:adelic-height}) unboundedly. Do not re-attempt this exact
strengthening expecting more than +1 lattice point per prime power; the
underlying reduced-denominator-implies-unit-numerator technique is reusable
for other denominator-exclusion arguments. \\
Fixed rank-3 finite-difference kernel squeezing more than one dyadic bit of
curvature information &
\(\mathrm{fixedRankSecondDifference}(2^{n+1}\cdot2,\,2^{n+1}\cdot3)=-2^{n+1}\)
exactly, for all \(n\): the two-adic valuation gained by the primitive
kernel factor \(2\varphi(j)\) is exactly tight.
\lean{fixedRankSecondDifference\_dyadic\_fixture\_exact}{TotientFixedRankLcmAsymptotic.lean:462}. \ev{Lean}
\scale{n/a} \coord{binary-digit} &
Route-pruning ammunition: rules out extracting more than one dyadic bit from
the bare 3-rank affine kernel \((1,-2,1)\); any further progress on the
fixed-rank curvature route needs new arithmetic input, not a sharper
valuation squeeze from the same kernel. The affine rank-3 kernel
classification itself
(\lean{affine\_rank\_three\_kernel\_classification}{TotientFixedRankLcmAsymptotic.lean:338}, pure linear algebra
over \(\mathbb Z\), zero number-theoretic content) is fully reusable for any
3-point finite-difference argument on either problem. \\
\end{longtable}
\end{landscape}

\section{Why the bounds are load-bearing}
\label{sec:promotion-audit}

The near-miss interface between the corpus's proved results and the four
open supply obligations — \(\#249\)-supply, \(\#257\)-reset,
\(\#257\)-cofinal-rows, and \(\#257\)-universal — was swept for
\emph{promotion candidates}: proved theorems whose \emph{stated} hypotheses
looked, on a first read of the \texttt{yields} clause, close enough to the
open obligation's shape that a routine strengthening (raising a scale bound,
widening a hypothesis, or reindexing a quantifier) might close the gap
without new mathematics. Fifty-six near-miss rows were catalogued across the
four obligations; eighteen of them were flagged \texttt{promotable\_claimed:
True} on this first pass. \textbf{Every one of the eighteen proofs was then
opened and read in full}, not merely re-inspected at the statement level,
and the promotion claim was checked against the actual proof body. The
audit table below reports the verdict for all eighteen.

Of the eighteen, \textbf{sixteen are conclusively NOT\_PROMOTABLE}: reading
the proof body exposes a documented arithmetic, logical, or coordinate
obstruction that a routine strengthening cannot cross, and the row's own
record states exactly what new input would be required.

\textbf{Two}
(\texttt{NM-02}, \texttt{NM-03} in the \(\#257\)-universal obligation) are
genuine exceptions: the auditor found the promotion is available with, in
its own words, ``no new mathematics'' — a pure widening of an already-general
proof to a wider stated target. Crucially, \emph{neither of these two closes
anything}: they widen an equivalence or a corollary from one fixed target
value to a family of targets, without touching the missing arithmetic
content (a cofinal certificate supply, a rank upper bound, an anti-concentration
estimate) that every genuine closure of \(\#249\) or \(\#257\) requires.

So
while the promotion audit is not literally eighteen-for-eighteen at the level
of ``can this exact statement be re-derived for a wider scope,'' it is
eighteen-for-eighteen at the level that matters for this paper's obligations:
\textbf{no promotion candidate in the audited set supplies missing arithmetic
content toward closing \#249 or \#257}. The two mechanically-widenable rows
are recorded honestly below rather than folded into the sixteen, because
collapsing that distinction would itself be exactly the kind of
quantifier-slippage error this paper's accuracy rules forbid.

\begin{landscape}
\small
\begin{longtable}{@{}p{0.12\linewidth}p{0.18\linewidth}p{0.13\linewidth}p{0.51\linewidth}@{}}
\toprule
\papertablehead
\textbf{Obligation} & \textbf{Row / site} & \textbf{Gap kind} &
\textbf{Verdict and exact blocker} \\
\midrule
\endfirsthead
\toprule
\papertablehead
\textbf{Obligation} & \textbf{Row / site} & \textbf{Gap kind} &
\textbf{Verdict and exact blocker} \\
\midrule
\endhead
\bottomrule
\endlastfoot
249-supply & e1-companion,
\lean{exists\_certifiedKill\_of\_first\_harmonic\_gap}{FirstHarmonicGap.lean:128} & hypothesis strength &
Not promotable. The consumer already has the obligation's exact quantifier
shape; the missing input is a single unproved arithmetic fact — a
constant-saving first-harmonic (Weyl-sum) cancellation bound
\(\sum_{N\in[X,2X)}\cos(2\pi\cdot(\ldots)/2^L)\le(9/10)X\) — not proved at
any \(X,h,L\) anywhere in the corpus. \\
249-supply & SGN-01,
\lean{actualLcmTailDiff\_shift\_pos}{TotientActualLcmOrbitSign.lean:39} & hypothesis strength &
Not promotable. Positivity is exactly HALF of the needed certificate: the
companion theorem in the same file proves integrality forces the residue to
the TOP edge, not a kill — the lower half of the band is discharged
unconditionally and cofinally, the upper half is untouched. \\
249-supply & TE-04,
\lean{actualLcmTailDiff\_notMem\_int\_of\_topEdgeResidueGap}{TotientActualLcmTopEdgeStaircase.lean:1260} & hypothesis strength &
Not promotable. The consumer side is finished and cofinal (proved for
every \(a\ge8\)); the residual is purely the one-sided residue-gap
\emph{producer}, unconditionally unsupplied at even one large \(a\). \\
249-supply & TE-05-weakest,
\lean{PowerTwoFlexibleActualTerminalDominanceSupply}{TotientActualLcmTopEdgeStaircase.lean:1908} & hypothesis strength &
Not promotable. The exact identity pinning the target shows the dominance
inequality is equivalent to \(\mathrm{carryOrbit}\le0\), and the sign
machinery (SGN-02) already proves the true carry orbit is strictly
\emph{positive} under integrality — the corridor-escape branch the identity
naturally supplies is exactly the branch already eliminated. \\
249-supply & SEP-02,
\lean{abs\_actualLcmTailOrbit\_sub\_rawApprox\_lt}{TotientActualLcmOrbitSeparation.lean:141} & scale only &
Not promotable. The analytic tail is fully discharged with an explicit,
uniform error radius; only two exponents (\(a=4,6\)) have the resulting
finite distance-to-integer question verified unconditionally
(\lean{actualLcmTailOrbit\_four\_and\_six\_notMem\_int}{TotientActualLcmShortKill.lean:35}); nothing beyond
\(a=6\) is proved. \\
249-supply & b7,
\lean{upperHalfMersenneProduct\_lower\_bound}{MersenneShadowDenominatorGrowth.lean:60} & coordinate only &
Not promotable. The only proved unbounded-growth quantity in the
Möbius--Mersenne coordinate, uniform in \(t\), but there is no landed
transport from a denominator lower bound to \(\mathrm{certifiedKill}\) or to
totient-tail non-integrality — a different coordinate from the obligation. \\
249-supply & d4/d5,
\lean{linearIndependent\_canonicalTotientKernelFamily}{TotientMahlerDefect.lean:935}, \lean{not\_irrational\_totientSeries\_implies\_unbounded\_carryRank\_unconditional}{TotientCarryKernelRigidity.lean:300} &
coordinate only & Not promotable. The scale side of a second, independent
open obligation (unbounded carry-kernel rank forced by rationality) is
finished uniformly in \(e\); the opposite (rationality-side rank upper
bound) is missing, and one natural route to it is proved dead
(\lean{false\_of\_compressedAdjointCertificate}{TotientMahlerDefect.lean:1183}). \\
257-reset & rc-2 / F4-adjacent,
\lean{one\_third\_lt\_scaledMersenneTail\_floorWall}{HalfCylinderMiddleCarryLowerBound.lean:2308} & multiple &
Not promotable. The excess-side run-length law is only half-assembled: an
exact iterate identity exists, and an upper bound on the terminal excess
exists separately, but no theorem in Lean composes the two into an
excess-side run-length statement, even though every ingredient is present. \\
257-reset & E1,
\lean{SeamGreedyRemainderGapBound}{HalfCylinderSeamLimit.lean:192} & coordinate only &
Not promotable. Supplies an UPPER bound on the remainder (the ceiling half
of the picture); the obligation needs a LOWER bound on the deviation
magnitude — the two do not compose without an additional input. \\
257-reset & rc-8,
prose \emph{erdos257\_reset\_crossing\_unification\_2026\_07\_24.md}\,\S6 &
multiple & Not promotable. Two gaps at once: the sign law is
\ev{Cert}/empirical only (flagged \ev{Open} as an unconditional lemma, not
proved), and even if proved it delivers only the SIGN of the deviation, not
its magnitude, which is what the obligation needs. \\
257-cofinal-rows & TH-sharp-capacity-progress,
\lean{exists\_exactRowStrictUpperFill\_of\_skippedCoreSharpCapacity}{BooleanMobiusSkipRow.lean:238} & hypothesis strength &
Not promotable. Unconditional and strictly more general than every actual
consumer, but it needs a per-\(c\) input
(\(\mathrm{localBinarySuffix}\,D\,1\,(2c{-}2)<2^{c-2}\)) that the corpus
only ever supplies through a route which then rigidly collapses onto one
specific support family — the deficit hypothesis is a self-imposed
restriction of the consumers, and the general input itself is unsupplied. \\
257-cofinal-rows & precriticalSuffix\_lt\_of\_future\_skip\_after\_takenBlock,
\lean{halfGreedy\_precriticalSuffix\_lt\_of\_future\_skip\_after\_takenBlock}{BooleanMobiusCriticalCapacityCofinal.lean:603} &
hypothesis strength & Not promotable. Unconditional and uniform in both
parameters \((c,t)\); the sole missing input is an orbit-level skip-gap
bound on the rational half-greedy orbit, never proved and never previously
isolated as a target anywhere in the banks. \\
257-cofinal-rows & greedyHalfFrozenMargin\_nonneg,
\lean{exists\_greedyHalfFrozenMargin\_nonneg\_of\_excess\_neg}{HalfCylinderFiniteShadow.lean:1198} & scale only &
Not promotable. The margin-nonnegativity horizon is produced by a
\emph{non-effective} limit argument (an existential horizon from a
convergence fact); the exact-row route needs the SPECIFIC effective horizon
\(J=k-2\), and effectivising it is explicitly recorded as unfinished
bookkeeping, not new mathematics — but it is still unfinished. \\
257-cofinal-rows & C4a,
\lean{HalfCarryCofinalTerminalOnlyStrip}{TerminalOnlyCofinal.lean:34} & multiple &
Not promotable. Trades off against the obligation in the opposite
direction on \emph{support locality}: tolerates any support inside a depth
window but demands a carry inside a tight \(\sim2\sqrt M\) strip, where the
obligation tolerates an exponentially looser carry but demands the support
be confined to a lower window — the gap is support locality, not carry
size, and neither socket implies the other. \\
257-universal & NM-02,
\lean{infinite\_greedyMersenneSkippedSupport\_of\_half\_mem}{GreedyAchievementSet.lean:2547} & scale only &
\textbf{PROMOTABLE (no new mathematics), but does not close anything.}
The proof was opened and found already uniform in the target in every
internal step; only the statement's own quantifier is artificially pinned
to \(t=1/2\). Promoting it converts the whole half-greedy refutation machine
from a one-target equivalence to one covering every rational target — the
highest-value mechanical promotion in the bank — but it supplies no new
arithmetic content: it is a restatement, not a producer. \\
257-universal & NM-03,
\lean{finite\_boolSupport\_ne\_half}{HalfCarryReachability.lean:589} & scale only &
\textbf{PROMOTABLE (no new mathematics), but does not close anything.} The
headline is pinned to target \(1/2\) while its engine
(\lean{finiteErdosSum\_den\_odd}{DyadicPrefixCompression.lean:1513}) is already fully general over
finite supports; verbatim re-derivation at any even-denominator target is
available with zero new arithmetic. Paired with NM-02: universal \#257 at
\(b=2\) is false iff some rational \(t\) with even reduced denominator has
an infinitely-skipping greedy Mersenne orbit — a genuine reformulation, not
a resolution. \\
257-universal & NM-04,
\lean{dyadic\_support\_fraction\_reciprocalMass\_diverges\_or\_gt\_one}{RationalSupportCarrySkeleton.lean:2210} & coordinate only &
Not promotable (partial mechanical content only). The dyadic restriction
\(v=1\) is a call-site choice, not a proof constraint — the general-\(v\)
conclusion is already landed for arbitrary \((v,h,L,F)\) at
\lean{shiftedBooleanCollision\_wrap\_bound\_of\_common\_multiple}{RationalSupportCarrySkeleton.lean:1968} — but
the strict form needed for non-dyadic targets requires one genuinely new
auxiliary bound (a mean-least-residue estimate) not present in the corpus. \\
257-universal & NM-12,
\lean{half\_mem\_mersenneAchievementSet\_or\_exists\_fatal\_gap}{HalfCutLocator.lean:623} & scale only &
Not promotable (unverified, not merely blocked). The killer and the
dichotomy scaffold are already target-generic; the two endpoint kills that
pin the dichotomy to \(1/2\) generalise on their face, but their proof
bodies were explicitly \emph{not} read in this audit pass — recorded as a
strong promotability hypothesis, not a verified promotion, since a
\texttt{norm\_num}-style numeric step could hide a \(1/2\)-specific bound. \\
\end{longtable}
\end{landscape}

\subsection*{Classification of the sixteen genuine blockers}

Reading all eighteen proof bodies rather than trusting the \texttt{yields}
clause exposes a small, recurring taxonomy of blocker shapes, tagged by
\texttt{gap\_kind} in the table above:

\begin{itemize}[nosep]
\item \textbf{hypothesis\_strength} (6 rows: e1-companion, SGN-01, TE-04,
TE-05-weakest, TH-sharp-capacity-progress,
precriticalSuffix\_lt\_of\_future\_skip\_after\_takenBlock). The consumer
theorem is already at the obligation's exact quantifier shape; a single
named arithmetic fact — a Weyl-sum cancellation bound, a one-sided residue
gap, an orbit-level skip-gap bound — is missing and is not a re-derivation
of anything on disk.
\item \textbf{scale\_only} (5 rows: SEP-02, greedyHalfFrozenMargin\_nonneg,
NM-02, NM-03, NM-12). The mathematics is either already general or reduces
to a finite verified range; what is missing is either extending a finite
census (SEP-02) or effectivising a non-constructive existence bound
(greedyHalfFrozenMargin\_nonneg). NM-02 and NM-03 are the two genuine
exceptions where scale widening costs nothing — and, being restatements, buy
nothing toward closure either.
\item \textbf{coordinate\_only} (4 rows: b7, d4/d5, E1, NM-04). A quantity
is proved unconditionally in one coordinate (Möbius--Mersenne denominator
growth, dyadic-kernel rank, remainder upper bound, dyadic support-fraction
mass) with no landed transport into the coordinate the obligation is stated
in.
\item \textbf{multiple} (3 rows: rc-2/F4-adjacent, rc-8, C4a). More than
one of the above simultaneously — typically an unassembled composition of
two otherwise-landed halves, or two orthogonal deficiencies (empirical-only
plus sign-not-magnitude) stacked on the same row.
\end{itemize}

No row in the audited eighteen falls under \texttt{quantifier\_order} — every
row with a genuine quantifier-order mismatch (e.g. \texttt{a12}'s
\(\exists N\,\exists L\,\forall h\) versus the obligation's
\(\forall h\,\forall N_0\,\exists N\,\exists L\)) was judged NOT promotable
at first read and excluded from the eighteen entirely, rather than surviving
to a full proof-body audit. The eighteen audited here are precisely the rows
whose \emph{stated} hypotheses looked closest to the obligation; that even
this most-favourable subset yields sixteen genuine blockers and only two
costless-but-empty restatements is the paper's evidence that the corpus's
bounds are load-bearing in the strict sense: no amount of routine
strengthening of what is already proved, at the level the audit checked,
supplies the missing arithmetic content that \(\#249\) or \(\#257\)
actually needs.
% ---- part a249_invent ----
\section{The mathematics this problem still needs}

Everything in Parts I--IV of this document is either a theorem about
$\varphi$ or a theorem about proof methods for $\varphi$. Nothing in it
decides Erd\H{o}s \#249, and this section does not either. What this section
does is different in kind from the rest: it takes each route that survives
the barrier classification, states the missing mathematics as an exact
sentence, identifies the shape of argument that could supply it, names the
precise point at which the nearest existing technique fails, and --- where
the evidence permits --- attempts the construction rather than describing it.

Four things are new here and are marked as such. (i) An exact normal form
that turns every surviving \#249 route into a statement about the binary
expansion of one explicit real number, with the thresholds computed
(Lemma~\ref{lem:orbit}, Proposition~\ref{prop:transfer},
Corollary~\ref{cor:digitform}). (ii) A proof that the first-harmonic gap is
strictly stronger than irrationality, by an explicit witness inside the same
coefficient class (Theorem~\ref{thm:lacunary}); the corpus had this only as
a plausibility argument. (iii) A one-sided smooth-number majorant for the
non-supplier budget and a bad-cofactor estimate, with their remaining
uniformity hypotheses exposed (Propositions~\ref{prop:dickman} and
\ref{prop:badcof}); the quantifier audit also shows that the proposed
shallow-modulus Siegel--Walfisz route does not fit the predicate. (iv) An
audit of the rationality-side rank route: its proposed generic counterexample
loses one zero-residue section per level, so it neither refutes nor retires
the route.

Where a claim is proved it is marked \ev{Math}; where it is checked by exact
or floating-point computation, \ev{Cert}; where it is quoted, \ev{Cited};
where it is a proposal, \ev{Open}. No proposal below is offered as progress
towards a solution. Two of them are reductions, and a reduction is not a
result.

\begin{rem}
Declarations in the shared tree are cited as
\texttt{Erdos249257.name} with file and line. Declarations in the newer
per-problem tree are cited with their path from the repository root,
\texttt{ErdosProblems/Erdos249/File.lean:line}; the hyperlink target for
those is the per-problem directory, not \texttt{Erdos249257/}.
\end{rem}

\subsection{One identity, and what it does to every surviving route}

Write $S=\sum_{n\ge 1}\varphi(n)/2^n$, $R_N=\sum_{m\ge 1}\varphi(N+m)/2^m$
for the local tail
(\lean{Erdos249257.\allowbreak TotientTailPeriodKiller.\allowbreak totientTail}{TotientTailPeriodKiller.lean:57}),
and $\Phi_N=\sum_{n\le N}\varphi(n)2^{N-n}\in\Z$ for the integer prefix, so
that $2^N S=\Phi_N+R_N$. For $h\ge 1$ set
\[
  \alpha_h \;:=\; (2^h-1)\,S .
\]

\begin{lem}[Doubling normal form]\label{lem:orbit}
For all $N\ge 0$ and $h\ge 1$,
\[
  R_{N+1}=2R_N-\varphi(N+1),
  \qquad
  R_{N+h}-R_N \;=\; 2^N\alpha_h-\bigl(\Phi_{N+h}-\Phi_N\bigr),
\]
so $R_{N+h}-R_N \equiv 2^N\alpha_h \pmod 1$, and consequently the exact
first character of the tail difference is the $\times 2$ orbit of one real
number:
\[
  \mathrm{tailOrbitFirstExp}(h,N)
  \;=\; e\bigl(2^N\alpha_h\bigr),
  \qquad e(x):=\exp(2\pi i x).
\]
Hence, for fixed $h$, the phases $\{\,R_{N+h}-R_N \bmod 1\,\}_{N\ge 0}$ are
the forward orbit of $\alpha_h \bmod 1$ under $x\mapsto 2x$.
\end{lem}

\noindent
The second identity and the character form are landed:
\lean{Erdos249257.\allowbreak TotientTailPeriodKiller.\allowbreak tail_diff_eq_scaled_\allowbreak totient_series_sub_prefix}{PivotAntiReconstruction.lean:54}
and
\lean{Erdos249257.\allowbreak TotientTailPeriodKiller.\allowbreak tailOrbitFirstExp_eq_scaledTotientSeriesFirstExp}{PivotAntiReconstruction.lean:66}
\quad \ev{Lean} \quad \scale{uniform} \quad \coord{other:binary-window}.
The recurrence $R_{N+1}=2R_N-\varphi(N+1)$ is one line from the definition
and is used below only for intuition \ev{Math}.

This is not a new fact --- it is two landed Lean theorems read together ---
but reading them together has a consequence the corpus never draws. The
first-harmonic block sum, which is the analytic frontier of \#249, is a
\emph{lacunary exponential sum in a single real variable}. The following
proposition makes the transfer exact, including the truncation budget.

\begin{prop}[Orbit transfer]\label{prop:transfer}
Define
\[
  \mathrm{OrbitBlockGap} :\equiv
  \forall h\ge 1\ \forall X_0\ \exists X\ge \max(X_0,1):\quad
  \sum_{N=X}^{2X-1}\cos\bigl(2\pi\,2^{N}\alpha_h\bigr)\;\le\;\tfrac{89}{100}\,X .
\]
Then $\mathrm{OrbitBlockGap}$ implies that $S$ is irrational.
\end{prop}

\begin{proof}
Fix $h\ge 1$ and $X_0$, and take $X$ as supplied. Choose $L$ with
$2^{L}\ge 1024\,(2X+h+L+2)$; such $L$ exists because $2^L/L\to\infty$, and it
satisfies the room inequality $16(2X+h+L+2)\le 2^L$ a fortiori. For every
$N<2X$ the landed truncation bound
\lean{Erdos249257.\allowbreak TotientTailPeriodKiller.\allowbreak norm_windowFirstExp_\allowbreak sub_tailOrbitFirstExp_lt}{PivotAntiReconstruction.lean:202}
gives
$\bigl\|\,\mathrm{windowFirstExp}(h,N,L)-e(2^N\alpha_h)\,\bigr\|
 < 2\pi(N+L+h+2)/2^{L} \le 2\pi/1024 < 1/100$,
so by
\lean{Erdos249257.\allowbreak TotientTailPeriodKiller.\allowbreak windowFirstCos_le_\allowbreak add_of_tailOrbitGap}{PivotAntiReconstruction.lean:220},
$\mathrm{windowFirstCos}(h,N,L)\le \cos(2\pi 2^N\alpha_h)+1/100$. Summing over
the $X$ values $N\in[X,2X)$ gives
$\sum_N \mathrm{windowFirstCos}(h,N,L)\le \tfrac{89}{100}X+\tfrac{1}{100}X
=\tfrac{9}{10}X$, so
\lean{Erdos249257.\allowbreak TotientTailPeriodKiller.\allowbreak exists_certifiedKill_of_\allowbreak first_harmonic_gap}{FirstHarmonicGap.lean:128}
produces $N\in[X,2X)$ with $\mathrm{certifiedKill}(h,N,L)$, and $N\ge X\ge X_0$.
This is exactly the certificate supply consumed by
\lean{Erdos249257.\allowbreak TotientTailPeriodKiller.\allowbreak irrational_totient_series_\allowbreak of_certificate_supply}{TotientTailPeriodKiller.lean:394}.
\end{proof}

\noindent
\ev{Math} \quad \scale{cofinal} \quad \coord{other:binary-window}. The proof
uses only landed lemmas; formalising it is a variant of
\lean{Erdos249257.\allowbreak irrational_totient_series_of_\allowbreak first_harmonic_norm_gap}{FirstHarmonicPivot.lean:83}
with the constant $21/25$ replaced by the real-part constant $89/100$ and the
room factor $16$ by $1024$. Nothing here is an improvement of the analytic
requirement; it is a change of coordinate that makes the requirement legible.

\begin{cor}[Digit form of the analytic requirement]\label{cor:digitform}
Let $\rho_h(X)$ denote the proportion of $N\in[X,2X)$ with
$\|2^N\alpha_h\|_{\R/\Z}\ge 1/4$. If for every $h\ge 1$ there are cofinally
many $X$ with $\rho_h(X)\ge 11/100$, then $S$ is irrational.
Equivalently, when $\alpha_h$ is not a dyadic rational: if for every $h$
there are cofinally many dyadic blocks $[X,2X)$ in which the binary
expansion of $\alpha_h$ contains at least $0.11X$ digit changes --- that is,
in which the mean binary run length is at most $9.1$ --- then $S$ is
irrational.
\end{cor}

\begin{proof}
$\|x\|\ge 1/4$ forces $\cos 2\pi x\le 0$, and the remaining terms are at most
$1$, so $\sum_{N}\cos(2\pi 2^N\alpha_h)\le (1-\rho_h(X))X$; require
$1-\rho\le 89/100$ and apply Proposition~\ref{prop:transfer}. For the digit
reading, $\|2^N\alpha_h\|\ge 1/4$ holds exactly when the binary digits of
$\alpha_h$ in positions $N+1,N+2$ differ.
\end{proof}

\noindent
\ev{Math}. This threshold is worth staring at. A single digit change in
every ninth position, in one block per scale, for each fixed $h$, would
settle Erd\H{o}s \#249. What is actually true, numerically, is far stronger.

\begin{table}[h]\centering
\begin{tabular}{r r r r r}
\hline
$h$ & $X$ & $\rho_h(X)$ & $X^{-1}\sum_{N\in[X,2X)}\cos(2\pi 2^N\alpha_h)$ & required \\
\hline
$1$ & $256$  & $0.4766$ & $\phantom{-}0.0037$ & $\le 0.89$ \\
$1$ & $2048$ & $0.5073$ & $-0.0043$ & $\le 0.89$ \\
$3$ & $256$  & $0.4648$ & $\phantom{-}0.0296$ & $\le 0.89$ \\
$3$ & $2048$ & $0.5020$ & $\phantom{-}0.0027$ & $\le 0.89$ \\
$8$ & $2048$ & $0.5112$ & $-0.0088$ & $\le 0.89$ \\
$13$& $2048$ & $0.5049$ & $-0.0195$ & $\le 0.89$ \\
\hline
\end{tabular}
\caption{Block statistics for $\alpha_h=(2^h-1)S$, computed from the exact
integer $\lfloor 2^{6000}S\rfloor$ (error $<1$ in the last bit, tail bound
$(M+2)2^{-M}$ at $M=6080$). Over positions $1..5800$ the binary expansion of
$\alpha_1$ has $2864$ runs, mean run length $2.0251$, longest run $13$; for
$\alpha_3$, $2964$ runs, mean $1.9568$, longest $12$; for $\alpha_8$, $2927$
runs, mean $1.9816$, longest $14$. A uniform digit model predicts mean run
length $2$ and longest run $\approx\log_2 5800\approx 12.5$. \ev{Cert}}
\label{tab:blocks}
\end{table}

\begin{obs}\label{obs:allroutes}
Under Lemma~\ref{lem:orbit} every route that survives the barrier
classification becomes a statement about the binary expansion of $S$ or of
some $\alpha_h$. Routes 1 and 2 ask for a positive proportion of digit
changes in a block; Route 3 asks for one short run at a prime-indexed
position; Route 4 asks that two blocks of digits of $S$ at distance $h$
disagree early. The routes differ in \emph{which positions} they interrogate
and \emph{how uniform} the margin must be, not in what they are about. This
is developed in \S\ref{sub:shape}.
\end{obs}

\subsection{Route 1: a constant-saving cancellation over a dyadic block}

\paragraph{The exact statement needed.}
\[
  \forall h\ge 1\ \forall X_0\ \exists X,L:\quad
  \max(X_0,1)\le X,\quad 16(2X+h+L+2)\le 2^{L},\quad
  \Bigl\|\sum_{N=X}^{2X-1} e\bigl(D(h,N,L)/2^{L}\bigr)\Bigr\|\le \tfrac{21}{25}X,
\]
where $D(h,N,L)=\sum_{j<L}\bigl(\varphi(N+h+1+j)-\varphi(N+1+j)\bigr)2^{L-1-j}$.
\lean{Erdos249257.\allowbreak TotientTailPeriodKiller.\allowbreak DTWFirstHarmonicNormGap}{FirstHarmonicPivot.lean:75},
consumer
\lean{Erdos249257.\allowbreak irrational_totient_series_of_\allowbreak first_harmonic_norm_gap}{FirstHarmonicPivot.lean:83}
\quad \ev{Open} \quad \scale{cofinal}. By Proposition~\ref{prop:transfer} the
weaker real-part form with constant $89/100$ suffices.

\paragraph{What $D$ is, as an arithmetic object.}
Dividing by $2^L$ and reindexing the two ranges onto a common variable
$t$, with $m=N+t$,
\begin{equation}\label{eq:window}
  \frac{D(h,N,L)}{2^{L}}
  =\underbrace{-\sum_{t=1}^{h}\frac{\varphi(N+t)}{2^{t}}}_{\text{head}}
  \;+\;\underbrace{(2^{h}-1)\sum_{t=h+1}^{L}\frac{\varphi(N+t)}{2^{t}}}_{\text{body}}
  \;+\;\underbrace{\sum_{t=L+1}^{L+h}\frac{\varphi(N+t)}{2^{t-h}}}_{\text{tail}} .
\end{equation}

So the phase is a \emph{geometrically weighted linear form in $\varphi$ at
$L+h$ consecutive arguments}, with the value at offset $t$ entering modulo
$2^{t}$ and with an odd multiplier $2^h-1$ throughout the body.

Two features control everything that follows. First, the weight at offset $t$ has
denominator exactly $2^{t}$, so offset $t$ contributes nothing unless
$v_2(\varphi(N+t))<t$: the $2$-adic valuation of the totient is the gate.

Second, $2^h-1$ is odd, so the multiplier never closes a gate that
$\varphi$ leaves open.

\paragraph{Which technique family is even the right shape.}
Four candidates, in decreasing order of relevance.

\emph{Weyl differencing and van der Corput} are the wrong shape and can be
dismissed exactly. Both require the phase to be a smooth or polynomial
function of the summation variable so that a difference operator lowers its
degree. Here $N\mapsto D(h,N,L)/2^{L}$ is, by Lemma~\ref{lem:orbit}, the
$\times 2$ orbit map: differencing in $N$ multiplies the phase by $2$ and
subtracts an integer, so the difference operator is an exact isometry of the
problem and lowers nothing. This is not a heuristic --- it is
$R_{N+1}=2R_N-\varphi(N+1)$ read modulo $1$.

\emph{The large sieve} needs a family of well-separated frequencies and a
sum over both the frequencies and the variable. Here there is one frequency
per $N$ and no family: the sum is over a single orbit. A large-sieve
inequality could be applied after introducing an artificial family (for
example, over the $h$ parameter), but the predicate quantifies $h$
universally, so an average over $h$ is not admissible.

\emph{Vaughan/Vinogradov bilinear decomposition} is the right shape, and is
what Route 2 already implements: decompose the argument $N+t$ at a large
prime factor, use $\varphi(mp)=\varphi(m)(p-1)$ to linearise the phase in
$p$, and sum over primes. \S\ref{sub:pivot} carries this as far as it goes.

\emph{Erd\H{o}s's 1948 digit method}, which proved irrationality of
$E=\sum_n 1/(2^n-1)=\sum_N d(N)/2^N$, is the right shape for the
\coord{mobius-mersenne} coordinate. That coordinate exists here and is
exact: from $\varphi=\mu * \mathrm{id}$,
\begin{equation}\label{eq:mobmers}
  S=\sum_{d\ge 1}\mu(d)\,\frac{2^{d}}{(2^{d}-1)^{2}}
  \qquad (\ev{Math};\ \text{verified to }16\ \text{digits},\ S=1.3676308019850223\ldots,\ \ev{Cert}).
\end{equation}

But the method does not transfer, for a reason that can be stated
quantitatively rather than vaguely. Erd\H{o}s's argument for $E$ exploits
that the coefficient $d(N)$ has enormous multiplicative fluctuation: over
$N\le Y$ the ratio of its maximum to its typical value is
$Y^{(\log 2+o(1))/\log\log Y}$, so certain highly composite $N$ deposit
identifiable spikes into the digit stream. The coefficient here is
$\varphi(N)$, which satisfies $\varphi(N)/N\in[c/\log\log N,1]$: its
multiplicative fluctuation over $N\le Y$ is a factor $e^{\gamma}\log\log Y$,
that is, $\log\log Y$ rather than $Y^{c/\log\log Y}$. There are no spikes to
find.

In the M\"obius--Mersenne coordinate \eqref{eq:mobmers} the situation is
worse rather than better: the coefficients $\mu(d)$ change sign, so the
digit blocks contributed by successive $d$ cancel rather than accumulate,
and there is no positivity to run a block argument on. This is, as far as
the evidence in this corpus goes, the specific reason \#249 is harder than
its Erd\H{o}s--Borwein ancestor.

\paragraph{The specific obstacle: what the estimate cannot be deduced from.}
The corpus records that no theorem proves the first-harmonic gap
inequivalent to irrationality, and argues the point from the structure of a
collapse proof. It can be settled outright, and the witness lives in the
same coefficient class that carries barriers B1 and B7.

\begin{thm}[The gap is strictly stronger than irrationality]\label{thm:lacunary}
Let $c(n)=1$ if $n=k!$ for some $k\ge 1$ and $c(n)=0$ otherwise, so
$0\le c(n)\le n$ for all $n\ge 1$, and let $\beta=\sum_{n\ge1}c(n)/2^{n}
=\sum_{k\ge 1}2^{-k!}$. Then $\beta$ is irrational, and for every $h\ge 1$
and every $X\ge 81(h+5)$,
\[
  \sum_{N=X}^{2X-1}\cos\bigl(2\pi\,2^{N}(2^{h}-1)\beta\bigr) \;>\; \tfrac{9}{10}X .
\]
Consequently the block-gap requirement fails at \emph{every} scale for
$\beta$, while $\beta$ is irrational. No proof of the block gap for $S$ can
therefore proceed from the irrationality of $S$ together with the growth
bound $c(n)\le n$; it must use arithmetic of $\varphi$.
\end{thm}

\begin{proof}
Irrationality of $\beta$ is Liouville's criterion. Write
$\gamma=(2^{h}-1)\beta=\sum_{k}(2^{h-k!}-2^{-k!})$, so for every $k$ with
$k!>2h$ the binary expansion of $\gamma$ carries ones exactly in the
positions $k!-h+1,\dots,k!$ and zeros between consecutive such blocks. If
$N\le k!-h-5$ for the least $k$ with $k!>N$, then
$\{2^{N}\gamma\}\le 2\cdot 2^{-5}=1/16$, whence $\|2^{N}\gamma\|\le 1/16$ and
$\cos(2\pi 2^{N}\gamma)\ge\cos(\pi/8)>0.9239$. The remaining $N$ lie in
$(k!-h-5,\,k!]$, at most $h+5$ values per factorial, and consecutive
factorials differ by a factor exceeding $2$ from $k\ge 3$, so the window
$[X,2X)$ meets at most one such interval. Therefore the sum exceeds
$0.9239\,(X-h-5)-(h+5)$, which exceeds $\tfrac{9}{10}X$ once
$X\ge 81(h+5)$.
\end{proof}

\noindent
\ev{Math} \quad \scale{cofinal} \quad \coord{other:binary-window}. This
upgrades the corpus's ``no collapse mechanism is known to apply'' to ``no
collapse mechanism can exist''. It also says precisely what the missing
input must do: it must rule out that $S$ behaves, in its doubling orbit,
like a Liouville number. That is the content of Corollary~\ref{cor:digitform},
and it is why the requirement is a digit-statistics statement and not an
irrationality statement.

\paragraph{Size of the quantity, honestly calibrated.}
Under any model in which the binary digits of $\alpha_h$ behave like fair
coin flips, $\rho_h(X)\to 1/2$ and the block sum is $O(\sqrt{X\log\log X})$
by the Erd\H{o}s--G\'al law of the iterated logarithm for lacunary series
\ev{Cited}. The requirement is $\rho_h\ge 0.11$ and a saving of a constant
factor. The measured values in Table~\ref{tab:blocks} are $\rho_h\approx 0.50$ and a block
average of order $X^{-1/2}$. So the analytic requirement is weaker than the
apparent truth by a factor $\sqrt X$, and the difficulty is entirely that no
technique produces \emph{any} nontrivial digit statistic for an explicit
constant of this kind. That is a statement about technology, and it is worth
separating from a statement about $S$.

\subsection{Route 2: the four-term pivot budget}\label{sub:pivot}

\paragraph{The exact statement needed.}
For every $h\ge 1$ there must exist $s\ge 1$ and $\eta\in(0,1)$ such that for
every $X_0$ there are $X\ge\max(X_0,1)$ and $L$ with $h\le L-s$,
$16(2X+h+L+2)\le 2^{L}$, and all four of
\[
\begin{aligned}
  \operatorname{Re}\bigl(\mathrm{pivotCenteredCorrelation}\bigr) &\le \tfrac{14}{25}X,
  &\qquad \bigl\|\mathrm{pivotFiberMeanContribution}\bigr\| &\le \tfrac{1}{100}X,\\
  \bigl\|\mathrm{pivotBadContribution}\bigr\| &\le \tfrac{1}{100}X,
  &\qquad \bigl\|\mathrm{pivotNonSupplierContribution}\bigr\| &\le \tfrac{8}{25}X .
\end{aligned}
\]
\lean{Erdos249257.\allowbreak TotientTailPeriodKiller.\allowbreak DTWPivotResidualDecorrelation}{FirstHarmonicPivot.lean:569},
consumer
\lean{Erdos249257.\allowbreak irrational_totient_\allowbreak series_of_pivotResidualDecorrelation}{FirstHarmonicPivot.lean:577}
\quad \ev{Open} \quad \scale{cofinal}. The decomposition is an exact identity,
\lean{Erdos249257.\allowbreak TotientTailPeriodKiller.\allowbreak windowFirstExp_sum_\allowbreak eq_pivot_decomposition}{FirstHarmonicPivot.lean:514}
\quad \ev{Lean}, and the supplier set at the canonical fibre is a membership
equality with a shifted dyadic interval of primes,
\lean{Erdos249257.\allowbreak TotientTailPeriodKiller.\allowbreak mem_pivotFiber_\allowbreak one_overlap_iff}{FirstHarmonicPivot.lean:423}
\quad \ev{Lean}.

\paragraph{One budget has a classical one-sided majorant.}
The non-supplier term is not identified exactly by smooth numbers.  What the
factorisation gives is the one-way containment needed for an upper bound.

\begin{prop}[Non-suppliers: the valid one-sided estimate]\label{prop:dickman}
Choose the admissible depth $L$ minimally for each large $X$, with the
predicate's previously fixed $h$ and $s$.  If
$n=N+L-s+1$ is not a supplier, then its largest prime factor satisfies
\[
  P(n)\le (4+o(1))\sqrt X .
\]
Therefore
\[
  \#\{N\in[X,2X):\neg\,\mathrm{pivotSupplier}(X,L,s,N)\}
  \le
  \Psi(2X+O(\log X),(4+o(1))\sqrt X)
  -\Psi(X+O(\log X),(4+o(1))\sqrt X).
\]
A uniform smooth-number asymptotic for this shifted interval would make the
right side $(1-\log2+o(1))X<\tfrac8{25}X$ and would discharge the budget.
\end{prop}

\begin{proof}
Let $p=P(n)$ and $m=n/p$.  Supplier failure means either
$p\le2\lfloor\sqrt X\rfloor$ or $m>\lfloor\sqrt X\rfloor/2$ (up to the
definition's exact floor and positivity clauses).  In the second case
$p=n/m\le(4+o(1))\sqrt X$, because minimal admissible $L$ gives
$L-s+1=O(\log X)$ and hence $n<2X+O(\log X)$.  Thus every non-supplier lies in
the displayed smooth-number set.  This proves only an upper bound, not an
asymptotic equality for the non-supplier count.
\end{proof}

\noindent
\ev{Math} for the containment; \ev{Cited} for the smooth-number estimate
needed to finish the numerical budget.  In particular the earlier identity
with Dickman density, and numerical extrapolations of an exact crossover
scale, are not claimed.  The implication is sufficient if the required
uniform shifted-interval estimate is supplied.

\begin{prop}[The bad-cofactor budget]\label{prop:badcof}
For $\eta\in(0,1)$ let $B(\eta)=\{m:\varphi(m)<\eta m\}$, with natural density
$D(\eta)$; by Schoenberg's theorem $D$ exists, is continuous, and
$D(0+)=0$ \ev{Cited}. Then
\[
  \#\{N\in[X,2X): N\ \text{a supplier with cofactor}\ m\in B(\eta)\}
  \;\le\;\bigl(D(\eta)+o(1)\bigr)X ,
\]
so a single choice of $\eta$ with $D(\eta)<1/200$ meets the $\tfrac{1}{100}X$
budget for all large $X$. Since $\eta$ is quantified existentially before
$X_0$, and a smaller $\eta$ only enlarges the good set, this choice costs
nothing elsewhere in the budget.
\end{prop}

\begin{proof}[Proof sketch]
Suppliers with cofactor $m$ number $\pi(2X/m)-\pi(X/m)$, which is at most
$(2X/m)/\log X$ up to $1+o(1)$ because $m\le\sqrt X/2$ forces
$\log(X/m)\ge\tfrac12\log X$. Sum over $m\in B(\eta)$ with $m\le\sqrt X/2$
and use that the logarithmic density of $B(\eta)$ tends to $D(\eta)$.
\end{proof}

\noindent
\ev{Math} \quad \scale{cofinal}. The choice is not tight: the least element
of $B(1/5)$ is $2\cdot3\cdot5\cdot7\cdot11\cdot13=30030$, with
$\varphi(m)/m=5760/30030=0.19181$, and the logarithmic density of $B(1/5)$ up
to $3\cdot10^{5}$ is $3\cdot 10^{-5}$ \ev{Cert}. So $\eta=1/5$ is already
defensible pending an explicit numerical bound on $D(1/5)$, and that bound is
itself a finite computation a reader could start today.

\paragraph{The pivot quantifiers force a deep modulus.}
The pivot sits at offset $t=L-s+1$ in \eqref{eq:window}.  Crucially, the
definition of \texttt{DTWPivotResidualDecorrelation} chooses $s$ once, after
$h$ and before the universal threshold $X_0$.  Thus $s$ is fixed as
$X\to\infty$; it cannot be increased with $X$ to make $t$ shallow.  The room
condition gives
\[
  2^L\ge16(2X+h+L+2),\qquad
  t=L-s+1\ge\log_2X-O_{h,s}(1).
\]
After removing the $2$-part of the coefficient, the prime-progression modulus
is still typically of order $X$ up to subpolynomial factors, not
$(\log X)^A$.  Standard Siegel--Walfisz therefore does not reach the
quantifier order of the stated predicate; choosing a hypothetical
$t=O(\log\log X)$ would amount to choosing $s$ after $X$, which the predicate
forbids.

The exact $2$-adic triviality observation remains useful: if
$v_2(\varphi(m))\ge t$, then the pivot phase is $1$.  It does not create a
nonempty analytic window under the actual quantifiers.  Hence the claimed
shallow-modulus route is withdrawn.  This does not disprove the pivot
predicate; it says only that the proposed prime-distribution argument does
not supply it. \ev{Math} \scale{cofinal}.

\paragraph{The specific obstacle: the weight is neither Type I nor Type II.}
Propositions~\ref{prop:dickman} and \ref{prop:badcof} isolate possible
bookkeeping bounds, subject to their stated uniform estimates.  The
fixed-$s$ quantifier leaves the fibre-mean modulus deep, and the centred term
has a separate correlation obstruction.

Factoring the pivot argument as
$n=mp$ and dividing out the pivot phase leaves
\[
  \operatorname{Re}\sum_{m}\ \sum_{p\,:\,mp\in[X,2X)+t}
    w(mp)\,\bigl(e\bigl(a_m(p-1)/2^{t}\bigr)-\overline{e}_m\bigr)
  \;\le\;\tfrac{14}{25}X ,
  \qquad |w|\equiv 1 ,
\]
where $\overline{e}_m$ is the fibre mean and, by \eqref{eq:window}, $w(n)$ is
the character of the same weighted totient sum with the single offset $t$
deleted.

Vinogradov's bilinear method requires that, after the decomposition
$n=mp$, the summand split as $\alpha_m\beta_p$ (Type II) or be independent of
$p$ given $m$ (Type I), up to boundedly many pieces. The weight $w(mp)$ is
neither: it is a function of the totients of the $L+h$ integers
\emph{adjacent} to $mp$, and $\varphi(mp+j)$ for $j\ne0$ bears no relation to
the factorisation $mp$ --- it is not a function of $m$, not a function of
$p$, and not a product of the two. No device is known that converts a weight
of this shape into bilinear form.

At $h=1$ and the canonical pivot the required estimate becomes fully
explicit. There $e\bigl((p-1)/4\bigr)=i^{\,p-1}=\chi_{-4}(p)$, so what is
needed is
\[
  \operatorname{Re}\sum_{X<p\le 2X}\chi_{-4}(p)\,
    e\Bigl(\frac{\varphi(p+1)}{8}+\frac{\varphi(p+2)}{16}+\cdots\Bigr)
  \;\le\;\Bigl(\frac{9}{10}-\delta\Bigr)\,\bigl(\pi(2X)-\pi(X)\bigr).
\]

The missing input is thus a \emph{non-correlation between a fixed quadratic
character at $p$ and the $2$-adic behaviour of $\varphi$ at the shifts
$p+1,p+2,\dots$}: a correlation statement for two arithmetic functions at
shifted arguments, of Chowla--Elliott type. The unconditional results in that
family --- Matom\"aki--Radziwi\l{}\l{} in almost all short intervals, Tao's
logarithmically averaged Chowla, Tao--Ter\"av\"ainen for odd order
\ev{Cited} --- are averaged, never pointwise at a single scale.

\paragraph{A swing: the predicate needs only cofinally many $X$, so average
over $X$.} This is the one place where the shape of the obligation is a gift.
Both $\mathrm{DTWFirstHarmonicNormGap}$ and
$\mathrm{DTWPivotResidualDecorrelation}$ ask for \emph{some} $X$ beyond each
$X_0$, never for all $X$. Hence it suffices to prove a logarithmically
averaged bound: if
\[
  \sum_{k\le K}\ \frac{1}{2^{k}}
  \Bigl|\sum_{N\in[2^{k},2^{k+1})}e\bigl(D(h,N,L_k)/2^{L_k}\bigr)\Bigr|
  \;=\;o(K)
  \quad\text{for each }h,
\]
with $L_k$ any admissible depth sequence, then infinitely many blocks satisfy
the gap and Proposition~\ref{prop:transfer} applies. Logarithmic averaging is
exactly the regime in which the entropy-decrement method operates.

The first
concrete step is not to prove this but to decide whether the method reaches
it: the phase $e\bigl(D(h,N,L)/2^{L}\bigr)$ is not multiplicative, but it is
a bounded \emph{local function of a multiplicative function} --- a fixed
continuous function of $\bigl(\varphi(N+1)/2,\dots,\varphi(N+L)/2^{L}\bigr)$
read modulo $1$ --- whereas entropy decrement is stated for correlations of
bounded multiplicative functions along fixed shifts. The precise question to
settle, and it can be started on immediately, is whether the decrement
survives when the observable is a local function of $\varphi$ with a growing
number of coordinates $L\approx\log_2X$ rather than a product of boundedly
many multiplicative values. \ev{Open}

\subsection{Route 3: uniform quantitative escape at primes}

\paragraph{The exact statement needed.}
\[
  \forall h\ge 1\ \forall N_0\ \exists p\ \text{prime}:\quad
  \max(N_0+h+1,\,h+5)\le p
  \ \wedge\
  \operatorname{Re}\bigl(\mathrm{tailOrbitFirstExp}(h,\,p-h-1)\bigr)<\tfrac{9}{10}.
\]
\lean{ErdosProblems.\allowbreak Erdos249.\allowbreak DTWNaturalPrimeTailOrbitStrictGap}{ErdosProblems/Erdos249/TotientStrictPrimeEscape.lean:157},
consumer
\lean{ErdosProblems.\allowbreak Erdos249.\allowbreak irrational_totient_\allowbreak series_of_naturalPrimeTailOrbitStrictGap}{ErdosProblems/Erdos249/TotientStrictPrimeEscape.lean:524}
\quad \ev{Open} \quad \scale{cofinal}.

\paragraph{What it says, after Lemma~\ref{lem:orbit}.}
$\operatorname{Re}e(2^{p-h-1}\alpha_h)<9/10$ says
$\|2^{p-h-1}\alpha_h\|_{\R/\Z}>\arccos(9/10)/2\pi=0.0717831\ldots$, so the
requirement is exactly: \emph{for each $h$, the indicated prime-indexed
$\times2$ orbit point stays more than $0.0717831\ldots$ from the nearest
integer, for cofinally many primes $p$}.  A short-run condition on the binary
digits can be a sufficient proxy for this circle-distance inequality, but it
is not equivalent to it; a run-length description alone does not determine
the residual position inside the dyadic cylinder.

\paragraph{The specific obstacle.} Two, stacked. First, this is a
\emph{pointwise} demand: it names an index, and the promotion audit records
that no pointwise producer in this programme has ever been supplied at even
one large index, across seven independent attempts, while the single audited
row whose missing input is a block average is the first-harmonic row.

Second, and sharper: digit behaviour \emph{along the primes} is strictly
harder than digit behaviour on average even for constants where the average
case is settled. The only constants whose digits are understood at all along a sparse
subsequence are ones built for the purpose (Champernowne;
Copeland--Erd\H{o}s) \ev{Cited}. No technique reads the digits of an
explicit constant at prime positions.

\paragraph{Numerical calibration.} Among the $501$ primes in $[1000,5000)$,
the proportion with $\operatorname{Re}e(2^{p-h-1}\alpha_h)<9/10$ is $0.8623$
for $h=1$, $0.8503$ for $h=2$, $0.8583$ for $h=5$, with minima below
$-0.9999$ \ev{Cert}. The cofinal supply is thus met by about six primes in
seven at these scales. As with Route 1, the requirement is very weak and the
obstacle is that nothing sees it.

\paragraph{A swing: replace the pointwise demand by a canonical fibre.}
The corpus warns against the subset consumer
\lean{Erdos249257.\allowbreak TotientTailPeriodKiller.\allowbreak exists_certifiedKill_of_\allowbreak first_harmonic_gap_subset}{FirstHarmonicGap.lean:104}
because a predicate quantifying existentially over subsets
$T\subseteq[X,2X)$ is collapse-exposed: the proved collapse
\lean{Erdos249257.\allowbreak dtwWindowSeparatedPairs_iff_\allowbreak irrational_totient_series}{PivotAntiReconstruction.lean:1765}
works precisely by \emph{choosing} $T=\{N,N+1\}$. That warning applies to a
free $T$.

It does not apply to a $T$ named in advance as a function of
$(h,X,L)$, and the corpus already contains the right one:
\[
  T_{h,X,L}\;:=\;\mathrm{pivotFiber}(X,L,L-h,1)
  \;=\;\{\,N\in[X,2X)\ :\ N+h+1\ \text{prime}\,\},
\]
a membership equality rather than a sampled surrogate
(\lean{Erdos249257.\allowbreak TotientTailPeriodKiller.\allowbreak mem_pivotFiber_\allowbreak one_overlap_iff}{FirstHarmonicPivot.lean:423},
\ev{Lean}).

The proposal is to prove
\begin{equation}\label{eq:primefibre}
  \forall h\ge1\ \forall X_0\ \exists X\ge X_0,\ L:\quad
  16(2X+h+L+2)\le 2^{L}
  \ \wedge\
  \sum_{N\in T_{h,X,L}}\mathrm{windowFirstCos}(h,N,L)\ \le\ \tfrac{9}{10}\,\bigl|T_{h,X,L}\bigr| ,
\end{equation}
which by the landed subset consumer yields a certificate at some
$N\ge X\ge X_0$, hence irrationality. \ev{Open}

Three features make this the sharpest available target.

(i) On this fibre the
cofactor is $m=1$, so the pivot coefficient is $(2^{h}-1)\varphi(1)=2^{h}-1$,
which is \emph{odd}, so the $2$-adic triviality described above cannot occur.

(ii) The pivot modulus is
$2^{h+1}$, \emph{fixed} once $h$ is fixed, so the fibre-mean term is governed
by the prime number theorem in progressions to a fixed power of two, where
the Ramanujan main term $c_{2^{h+1}}(2^{h}-1)=\mu(2^{h+1})=0$ vanishes for
every $h\ge1$ and the error term is effective; the fixed-$s$ deep-modulus
obstruction does not apply to this separately defined canonical fibre.

(iii) The bookkeeping budgets of \S\ref{sub:pivot} disappear:
inside $T$ there are no non-suppliers and no bad cofactors. What remains is
exactly one estimate --- decorrelation of the fixed phase
$e\bigl((2^{h}-1)(p-1)/2^{h+1}\bigr)$ from the residual totient weight --- and
it is the same estimate isolated at the end of \S\ref{sub:pivot}. This does
not make that estimate easier. It removes everything that is not that
estimate.

\begin{rem}
Honesty about \eqref{eq:primefibre}: it is not proved to be strictly stronger
than irrationality, and no collapse is proved either. What can be said is
that the one collapse mechanism that exists in this lane consumes the freedom
to choose the sample, and \eqref{eq:primefibre} has no such freedom. Deciding
this either way --- exhibiting a collapse, or a witness in the coefficient
class of Theorem~\ref{thm:lacunary} that separates it from irrationality ---
is itself a well-defined finite piece of work.
\end{rem}

\subsection{Route 4: depth-locked full-depth escape}

\paragraph{The exact statement needed.}
\[
  \mathrm{ApFullDepthEscape} :\equiv
  \forall d\ge 1\ \forall N\ \exists t\ge 1:\ \mathrm{certifiedKill}(td,\,N,\,td),
\]
unpacked, $(N+2td+2:\Z)<D(td,N,td)\bmod 2^{td}<2^{td}-(N+2td+2)$.
\lean{ErdosProblems.\allowbreak Erdos249.\allowbreak PeriodMultipleEscape.\allowbreak ApFullDepthEscape}{ErdosProblems/Erdos249/PeriodMultipleEscape.lean:441},
consumer
\lean{ErdosProblems.\allowbreak Erdos249.\allowbreak PeriodMultipleEscape.\allowbreak irrational_totient_\allowbreak series_of_apFullDepthEscape}{ErdosProblems/Erdos249/PeriodMultipleEscape.lean:450},
ambient equivalence
\lean{ErdosProblems.\allowbreak Erdos249.\allowbreak PeriodMultipleEscape.\allowbreak periodMultipleKillSupply_iff_irrational}{ErdosProblems/Erdos249/PeriodMultipleEscape.lean:432}
\quad \ev{Open} \quad \scale{cofinal}. It is the shortest fully stated open
inequality the programme has produced.

\paragraph{What it says: an anti-self-similarity statement about the digits
of $S$ alone.}
Writing $h=td$ and using
$D(h,N,L)/2^{L}=(R_{N+h}-R_N)-(R_{N+L+h}-R_{N+L})/2^{L}$ at $L=h$, together
with the landed strip $|R_{M+h}-R_M|<M+h+2$
(\lean{Erdos249257.\allowbreak TotientTailPeriodKiller.\allowbreak abs_tail_\allowbreak diff_lt}{CarrySurvivorExtinction.lean:79},
\ev{Lean}), applied at $M=N+h$, one gets:

\begin{prop}\label{prop:route4}
$\mathrm{certifiedKill}(h,N,h)$ holds whenever
$\bigl\|2^{N+h}S-2^{N}S\bigr\|_{\R/\Z}>2(N+2h+2)/2^{h}$.
\end{prop}

\begin{proof}
$\|D(h,N,h)/2^{h}\|\ge\|R_{N+h}-R_N\|-|R_{N+2h}-R_{N+h}|/2^{h}
 >2(N+2h+2)/2^{h}-(N+2h+2)/2^{h}=(N+2h+2)/2^{h}$, and
$R_{N+h}-R_N\equiv 2^{N}(2^{h}-1)S=2^{N+h}S-2^{N}S\pmod 1$ by
Lemma~\ref{lem:orbit}. A residue at distance more than $(N+2h+2)/2^{h}$ from
$\Z$ is exactly a certified kill at depth $h$.
\end{proof}

\noindent
\ev{Math}. So Route 4 asks: \emph{for every ray $d$ and every basepoint $N$,
some multiple $h=td$ is such that the tail of the binary expansion of $S$
beginning at position $N+1$ and the tail beginning at position $N+h+1$, read
as reals in $[0,1)$, are more than $2(N+2h+2)/2^{h}$ apart in $\R/\Z$.} Two
things follow. First, unlike Routes 1--3, this involves only $S$ itself, not
the multiples $\alpha_h$. Second, the required precision grows only like
$\log_2(N+2h)$ while the available depth grows linearly in $t$: the demand
weakens rapidly along each ray.

\paragraph{Which technique family fits, and why it is unavailable.}
A statement that a sequence does not almost repeat at some multiple of every
gap, to logarithmic precision, is a \emph{subword-complexity} statement, and
the mature technology for such statements about a specific real number is the
Adamczewski--Bugeaud method: the Schmidt subspace theorem applied to the
$b$-ary expansion, which shows that an algebraic irrational cannot have too
many long repetitions \ev{Cited}. The shape fits exactly. The hypothesis does
not: the method conditions on algebraicity, and $S$ is not known to be
algebraic --- indeed it is not known to be irrational, which is the whole
problem. There is no version of the method that runs from an analytic
description of the constant. That is the entire obstruction, and it is worth
stating plainly because it explains why Route 4 looks so much easier than it
is: the required combinatorial property is weak, but the only machine that
produces such properties needs an input we do not have.

\paragraph{What the data says.}
At $N=300$, the least $t$ with $\mathrm{certifiedKill}(td,300,td)$ exists
and is small for every ray $d\le 24$: $t=10$ for $d=1$, $t=5$ for $d=2$,
$t=4$ for $d=3$, $t=3$ for $d=4$, $t=2$ for $5\le d\le 9$, and $t=1$ for
$10\le d\le 24$ --- in every case at or near the first depth admitted by the
room floor. The residues are not marginal: $r/2^{h}=0.3594$ at $d=1$,
$0.5977$ at $d=3$, $0.6853$ at $d=7$, $0.2620$ at $d=17$, $0.5887$ at
$d=23$, against edge radii $3.1\cdot10^{-1}$, $8.0\cdot10^{-2}$,
$2.0\cdot10^{-2}$, $2.6\cdot10^{-3}$, $4.1\cdot10^{-5}$ \ev{Cert}. This is a
floor, not a trend, exactly as barrier B1 requires; it decides nothing.

\paragraph{A heuristic size computation.}
Under a uniform model for $\|2^{N}(2^{h}-1)S\|$ the failure probability at
depth $h=td$ is $\approx 4(N+2h+2)/2^{h}$, which is $\approx 1$ at the first
admissible $t$ and falls by a factor $2^{d}$ at each subsequent $t$. The
probability that a given $(d,N)$ fails at every $t$ is therefore a product
$\prod_{j\ge0}\min(1,\,c\,2^{-jd})$, super-exponentially small, and the
expected number of failing pairs $(d,N)$ over all $N$ converges. At $d=1$,
$N=300$ the product evaluates to about $10^{-5}$, and the observed value is a
success at the very first admissible depth. The model therefore predicts
Route 4 holds with enormous margin --- which is precisely why no finite
computation will ever be evidence for it. \ev{Open} (heuristic).

\subsection{Route 5: the rationality-side rank bound remains open}

\paragraph{The exact statement that was wanted.}
\[
  \exists C\ \forall c:\N\to\N\ \bigl(c(n)\le n,\ \neg\,\mathrm{Irrational}(X_c)\bigr)
  \ \forall v>0\ \forall u\ \mathrm{IsTemperedBinaryOrbit}(c,v,u)\ \forall e:\quad
  \operatorname{rk}_e(u)\le C ,
\]
where $X_c=\sum_{n\ge1}c(n)/2^{n}$ and $\operatorname{rk}_e(u)$ is the
dimension of the span of the dyadic sections of $u$ through level $e$.
A suitable subexponential upper bound would also contradict the landed
$\varphi$-specific floor
\lean{Erdos249257.\allowbreak not_irrational_totientSeries_\allowbreak implies_unbounded_carryRank_unconditional}{TotientCarryKernelRigidity.lean:300}
\quad \ev{Lean}.  No such rationality-side upper bound is proved.

\paragraph{The right vocabulary, and the audited gap.}
Finite-dimensional span of all dyadic sections is the definition of a
$2$-regular sequence in the sense of Allouche and Shallit \ev{Cited}.  The
unconditional theorem
\lean{Erdos249257.\allowbreak not_finiteDimensional_\allowbreak span_fullTotientKernel}{TotientMahlerDefect.lean:1145}
says that $\varphi$ is not $2$-regular.  An attempted generic counterexample
chose a coefficient sequence $c(n)\le n$ with rational binary series and
non-$2$-regular dyadic kernel.  Those facts do not refute the desired bound
for its carry orbit.

The generic recurrence
\[
  v\,c(N+1)=2u(N)-u(N+1)
\]
does express each \emph{positive-residue} level-$j$ section of $c$ as a
linear combination of two level-$j$ sections of $u$.  It does not control the
zero-residue section at each level.  Across levels $1,\ldots,e$ those omitted
sections can contribute up to $e$ new directions, not a fixed $O(1)$ error.
Consequently non-$2$-regularity of $c$ alone does not imply unbounded
$\operatorname{rk}_e(u)$, and the previously claimed bound
$\operatorname{rk}_e(u)\ge2^{e-1}-3$ has no valid proof here.  The proposed
greedy set construction therefore supplies neither a counterexample nor a
reason to retire the route. \ev{Gap} \scale{uniform}
\coord{other:carry-kernel}.

The Lean theorem for $\varphi$ remains valid: it uses the special structure
and independently proved linear independence of the totient kernel, not the
invalid generic $O(1)$ bridge.  The honest status is thus asymmetric.  The
large $\varphi$-specific rank floor is proved; a rationality-driven ceiling is
open; and the generic countermodel attempt is inconclusive.

\subsection{What the shape of the wall suggests about the object}\label{sub:shape}

Everything in this subsection is inference from the assembled evidence, not
theorem. It is written because the assembly makes one inference available
that no individual reformulation does, and because refusing to draw it would
be a different kind of dishonesty than drawing it carelessly.

\paragraph{First, a fact rather than an inference: the routes are one route
in five coordinates.}
Lemma~\ref{lem:orbit} collapses the coordinate spread that made the corpus
look like many independent attacks.

\begin{center}
\begin{tabular}{p{0.16\linewidth} p{0.44\linewidth} p{0.32\linewidth}}
\hline
Route & What it asks of the binary expansion & Positions interrogated \\
\hline
1, 2 & at least $11\%$ of positions carry a digit change; mean run length
$\le 9.1$ in one block per scale & a full dyadic block of $\alpha_h$ \\
3 & one run of length $\le 3$ & positions $p-h$, $p$ prime, in $\alpha_h$ \\
4 & the tails at gap $h$ differ by more than $2(N+2h+2)2^{-h}$ & every
basepoint, some multiple of every ray, in $S$ \\
5 & an open rationality-side upper bound on the carry-orbit kernel rank &
all dyadic sections through level $e$ \\
\hline
\end{tabular}
\end{center}

\noindent
This is the honest content of the operator's ``hundred measurements of one
wall''. The measurements are of one thing, and the thing is the binary digit
sequence of $S$ and of its multiples $(2^{h}-1)S$. The routes differ in which
positions they interrogate and how uniform a margin they demand; they do not
differ in subject.

\paragraph{Second: the wall is a wall about explicit constants, not about
$\varphi$.}
Of the seven barriers, six --- B1, B2, B3, B5, B6, B7 --- are statements that
a class of proof cannot see the digits: bounded inspection, retargeting,
strengthening, bounded state, finite linear compression, coarse invariants.
None asserts that the digits misbehave. Only B4 is internal to the object,
and even B4 says that one engineering scheme self-cancels, not that the
quantity being engineered is unattainable. Meanwhile the two object-level
witnesses the corpus can produce --- the $\gamma$-splice of B1, the parity
coboundary of B7 --- and the one produced here
(Theorem~\ref{thm:lacunary}) are all lacunary or eventually periodic:
constructed, measure-zero, and highly structured. In two centuries nobody has
exhibited a naturally occurring constant with Liouville-type digit behaviour;
every known example is built to have it. So the shape of the barrier set
points at our technology, not at $S$.

\paragraph{Third: is the object behind the measurements simple or complex?}
The evidence available leans one way, and I will say how far it leans.

\emph{For ``generic, and the difficulty is entirely ours''.} (a) Every
statistic measurable is indistinguishable from a fair-coin digit model:
$\rho_h(X)\approx0.50$ against a required $0.11$; mean run length
$1.96$--$2.03$ against a permitted $9.1$; longest run $12$--$14$ over $5800$
positions against a model prediction of $\approx12.5$; block sums of order
$X^{-1/2}$ against a permitted $0.89$ (Table~\ref{tab:blocks}, \ev{Cert}). (b) The
requirements are weaker than the apparent truth by a factor $\sqrt X$; a
proof does not need to understand the digits, only to exclude a pathology.
(c) All countermodels are constructed lacunary objects. (d) The one
coordinate in which $S$ has classical structure --- the M\"obius--Mersenne
form \eqref{eq:mobmers} --- is an alternating sum of rational functions of
$2^{d}$ with no functional equation, no modularity, no continued-fraction
structure, and no algebraicity: there is no hidden object to find, which is
consistent with genericity.

\emph{Against, or at least complicating.} (a) B4's self-cancellation is
exact, unconditional, and holds at every depth $K$: the cost of forcing a
depth-$K$ residue by a Dirichlet prime is $p\ge1+2^{K-1}$ while the payoff is
amplitude $2^{K-1}$, so the trade is precisely null. An exact null of that
form is a structural fact and not an accident of parametrisation, and it is
the one place where the object itself pushes back. (b) The exact census
records residues approaching the forbidden edge --- a closest central margin
of about $2.2\cdot10^{-4}$ of the modulus at $t=100$ --- which is what a
uniform model predicts and which rules out, permanently, any argument that
proceeds by a crude uniform margin. (c) Proposition~\ref{prop:dickman} shows
that even the pieces of the frontier that are provable become true only past
$X\approx10^{30}$--$10^{40}$; a genuinely simple object would not usually
require asymptotics that begin so late.

\emph{Verdict, flagged as inference.} The evidence supports ``the digits of
$S$ are generic and the wall is technological'' over ``the object is
subtle'', but it does not determine it, and the reason it cannot is
structural: every test performed is a test a generic object passes, so
passing them is weak evidence, and the only tests that would discriminate are
exactly the ones no technique can run. What the evidence \emph{does} determine
is narrower and firmer: the difficulty of \#249 is not located in $\varphi$'s
irregularity --- $\varphi$ is the smoothest interesting multiplicative
function, with multiplicative fluctuation $\log\log Y$ where $d(n)$ has
$Y^{c/\log\log Y}$ --- but in the fact that we possess no method whatsoever
for lower-bounding digit changes of a constant that was not designed to have
them. That is why the Erd\H{o}s--Borwein constant fell in 1948 and this one
has not: $d(n)$ has spikes to deposit into the digit stream and $\varphi(n)$
does not, and in the M\"obius coordinate the signs cancel.

\paragraph{What would actually change the picture.}
One theorem, of any strength, giving a nontrivial digit statistic for $S$ or
for some $\alpha_h$: an upper bound $o(N)$ on the longest binary run in the
first $N$ positions would not by itself decide \#249, but it would be the
first evidence that the object is legible at all, and by
Corollary~\ref{cor:digitform} any bound strong enough to give a positive
proportion of digit changes in one block per scale would decide it. Failing
that, the two concrete openings identified above are: the logarithmically
averaged form of the block gap, where the entropy-decrement method is at
least in the right regime (\S\ref{sub:pivot}); and the canonical prime fibre
\eqref{eq:primefibre}, where every budget except one decorrelation estimate
has been discharged. Neither is close. Both are well posed, and both can be
started on today.
% ---- part a249_p5 ----
\section{What is open, stated exactly}

Erdős \#249 asks whether $S = \sum_{n\ge 0} \varphi(n)/2^n$ is irrational. It is
\ev{Open}: nothing in this corpus decides it.

Every attack recorded in Parts
I--IV terminates, in one of exactly two ways, at a small number of
\emph{cofinal supply obligations} on the totient tail
$R_N := \sum_{j\ge 0}\varphi(N+1+j)/2^{j+1}$
(\lean{Erdos249257.TotientTailPeriodKiller.totientTail}{TotientTailPeriodKiller.lean:57}):
either it proves a genuine theorem of the shape ``\emph{if} this cofinal
predicate holds, \emph{then} $S$ is irrational,'' with the predicate itself
left completely unproved, or it prunes one specific proof shape and shows it
cannot supply the predicate.

This section renders every such obligation
exactly as Lean states it, gives its site, records which of them are proved
\emph{equivalent} to one another (iff) and which are proved only
\emph{sufficient} (one-directional, an easier target implied by a harder
one), and closes by naming, without overclaiming, the piece of new
mathematics each surviving form would need. Quantifier order is preserved
exactly as written in Lean; nowhere below does a finite verified list stand
in for a cofinal supply, and nowhere does an implication get reported as a
solved case.

Throughout, $H_t := \mathrm{periodLcm}(t) = \operatorname{lcm}(1,\dots,t)$
(\lean{Erdos249257.CarrySurvivorExtinction.periodLcm}{CarrySurvivorExtinction.lean:515}),
$D(h,N,L) := \sum_{j<L}\bigl(\varphi(N{+}h{+}1{+}j)-\varphi(N{+}1{+}j)\bigr)\cdot 2^{L-1-j}\in\mathbb{Z}$
is \texttt{windowDiscrepancy}
(\lean{Erdos249257.TotientTailPeriodKiller.windowDiscrepancy}{TotientTailPeriodKiller.lean:66}),
and the base predicate named ``$\mathrm{Sep}(h,N,L)$'' in the task brief is
literally \texttt{certifiedKill}:
\[
  \mathrm{certifiedKill}(h,N,L) :\equiv
  (N+h+L+2 : \mathbb{Z}) < D(h,N,L) \bmod 2^L
  \ \wedge\
  D(h,N,L) \bmod 2^L < 2^L - (N+h+L+2).
\]
\lean{Erdos249257.TotientTailPeriodKiller.certifiedKill}{TotientTailPeriodKiller.lean:72}
\quad \ev{Lean} (the definition; decidable, instance at line 84)
\quad \scale{n/a} \quad \coord{other:binary-window}. In words: the residue
of the window discrepancy modulo $2^L$ avoids a shrinking radius-$(N+h+L+2)$
neighbourhood of $0$ inside the growing modulus $2^L$.

\subsection{The certificate-completeness converter}

Every reformulation below ultimately measures the same underlying real
quantity, because certificates are \emph{complete} receipts of
non-integrality, not merely sufficient ones:
\[
  \bigl(\exists L,\ \mathrm{certifiedKill}(h,N,L)\bigr)
  \iff
  R_{N+h} - R_N \notin \operatorname{range}\bigl((\uparrow)\colon \mathbb{Z}\to\mathbb{R}\bigr).
\]
\lean{Erdos249257.exists_certifiedKill_iff_tail_diff_notMem_int}{LcmConeFlatness.lean:316}
(wrapped as \texttt{totient\_tail\_window\_kill\_exists\_iff\_tail\_diff\allowbreak\_not\_int} at
\lean{Erdos249257.totient\_tail\_window\_kill\_exists\_iff\_tail\_diff\_not\_int}{CertificateKernel.lean:19025})
\quad \ev{Lean} \quad \scale{uniform} (holds for all $h,N$)
\quad \coord{other:binary-window}.

\emph{Note.} This iff is the reason every ``certificate supply'' form below
can be restated, without loss, in ``pure non-integrality'' language with no
certificate vocabulary at all (the B8/TE forms). It does \emph{not} assert
that the cofinal supply exists; it only says the certificate route and the
real-analytic route are the same route.

\subsection{The base cofinal obligation — the wall}

\begin{quote}
\textbf{Obligation 1 (certificate-supply normal form).}
\[
  \mathrm{Irrational}(S)
  \quad\Longleftrightarrow\quad
  \forall h\ge 1,\ \forall N_0,\ \exists N\ge N_0,\ \exists L,\
  \mathrm{certifiedKill}(h,N,L).
\]
\end{quote}
This is exactly the quantifier structure named ``$\mathrm{Sep}(h,N,L)$'' in
the task brief: $\forall h\ge 1\ \forall N_0\ge 0\ \exists N\ge N_0\ \exists L\ \mathrm{Sep}(h,N,L)$.
\lean{Erdos249257.irrational_totient_series_iff_certificate_supply}{LcmConeFlatness.lean:412}
\quad \ev{Lean} (the equivalence is proved; the supply itself is unsupplied)
\quad \scale{cofinal} \quad \coord{other:binary-window}.

The reverse direction is by contradiction against the tail-period law: if $S$ were
rational, \texttt{eventual\_period\_of\_not\_irrational}
(\lean{Erdos249257.eventual_period_of_not_irrational}{TotientTailPeriodKiller.lean:358},
\ev{Lean}, unconditional) supplies a period $h>0$ and pre-period $N_0$ with
$R_{N+h}-R_N\in\mathbb{Z}$ for all $N\ge N_0$; the certificate hypothesis
instantiated at that $(h,N_0)$ produces an $N\ge N_0$ where
$R_{N+h}-R_N\notin\mathbb{Z}$, a contradiction via \texttt{tail\_diff\_notMem\_int\_of\_certifiedKill}
(\lean{Erdos249257.tail_diff_notMem_int_of_certifiedKill}{TotientTailPeriodKiller.lean:262}).

For the forward direction, irrationality gives pointwise non-integrality of
every positive tail shift and certificate completeness supplies a witness
already at $N=N_0$.

\emph{What supplying Obligation 1 would take:} an arithmetic (not merely
existential) argument that, for every period length $h$ and past every
threshold $N_0$, the totient window discrepancy's residue mod $2^L$ can be
pushed off $0$ by a margin growing with the modulus — i.e., genuine
cancellation/anti-concentration information about $\varphi$ on a sliding
window, at every scale, for every $h$. No such argument exists in the
corpus; the rows below are exact normal forms or sufficient producers for
this missing input, never a proof that the input holds.

\subsection{Collapsing free parameters: the multiple / diagonal / cone chain}

Three further reformulations replace Obligation 1's two free parameters
$(h,N_0)$ with successively less information.  The multiple and cone
predicates have explicit reductions to irrationality; the lcm-diagonal
predicate has a registered iff.  Since Obligation 1 itself is an iff, every
predicate both implied by Obligation 1 and sufficient for irrationality is
also propositionally equivalent to it, even when the corpus exposes only
the two directed constructions rather than a separately named iff.

\paragraph{Multiple-period collapse (formally looser, but endpoint-equivalent).}
\[
  \forall h_0>0,\ \forall N_0,\ \exists m>0,\ \exists N\ge N_0,\ \exists L,\
    \mathrm{certifiedKill}(m\cdot h_0, N, L)
  \ \Longrightarrow\ \mathrm{Irrational}(S).
\]
\lean{Erdos249257.irrational_totient_series_of_multiple_certificate_supply}{CarrySurvivorExtinction.lean:502}
\quad \ev{Lean} \quad \scale{cofinal} \quad \coord{other:lcm-period-multiple}.
Obligation 1's hypothesis trivially implies this one (take $m=1$), while the
displayed theorem sends it back to irrationality and hence, by the base iff,
back to Obligation 1.  It is an easier-looking target, not progress; it
remains exactly as open as \#249.

\paragraph{Diagonal collapse — one free parameter.}
\[
  \mathrm{Irrational}(S)
  \ \Longleftrightarrow\
  \forall t_0,\ \exists t\ge t_0,\ \exists L,\
    \mathrm{certifiedKill}(H_t, H_t, L).
\]
\lean{Erdos249257.irrational_totient_series_iff_lcm_diagonal_certificate_supply}{LcmConeFlatness.lean:426}
\quad \ev{Lean} \quad \scale{cofinal} \quad \coord{other:lcm-diagonal}.
This is the canonical single-quantifier restatement: write
$P(t) :\equiv \exists L,\ \mathrm{certifiedKill}(H_t,H_t,L)$; the obligation
is $\forall t_0\,\exists t\ge t_0,\ P(t)$. The reduction is a genuine proof,
not a relabelling: given $t\ge\max(h_0,N_0)$, $h_0\mid H_t$ and $H_t\ge
t\ge N_0$ simultaneously, so a single diagonal witness at $t$ discharges
\emph{both} of Obligation 1's free parameters at once, via the intermediate
lemma \lean{Erdos249257.irrational_totient_series_of_lcm_certificate_supply}{LcmDiagonalReduction.lean:92}
which itself reduces to the multiple-period form above.  Conversely,
irrationality and pointwise certificate completeness give the diagonal
witness already at $t=t_0$.

$P(t)$ is verified for every $t\le82$
(\lean{ErdosProblems.Skip.LadderT67.exists_diagonalKill_le_82}{ErdosProblems/Skip/LadderT67.lean:71264}).
The historical bank contained 28 explicit values through $t=64$
(\lean{Erdos249257.certifiedKill_diagonal_all_imported}{DiagonalPincerCertificates.lean:2870},
depths $[6,5,7,7,9,14,15,14,21,22,23,26,\dots]$ for
$t\in\{1,2,3,4,5,7,8,9,11,13,16,17,\dots\}$; endpoint
\lean{Erdos249257.certifiedKill_diagonal_t64}{DiagonalPincerCertificateT64Endpoint.lean:1928})
\ev{Cert} \scale{fixed} \coord{other:lcm-diagonal}).  The current contiguous
band is still a finite floor, not a cofinal supply; no certificate at $t=83$
is claimed.

\paragraph{Cone collapse — two multipliers, still the same wall.}
\[
  \forall t_0,\ \exists t\ge t_0,\ \exists q,m,L,\ 0<q\ \wedge\
    \mathrm{certifiedKill}(m\cdot H_t,\, q\cdot H_t,\, L)
  \ \Longrightarrow\ \mathrm{Irrational}(S).
\]
\lean{Erdos249257.irrational_totient_series_of_lcm_cone_window_kill_supply}{CertificateKernel.lean:19014}
\quad \ev{Lean} \quad \scale{cofinal} \quad \coord{other:lcm-cone}. The
diagonal form above is exactly the cell $q=m=1$, so any diagonal witness
trivially witnesses the cone form.  The cone predicate is therefore a
formally looser target, but its sufficiency theorem and the base iff make it
propositionally equivalent to \#249 rather than an unconditional advance. It
rests on a genuine
strengthening of the tail-period law, \emph{lcm-cone flatness}: if $S$ is
rational there is $t_1$ such that for every $t\ge t_1$ and every $q>0,m$,
$R_{qH_t+mH_t}-R_{qH_t}\in\mathbb{Z}$ — rationality flattens the
\emph{whole} cone $\{k H_t : k\ge 1\}$, not just one difference
(\lean{Erdos249257.rational_totient_series_forces_lcm_cone_flatness}{CertificateKernel.lean:19002},
\ev{Lean}, unconditional).

A sharper cone-form producer,
\texttt{coneNonflatCert}, needs only a \emph{one-sided} radius per vertex —
information-theoretically half of \texttt{certifiedKill}'s pairwise floor —
and proves that some pair in a finite multiplier menu $Q$ is non-integral;
the corresponding supply,
$\exists$ an unbounded-scale menu $Q$ with \texttt{coneNonflatCert} firing,
is \ev{Open}, \scale{cofinal},
\lean{Erdos249257.irrational_totient_series_of_lcm_cone_nonflat_supply}{CertificateKernel.lean:19140}.

\paragraph{Pure non-integrality form — certificate vocabulary stripped out.}
Via the completeness iff above, the diagonal and cone forms restate,
\emph{exactly}, with no reference to $\mathrm{certifiedKill}$ at all:
\[
  \forall t_0,\ \exists t\ge t_0,\ R_{2H_t}-R_{H_t}\notin\mathbb{Z}
  \ \Longleftrightarrow\ \text{diagonal obligation above}
  \ \Longrightarrow\ \mathrm{Irrational}(S),
\]
\lean{Erdos249257.irrational_totient_series_of_lcm_diagonal_nonintegrality_supply}{CertificateKernel.lean:19034}
(cone analogue at
\lean{Erdos249257.irrational\_totient\_series\_of\_lcm\_cone\_nonintegrality\_supply}{CertificateKernel.lean:19046})
\quad \ev{Lean} \quad \scale{cofinal} \quad \coord{other:real-analytic-nonintegrality}.
This is the frontier of \#249 with every piece of certificate machinery
removed: \emph{does $R_{2H_t}-R_{H_t}\notin\mathbb{Z}$ for infinitely many
$t$?} Nothing decides this either.

\subsection{The exponential-sum form}

A first-harmonic (Weyl-sum) cancellation statement is proved sufficient for
Obligation 1 by an elementary pigeonhole argument, with no case analysis and
no scale-degrading constants:
\[
  \mathrm{DTWFirstHarmonicNormGap} :\equiv\
  \forall h>0,\ \forall X_0,\ \exists X,L,\quad
  \max(X_0,1)\le X\ \wedge\
  16\,(2X+h+L+2)\le 2^L\ \wedge\
  \Bigl\|\sum_{N\in[X,2X)} e\bigl(D(h,N,L)/2^L\bigr)\Bigr\| \le \tfrac{21}{25}X.
\]
\lean{Erdos249257.DTWFirstHarmonicNormGap}{FirstHarmonicPivot.lean:75}
\quad \ev{Open} (the predicate; unproved at every $h,X_0$)
\quad \scale{cofinal} \quad \coord{other:first-harmonic}.

\[
  \mathrm{DTWFirstHarmonicNormGap} \ \Longrightarrow\ \mathrm{Irrational}(S).
\]
\lean{Erdos249257.irrational_totient_series_of_first_harmonic_norm_gap}{FirstHarmonicPivot.lean:83}
\quad \ev{Lean} \quad \scale{cofinal} \quad \coord{other:first-harmonic}. The
consumer already has Obligation 1's exact free-parameter shape ($X$ is a
completely free threshold, so applying the gap at $X\ge N_0$ gives
$\exists N\ge N_0\ \exists L$ directly); what is missing is not scale and not
coordinate, it is the arithmetic input itself — not one instance of a
constant-saving cancellation bound for the complex exponential sum
$\sum_{N\in[X,2X)} e(D(h,N,L)/2^L)$ is proved anywhere in the corpus, at any
$h,X,L$.

A strictly stronger \emph{subset} form (a saving on any nonempty
finite $T\subseteq[0,2X)$, no density or partition hypothesis) is also
proved sufficient and unconditional as an implication:
\lean{Erdos249257.exists_certifiedKill_of_first_harmonic_norm_gap}{FirstHarmonicPivot.lean:53}
consumes the elementary block lemma
\lean{Erdos249257.exists_certifiedKill_of_first_harmonic_gap}{FirstHarmonicGap.lean:128}
(\ev{Lean}, unconditional: \emph{any} constant-saving first-harmonic gap on
one dyadic block forces a finite kill certificate, via $\cos(\pi/8)>9/10$
and averaging).

\emph{What supplying this form would take:} genuine
cancellation for the first additive character of the totient window
discrepancy — a Weyl-type bound, uniform in $h$ and growing $L$, on a
character sum built from $\varphi$ differences, not merely an existence
statement. \lean{Erdos249257.ResidualGaugeObstruction.det_residualMonomialMatrix}{ResidualGaugeObstruction.lean:47}
rules out one entire class of candidate proofs: a residual-blind
determinant/rank certificate cannot see this cancellation, because an
explicit ``locked'' unit-norm gauge reconstructs the exact bad configuration
while keeping any Vandermonde-shaped minor nonzero — any real argument here
must couple the rows arithmetically, not merely normalise columns.

\subsection{The actual-orbit reformulation and the top-edge staircase}

A second, independently developed lane (the public pinned modules
\texttt{TotientActual*}/\texttt{TotientFixedRank*}) restates \#249 in terms of the \emph{actual}
power-of-two LCM diagonal, $H = H_{2^a}$, and proves a genuine \textbf{iff}
with the totient series itself — the strongest form of ``exact statement of
what remains'' anywhere in the corpus.

\begin{quote}
\textbf{The exact equivalence.}
\[
  \mathrm{Irrational}(S)
  \iff
  \forall a_0,\ \exists a\ge a_0,\ \mathrm{actualLcmTailOrbit}(a)\notin\mathbb{Z}.
\]
\end{quote}
\lean{Erdos257PeriodNoncollapse.DiagonalFreshLossBridge.PowerTwoOddWindowAffine.irrational_totientSeries_iff_actualLcmOrbitNonintegralitySupply}{TotientActualLcmOrbitNonintegrality.lean:37}
\quad \ev{Lean} \quad \scale{n/a} (an unconditional equivalence, not a
supply) \quad \coord{mobius-mersenne}, where
$\mathrm{actualLcmTailOrbit}(a) := 2^H(2^H-1)S - (\mathrm{totientPrefix}(2H)-\mathrm{totientPrefix}(H))$,
$H=H_{2^a}$
(\lean{Erdos257PeriodNoncollapse.DiagonalFreshLossBridge.PowerTwoOddWindowAffine.actualLcmTailOrbit_eq_scaled_totientSeries_sub_prefix}{TotientActualLcmOrbitSeparation.lean:68}).
\emph{This is \#249 restated with nothing left over}: no auxiliary
hypothesis, no scale caveat — cofinal non-integrality of this one sparse
sequence at power-of-two heights is exactly Erdős \#249. Everything else in
this subsection is an attempt to reach the right-hand side.

\paragraph{Route-pruning: total staircase collapse is impossible.}
Before any positive target, the file proves one entire proof shape is
empty. For $a\ge 8$ and room bound $J+K+(a{+}6)<2\cdot 2^a$ with a wide
enough modulus $2H+J+K+2<2^m$, a terminal window where \emph{every} one of
the last $m$ letters vanishes mod its own growing power of two is
impossible:
\[
  \neg\,\mathrm{ActualLcmTerminalDyadicStaircase}(a,J,K,m).
\]
\lean{Erdos257PeriodNoncollapse.DiagonalFreshLossBridge.PowerTwoOddWindowAffine.not_actualLcmTerminalDyadicStaircase_of_room}{TotientActualLcmTopEdgeStaircase.lean:251}
\quad \ev{Lean} \quad \scale{bounded} \quad \coord{mobius-mersenne}. The
mechanism is a bare positivity/divisibility contradiction (a positive
quantity below a modulus, divisible by that modulus, must be $0$) against
the unconditional positivity of every short-window letter,
$0<\mathrm{lcmRayArithmeticLetter}(2^a,j)$ for $a\ge8$, $0<j<2\cdot2^a$
(\lean{Erdos257PeriodNoncollapse.DiagonalFreshLossBridge.PowerTwoOddWindowAffine.lcmRayArithmeticLetter_pos_of_lt_two_mul}{TotientActualLcmOrbitArithmetic.lean:1703},
\ev{Lean}, unconditional, no rationality hypothesis).

\emph{Do not attempt a
full terminal-staircase producer}: it is provably empty. The corpus's own
route past this dead end is the \emph{punctured} staircase (all but the
last letter vanish; the last is retained and pinned exactly to the half-turn
$2^{m-1}$,
\lean{Erdos257PeriodNoncollapse.DiagonalFreshLossBridge.PowerTwoOddWindowAffine.puncturedDyadicStaircase_penultimate_eq_half}{TotientActualLcmTopEdgeStaircase.lean:1187},
\ev{Lean}, \scale{bounded}) and the one-sided residue-gap family below.

\paragraph{The surviving one-sided target.}
\[
  \mathrm{ActualLcmTopEdgeResidueGap}(a,J,K,m) :\equiv
  m\le K\ \wedge\
  2H+J+K+2 < 2^m\ \wedge\
  D(H,\,H+J,\,K) \bmod 2^m \le 2^m - (2H+J+K+2).
\]
\lean{Erdos257PeriodNoncollapse.DiagonalFreshLossBridge.PowerTwoOddWindowAffine.ActualLcmTopEdgeResidueGap}{TotientActualLcmTopEdgeStaircase.lean:900}
\quad \ev{Open} (a definition, unproved at any $a$)
\quad \scale{bounded} \quad \coord{mobius-mersenne}.

Note this is a
\emph{one-sided} inequality — only the upper (positive) carry arc is
excluded, not a symmetric two-sided band — because the corridor's lower half
is already discharged unconditionally: for $a\ge8$ and $J+(a{+}6)<2\cdot2^a$,
\[
  0 < R_{2H+J}-R_{H+J}
\]
\lean{Erdos257PeriodNoncollapse.DiagonalFreshLossBridge.PowerTwoOddWindowAffine.actualLcmTailDiff_shift_pos}{TotientActualLcmOrbitSign.lean:39}
\quad \ev{Lean} \quad \scale{bounded} \quad \coord{mobius-mersenne} — a rare
genuinely unconditional, non-hypothetical real-analytic theorem in this
corpus, needing no rationality assumption at all.

Its companion shows that
\emph{if} the orbit is integral, the residue is forced to the exact top
edge $2^K-e$ ($e>0$ small), \emph{outside} the arc a symmetric certificate
would need:
\lean{Erdos257PeriodNoncollapse.DiagonalFreshLossBridge.PowerTwoOddWindowAffine.actualLcm_integral_forces_topEdgeResidue}{TotientActualLcmOrbitSign.lean:211}
(\ev{Lean}, \scale{bounded}) — this is the exact statement of ``here is
what remains,'' pinning the missing exclusion to one named residue class
rather than leaving it implicit.

\[
  \mathrm{ActualLcmTopEdgeResidueGap}(a,J,K,m)\ \Longrightarrow\
  R_{2H+J}-R_{H+J}\notin\mathbb{Z}.
\]
\lean{Erdos257PeriodNoncollapse.DiagonalFreshLossBridge.PowerTwoOddWindowAffine.actualLcmTailDiff_notMem_int_of_topEdgeResidueGap}{TotientActualLcmTopEdgeStaircase.lean:1260}
\quad \ev{Lean} \quad \scale{bounded} \quad \coord{mobius-mersenne}.

The
cofinal target built from it:
\[
  \mathrm{PowerTwoActualLcmTopEdgeResidueGapSupply} :\equiv\
  \forall a_0,\ \exists a,K,m,\quad
  a_0\le a\ \wedge\ 8\le a\ \wedge\ K+(a{+}6)<2\cdot 2^a\ \wedge\
  \mathrm{ActualLcmTopEdgeResidueGap}(a,0,K,m).
\]
\lean{Erdos257PeriodNoncollapse.DiagonalFreshLossBridge.PowerTwoOddWindowAffine.PowerTwoActualLcmTopEdgeResidueGapSupply}{TotientActualLcmTopEdgeStaircase.lean:1325}
\quad \ev{Open} \quad \scale{cofinal} \quad \coord{mobius-mersenne}.
\[
  \mathrm{PowerTwoActualLcmTopEdgeResidueGapSupply} \Longrightarrow \mathrm{Irrational}(S).
\]
\lean{Erdos257PeriodNoncollapse.DiagonalFreshLossBridge.PowerTwoOddWindowAffine.irrational_totientSeries_of_topEdgeResidueGapSupply}{TotientActualLcmTopEdgeStaircase.lean:2184}
\quad \ev{Lean} \quad \scale{cofinal} \quad \coord{mobius-mersenne}.

\paragraph{Five strictly weaker links, each independently proved sufficient.}
The same file proves five further cofinal predicates, each strictly weaker
than \texttt{PowerTwoActualLcmTopEdgeResidueGapSupply} (each implies it, so
each is an easier target), with a direct \#249 endpoint of its own —
proving the \emph{weakest} of the six closes the entire cluster.

\begin{enumerate}
\item[(i)] $\mathrm{PowerTwoAdjacentSuffixMidbandSupply}$: replaces the
$m$-bit residue test with a symmetric two-sided band on the
\emph{adjacent-suffix} residue at depth $m$, buffered by the larger depth
$m{+}1$ so either branch stays inside the sign corridor. \lean{Erdos257PeriodNoncollapse.DiagonalFreshLossBridge.PowerTwoOddWindowAffine.PowerTwoAdjacentSuffixMidbandSupply}{TotientActualLcmTopEdgeStaircase.lean:1334}
\ev{Open} \scale{cofinal} \coord{mobius-mersenne}. Sufficiency:
\lean{Erdos257PeriodNoncollapse.DiagonalFreshLossBridge.PowerTwoOddWindowAffine.powerTwoActualLcmTopEdgeResidueGapSupply_of_adjacentSuffixMidband}{TotientActualLcmTopEdgeStaircase.lean:2108}, direct endpoint
\lean{Erdos257PeriodNoncollapse.DiagonalFreshLossBridge.PowerTwoOddWindowAffine.irrational_totientSeries_of_adjacentSuffixMidbandSupply}{TotientActualLcmTopEdgeStaircase.lean:2191}.

\item[(ii)] $\mathrm{PowerTwoOddGuardTopEdgeHalfWordBandSupply}$: at odd
depth $m=2q+1$ the adjacent-suffix residue is exactly twice a half-word
residue, and both directed edge widths halve to the same threshold
$H+q+2$:
\[
  H+q+2 \le \bigl(\text{half-word correction word}\bigr) \bmod 4^q
  \le 4^q - (H+q+2).
\]
\lean{Erdos257PeriodNoncollapse.DiagonalFreshLossBridge.PowerTwoOddWindowAffine.PowerTwoOddGuardTopEdgeHalfWordBandSupply}{TotientActualLcmTopEdgeStaircase.lean:1349}
\ev{Open} \scale{cofinal} \coord{mobius-mersenne}. This is a substantially
weaker demand than the earlier fixed $1/32$ central band elsewhere in the
corpus. Sufficiency:
\lean{Erdos257PeriodNoncollapse.DiagonalFreshLossBridge.PowerTwoOddWindowAffine.powerTwoAdjacentSuffixMidbandSupply_of_oddGuardTopEdgeHalfWordBand}{TotientActualLcmTopEdgeStaircase.lean:2057}.

\item[(iii)] $\mathrm{PowerTwoActualFinalTopEdgeMagnitudeSupply}$: the exact
centered-lift restatement of (ii),
\[
  H+q+2 \le |\mathrm{actualOddHalfCenteredLift}(a,q)|,
\]
\lean{Erdos257PeriodNoncollapse.DiagonalFreshLossBridge.PowerTwoOddWindowAffine.PowerTwoActualFinalTopEdgeMagnitudeSupply}{TotientActualLcmTopEdgeStaircase.lean:1471}
\ev{Open} \scale{cofinal} \coord{mobius-mersenne}. Proved
\textbf{equivalent} (an iff, not merely sufficient) to (ii):
\[
  \mathrm{PowerTwoOddGuardTopEdgeHalfWordBandSupply}
  \iff
  \mathrm{PowerTwoActualFinalTopEdgeMagnitudeSupply}.
\]
\lean{Erdos257PeriodNoncollapse.DiagonalFreshLossBridge.PowerTwoOddWindowAffine.topEdgeHalfWordBandSupply_iff_actualFinalCenteredMagnitudeSupply}{TotientActualLcmTopEdgeStaircase.lean:1479}
\quad \ev{Lean}.

\item[(iv)] $\mathrm{PowerTwoFlexibleActualTopEdgeMagnitudeSupply}$: the
same magnitude test at an \emph{arbitrary} odd rank whose adjacent depth
still fits the corridor (not forced to the canonical guarded depth):
\lean{Erdos257PeriodNoncollapse.DiagonalFreshLossBridge.PowerTwoOddWindowAffine.PowerTwoFlexibleActualTopEdgeMagnitudeSupply}{TotientActualLcmTopEdgeStaircase.lean:1927}
\ev{Open} \scale{cofinal} \coord{mobius-mersenne}. Sufficiency routes back
through (i), not through the corridor-escape chain below:
\lean{Erdos257PeriodNoncollapse.DiagonalFreshLossBridge.PowerTwoOddWindowAffine.powerTwoAdjacentSuffixMidbandSupply_of_flexibleActualTopEdgeMagnitude}{TotientActualLcmTopEdgeStaircase.lean:2007}, direct endpoint
\lean{Erdos257PeriodNoncollapse.DiagonalFreshLossBridge.PowerTwoOddWindowAffine.irrational_totientSeries_of_flexibleActualTopEdgeMagnitudeSupply}{TotientActualLcmTopEdgeStaircase.lean:2213}.

\item[(v)] $\mathrm{PowerTwoFlexibleActualTerminalDominanceSupply}$: a
strictly \emph{one-sided} version of (iv) — only the upper comparison, no
absolute value:
\[
  \mathrm{diagonalWindowIncrement}(2^a,\,2q{+}2) \le 2\cdot\mathrm{actualOddHalfCenteredLift}(a,q).
\]
\lean{Erdos257PeriodNoncollapse.DiagonalFreshLossBridge.PowerTwoOddWindowAffine.PowerTwoFlexibleActualTerminalDominanceSupply}{TotientActualLcmTopEdgeStaircase.lean:1908}
\ev{Open} \scale{cofinal} \coord{mobius-mersenne}. Direct endpoint
\lean{Erdos257PeriodNoncollapse.DiagonalFreshLossBridge.PowerTwoOddWindowAffine.irrational_totientSeries_of_flexibleActualTerminalDominanceSupply}{TotientActualLcmTopEdgeStaircase.lean:2221}.
\end{enumerate}

\paragraph{The weakest known link, and the exact identity pinning it.}
A sixth predicate, strictly weaker again than (v) — it is the disjunction of
(v)'s inequality with the opposite-direction escape — is the true minimum
of the whole cluster:
\[
  \mathrm{PowerTwoFlexibleActualTerminalCarryCorridorEscapeSupply} :\equiv
  \forall a_0,\ \exists a,q,\quad \cdots\ \wedge\
  \Bigl(2u \le d - B\ \vee\ d \le 2u\Bigr),
\]
where $d=\mathrm{diagonalWindowIncrement}(2^a,2q{+}2)$,
$u=\mathrm{actualOddHalfCenteredLift}(a,q)$,
$B=2H+(2q{+}1)+2$.
\lean{Erdos257PeriodNoncollapse.DiagonalFreshLossBridge.PowerTwoOddWindowAffine.PowerTwoFlexibleActualTerminalCarryCorridorEscapeSupply}{TotientActualLcmTopEdgeStaircase.lean:1895}
\ev{Open} \scale{cofinal} \coord{mobius-mersenne}. Direct endpoint
\lean{Erdos257PeriodNoncollapse.DiagonalFreshLossBridge.PowerTwoOddWindowAffine.irrational_totientSeries_of_flexibleActualTerminalCarryCorridorEscapeSupply}{TotientActualLcmTopEdgeStaircase.lean:2230}.

This form is not arbitrary: under integrality it is exactly one side of a
proved identity,
\[
  2\cdot\mathrm{actualOddHalfCenteredLift}(a,q)
  = \mathrm{diagonalWindowIncrement}(2^a,2q{+}2) - \mathrm{carryOrbit}(H,H,z,2q{+}1),
\]
\lean{Erdos257PeriodNoncollapse.DiagonalFreshLossBridge.PowerTwoOddWindowAffine.two_mul_actualOddHalfCenteredLift_eq_terminal_sub_trueCarry}{TotientActualLcmTopEdgeStaircase.lean:1678}
\quad \ev{Lean} \quad \scale{bounded} \quad \coord{mobius-mersenne} — an
\emph{equality}, not an inequality: the doubled centered-lift state is
exactly (terminal arithmetic letter) minus (true carry). Escaping either
side of this named open interval is exactly the corridor-escape supply.

Crucially, the sign machinery has already eliminated the branch this
identity most naturally supplies:
\lean{Erdos257PeriodNoncollapse.DiagonalFreshLossBridge.PowerTwoOddWindowAffine.actualLcm_trueEndpointSurvivor_neg}{TotientActualLcmOrbitSign.lean:172}
(\ev{Lean}, \scale{bounded}) proves that \emph{under integrality} the true
carry orbit is strictly \emph{positive} throughout the corridor, which
means the dominance branch (v) is the one the identity naturally produces
positive evidence \emph{against}, and the surviving open target is really
the \emph{lower}-escape branch,
$2u \le d - B$ — the dominance branch itself is stated separately as
\lean{Erdos257PeriodNoncollapse.DiagonalFreshLossBridge.PowerTwoOddWindowAffine.actualLcmTailOrbit_notMem_int_of_actualTerminalDominance}{TotientActualLcmTopEdgeStaircase.lean:1880}. This is the weakest exact
finite arithmetic normal form currently formalised in this corridor cluster,
not the programme's leading theorem interface: the direct fixed-full-block
first-harmonic gap above remains the constitutional outward socket.  The
corpus: a fully closed-form real-integer quantity whose escape from a named
interval is necessary and sufficient (mod a fit hypothesis
$2(H{+}q{+}2)\le 4^q$) for non-integrality at that rank.

\paragraph{A second, independent finite-window target: the short-arithmetic-kill supply.}
A separate reduction ties the diagonal form (the collapsing-free-parameters chain above) to this actual-orbit
lane via an exact digit formula with no cleanliness hypothesis at all,
$\mathrm{lcmRayArithmeticLetter}(t,j) = \varphi(H_t+j)-\varphi(H_t)$ for
\emph{every} offset $j$
(\lean{Erdos257PeriodNoncollapse.DiagonalFreshLossBridge.PowerTwoOddWindowAffine.lcmRayArithmeticLetter_eq_deltaTotient}{TotientActualLcmOrbitArithmetic.lean:298},
\ev{Lean}), and
$\mathrm{LcmDiagonalArithmeticKill}(t,L)\iff\mathrm{certifiedKill}(H_t,H_t,L)$
(\lean{Erdos257PeriodNoncollapse.DiagonalFreshLossBridge.PowerTwoOddWindowAffine.lcmDiagonalArithmeticKill_iff_certifiedKill}{TotientActualLcmOrbitArithmetic.lean:2045},
\ev{Lean}).

Consequently
\[
  \mathrm{PowerTwoActualLcmShortArithmeticKillSupply} :\equiv
  \forall a_0,\ \exists a,L,\quad a_0\le a\ \wedge\ L<2\cdot2^a\ \wedge\
  \mathrm{certifiedKill}(H_{2^a}, H_{2^a}, L)
\]
\lean{Erdos257PeriodNoncollapse.DiagonalFreshLossBridge.PowerTwoOddWindowAffine.PowerTwoActualLcmShortArithmeticKillSupply}{TotientActualLcmOrbitArithmetic.lean:2107}
\ev{Open} \scale{cofinal} \coord{mobius-mersenne} is \emph{literally the
diagonal obligation $P(t)$ of the collapsing-free-parameters chain above, restricted to powers of two and to a
short window $L<2\cdot2^a$}. It is proved sufficient for the actual-orbit
iff above via
\lean{Erdos257PeriodNoncollapse.DiagonalFreshLossBridge.PowerTwoOddWindowAffine.actualLcmOrbitNonintegralitySupply_of_shortArithmeticKill}{TotientActualLcmOrbitArithmetic.lean:2152},
and it is verified at exactly two exponents, $a=4$ (depth $L=23$) and
$a=6$ (depth $L=93$), each traced to a pre-existing compressed diagonal
certificate:
\lean{Erdos257PeriodNoncollapse.DiagonalFreshLossBridge.PowerTwoOddWindowAffine.actualLcmTailOrbit_four_and_six_notMem_int}{TotientActualLcmShortKill.lean:35}
\ev{Cert} \scale{fixed}.

The bounded stub
\lean{Erdos257PeriodNoncollapse.DiagonalFreshLossBridge.PowerTwoOddWindowAffine.powerTwoActualLcmShortArithmeticKillSupply_through_six}{TotientActualLcmShortKill.lean:54}
(\ev{Lean}, \scale{bounded}, proof is a single hard-coded witness
$(a,L)=(6,93)$, with no argument in $a$ at all) proves the supply only for
$a_0\le 6$; extending it to a single further exponent (say $a=7$) is a
self-contained, independently interesting target, not a re-run.

\paragraph{A Diophantine-flavoured alternative: the raw-approximant separation supply.}
Unconditionally, for every $a,q$:
\[
  \bigl|\mathrm{actualLcmTailOrbit}(a) - \mathrm{actualLcmRawApprox}(a,q)\bigr|
  < \frac{4H+2(2q{+}1)+4}{2^{2q+2}},
\]
\lean{Erdos257PeriodNoncollapse.DiagonalFreshLossBridge.PowerTwoOddWindowAffine.abs_actualLcmTailOrbit_sub_rawApprox_lt}{TotientActualLcmOrbitSeparation.lean:141}
\quad \ev{Lean} \quad \scale{uniform} \quad \coord{mobius-mersenne}, where
$\mathrm{actualLcmRawApprox}(a,q)$ is an explicit finite computable rational.
This reduces the whole analytic problem to a finite question: a cofinal
$1/32$-separation of the raw approximant from every integer,
\[
  \mathrm{PowerTwoActualLcmOrbitSeparationSupply} :\equiv
  \forall a_0,\ \exists a\ge\max(2,a_0),\ \exists q,\quad
  \mathrm{oddGuardedCanonicalAdjacentSuffixDepth}(2^a)=2q{+}1\ \wedge\
  \forall z\in\mathbb{Z},\ \tfrac{1}{32}+\mathrm{errorRadius}(a,q)\le|\mathrm{actualLcmTailOrbit}(a)-z|,
\]
\ev{Open} \scale{cofinal} \coord{mobius-mersenne}, sufficient via
\lean{Erdos257PeriodNoncollapse.DiagonalFreshLossBridge.PowerTwoOddWindowAffine.irrational_totientSeries_of_actualLcmOrbitSeparationSupply}{TotientActualLcmOrbitSeparation.lean:305}
(\ev{Lean}). Note the depth $q$ is not a free search parameter: the supply
pins exactly one admissible depth per scale $a$ via
$\mathrm{oddGuardedCanonicalAdjacentSuffixDepth}$.

\subsection{A problem-agnostic normal form: the guard-cylinder compression}

Independently of which coordinate is used, \emph{any} tail-difference
certificate — for arbitrary $h,N,L$, not specific to \#249 — compresses to
a two-bit test at logarithmic depth:
\[
  \bigl(\exists L,\ \mathrm{certifiedKill}(h,N,L)\bigr)
  \iff
  \mathrm{GuardCylinderWitness}(h,N).
\]
\lean{Erdos257PeriodNoncollapse.DiagonalFreshLossBridge.PowerTwoOddWindowAffine.exists_certifiedKill_iff_guardCylinderWitness}{TotientActualLcmTopEdgeStaircase.lean:844}
\quad \ev{Lean} \quad \scale{uniform} \quad \coord{seam-integer}, where the
witness is $\exists s,b,\ \mathrm{certifiedKill}(h,N{+}s,b{+}1) \vee
(\text{room} \wedge \mathrm{DyadicMixedGuard}(D(h,N{+}s,b{+}2),b))$ at
$b=\lfloor\log_2(N{+}h{+}L{+}2)\rfloor+1$. This statement mentions no
Mersenne or totient structure whatsoever and is stated for arbitrary
$h,N,L$; it compresses the search space for \emph{any} tail-difference
non-integrality certificate — of any depth $L$, however large — down to a
socket at a logarithmic scale plus a two-bit mixed-guard cylinder
($01$ or $10$). This is directly reusable, unchanged, for any binary-series
tail-difference problem, including \#257's own denominators.

\subsection{One open Farey problem and one failed rank shortcut}

The Farey growth law below is a genuinely independent live obligation.  The
rank statement is retained only to record a counterfactual sufficient input
and the reason that the generic compression route to it is closed; it is not a
second open problem supported by the present argument.

\paragraph{A counterfactual rank criterion.}
Unconditionally, for every level $e$, the dyadic totient-kernel family of
$2^e+1$ channels is linearly independent over $\mathbb{Q}$ (via CRT and
Dirichlet's theorem on primes in arithmetic progression):
\lean{Erdos249257.linearIndependent_canonicalTotientKernelFamily}{TotientMahlerDefect.lean:935}
\quad \ev{Lean} \quad \scale{uniform} \quad \coord{other:carry-kernel-rank}.
Consequently, if $S$ is rational there is a tempered integral binary carry
orbit $u$ with
\[
  \forall e,\quad 2^e-1 \le \operatorname{rank}_{\mathbb{Q}}\operatorname{span}\bigl(\mathrm{canonicalCarryKernelFamily}(u,e)\bigr).
\]
\lean{Erdos249257.not_irrational_totientSeries_implies_unbounded_carryRank_unconditional}{TotientCarryKernelRigidity.lean:300}
\quad \ev{Lean} \quad \scale{uniform} \quad \coord{other:carry-kernel-rank}.

An opposite inequality — a rationality-side rank \emph{upper} bound — would
of course contradict the displayed lower bound:
\[
  \exists C,\quad \forall\text{ tempered integral binary carry orbit } u
  \text{ for a rational } S,\quad
  \forall e,\quad \operatorname{rank}_{\mathbb{Q}}\operatorname{span}\bigl(\mathrm{canonicalCarryKernelFamily}(u,e)\bigr) \le C
\]
(or any bound growing slower than $2^e-1$).  This is a logically sufficient
new theorem schema, not an open proposition isolated by the preceding
mathematics: the corpus gives no
reason rationality should force it, and explicit finite-rank
shift-polynomial countermodels show that periodic denominator data alone do
not.

The most direct generic construction is also impossible: no
\texttt{CompressedAdjointCertificate} ($Q\cdot v\cdot A =
\mathrm{boundary}$ with $|\mathrm{boundary}|<Q\cdot v$ and $A\ne0$) can
exist —
\lean{Erdos249257.false_of_compressedAdjointCertificate}{TotientMahlerDefect.lean:1183}
\quad \ev{Lean} — and the full family is proved infinite-dimensional in
span:
\lean{Erdos249257.not_finiteDimensional_span_fullTotientKernel}{TotientMahlerDefect.lean:1145}
\quad \ev{Lean}.

What rationality \emph{does} buy — uniform eventual
periodicity of $u$'s dyadic sections modulo some $v>0$ — is proved to
\emph{not} promote to any finite rank bound without further arithmetic
input:
\lean{Erdos249257.not_irrational_totientSeries_implies_mod_period_and_unbounded_rank}{TotientTailCarryPeriod.lean:224}
\quad \ev{Lean}. This pins the obstacle precisely, but supplies no rank
upper bound.

\paragraph{Obligation 3 (Farey growth law).}
The corpus's strongest \emph{unconditional} statement about $S$ comes from
a Farey-gap denominator exclusion at a single fixed window $K=240$:
\[
  \forall p\in\mathbb{Q},\quad p.\mathrm{den} \le 79639646646701375323355774875831053
  \ \Longrightarrow\ S\ne p.
\]
\lean{Erdos249257.tsum_totient_div_pow_two_ne_ratCast_of_den_le_79639646646701375323355774875831053}{CertificateKernel.lean:18384}
\quad \ev{Lean} \quad \scale{bounded} \quad \coord{farey}. This bound is
\emph{sharp} at $K=240$ — the mediant $q=79639646646701375323355774875831054$
is proved to be the exact first failing denominator
(\lean{Erdos249257.gap_check_window_1_240_first_failure}{GapFareyBound.lean:225},
\ev{Lean}) — so re-running the same window buys nothing further; a new $K$
needs a freshly committed totient residue and freshly computed
continued-fraction convergents, both hard-coded numerals.

The cofinal
upgrade would be: a proved growth law $g(K)\to\infty$ such that for every
$K$ the $(N{=}1,K)$ gap check passes for all $q\le g(K)$ — equivalently, a
lower bound on the convergent denominators of the underlying totient-window
constant, uniform in $K$. \ev{Open} \scale{cofinal} \coord{farey}.

One
natural strengthening is proved to be a hard ceiling, not a route to
unboundedness: at a prime-power reduced denominator, the ``unit-gap''
refinement can rescue \emph{at most one} additional lattice point beyond
the ordinary gap certificate —
\lean{Erdos249257.reducedDenominatorCandidateCount_prime_pow_le_one}{PrimitiveWeightCertificate.lean:111}
\quad \ev{Lean}. This route is closed; growing $K$ genuinely requires new
per-window computation, not a refinement of the existing one.

\subsection{Equivalence map}

\begin{center}
\begin{longtable}{p{0.30\linewidth} p{0.20\linewidth} p{0.42\linewidth}}
\textbf{Form} & \textbf{Relation to Obligation 1} & \textbf{Site} \\
\hline
\endhead
Sep$(h,N,L)$ supply (Obl.\ 1, base form) & \textbf{equivalent} to Irrational$(S)$ & LcmConeFlatness.lean:412 \\
multiple-period supply & propositionally equivalent via Obl.\ 1; explicit directions are Obl.\ 1 $\Rightarrow$ multiple and multiple $\Rightarrow$ irrationality & CarrySurvivorExtinction.lean:502 \\
lcm-diagonal supply $P(t)$ & \textbf{equivalent} to Irrational$(S)$; single free parameter & LcmConeFlatness.lean:426 \\
lcm-cone supply & propositionally equivalent via Obl.\ 1; $P(t)$ is the cell $q{=}m{=}1$ & CertificateKernel.lean:18686 \\
cone-menu (\texttt{coneNonflatCert}) supply & weaker; half the pairwise radius & CertificateKernel.lean:18812 \\
pure non-integrality (diagonal/cone) & \textbf{equivalent} to the corresp.\ certificate form, via the completeness iff & CertificateKernel.lean:18706, :18718 \\
DTWFirstHarmonicNormGap & sufficient for Obl.\ 1 (not shown weaker/stronger) & FirstHarmonicPivot.lean:83 \\
actual-orbit nonintegrality supply & \textbf{equivalent} to Irrational$(S)$ itself & TotientActualLcmOrbitNonintegrality.lean:37 \\
short-arithmetic-kill supply & special case of $P(t)$: powers of two, short window & TotientActualLcmOrbitArithmetic.lean:2107 \\
top-edge residue-gap supply & sufficient for actual-orbit supply; one-sided & TotientActualLcmTopEdgeStaircase.lean:1325 \\
adjacent-suffix midband supply & weaker than top-edge residue-gap & TotientActualLcmTopEdgeStaircase.lean:1334 \\
odd-guard half-word band supply & weaker again; \textbf{equivalent} to final-magnitude form & TotientActualLcmTopEdgeStaircase.lean:1349 \\
actual final top-edge magnitude supply & \textbf{equivalent} to odd-guard half-word band & TotientActualLcmTopEdgeStaircase.lean:1471 \\
flexible top-edge magnitude supply & weaker than midband; feeds midband, not corridor-escape & TotientActualLcmTopEdgeStaircase.lean:1927 \\
flexible terminal-dominance supply & weaker again; one-sided & TotientActualLcmTopEdgeStaircase.lean:1908 \\
flexible terminal carry-corridor-escape supply & \textbf{weakest known} in this cluster & TotientActualLcmTopEdgeStaircase.lean:1895 \\
raw-approximant separation supply & alternative (Diophantine) sufficient form, same target & TotientActualLcmOrbitSeparation.lean:305 \\
rank-compression upper bound (Obl.\ 2) & independent obligation, different coordinate & TotientCarryKernelRigidity.lean:284 \\
Farey growth law (Obl.\ 3) & independent obligation, different coordinate & GapFareyBound.lean, CertificateKernel.lean:18056 \\
guard-cylinder witness & problem-agnostic normal form of \emph{any} certificate, not scale-comparable & TotientActualLcmTopEdgeStaircase.lean:844 \\
\end{longtable}
\end{center}

\emph{Reading the table.} \emph{Weaker} in rows not marked
equivalent describes an explicit implication between producer predicates:
the stronger hypothesis constructs the easier-looking one.
\emph{Equivalent} marks a registered \texttt{iff};
\emph{propositionally equivalent} marks two proved directions obtained by
composing the displayed implications with the base iff.  In particular,
Obligation~1 and the actual-orbit supply are logically equivalent through
$\mathrm{Irrational}(S)$, although the corpus does not claim a direct
witness-to-witness converter between their predicates.  The
short-arithmetic-kill supply is one concrete special case of Obligation~1's
diagonal form.

\subsection{What new mathematics each surviving form would need}

None of the forms above is close to complete; every one reduces the open
content to a named finite-flavoured arithmetic statement, never removes it.
Two forms above have their missing ingredient identified precisely enough
to name what field of technique would supply it, without any claim that the
technique exists or is easy to apply.

\paragraph{The exponential-sum form.} \texttt{DTWFirstHarmonicNormGap} asks
for cancellation in the first additive character of a totient-driven
discrepancy: a constant-saving bound, uniform in the period $h$ and growing
with the block size $X$ and depth $L$, on
$\sum_{N\in[X,2X)} e\bigl(D(h,N,L)/2^L\bigr)$, where $D$ is built from
consecutive differences of $\varphi$. This is a Weyl-sum-shaped statement
about additive character sums of an arithmetic function on a window; no
argument of this shape is attempted or proved anywhere in the corpus.

\paragraph{The staircase form.} The actual-orbit top-edge cluster's surviving
target, in every one of its six equivalent-or-sufficient guises, reduces
under the exact identity of the actual-orbit lane above to control of a single \emph{terminal
arithmetic letter}: the value $\mathrm{diagonalWindowIncrement}(2^a,2q{+}2)
= \varphi(2H{+}2q{+}2)-\varphi(2H)$ measured against twice a centred
carry-lift state. What is needed is not a new certificate mechanism — the
whole apparatus of arithmetic letters, centred lifts, and carry orbits is
already exact and unconditional — but a genuine size or residue statement
about this one totient value at cofinally many scales $a$, sufficient to
force it off the interval the identity names. The corpus's own finite
census (through every $t\le82$ on the diagonal form, and $a=4,6$ on the short
arithmetic-kill form) gives no asymptotic evidence either way; it is a
floor, not a trend.

No claim is made that either obstruction is close to resolution, and none
of the material above decides Erdős \#249, \#257, or any weakening of them.

\clearpage
\section*{Statements and declarations}

\paragraph{Artefact and data availability.} The linked Lean declarations,
fixed toolchain, and library manifest are published in the companion
repository. This manuscript provides navigation and exposition rather than
proof authority.

\paragraph{Authority boundary.} Kernel checking establishes that a proposition
was proved; it does not authorise the exposition, literature judgements, or any
claim that Problem 249 is solved.

\paragraph{Funding and competing interests.} This work received no external
funding. The author declares no competing interests.

\end{document}
